\section{I.1.1. The ${\mathcal F}$-functional and its
monotonicity} \label{I.1.1}
\label{Derivation of first equation of Section 1.1}

The goal of this section is to show that in an appropriate sense, Ricci
flow is a gradient flow on the space of metrics.
We introduce the entropy functional ${\mathcal F}$.
We compute its formal variation and show that the corresponding
gradient flow is a modified Ricci flow.

In Sections \ref{I.1.1} through \ref{I.5} of these notes,  $M$
is a closed manifold.
We will use the Einstein summation convention freely. We also follow
Perelman's convention that a condition like $a > 0$ means that
$a$ should be considered to be a small parameter, while a condition
like $A < \infty$ means that $A$ should be considered to be a large
parameter. This convention is only for pedagogical purposes and
may be ignored by a logically minded reader.


Let ${\mathcal M}$ denote the space of smooth Riemannian metrics $g$ on
$M$. We think of ${\mathcal M}$ formally as an infinite-dimensional
manifold.  The tangent space $T_g {\mathcal M}$ consists of the
symmetric covariant $2$-tensors $v_{ij}$ on $M$. Similarly, $C^\infty(M)$
is an infinite-dimensional manifold with $T_f C^\infty(M) =
C^\infty(M)$.
The diffeomorphism group
$\Diff(M)$ acts on ${\mathcal M}$ and $C^\infty(M)$ by pullback.

Let $dV$ denote the Riemannian volume density associated to a metric $g$.
We use the convention that
$\triangle \: = \: \Div \: \grad$.

\begin{definition}[F-entropy definition]
  \label{def:kl-f-entropy-f-entropy-definition}
  \dcref{Definition 5.1}
  \group{Kleiner--Lott Entropy Functionals}
  \source{kleiner-lott}

The ${\mathcal F}$-functional ${\mathcal F} \: : \: {\mathcal M} \times C^\infty(M)
\rightarrow \R$ is given by
\begin{equation}
{\mathcal F}(g, f) \: = \:
\int_M \left( R \: + \: |\nabla f|^2 \right) \: e^{-f} \: dV .
\end{equation}
\end{definition}

Given $v_{ij} \in T_g{\mathcal M}$ and $h \in T_f C^\infty(M)$, the
evaluation of the differential $d{\mathcal F}$ on $(v_{ij}, h)$ is
written as $\delta {\mathcal F}(v_{ij}, h)$.
Put $v \: = \: g^{ij} \: v_{ij}$.
\begin{proposition}[first variation of F-entropy]
  \label{prop:kl-f-entropy-first-variation-f-entropy}
  \dcref{Proposition 5.3}
  \group{Kleiner--Lott Entropy Functionals}
  \source{kleiner-lott}
 (cf. I.1.1)
We have
\begin{equation}
\delta {\mathcal F}(v_{ij}, h) \: = \:
\int_M e^{-f} \: \left[ - \: v_{ij}
(R_{ij} \: + \: \nabla_i \nabla_j f) \: + \:
\left( \frac{v}{2} \: - \: h \right) \: (2 \triangle f \: - \:
|\nabla f|^2 \: + \: R) \right] \: dV.
\end{equation}
\end{proposition}
\par\medskip\noindent\textit{Source argument (not certified as complete).}\ \begingroup\itshape
From a standard formula,
\begin{equation}
\delta R \: = \: - \: \triangle v \: + \: \nabla_i \nabla_j v_{ij}
\: - \: R_{ij} v_{ij}.
\end{equation}
As
\begin{equation}
|\nabla f|^2 \: = \: g^{ij} \nabla_i f \nabla_j f,
\end{equation}
we have
\begin{equation}
\delta |\nabla f|^2 \: = \: - \: v^{ij} \nabla_i f \nabla_j f
\: + \: 2 \: \langle \nabla f, \nabla h \rangle.
\end{equation}
As $dV \: = \: \sqrt{\det(g)} \: dx_1 \ldots dx_n$, we have
$\delta (dV) \: = \: \frac{v}{2} \: dV$, so
\begin{equation} \label{measvar}
\delta \left( e^{-f} \: dV \right) \: = \: \left( \frac{v}{2} \: - \: h
\right) \: e^{-f} \: dV.
\end{equation}
Putting this together gives
\begin{align} \label{putting}
\delta {\mathcal F} \: = \:
\int_M \:  e^{-f} \: & \left[  - \: \triangle v \: + \:
\nabla_i \nabla_j v_{ij} \: - \: R_{ij} v_{ij} \: - \:
v_{ij} \nabla_i f \nabla_j f \: + \right. \\
& \left.  2 \: \langle \nabla f, \nabla h \rangle
\: + \: (R \: + \: |\nabla f|^2 ) \: \left( \frac{v}{2} \: - \: h \right)
\right] \:  dV. \notag
\end{align}

The goal now is to rewrite the right-hand side of
(\ref{putting}) so that $v_{ij}$ and $h$ appear
algebraically, i.e. without derivatives.
As
\begin{equation} \label{id}
\triangle e^{-f} \: = \: (|\nabla f|^2 \: - \: \triangle f) \: e^{-f},
\end{equation}
we have
\begin{equation}
\int_M e^{-f} \: [ - \: \triangle v ] \: dV \: = \:
- \int_M (\triangle e^{-f} )\: v \: dV \: = \:
\int_M e^{-f} \: (\triangle f \: - \: |\nabla f|^2 ) \: v \: dV.
\end{equation}
Next,
\begin{align}
\int_M e^{-f} \nabla_i \nabla_j v_{ij} \: dV \: & = \:
\int_M (\nabla_i \nabla_j e^{-f}) \: v_{ij} \: dV \: = \:
- \: \int_M \nabla_i ( e^{-f} \nabla_j f) \: v_{ij} \: dV \\
& = \:
\int_M e^{-f} (\nabla_i f \nabla_j f \: - \: \nabla_i \nabla_j f)
\: v_{ij} \: dV. \notag
\end{align}
Finally,
\begin{align}
2 \: \int_M e^{-f} \langle \nabla f, \nabla h \rangle \: dV \: & = \:
- \: 2 \: \int_M  \langle \nabla e^{-f}, \nabla h \rangle \: dV \: = \:
2 \: \int_M  (\triangle e^{-f}) \: h \: dV \\
& = \:
2 \: \int_M e^{-f} \:
(|\nabla f|^2 \: - \: \triangle f) \: h \: dV. \notag
\end{align}
Then
\begin{align} \label{Fvariation}
\delta {\mathcal F} \: = \:
\int_M e^{-f} \: & \left[
\left( \frac{v}{2} \: - \: h \right) (2 \triangle f \: - \:
2 |\nabla f|^2 ) \: - \: v_{ij} (R_{ij} \: + \: \nabla_i \nabla_j f)
\right. \\
&
\left.  + \: \left( \frac{v}{2} \: - \: h \right) \:
(R \: + \: |\nabla f|^2 )
\right] \: dV \notag \\
=  \: \int_M e^{-f} \: & \left[ - \: v_{ij}
(R_{ij} \: + \: \nabla_i \nabla_j f) \: + \:
\left( \frac{v}{2} \: - \: h \right) \: (2 \triangle f \: - \:
|\nabla f|^2 \: + \: R) \right] \: dV. \notag
\end{align}
This proves the proposition.
\endgroup\par\medskip

We would like to get rid of the $\left( \frac{v}{2} \: - \: h \right) \: (2 \triangle f \: - \:
|\nabla f|^2 \: + \: R)$ term in (\ref{Fvariation}).  We can do this by restricting
our variations so that $\frac{v}{2} \: - \: h \: = \: 0$. From
(\ref{measvar}), this amounts to assuming that assuming $e^{-f} \: dV$ is fixed.
We now fix a smooth measure $dm$ on $M$ and relate $f$ to $g$ by
requiring that $e^{-f} \: dV \: = \: dm$. Equivalently, we define
a section $s \: : \: {\mathcal M} \rightarrow {\mathcal M} \times C^\infty(M)$
by $s(g) \: = \: \left( g, \ln \left( \frac{dV}{dm}\right) \right)$.
Then the composition ${\mathcal F}^m \: = \:
{\mathcal F}\circ s$ is a function on ${\mathcal M}$ and its differential
is given by
\begin{equation}
d{\mathcal F}^m (v_{ij}) \: = \: \int_M e^{-f} \: \left[ - \: v_{ij}
(R_{ij} \: + \: \nabla_i \nabla_j f) \right] \: dV.
\end{equation}

Defining a formal Riemannian metric on ${\mathcal M}$ by
\begin{equation}
\langle v_{ij}, v_{ij} \rangle_{g} \: = \: \frac12
\int_M v^{ij} \: v_{ij} \: dm,
\end{equation}
the gradient flow of ${\mathcal F}^m$ on ${\mathcal M}$
is given by
\begin{equation} \label{Fflow}
(g_{ij})_t \: = \: -2 \: (R_{ij} \: + \: \nabla_i \nabla_j f).
\end{equation}
The induced flow equation for $f$ is
\begin{equation}
f_t \: = \: \frac{\partial}{\partial t} \: \ln \left( \frac{dV}{dm}\right)
\: = \: \frac12 \: g^{ij} \: (g_{ij})_t \: = \: - \: \triangle f \: - \: R.
\end{equation}
As with any gradient flow, the function ${\mathcal F}^m$ is nondecreasing
along the flow line with its derivative being given by the
length squared of the gradient, i.e.
\begin{equation} \label{Fchange}
{\mathcal F}^m_t \: = \: 2 \int_M |R_{ij} \: + \: \nabla_i \nabla_j f|^2 \: dm,
\end{equation}
as follows from (\ref{Fvariation}) and (\ref{Fflow})

We now perform time-dependent
diffeomorphisms to transform (\ref{Fflow}) into the
Ricci flow equation. If $V(t)$ is the time-dependent generating vector field
of the diffeomorphisms then the new equations for $g$ and $f$ become
\begin{align}
(g_{ij})_t \: & = \: -2 \: (R_{ij} \: + \: \nabla_i \nabla_j f) \: + \:
{\mathcal L}_V g, \\
f_t & \: = \: - \: \triangle f \: - \: R \: + \: {\mathcal L}_V f. \notag
\end{align}
Taking $V \: = \: \nabla f$ gives
\begin{align} \label{newevolution}
(g_{ij})_t \: & = \: -2 \: R_{ij}, \\
f_t & \: = \: - \: \triangle f \: - \: R \: + \: |\nabla f|^2. \notag
\end{align}
As ${\mathcal F}(g, f)$ is unchanged by a simultaneous pullback of $g$ and $f$,
and the right-hand side of (\ref{Fchange}) is also unchanged under a
simultaneous pullback, it follows that
under the new evolution equations (\ref{newevolution})
we still have
\begin{equation} \label{Fder}
\frac{d}{dt} \: {\mathcal F}(g(t), f(t)) \: = \:
2 \int_M |R_{ij} \: + \: \nabla_i \nabla_j f|^2 \: e^{-f} \: dV
\end{equation}
(This can also be checked directly).

Because of the diffeomorphisms that we applied,
$g$ and $f$ are no longer related by
$e^{-f} \: dV \: = \: dm$. We do have that
$\int_M e^{-f} \: dV$ is constant in $t$, as
$e^{-f} \: dV$ is related to $dm$ by a diffeomorphism.

The relation between $g$ and $f$ is as follows:
we solve (or assume that we have a solution for)
the first equation in (\ref{newevolution}), with some initial
metric.
Then given the solution $g(t)$, we require that $f$ satisfy
the second equation in (\ref{newevolution}) (which is in
terms of $g(t)$).

The second equation in
(\ref{newevolution}) can be written as
\begin{equation} \label{backwards}
\frac{\partial}{\partial t} \: e^{-f} \: = \: - \: \triangle e^{-f}
\: + \: R \: e^{-f}.
\end{equation}
As this is a backward heat equation, we cannot solve for $f$
forward in time starting with an arbitrary smooth function.
Instead, (\ref{newevolution}) will be applied by
starting with a solution for $(g_{ij})_t \: = \: -2 \: R_{ij}$ on
some time interval $[t_1, t_2]$ and then solving
(\ref{backwards}) backwards in time on $[t_1, t_2]$
(which can always be done) starting with some initial
$f(t_2)$. Having done this, the solution $(g(t), f(t))$ on $[t_1, t_2]$
will satisfy (\ref{Fder}).

\section{Basic example for I.1}

In this section we compute ${\mathcal F}$ in a Euclidean example.

Consider $\R^n$ with the standard metric, constant in time.
Fix $t_0 \: > \: 0$. Put $\tau \: = \: t_0 \: - \: t$ and
\begin{equation}
f(t, x) \: = \: \frac{|x|^2}{4\tau} \: + \: \frac{n}{2} \ln(4 \pi \tau),
\end{equation}
so
\begin{equation}
e^{-f} \: = \: (4 \pi \tau)^{-n/2} \: e^{- \frac{|x|^2}{4\tau}}.
\end{equation}
This is the standard heat kernel when considered for $\tau$ going from
$0$ to $t_0$, i.e. for $t$ going from $t_0$ to $0$.
One can check that $(g,f)$ solves (\ref{newevolution}).
As
\begin{equation} \label{Gaussianintegral}
\int_{\R^n} e^{- \frac{|x|^2}{4\tau}} \: dV \: = \: (4 \pi \tau)^{n/2},
\end{equation}
$f$ is properly normalized.
Then $\nabla f \: = \: \frac{{\bf x}}{2 \tau}$ and
$|\nabla f|^2 \: = \: \frac{|x|^2}{4 \tau^2}$.
Differentiating (\ref{Gaussianintegral}) with respect to $\tau$ gives
\begin{equation} \label{Gaussianintegral2}
\int_{\R^n} \frac{|x|^2}{4\tau^2} \:
e^{- \frac{|x|^2}{4\tau}} \: dV \: = \: (4 \pi \tau)^{n/2} \:
\frac{n}{2\tau},
\end{equation}
so
\begin{equation}
\int_{\R^n} |\nabla f|^2 \: e^{- f} \: dV \: = \: \frac{n}{2\tau}.
\end{equation}
Then ${\mathcal F}(t) \: = \: \frac{n}{2\tau} \: = \: \frac{n}{2(t_0-t)}$.
In particular, this is nondecreasing as a function of
$t \in [0, t_0)$.

\section{I.2.2. The $\lambda$-invariant and its applications}
\label{I.2.2}

In this section we define $\lambda(g)$ and show that it is nondecreasing under
Ricci flow.  We use this to show that a steady breather on a compact manifold
is a gradient steady soliton.

\begin{proposition}[uniqueness for F-entropy]
  \label{prop:kl-invariant-its-applications-uniqueness-f-entropy}
  \dcref{Proposition 7.1}
  \group{Kleiner--Lott Entropy Functionals}
  \source{kleiner-lott}

Given a metric $g$, there is a unique minimizer
$\overline{f}$ of ${\mathcal F}(g, f)$ under the constraint
$\int_M e^{-f} \: dV \: = \: 1$.
\end{proposition}
\begin{proof}
Write
\begin{equation}
{\mathcal F} \: = \: \int_M \left( R e^{-f} \: + \: 4 |\nabla e^{-f/2}|^2
\right) \: dV.
\end{equation}
Putting $\Phi \: = \: e^{- f/2}$,
\begin{equation}
{\mathcal F} \: = \: \int_M \left( 4 |\nabla \Phi|^2 \: + \: R \: \Phi^2
\right) \: dV \: = \:
\int_M \Phi \left( - \: 4 \triangle \Phi \: + \: R \Phi
\right) \: dV.
\end{equation}
The constraint equation becomes $\int_M \Phi^2 \: dV \: = \: 1$.
Then $\lambda$ is the smallest eigenvalue of
$- \: 4 \triangle \: + \: R$ and $e^{-\overline{f}/2}$
is a corresponding normalized
eigenvector.  As the operator is a Schr\"odinger operator, there is a
unique normalized positive eigenvector \cite[Chapter XIII.12]{Reed-Simon}.
\end{proof}

\begin{definition}[lambda invariant]
  \label{def:kl-invariant-its-applications-lambda-invariant}
  \dcref{Definition 7.4}
  \group{Kleiner--Lott Entropy Functionals}
  \source{kleiner-lott}

The $\lambda$-functional is given by
$\lambda(g) \: =\: {\mathcal F}(g, \overline{f})$.
\end{definition}

If $g(t)$ is a smooth family of metrics then it follows from eigenvalue
perturbation theory that $\lambda(g(t))$ and $\overline{f}(t)$ are
smooth in $t$ \cite[Chapter XII]{Reed-Simon}.

\begin{proposition}[monotonicity for entropy]
  \label{prop:kl-invariant-its-applications-monotonicity-entropy}
  \dcref{Proposition 7.5}
  \group{Kleiner--Lott Entropy Functionals}
  \source{kleiner-lott}
 (cf. I.2.2) \label{mono}
If $g(\cdot)$ is a Ricci flow solution then $\lambda(g(t))$ is
nondecreasing in $t$.
\end{proposition}
\begin{proof}
Consider a time
interval $[t_1, t_2]$, and the minimizer $\overline{f}(t_2)$. In
particular, $\lambda(t_2) \: = \: {\mathcal F}(g(t_2),
\overline{f}(t_2))$. Put $u(t_2) \: = \: e^{-\overline{f}(t_2)}$. Solve the
backward heat equation
\begin{equation}
\frac{\partial u}{\partial t} \: = \: - \: \triangle u \: + \: Ru
\end{equation}
backward on $[t_1, t_2]$.

We claim that
$u(x^\prime,t^\prime) > 0$
for all $x^\prime \in M$ and $t^\prime \in [t_1, t_2]$. To see this,
we take $t^\prime \in [t_1, t_2)$,
and let $h$ be the solution to the forward heat equation
$\frac{\partial h}{\partial t} \: = \: \triangle h$ on $(t^\prime, t_2]$
with $\lim_{t \rightarrow t^\prime} h(t) \: = \: \delta_{x^\prime}$.
We have
\begin{equation}
\frac{d}{dt} \: \int_M u(t) \: h(t) \: dV \: = \:
\int_M \left[ (\partial_t u \: + \: \triangle u \: - \: Ru) \: v \: + \: u \:
(\partial_t h \: - \: \triangle h) \right] \: dV \: = \: 0.
\end{equation}
One knows that $h(t) > 0$ for all $t \in (t^\prime, t_2]$. Then
\begin{equation}
u(x^\prime, t^\prime) \: = \: \int_M u(x, t^\prime) \:\delta_{x^\prime}(x) \:
dV(x) \: = \: \lim_{t \rightarrow t^\prime}
\int_M u(t) \: h(t) \:
dV \: = \:
\int_M u(t_2) \: h(t_2) \: dV \: > \: 0.
\end{equation}

For $t \in [t_1, t_2]$, define $f(t)$ by $u(t) \: = \: e^{-f(t)}$.
By (\ref{Fder}),
$
{\mathcal F}(g(t_1), f(t_1)) \: \le \:
{\mathcal F}(g(t_2), f(t_2))$. By the definition of
$\lambda$, $\lambda(t_1) \: = \: {\mathcal F}(g(t_1), \overline{f}(t_1))
\: \le \: {\mathcal F}(g(t_1), f(t_1))$. (We are using the fact that
$\int_M e^{-f(t_1)} \: dV(t_1) \: = \:
\int_M e^{-f(t_2)} \: dV(t_2) \: = \: 1$.)
Thus $\lambda(t_1) \: \le \: \lambda(t_2)$.
\end{proof}

\begin{definition}[steady breather definition]
  \label{def:kl-invariant-its-applications-steady-breather-definition}
  \dcref{Definition 7.9}
  \group{Kleiner--Lott Entropy Functionals}
  \source{kleiner-lott}

A {\em steady breather} is a Ricci flow solution on an interval
$[t_1, t_2]$ that satisfies the equation
$g(t_2) \: = \: \phi^* g(t_1)$ for some $\phi \in \Diff(M)$.
\end{definition}

Steady soliton solutions are steady breathers.

Again, we are assuming that $M$ is compact. The next result is not
essential for the sequel, but gives a good illustration of how a
monotonicity formula is used.

\begin{proposition}[rigidity of steady breathers]
  \label{prop:kl-invariant-its-applications-rigidity-steady-breathers}
  \dcref{Proposition 7.10}
  \group{Kleiner--Lott Entropy Functionals}
  \source{kleiner-lott}
  (cf. I.2.2) \label{steadybreather}
A steady breather is a gradient steady soliton.
\end{proposition}
\begin{proof}
  \uses{prop:kl-invariant-its-applications-monotonicity-entropy}
We have $\lambda(g(t_2)) \: = \: \lambda(\phi^* g(t_1)) \: = \:
\lambda(g(t_1))$. Thus we have equality in Proposition
\ref{mono}. Tracing through the proof,
${\mathcal F}(g(t), f(t))$ must be constant in $t$. From
(\ref{Fder}), $R_{ij} \: + \: \nabla_i \nabla_j f \: = \: 0$.
Then $R \: + \: \triangle f \: = \: 0$ and so
(\ref{newevolution}) becomes (\ref{steady}).
\end{proof}

One can sharpen Proposition \ref{mono}.

\begin{lemma}[differential inequality for entropy]
  \label{lem:kl-invariant-its-applications-differential-inequality-entropy}
  \dcref{Lemma 7.11}
  \group{Kleiner--Lott Entropy Functionals}
  \source{kleiner-lott}
 (cf. Proposition I.1.2) \label{2/3}
\begin{equation}
\frac{d\lambda}{dt} \: \ge \: \frac{2}{n} \: \lambda^2(t).
\end{equation}
\end{lemma}
\begin{proof}
  \uses{prop:kl-invariant-its-applications-monotonicity-entropy}
Given a time interval $[t_1, t_2]$, with the notation of the proof of Proposition \ref{mono} we have
\begin{align}
\lambda(t_1) \: & \le \: {\mathcal F}(g(t_1), f(t_1)) \: = \:
{\mathcal F}(g(t_2), f(t_2)) \: - \: 2 \: \int_{t_1}^{t_2} \int_M |R_{ij} + \nabla_i \nabla_j f|^2 \:
e^{-f} \: dV \: dt \\
& = \: \lambda(t_2)  \: - \: 2 \: \int_{t_1}^{t_2} \int_M |R_{ij} + \nabla_i \nabla_j f|^2 \:
e^{-f} \: dV \: dt. \notag
\end{align}
Then
\begin{equation}
\frac{d\lambda}{dt} \Big|_{t=t_2}
\: = \: \lim_{t_1 \rightarrow t_2^-} \frac{\lambda(t_2) - \lambda(t_1)}{t_2 - t_1} \: \ge \:
2 \int_M |R_{ij} \: + \: \nabla_i \nabla_j {f}|^2 \: e^{- {f}}
\: dV,
\end{equation}
where the right-hand side is evaluated at time $t_2$.

Hence for all $t$,
\begin{equation} \label{dlambdadt}
\frac{d\lambda}{dt} \: \ge \:
2 \int_M |R_{ij} \: + \: \nabla_i \nabla_j \overline{f}|^2 \: e^{-\overline{f}}
\: dV
\end{equation}
and so
\begin{equation}
\frac{d\lambda}{dt} \: \ge \: \frac{2}{n} \int_M (R + \triangle \overline{f})^2 \:
e^{- \overline{f}} \: dV.
\end{equation}
From the Cauchy-Schwarz inequality and the fact that
$\int_M e^{- \overline{f}} \: dV \: = \: 1$,
\begin{equation}
\left( \int_M (R + \triangle \overline{f}) \: e^{- \overline{f}} \: dV \right)^2
\: \le \: \int_M (R + \triangle \overline{f})^2 \:
e^{- \overline{f}} \: dV.
\end{equation}
Finally, (\ref{id}) gives
\begin{equation} \label{using}
\int_M (R + \triangle \overline{f}) \: e^{- \overline{f}} \: dV \: = \:
\int_M (R + | \nabla \overline{f}|^2 ) \: e^{- \overline{f}} \: dV \: = \:
{\mathcal F}(g(t), \overline{f}(t)) \: = \: \lambda(t).
\end{equation}
This proves the lemma.
\end{proof}
\section{I.2.3. The rescaled $\lambda$-invariant}

In this section we show the monotonicity of a scale-invariant
version of $\lambda$.  This will be used in Section \ref{II.8}.
We then show that an expanding breather on a compact manifold
is a gradient expanding soliton.

Put $\overline{\lambda}(g) \: = \: \lambda(g) \: V(g)^{\frac{2}{n}}$.
As $\overline{\lambda}$ is scale-invariant, it is constant in $t$
along a steady, shrinking or expanding soliton solution.

\begin{proposition}[monotonicity of the rescaled lambda invariant]
  \label{prop:kl-rescaled-invariant-monotonicity-rescaled-lambda-invariant}
  \dcref{Proposition 8.1}
  \group{Kleiner--Lott Entropy Functionals}
  \source{kleiner-lott}
 \label{expandingprop}  (cf. Claim of I.2.3)
If $g(\cdot)$ is a Ricci flow solution and
$\overline{\lambda}(g(t)) \: \le \: 0$ for some $t$ then
$\frac{d}{dt} \overline{\lambda}(g(t)) \: \ge \: 0$.
\end{proposition}
\begin{proof}
We have
\begin{equation}
\frac{d\overline{\lambda}}{dt} \: = \:
\frac{d\lambda}{dt} \: V(t)^{\frac{2}{n}} \: - \:
\frac{2}{n} \: V(t)^{\frac{2-n}{n}} \: \lambda(t) \: \int_M R \: dV.
\end{equation}
From (\ref{dlambdadt}),
\begin{equation}
\frac{d\overline{\lambda}}{dt} \: \ge \:
V(t)^{\frac{2}{n}} \: \left[
2 \int_M |R_{ij} \: + \:
\nabla_i \nabla_j \overline{f}|^2 \: e^{- \overline{f}} \:  dV
 \: - \:
\frac{2}{n} \: V(t)^{-1} \: \lambda(t) \: \int_M R \: dV \right].
\end{equation}
Using the spatially-constant function $\ln V(t)$ as a test function for ${\mathcal F}$
gives
\begin{equation} \label{ttest}
\lambda(t) \: \le \: V(t)^{-1} \: \int_M R \: dV.
\end{equation}
The assumption that $\lambda(t) \: \le \: 0$ gives
\begin{equation}
- \: \lambda(t)^2 \: \le \: - \: V(t)^{-1} \: \lambda(t) \: \int_M R \: dV
\end{equation}
and so
\begin{equation}
\frac{d\overline{\lambda}}{dt} \: \ge \:
V(t)^{\frac{2}{n}} \: \left[
2 \int_M |R_{ij} \: + \:
\nabla_i \nabla_j \overline{f}|^2 \: e^{- \overline{f}} \: dV
 \: - \:
\frac{2}{n} \: \lambda(t)^2 \right].
\end{equation}
Next,
\begin{equation}
|R_{ij} \: + \:
\nabla_i \nabla_j \overline{f}|^2 \: \ge \:
\left| R_{ij} \: + \:
\nabla_i \nabla_j \overline{f} \: - \: \frac{1}{n} \: (R + \triangle \overline{f}) \: g_{ij}
\right|^2 \: + \: \frac{1}{n} \: (R + \triangle \overline{f})^2.
\end{equation}
Using (\ref{using}), one obtains
\begin{align} \label{cs}
\frac{d\overline{\lambda}}{dt} \: \ge \: 2 \:
V(t)^{\frac{2}{n}} \: & \left[
\int_M \left| R_{ij} \: + \:
\nabla_i \nabla_j \overline{f} \: - \: \frac{1}{n} \: (R + \triangle \overline{f}) \:
g_{ij} \right|^2 \: e^{- \overline{f}} \: dV
\: + \: \frac{1}{n} \: \int_M (R + \triangle \overline{f})^2 \: e^{- \overline{f}} \: dV \right. \\
& \left. - \: \frac{1}{n} \: \left( \int_M (R + \triangle \overline{f}) \: e^{- \overline{f}} \: dV
\right)^2 \right]. \notag
\end{align}
As $\int_M e^{- \overline{f}} \: dV \: = \: 1$, the Cauchy-Schwarz inequality implies
that the right-hand side of (\ref{cs}) is nonnegative.
\end{proof}

\begin{corollary}[rigidity at constant rescaled lambda]
  \label{cor:kl-rescaled-invariant-rigidity-at-constant-rescaled-lambda}
  \dcref{Corollary 8.9}
  \group{Kleiner--Lott Entropy Functionals}
  \source{kleiner-lott}
 \label{expp}
If $\overline{\lambda}$ is a
constant nonpositive number on an interval $[t_1, t_2]$ then the
Ricci flow solution is a
gradient soliton.
\end{corollary}
\begin{proof}
From equation (\ref{cs}), we obtain that $R + \triangle \overline{f} \: = \: \alpha(t)$ for some
function $\alpha$ that is spatially constant, and
$R_{ij} \: + \: \nabla_i \nabla_j \overline{f} \: = \: \frac{\alpha(t)}{n} \: g_{ij}$.
Thus $g$ evolves by diffeomorphisms and dilations.
After a shift of the time parameter,
$\alpha(t)$ is proportionate to $t$
cf. \cite[Lemma 2.4]{Chow-Knopf}.
This is the gradient soliton equation of Appendix \ref{Riccisolitons}.
\end{proof}

\begin{definition}[expanding breather definition]
  \label{def:kl-rescaled-invariant-expanding-breather-definition}
  \dcref{Definition 8.10}
  \group{Kleiner--Lott Entropy Functionals}
  \source{kleiner-lott}

An {\em expanding breather} is a Ricci flow solution on $[t_1, t_2]$ that
satisfies $g(t_2) \: = \: c \: \phi^* g(t_1)$ for some
$c > 1$ and $\phi \in \Diff(M)$.
\end{definition}

Expanding soliton solutions are
expanding breathers.

Again, we are assuming that $M$ is
compact.

\begin{proposition}[rigidity of expanding breathers]
  \label{prop:kl-rescaled-invariant-rigidity-expanding-breathers}
  \dcref{Proposition 8.11}
  \group{Kleiner--Lott Entropy Functionals}
  \source{kleiner-lott}

An expanding breather is a gradient expanding soliton.
\end{proposition}
\begin{proof}
  \uses{cor:kl-rescaled-invariant-rigidity-at-constant-rescaled-lambda,prop:kl-rescaled-invariant-monotonicity-rescaled-lambda-invariant}
First, $\overline{\lambda}(t_2) \: = \: \overline{\lambda}(t_1)$.
As $V(t_2) \: > \: V(t_1)$, we must have
$\frac{dV}{dt} \: > \: 0$ for some $t \in [t_1, t_2]$.
From (\ref{ttest}), $\frac{dV}{dt} \: = \:
-  \int_M R \: dV \: \le \: - \: \lambda(t) \: V(t)$, so
$\overline{\lambda}(t)$ must be negative for some
$t \in [t_1, t_2]$.
Proposition \ref{expandingprop} implies that
$\overline{\lambda}(t_1) \: < \: 0$. Then as
$\overline{\lambda}(t_2) \: = \: \overline{\lambda}(t_1)$,
it follows that $\overline{\lambda}$ is a negative
constant on $[t_1, t_2]$. From Corollary \ref{expp}, the solution is a gradient expanding
soliton.
\end{proof}

\section{I.2.4. Gradient steady solitons on compact manifolds}

It was shown in Section \ref{I.2.2} that a steady breather on a compact manifold
is a gradient steady soliton.  We now show that it is in fact Ricci flat.
This was previously shown in \cite[Theorem 20.1]{Hamilton}.

\begin{proposition}[compactness for Ricci solitons]
  \label{prop:kl-gradient-steady-solitons-compact-manifolds-compactness-ricci-solitons}
  \dcref{Proposition 9.1}
  \group{Kleiner--Lott Entropy Functionals}
  \source{kleiner-lott}

A gradient steady solution on a compact manifold is Ricci flat.
\end{proposition}
\begin{proof}
  \uses{prop:kl-invariant-its-applications-monotonicity-entropy}
As we are in the equality case of Proposition \ref{mono},
the function $f(t)$ must be the minimizer of ${\mathcal F}(g(t), \cdot)$
for all $t$. That is,
\begin{equation}
(- \: 4 \triangle \: + \: R) \: e^{-\frac{f}{2}} \: = \: \lambda \:
e^{-\frac{f}{2}}
\end{equation}
for all $t$, where $\lambda$ is constant in $t$.
Equivalently, $2 \triangle f \: - \: |\nabla f|^2 \: + \: R \: = \:
\lambda$. As $R \: + \: \triangle f \: = \: 0$, we have
$\triangle f \: - \: |\nabla f|^2 \: = \: \lambda$. Then
$\triangle e^{-f} \: = \: - \lambda \: e^{-f}$. Integrating gives
$0 \: = \: \int_M \triangle e^{-f} \: dV \: = \: - \: \lambda \:
\int_M e^{-f} \: dV$, so $\lambda \: = \: 0$. Then
$0 \: = \: - \: \int_M e^{-f} \: \triangle e^{-f} \: dV \: = \:
\int_M \left| \nabla e^{-f} \right|^2 \: dV$,
so $f$ is constant and $g$ is Ricci flat.
\end{proof}

A similar argument shows that a gradient expanding soliton on a
compact manifold comes from an Einstein metric with negative
Ricci curvature.

\section{Ricci flow as a gradient flow}

We have shown in Section \ref{I.1.1} that the modified Ricci flow is the gradient flow
for the functional ${\mathcal F}^m$ on the space of metrics ${\mathcal M}$.
One can ask if the unmodified Ricci flow is a gradient flow. This turns out to be true
provided that one
considers it as a flow on the space ${\mathcal M}/\Diff(M)$.

As mentioned in II.8, Ricci flow is the gradient flow
for the function $\lambda$. More precisely, this statement  is valid on ${\mathcal M}/\Diff(M)$,
with the latter being equipped with an appropriate metric.
To see this, we first
consider $\lambda$ as a function on the space of metrics ${\mathcal M}$.
Here the formal Riemannian metric on ${\mathcal M}$ comes from saying
that for $v_{ij} \in T_g {\mathcal M}$,
\begin{equation} \label{riemmetric}
\langle v_{ij}, v_{ij} \rangle \: = \: \frac12 \:
\int_M v^{ij} \: v_{ij} \: \Phi^2 \: dV(g),
\end{equation}
where $\Phi \: = \: \Phi(g)$ is the unique
normalized positive eigenvector corresponding to $\lambda(g)$.

\begin{lemma}[gradient-flow formula for the lambda invariant]
  \label{lem:kl-ricci-flow-as-gradient-flow-gradient-flow-formula-lambda-invariant}
  \dcref{Lemma 10.2}
  \group{Kleiner--Lott Entropy Functionals}
  \source{kleiner-lott}

The formal gradient flow of $\lambda$
is
\begin{equation}
\frac{\partial g_{ij}}{\partial t} \: = \:
- \: 2 \: \left( R_{ij} \: - \: 2 \: \nabla_i \nabla_j \ln \Phi
\right).
\end{equation}
\end{lemma}
\begin{proof}
We set
\begin{equation}
\label{eqnlambdadef}
\lambda(g) \: = \: \inf_{f \in C^\infty(M)
\: : \: \int_M e^{-f} \: dV \: = \: 1}
{\mathcal F}(g, f).
\end{equation}
To calculate the variation in $\lambda$
due to a variation $\delta g_{ij} \: = \: v_{ij}$,
we let $h=\delta f$ be the variation induced by letting $f$ be the
minimizer in (\ref{eqnlambdadef}). Then
\begin{equation}
0=\delta \left(\int_M e^{-f}\: dV\right) = \int_M \left(\frac{v}{2}-h\right)\:e^{-f}\:dV.
\end{equation}
Now  equation
(\ref{Fvariation}) gives
\begin{align} \label{lambdavar}
\delta \lambda(v_{ij})
\: = \:
\int_M e^{-f} \: & \left[  - \: v_{ij}
(R_{ij} \: + \: \nabla_i \nabla_j f) \: + \right. \\
& \left. \left( \frac{v}{2} \: - \: h \right) \: (2 \triangle f \: - \:
|\nabla f|^2 \: + \: R) \right] \: dV. \notag
\end{align}
As $\Phi \: = \: e^{-\frac{f}{2}}$ satisfies
\begin{equation}
- \: 4 \: \triangle \Phi \: + \: R \Phi \: = \: \lambda \Phi,
\end{equation}
it follows that
\begin{equation} \label{feqn}
2 \: \triangle f \: - \: | \nabla f|^2 \: + \: R \: = \: \lambda.
\end{equation}
Hence the last term in (\ref{lambdavar}) vanishes, and
(\ref{lambdavar}) becomes
\begin{equation}
\delta \lambda(v_{ij})
\: = \:
- \: \int_M e^{-f} \: v_{ij}
(R_{ij} \: + \: \nabla_i \nabla_j f) \: dV.
\end{equation}
The corresponding gradient flow is
\begin{equation} \label{gradflow}
\frac{\partial g_{ij}}{\partial t} \: = \:
- \: 2 \: \left( R_{ij} \: + \: \nabla_i \nabla_j f
\right)\: = \:
- \: 2 \: \left( R_{ij} \: - \: 2 \: \nabla_i \nabla_j \ln \Phi
\right).
\end{equation}
\end{proof}

We note that it follows from (\ref{feqn}) that
\begin{equation}
\nabla_j \left( (R_{ij} \: + \: \nabla_i \nabla_j f) \: e^{-f} \right)
\: = \: 0.
\end{equation}
This implies that the gradient vector field of $\lambda$ is perpendicular
to the infinitesimal diffeomorphisms at $g$, as one would expect.

In the sense of  \cite{Bourguignon}, the quotient space ${\mathcal M}/\Diff(M)$
is a stratified
infinite-dimensional Riemannian manifold, with the strata corresponding
to the possible isometry groups $\Isom(M, g)$.
We give it the quotient Riemannian metric coming from
(\ref{riemmetric}).
The modified Ricci flow (\ref{gradflow}) on
${\mathcal M}$ projects to a flow on ${\mathcal M}/\Diff(M)$
that coincides with the projection of the unmodified Ricci flow
$\frac{dg}{dt} \: = \: - \: 2 \Ric$. The upshot is that  the
Ricci flow, as a flow on ${\mathcal M}/\Diff(M)$, is the gradient flow of
$\lambda$, the latter now being considered as a function on
${\mathcal M}/\Diff(M)$.

One sees an intuitive explanation for Proposition \ref{steadybreather}.
If a gradient flow on a finite-dimensional manifold has a periodic
orbit then it must be a fixed-point.  Applying this principle formally
to the Ricci flow on ${\mathcal M}/\Diff(M)$, one infers that a steady
breather only evolves by diffeomorphisms.

\section{The ${\mathcal W}$-functional}
\label{Basic Example for Section 3.1}

\begin{definition}[W-entropy definition]
  \label{def:kl-w-entropy-w-entropy-definition}
  \dcref{Definition 11.1}
  \group{Kleiner--Lott Entropy Functionals}
  \source{kleiner-lott}

The {\em ${\mathcal W}$-functional} ${\mathcal W} \: : \: {\mathcal M} \times C^\infty(M) \times
\R^+ \rightarrow \R$ is given by
\begin{equation}
{\mathcal W}(g, f, \tau) \: = \: \int_M \left[ \tau \left(
|\nabla f|^2 + R \right) \: + \: f \: - \: n \right] \:
(4\pi \tau)^{-\frac{n}{2}} \: e^{-f} \: dV.
\end{equation}
\end{definition}

The ${\mathcal W}$-functional
is a scale-invariant variant of ${\mathcal F}$. It
has the symmetries ${\mathcal W}(\phi^* g, \phi^* f, \tau) \: = \:
{\mathcal W}(g, f, \tau)$ for $\phi \in \Diff(M)$, and
${\mathcal W}(cg, f, c\tau) \: = \: {\mathcal W}(g, f, \tau)$
for $c > 0$. Hence it is constant in $t \: = \: - \: \tau$ along
a gradient shrinking soliton defined for $t \in (-\infty, 0)$,
as in Appendix \ref{Riccisolitons}. In this sense, ${\mathcal W}$ is
constant on gradient shrinking solitons just as ${\mathcal F}$ is
constant on gradient steady solitons.

As an example of a gradient shrinking soliton,
consider $\R^n$ with the flat metric, constant in time
$t \in (-\infty, 0)$. Put $\tau \: = \: - \: t$ and
\begin{equation}
f(t, x) \: = \: \frac{|x|^2}{4\tau},
\end{equation}
so
\begin{equation}
e^{-f} \: = \: e^{- \frac{|x|^2}{4\tau}}.
\end{equation}
One can check that $(g(t), f(t), \tau(t))$ satisfies
(\ref{neweqn}) and (\ref{Wnorm}).
Now
\begin{equation}
\tau ( |\nabla f|^2 \: + \: R) \: + \: f \: - \: n \: = \:
\tau \cdot \frac{|x|^2}{4 \tau^2} \: + \: \frac{|x|^2}{4\tau} \: - \: n
\: = \:  \frac{|x|^2}{2\tau} \: - \: n.
\end{equation}
It follows from (\ref{Gaussianintegral}) and (\ref{Gaussianintegral2})
that ${\mathcal W}(t) \: = \: 0$ for all $t$.

\section{I.3.1. Monotonicity of the ${\mathcal W}$-functional}

In this section
we compute the variation of ${\mathcal W}$, in analogy with
the computation in Section \ref{I.1.1} of the variation of ${\mathcal F}$.
We then show that a shrinking breather on a compact manifold
is a gradient shrinking soliton.

As in Section \ref{I.1.1}, we write $\delta g_{ij} \: = \: v_{ij}$ and $\delta f \: = \: h$.
Put $\sigma \: = \: \delta \tau$.

\begin{proposition}[first variation of W-entropy]
  \label{prop:kl-w-entropy-first-variation-w-entropy}
  \dcref{Proposition 12.1}
  \group{Kleiner--Lott Entropy Functionals}
  \source{kleiner-lott}

We have
\begin{align}
\delta {\mathcal W}(v_{ij}, h, \sigma) \: = \:
\int_M & \left[ \sigma  (R \: + \: |\nabla f|^2) \:
- \: \tau v_{ij} (R_{ij} \: + \: \nabla_i \nabla_j f) \: + \: h \: +
\right.  \\
& \left.
\left[ \tau (2 \triangle f \: - \: |\nabla f|^2 \: + \: R) \: + \: f
\: - \: n \right] \:
\left( \frac{v}{2} \: - \: h \: - \: \frac{n\sigma}{2\tau} \right) \right] \:
(4 \pi \tau)^{-n/2} \: e^{-f} \: dV. \notag
\end{align}
\end{proposition}
\begin{proof}
One finds
\begin{equation}
\delta \left( (4 \pi \tau)^{-n/2} \: e^{-f} \: dV \right) \: = \:
\left( \frac{v}{2} \: - \: h \: - \: \frac{n\sigma}{2\tau} \right) \:
(4 \pi \tau)^{-n/2} \: e^{-f} \: dV.
\end{equation}
Writing
\begin{equation}
{\mathcal W} \: = \: \int_M \left[ \tau (R \: + \: |\nabla f|^2) \: + \: f
\: - \: n \right] \: (4 \pi \tau)^{-n/2} \: e^{-f} \: dV
\end{equation}
we can use
(\ref{Fvariation}) to obtain
\begin{align}
\delta {\mathcal W} \: = \:
\int_M & \left[ \sigma  (R \: + \: |\nabla f|^2) \: + \:
\tau  \left( \frac{v}{2} \:  - \: h \: \right) (2 \triangle f \: - \:
2|\nabla f|^2) \: - \: \tau v_{ij} (R_{ij} \: + \: \nabla_i \nabla_j f) \: +
\right. \\
& \left. h \: + \:
\left[ \tau (R \: + \: |\nabla f|^2) \: + \: f
\: - \: n \right] \:
\left( \frac{v}{2} \: - \: h \: - \: \frac{n\sigma}{2\tau} \right) \right] \:
(4 \pi \tau)^{-n/2} \: e^{-f} \: dV.
\end{align}
Then (\ref{id}) gives
\begin{align}
\delta {\mathcal W} \:
= \:
\int_M & \left[ \sigma  (R \: + \: |\nabla f|^2) \:
- \: \tau v_{ij} (R_{ij} \: + \: \nabla_i \nabla_j f) \: + \: h \: +
\right. \\
& \left.
\left[ \tau (2 \triangle f \: - \: |\nabla f|^2 \: + \: R) \: + \: f
\: - \: n \right] \:
\left( \frac{v}{2} \: - \: h \: - \: \frac{n\sigma}{2\tau} \right) \right] \:
(4 \pi \tau)^{-n/2} \: e^{-f} \: dV. \notag
\end{align}
This proves the proposition.
\end{proof}

We now fix a smooth measure $dm$ on $M$ with mass $1$ and
relate $f$ to $g$ and $\tau$ by
requiring that $(4 \pi \tau)^{-n/2} \: e^{-f} \: dV \: =\: dm$. Then
$\frac{v}{2} \: - \: h \: - \: \frac{n\sigma}{2\tau} \: = \: 0$ and
\begin{equation} \label{Wvar}
\delta {\mathcal W} \: = \: \int_M  \left[ \sigma  (R \: + \: |\nabla f|^2) \:
- \: \tau v_{ij} (R_{ij} \: + \: \nabla_i \nabla_j f) \: + \: h \right] \:
(4 \pi \tau)^{-n/2} \: e^{-f} \: dV.
\end{equation}

We now consider $\frac{d{\mathcal W}}{dt}$ when
\begin{align} \label{oldeqn}
(g_{ij})_t \: & = \: - \: 2 \: (R_{ij} \: + \: \nabla_i \nabla_j f ), \\
f_t \: & = \: - \: \triangle f \: - \: R \: + \: \frac{n}{2\tau}, \notag \\
\tau_t \: & = \: -1. \notag
\end{align}
To apply (\ref{Wvar}), we put
\begin{align}
v_{ij} \: & = \: - \: 2 \: (R_{ij} \: + \: \nabla_i \nabla_j f ), \\
h \: & = \: - \: \triangle f \: - \: R \: + \: \frac{n}{2\tau}, \notag \\
\sigma \: & = \: -1. \notag
\end{align}
We do have
$\frac{v}{2} \: - \: h \: - \: \frac{n\sigma}{2\tau} \: = \: 0$.
Then from (\ref{Wvar}),
\begin{align} \label{Wtimeder}
\frac{d{\mathcal W}}{dt} \: = \:
\int_M & \left[- \:  (R \: + \: |\nabla f|^2)
\: + \: 2 \: \tau |R_{ij} \: + \: \nabla_i \nabla_j f|^2 \:
\right. \\
& -\left. \triangle f \: - \: R \: + \: \frac{n}{2\tau}
 \right] \:
(4 \pi \tau)^{-n/2} \: e^{-f} \: dV \notag \\
= \:
\int_M & \left[- \:  2(R \: + \: \triangle f)
\: + \: 2 \: \tau |R_{ij} \: + \: \nabla_i \nabla_j f|^2
\: + \: \frac{n}{2\tau}
 \right] \:
(4 \pi \tau)^{-n/2} \: e^{-f} \: dV \notag \\
= \: \int_M &
2 \: \tau |R_{ij} \: + \: \nabla_i \nabla_j f
\: - \: \frac{1}{2\tau} \: g_{ij}|^2 \:
(4 \pi \tau)^{-n/2} \: e^{-f} \: dV. \notag
\end{align}
Adding a Lie derivative to the right-hand side of
(\ref{oldeqn}) gives the new flow equations
\begin{align} \label{neweqn}
(g_{ij})_t \: & = \: - \: 2 \: R_{ij}, \\
f_t \: & = \: - \: \triangle f \: + \: |\nabla f|^2
\: - \: R \: + \: \frac{n}{2\tau}, \notag \\
\tau_t \: & = \: -1, \notag
\end{align}
with (\ref{Wtimeder}) still holding. We no longer have
$(4 \pi \tau)^{-n/2} \: e^{-f} \: dV \: = \: dm$, but we do have
\begin{equation} \label{Wnorm}
\int_M (4 \pi \tau)^{-n/2} \: e^{-f} \: dV \: = \: 1.
\end{equation}


We now want to look at the variational problem of
minimizing ${\mathcal W}(g,f,\tau)$ under the constraint
that $\int_M (4 \pi \tau)^{-n/2} \: e^{-f} \: dV \: = \: 1$.
We write
\begin{equation}
\mu(g, \tau) \: = \: \inf_f \{ {\mathcal W}(g,f,\tau) \: : \:
\int_M (4 \pi \tau)^{-n/2} \: e^{-f} \: dV \: = \: 1\}.
\end{equation}
Making the change of variable $\Phi \: = \: e^{- \frac{f}{2}}$, we are
minimizing
\begin{equation}
(4 \pi \tau)^{-n/2} \: \int_M \left[ \tau (4 |\nabla \Phi|^2 \: + \: R \Phi^2) \: - \:
2 \Phi^2 \log \Phi \: - \: n \: \Phi^2 \right] \: dV
\end{equation}
under the constraint $(4 \pi \tau)^{-n/2} \: \int_M \Phi^2 \: dV \: = \: 1$.
From  \cite[Section 1]{Rothaus} the infimum is finite and
there is a positive
continuous minimizer $\Phi$. It will be a weak solution of the
variational equation
\begin{equation}
\tau (- 4 \triangle + R) \Phi \: = \: 2\Phi \log \Phi \: + \:
(\mu(g, \tau) + n) \Phi.
\end{equation}
From elliptic theory, $\Phi$ is smooth. Then $f \: = \: - \: 2 \log \Phi$
is also smooth.

As in Section \ref{I.2.2}, it follows that
$\mu(g(t), t_0 - t)$ is nondecreasing in $t$ for a Ricci flow
solution, where $t_0$ is any fixed number and $t < t_0$.
If it is constant in $t$ then the solution must be
a gradient shrinking soliton that goes singular at time $t_0$.

\begin{definition}[shrinking breather definition]
  \label{def:kl-w-entropy-shrinking-breather-definition}
  \dcref{Definition 12.17}
  \group{Kleiner--Lott Entropy Functionals}
  \source{kleiner-lott}

A {\em shrinking breather} is a Ricci flow solution on $[t_1, t_2]$ that
satisfies $g(t_2) \: = \: c \: \phi^* g(t_1)$ for some
$c < 1$ and $\phi \in \Diff(M)$.
\end{definition}

Gradient shrinking soliton solutions
are shrinking breathers.

Again, we are assuming that $M$ is
compact.

\begin{proposition}[rigidity of shrinking breathers]
  \label{prop:kl-w-entropy-rigidity-shrinking-breathers}
  \dcref{Proposition 12.18}
  \group{Kleiner--Lott Entropy Functionals}
  \source{kleiner-lott}

A shrinking breather is a gradient shrinking soliton.
\end{proposition}
\begin{proof}
Put $t_0 \: = \: \frac{t_2 - c t_1}{1-c}$. Then if
$\tau_1 \: = \: t_0 - t_1$ and $\tau_2 \: = \: t_0 - t_2$, we
have $\tau_2 \: = \: c\tau_1$. Hence
\begin{equation}
\mu(g(t_2), \tau_2) \: = \: \mu \left( \frac{\tau_2}{\tau_1} \:
\phi^* g(t_1), \tau_2 \right) \: = \:
\mu \left(\phi^* g(t_1), \tau_1 \right)
 \: = \:
\mu \left(g(t_1), \tau_1 \right).
\end{equation}
It follows that the solution is a gradient shrinking soliton.
\end{proof}

\section{I.4. The no local collapsing theorem I}

In this section we prove the no local collapsing theorem.

\begin{definition}[locally collapsing definition]
  \label{def:kl-no-local-collapsing-theorem-i-locally-collapsing-definition}
  \dcref{Definition 13.1}
  \group{Kleiner--Lott Entropy Functionals}
  \source{kleiner-lott}

A smooth Ricci flow solution $g(\cdot)$ on a time interval
$[0, T)$ is said to be {\em locally collapsing} at $T$ if there
is a sequence of times $t_k \rightarrow T$ and a sequence
of metric balls $B_k = B(p_k, r_k)$ at times $t_k$ such that
$r_k^2/t_k$ is bounded, $|\Rm|(g(t_k)) \: \le \: r_k^{-2}$
in $B_k$ and $\lim_{k \rightarrow \infty}
r_k^{-n} \: \vol(B_k) \: = \: 0$.
\end{definition}

\begin{remark}[estimates for noncollapsing]
  \label{rem:kl-no-local-collapsing-theorem-i-estimates-noncollapsing}
  \dcref{Remark 13.2}
  \group{Kleiner--Lott Entropy Functionals}
  \source{kleiner-lott}

In the definition of noncollapsing, $T$ could be infinite.
This is why it is written that $r_k^2/t_k$ stays bounded, while
if $T \: < \: \infty$ then this is obviously the same as saying that
$r_k$ stays bounded.
\end{remark}

\begin{theorem}[noncollapsing theorem]
  \label{thm:kl-no-local-collapsing-theorem-i-noncollapsing-theorem}
  \dcref{Theorem 13.3}
  \group{Kleiner--Lott Entropy Functionals}
  \source{kleiner-lott}
 \label{noncollthm} (cf. Theorem I.4.1)
If $M$ is closed and $T < \infty$ then $g(\cdot)$ is not locally
collapsing at $T$.
\end{theorem}
\begin{proof}
We first sketch the idea of the proof.
In Section \ref{Basic Example for Section 3.1} we showed that in the
case of flat $\R^n$, taking $e^{-f}(x) \: = \:
e^{- \frac{|x|^2}{4\tau}}$, we get ${\mathcal W}(g, f, \tau) \: = \: 0$.
So putting $\tau \: = \: r_k^2$ and $e^{-f_k}(x) \: = \:
e^{- \frac{|x|^2}{4r_k^2}}$, we have ${\mathcal W}(g, f_k, r_k^2) \: = \: 0$.
In the collapsing case, the idea is to use a test function $f_k$ so that
\begin{equation}
e^{-f_k}(x) \: \sim \:
e^{- c_k} \: e^{- \frac{\dist_{t_k}(x, p_k)^2}{4r_k^2}},
\end{equation}
where $c_k$ is determined by the normalization condition
\begin{equation} \label{normalization}
\int_M (4 \pi r_k^2)^{- n/2} \: e^{-f_k} \: dV \: = \: 1.
\end{equation}
The main difference between computing (\ref{normalization}) in $M$ and
in $\R^n$ comes from the difference in volumes, which means that
$e^{-c_k} \sim \frac{1}{r_k^{-n} \: vol(B_k)}$.
In particular, as $k \rightarrow
\infty$, we have $c_k \rightarrow - \infty$.

Now that $f_k$ is normalized correctly, the main difference between
computing ${\mathcal W}(g(t_k), f_k, r_k^2)$ in $M$, and the
analogous computation for the
Gaussian in $\R^n$,
comes from the $f$ term in the integrand of ${\mathcal W}$.  Since
$f_k \sim c_k$, this will drive ${\mathcal W}(g(t_k), f_k, r_k^2)$ to
$- \infty$ as $k \rightarrow \infty$, so $\mu(g(t_k), r_k^2)
\rightarrow - \infty$; by the monotonicity of $\mu(g(t),t_0-t)$
it follows that $\mu(g(0),t_k+r_k^2)\ra -\infty$ as $k\ra \infty$.
This contradicts the fact that $\mu(g(0),\tau)$ is a continuous function
of $\tau$.

To write this out precisely, let us put $\Phi \: = \: e^{- f/2}$, so
that
\begin{equation}
{\mathcal W}(g, f, \tau) \: = \: (4 \pi \tau)^{-n/2} \:
\int_M \left[ 4 \tau \: |\nabla \Phi|^2 \: + \:
( \tau R \: - \: 2 \ln \Phi
\: - \: n ) \: \Phi^2 \right] \: dV.
\end{equation}
For the argument, it is enough to obtain small values of ${\mathcal W}$ for
positive $\Phi$.   Since $\lim_{s\ra 0}(-2\ln s)\:s^2\:=\:0$,
by an approximation it is enough to obtain small values of ${\mathcal W}$ for
nonnegative $\Phi$, where the integrand is declared to be
$4 \tau \: |\nabla \Phi|^2$ at points where $\Phi$ vanishes. Take
\begin{equation}
\Phi_k(x) \: = \: e^{- c_k/2} \: \phi(\dist_{t_k}(x, p_k)/r_k),
\end{equation}
where $\phi \: : \: [0, \infty) \rightarrow [0, 1]$ is
a monotonically nonincreasing function such that
$\phi(s) \: = \: 1$ if $s \in [0, 1/2]$, $\phi(s) \: = \: 0$ if $s \: \ge \: 1$
and $|\phi^\prime(s)| \: \le \: 10$ for $s \in [1/2, 1]$.
The function $\Phi_k$ is {\it a priori} only Lipschitz, but by smoothing
it slightly we can use $\Phi_k$ in the variational formula
to bound ${\mathcal W}$ from above.

The constant $c_k$ is determined by
\begin{equation}
e^{c_k} \: = \: \int_M (4 \pi r_k^2)^{-n/2} \:
\phi^2(\dist_{t_k}(x, p_k)/r_k) \: dV \: \le \:
(4 \pi r_k^2)^{-n/2} \:
\vol(B_k).
\end{equation}
Thus $c_k \rightarrow - \infty$.
Next,
\begin{equation} \label{mess}
{\mathcal W}(g(t_k), f_k, r_k^2) \: = \: (4 \pi r_k^2)^{-n/2} \:
\int_M \left[ 4 r_k^2 \: |\nabla \Phi_k|^2 \: + \:
( r_k^2 R \: - \: 2 \ln \Phi_k
\: - \: n ) \: \Phi_k^2 \right] \: dV.
\end{equation}
Let $A_k(s)$ be the mass of the distance sphere $S(p_k, r_k s)$ around $p_k$.
Put
\begin{equation}
\overline{R}_k(s) \: = \: r_k^{2} \: A_k(s)^{-1}
\: \int_{S(p_k, r_k s)} R \: d\area.
\end{equation}
We can compute the integral in (\ref{mess}) radially to get
\begin{equation}
{\mathcal W}(g(t_k), f_k, r_k^2) \:
= \:  \frac{\int_0^1 \left[ 4 (\phi^\prime(s))^2 \: + \:
(\overline{R}_k(s) \: + \: c_k \: - \: 2 \: \ln \phi(s) \:
\: - \: n ) \: \phi^2(s) \right] \: A_k(s) \:
ds}{\int_0^1 \phi^2(s) \: A_k(s) \: ds}.
\end{equation}
The expression
$4 (\phi^\prime(s))^2 \: - \: 2 \: \ln \phi(s) \: \phi^2(s)$ vanishes
if $s \notin [1/2, 1]$, and is bounded above by $400 \: + \: e^{-1}$ if
$s \in [1/2, 1]$. Then the lower
bound on the Ricci curvature and the Bishop-Gromov
inequality give
\begin{align}
\frac{\int_0^1 \left[ 4 (\phi^\prime(s))^2 \: - \: 2 \: \ln \phi(s)
\: \phi^2(s) \right] \: A_k(s) \:
ds}{\int_0^1 \phi^2(s) \: A_k(s) \: ds} \: & \le \: 401 \:
\frac{\vol(B(p_k, r_k)) - \vol(B(p_k, r_k/2))}{\vol(B(p_k, r_k/2))} \\
& \le \: 401 \: \left(
\frac{\int_0^1 \sinh^{n-1}(s)
\: ds}{\int_0^{1/2} \sinh^{n-1}(s) \: ds} \: - \: 1 \right).
\notag
\end{align}
Next, from the upper bound on scalar curvature,
$\overline{R}_k(s) \: \le \: n(n-1)$ for $s \in [0,1]$.
Putting this together gives
${\mathcal W}(g(t_k), f_k, r_k^2) \: \le \: \const \: + \: c_k$ and
so ${\mathcal W}(g(t_k), f_k, r_k^2) \rightarrow -
\infty$ as $k \rightarrow \infty$.

Thus $\mu(g(t_k), r_k^2) \rightarrow - \infty$.
For any $t_0>t$, $\mu(g(t), t_0 - t)$ is
nondecreasing in $t$. Hence $\mu(g(0), t_k \: + \: r_k^2) \: \le \:
\mu(g(t_k), r_k^2)$, so $\mu(g(0), t_k \: + \: r_k^2) \rightarrow
- \infty$. Since $T$ is finite, $t_k$ and $r_k^2$ are
uniformly bounded,
and $t_k$ uniformly positive,
which
contradicts the fact that $\mu(g(0),\tau)$ is a continuous function of $\tau$.
\end{proof}

\begin{remark}[estimates for scalar curvature]
  \label{rem:kl-no-local-collapsing-theorem-i-estimates-scalar-curvature}
  \dcref{Remark 13.13}
  \group{Kleiner--Lott Entropy Functionals}
  \source{kleiner-lott}

In the preceding argument we only used the upper bound on scalar
curvature and the lower bound on Ricci
curvature, i.e. in the definition
of local collapsing one could have assumed that
$R(g_{ij}(t_k)) \: \le \: n(n-1) \: r_k^{-2}$ in $B_k$ and
$\Ric(g_{ij}(t_k)) \: \ge \: - \: (n-1) \: r_k^{-2}$ in $B_k$.
In fact, one can also remove the lower bound on Ricci curvature
(observation of Perelman, communicated by Gang Tian).
The necessary ingredients of the preceding argument were that \\
1. $r_k^{-n} \: \vol(B(p_k, r_k)) \rightarrow 0$, \\
2. $r_k^2 \: R$ is uniformly bounded above on $B(p_k, r_k)$ and \\
3. $\frac{\vol(B(p_k, r_k))}{\vol(B(p_k, r_k/2))}$ is uniformly bounded above.

Suppose only that $r_k^{-n} \: \vol(B(p_k, r_k)) \rightarrow 0$ and
for all $k$, $r_k^2 \: R \: \le \: n(n-1)$ on $B(p_k, r_k)$. If
$\frac{\vol(B(p_k, r_k))}{\vol(B(p_k, r_k/2))} \: < \: 3^n$ for all
$k$ then we are done.  If not, suppose that for a given $k$,
$\frac{\vol(B(p_k, r_k))}{\vol(B(p_k, r_k/2))} \: \ge \: 3^n$.
Putting $r_k^\prime \: = \: r_k/2$, we have that
$(r_k^\prime)^{-n} \: \vol(B(p_k, r_k^\prime)) \: \le \:
r_k^{-n} \: \vol(B(p_k, r_k))$ and
$(r_k^\prime)^2 \: R \: \le \: n(n-1)$ on $B(p_k, r_k^\prime)$.
We replace $r_k$ by $r_k^\prime$. If now
$\frac{\vol(B(p_k, r_k))}{\vol(B(p_k, r_k/2))} \: < \: 3^n$ then we
stop. If not then we repeat the process and replace $r_k$ by $r_k/2$.
Eventually we will achieve that
$\frac{\vol(B(p_k, r_k))}{\vol(B(p_k, r_k/2))} \: < \: 3^n$.
Then we can apply the preceding argument to this new sequence of pairs
$\{(p_k, r_k)\}_{k=1}^\infty$.
\end{remark}

\begin{definition}[kappa noncollapsing]
  \label{def:kl-no-local-collapsing-theorem-i-kappa-noncollapsing}
  \dcref{Definition 13.14}
  \group{Kleiner--Lott Entropy Functionals}
  \source{kleiner-lott}
 \label{noncollapsed1} (cf. Definition I.4.2)
We say that a metric $g$ is {\em $\kappa$-noncollapsed on the scale
$\rho$} if every metric ball $B$ of radius $r < \rho$, which satisfies
$|\Rm(x)| \: \le \: r^{-2}$ for every $x \in B$, has volume at least
$\kappa r^n$.
\end{definition}

\begin{remark}[structural remark on noncollapsing]
  \label{rem:kl-no-local-collapsing-theorem-i-structural-remark-noncollapsing}
  \dcref{Remark 13.15}
  \group{Kleiner--Lott Entropy Functionals}
  \source{kleiner-lott}

We caution the reader that this definition differs slightly from the
definition of noncollapsing that is used from section I.7 onwards.
\end{remark}

Note that except for the overall scale $\rho$, the $\kappa$-noncollapsed
condition is scale-invariant.
From the proof of Theorem \ref{noncollthm} we extract the following
statement. Given a Ricci flow defined on
an interval $[0, T)$, with $T < \infty$, and a scale $\rho$, there
is some number $\kappa = \kappa(g(0), T, \rho)$ so that the solution is
$\kappa$-noncollapsed on the scale $\rho$ for all $t \in [0, T)$.
We note that the estimate on $\kappa$ deteriorates as $T \rightarrow
\infty$, as there are Ricci flow solutions that collapse at
long time.

\section{I.5. The ${\mathcal W}$-functional as a time derivative} \label{I.5}

We will only discuss one formula from I.5, showing that along a Ricci flow,
${\mathcal W}$ is itself the time-derivative of an integral expression.

Again, we put $\tau \: = \: - \: t$. Consider the evolution equations
(\ref{oldeqn}), with
$(4 \pi \tau)^{-n/2} \: e^{-f} \: dV \: =\: dm$. Then
\begin{align}
\frac{d}{d\tau} \left( \tau \int_M \left( f \: - \: \frac{n}{2} \right)
\: dm \right) \: & = \:
\int_M \left( f \: - \: \frac{n}{2} \right)
\: dm \: + \: \tau \: \int_M  \left( \triangle f \: + \: R \: - \:
\frac{n}{2\tau} \right) \: dm \\
& = \:
\int_M \left( f \: - \: \frac{n}{2} \right)
\: dm \: + \: \tau \: \int_M  \left( |\nabla f|^2 \: + \: R \: - \:
\frac{n}{2\tau} \right) \: dm \notag \\
& = \: {\mathcal W}(g(t), f(t), \tau). \notag
\end{align}
With respect to the evolution equations
(\ref{neweqn}) obtained by performing diffeomorphisms, we get
\begin{equation}
\frac{d}{d\tau} \left( \tau \int_M \left( f \: - \: \frac{n}{2} \right)
\: (4 \pi \tau)^{-n/2} \: e^{-f} \: dV \right) \:
 = \: {\mathcal W}(g(t), f(t), \tau).
\end{equation}

Similarly, with respect to (\ref{newevolution}),
\begin{equation}
\frac{d}{dt} \left(- \int_M f
\: e^{-f} \: dV \right) \:
 = \: {\mathcal F}(g(t), f(t)).
\end{equation}

\section{I.7. Overview of reduced length and reduced volume} \label{I.7}

We first give a brief summary of I.7.
In I.7, the variable $\tau \: = \: t_0 \: - t$ is used and so the
corresponding Ricci flow equation is $(g_{ij})_\tau \: = \: 2 R_{ij}$.
The goal is to prove a no local collapsing theorem
by means of the ${\mathcal L}$-lengths of curves
$\gamma \: : \: [\tau_1, \tau_2] \rightarrow M$, defined by
\begin{equation} \label{defL}
{\mathcal L}(\gamma) \: = \:
\int_{\tau_1}^{\tau_2} \sqrt{\tau} \: \left( R(\gamma(\tau)) \: + \:
\big| \dot{\gamma}(\tau) \big|^2 \right) \: d\tau,
\end{equation}
where the scalar curvature $R(\ga(\tau))$ and the norm $|\dot{\gamma}(\tau)|$
are evaluated using the metric at time $t_0-\tau$.
Here $\tau_1 \ge 0$.
With $X \: = \: \frac{d\gamma}{d\tau}$,
the corresponding ${\mathcal L}$-geodesic equation is
\begin{equation}
\nabla_X X \: - \: \frac12 \: \nabla R \: + \: \frac{1}{2\tau} X \: + \:
2 \Ric(X, \cdot) \: = \: 0,
\end{equation}
where again the connection and curvature are taken at the corresponding time,
and the $1$-form $\Ric(X,\cdot)$ has been identified with the
corresponding dual vector field.

Fix $p \in M$.
Taking $\tau_1 \: = \: 0$ and $\gamma(0) \: = \: p$,
the vector $v \: = \: \lim_{\tau \rightarrow 0}
\sqrt{\tau} \: X(\tau)$ is well-defined in $T_pM$
and is called the initial vector of the geodesic. The
${\mathcal L}$-exponential map ${\mathcal L}exp_\tau \: : \: T_p M
\rightarrow M$ sends $v$ to $\gamma(\tau)$.

The function $L(q, \overline{\tau})$ is the infimal ${\mathcal L}$-length of
curves $\gamma$ with $\gamma(0) \: = \: p$ and $\gamma(\overline{\tau})
\: = \: q$. Defining the reduced length by
\begin{equation} \label{reducedl}
l(q, {\tau}) \: = \: \frac{L(q, \tau)}{2 \sqrt{\tau}}
\end{equation}
and the reduced volume by
\begin{equation} \label{reducedv}
\tilde{V}(\tau) \: = \: \int_M \tau^{-\frac{n}{2}} \: e^{-l(q, \tau)} \:
dq,
\end{equation}
the goal is to show that $\tilde{V}(\tau)$ is
nonincreasing in $\tau$, i.e. nondecreasing in $t$.
To do this, one uses the ${\mathcal L}$-exponential map to
write $\tilde{V}(\tau)$ as an integral over $T_pM$:
\begin{equation}
\tilde{V}(\tau) \: = \: \int_{T_pM} \tau^{-\frac{n}{2}} \:
e^{-l({\mathcal L}exp_\tau(v), \tau)} \: {\mathcal J}(v, \tau) \:
\chi_{\tau}(v) \:
dv,
\end{equation}
where ${\mathcal J}(v, \tau) \: = \:
\det \: d\left( {\mathcal L}exp_\tau \right)_v$ is the Jacobian factor
in the change of variable and
$\chi_\tau$ is a cutoff function related to the
${\mathcal L}$-cut locus of $p$.
To show that $\tilde{V}(\tau)$ is
nonincreasing in $\tau$ it suffices to show that
$\tau^{-\frac{n}{2}} \:
e^{-l({\mathcal L}exp_\tau(v), \tau)}
\: {\mathcal J}(v, \tau)$ is nonincreasing
in $\tau$ or, equivalently, that $- \: \frac{n}{2} \: \ln(\tau) \:
- \: l({\mathcal L}exp_\tau(v), \tau) \: + \: \ln \:{\mathcal J}(v, \tau)$ is
nonincreasing in $\tau$. Hence it is necessary to compute
$\frac{dl({\mathcal L}exp_\tau(v), \tau)}{d\tau}$  and
$\frac{d{\mathcal J}(v, \tau)}{d\tau}$. The
computation of the latter will involve the ${\mathcal L}$-Jacobi fields.

The fact that $\tilde{V}(\tau)$ is
nonincreasing in $\tau$ is then used to show that the Ricci flow
solution cannot be  collapsed near $p$.

\section{Basic example for I.7} \label{7Euclidean}

 In this section we say what the various expressions of I.7 become in
 the model case of a flat Euclidean Ricci solution.

If $M$ is flat $\R^n$ and $p \: = \: \vec{0}$ then the unique
${\mathcal L}$-geodesic $\gamma$ with $\gamma(0) \: =\: \vec{0}$ and
$\gamma(\overline{\tau}) \: = \: \vec{q}$ is
\begin{equation}
\gamma(\tau) \: = \: \left( \frac{\tau}{\overline{\tau}}
\right)^{\frac12} \: \vec{q} \: = \: 2 \: \tau^{\frac12} \: \vec{v}.
\end{equation}
The function $L$ is given by
\begin{equation}
L(q, \overline{\tau}) \: = \: \frac12 \: \overline{\tau}^{\: -
\frac12} \: |q|^2
\end{equation}
and the reduced length (\ref{reducedl}) is given by
\begin{equation}
l(q, \overline{\tau}) \: = \: \frac{|q|^2}{4 \overline{\tau}}.
\end{equation}
The function $\overline{L}(q, \overline{\tau}) \: = \:
2 \: \overline{\tau}^{\frac12} \: L(q, \overline{\tau})$
is
\begin{equation}
\overline{L}(q, \overline{\tau}) \: = \: |q|^2.
\end{equation}
Then
\begin{equation}
\tilde{V}(\tau) \: = \:
\int_{\R^n} \tau^{-\frac{n}{2}} \: e^{-\frac{|q|^2}{4 {\tau}}} \:
d^nq \: = \: (4\pi)^{\frac{n}{2}}
\end{equation}
is constant in $\tau$.

\section{Remarks about ${\mathcal L}$-Geodesics and
${\mathcal L}\exp$} \label{Lremarks}

In this section we discuss the variational equation corresponding to
(\ref{defL}).

To derive the ${\mathcal L}$-geodesic equation, as in Riemannian geometry
we consider a
$1$-parameter family of curves
$\gamma_s \: : \: [\tau_1, \tau_2] \rightarrow
M$, parametrized by $s \in (-\epsilon, \epsilon)$. Equivalently, we
have a map $\tilde\gamma(s, \tau)$ with $s \in (-\epsilon, \epsilon)$ and
$t \in [\tau_1, \tau_2]$. Putting $X \: = \:
\frac{\partial \tilde\gamma}{\partial \tau}$ and
$Y \: = \:
\frac{\partial \tilde\gamma}{\partial s}$, we have
$[X, Y] \: = \: 0$. Then $\nabla_X Y \: = \: \nabla_Y X$.
Restricting to the curve $\gamma(\tau) \: = \: \tilde\gamma(0, \tau)$ and
writing $\delta_Y$ as shorthand for $\frac{d}{ds} \big|_{s=0}$, we
have $(\delta_Y \gamma)(\tau) \:  =\: Y(\tau)$ and $(\delta_Y X)(\tau) \: = \:
(\nabla_X Y)(\tau)$.
Then
\begin{equation}
\delta_Y {\mathcal L} \: = \:
\int_{\tau_1}^{\tau_2} \sqrt{\tau} \: \left( \langle Y, \nabla R \rangle
\: + \: 2 \: \langle \nabla_X Y, X \rangle \right) \: d\tau.
\end{equation}
Using the fact that $\frac{dg_{ij}}{d\tau} \: = \: 2 R_{ij}$, we have
\begin{equation}
\frac{d\langle Y, X \rangle}{d\tau} \: = \:
\langle \nabla_X Y, X \rangle \: + \:
\langle Y, \nabla_X X \rangle \: + \:
2 \: \Ric(Y, X).
\end{equation}
Then
\begin{align}
\label{firstvariation}
& \int_{\tau_1}^{\tau_2} \sqrt{\tau} \: \left( \langle Y, \nabla R \rangle
\: + \: 2 \: \langle \nabla_X Y, X \rangle \right) \: d\tau \: = \\
& \int_{\tau_1}^{\tau_2} \sqrt{\tau} \: \left( \langle Y, \nabla R \rangle
\: + \: 2 \: \frac{d}{d\tau} \langle Y, X \rangle \: - \: 2
\langle Y, \nabla_X X \rangle \: - \: 4 \: \Ric(Y, X)
 \right) \: d\tau \: = \notag \\
& 2 \: \sqrt{{\tau}} \: \langle X, Y \rangle \big|_{\tau_1}^{\tau_2}
\: + \:
\int_{\tau_1}^{\tau_2} \sqrt{\tau} \: \left\langle Y, \nabla R
\: - \: 2 \nabla_X X  \: - \: 4 \: \Ric(X, \cdot) \: -\: \frac{1}{\tau}
\: X \right\rangle
 \: d\tau. \notag
\end{align}
Hence the ${\mathcal L}$-geodesic equation is
\begin{equation} \label{Lgeod}
\nabla_X X \: - \: \frac12 \: \nabla R \: + \: \frac{1}{2\tau} X \: + \:
2 \Ric(X, \cdot) \: = \: 0.
\end{equation}

We now discuss some technical issues about ${\mathcal L}$-geodesics and
the ${\mathcal L}$-exponential map.
We are assuming that $(M,g(\cdot))$ is a Ricci flow, where the curvature
operator
of $M$ is uniformly bounded on a $\tau$-interval $[\tau_1,\tau_2]$, and
each $\tau$-slice $(M,g(\tau))$ is complete for $\tau\in [\tau_1,\tau_2]$.
By Appendix \ref{applocalder},  for every $\tau'<\tau_2$ there
is a constant $D< \infty$ such that
\begin{equation}
\label{dr}
|\nabla R(x,\tau)|< \frac{D}{\sqrt{\tau_2 - \tau}}
\end{equation}
for all $x\in M$, $\tau \in [\tau_1, \tau_2)$.


Making the change of variable $s \: = \: \sqrt{\tau}$ in the formula
for ${\mathcal L}$-length, we get
\begin{equation}
\label{changetos}
{\mathcal L}(\gamma) \: = \:
2 \: \int_{s_1}^{s_2} \: \left( \frac14
\Big| \frac{d\gamma}{ds} \Big|^2 \: + \:
s^2 \: R(\gamma(s)) \right) \: ds.
\end{equation}
The Euler-Lagrange equation becomes
\begin{equation}
\label{sE-L}
\nabla_{\hat X}\hat X-2s^2\nabla R+4s\Ric(\hat X,\cdot)=0,
\end{equation}
where $\hat X \: = \: \frac{d\gamma}{ds}=2sX$.
Putting $s_1 \: = \: \sqrt{\tau_1}$,
it follows from standard existence theory
for ODE's that for each $p \in M$ and $v\in T_pM$,
there is a unique solution $\gamma(s)$ to
(\ref{sE-L}), defined on an interval $[s_1,s_1+\eps)$, with $\gamma(s_1) = p$
and
\begin{equation}
\frac{1}{2} \: \gamma^\prime(s_1) \: = \:
\lim_{\tau\ra \tau_1}\sqrt{\tau}\frac{d\gamma}{d\tau} \: = \: v.
\end{equation}
If $\gamma(s)$ is defined for $s \in [s_1,s']$ then
\begin{align}
\frac{d}{ds}|\hat X|^2=\frac{d}{ds} \: \langle \hat X,\hat X\rangle
&= 4s\Ric(\hat X,\hat X)+2\langle\nabla_{\hat X}\hat X,\hat X\rangle \\
&= - 4s\Ric(\hat X,\hat X) + 4s^2 \langle \nabla R, \hat X \rangle \notag
\end{align}
and so if $\hat X (s) \neq 0$ then
\begin{equation}
\frac{d}{ds}|\hat X|=
\frac{1}{2|\hat X|} \: \frac{d}{ds}|\hat X|^2 \: = \:
- 2s |\hat X|
 \Ric \left( \frac{\hat X}{|\hat X|}, \frac{\hat X}{|\hat X|} \right)
+ 2s^2 \left\langle \nabla R, \frac{\hat X}{|\hat X|} \right\rangle.
\end{equation}
By (\ref{dr}),
\begin{equation}
\frac{d}{ds}|\hat X| \: \le \: C_1 |\hat X| +
\frac{C_2}{\sqrt{s_2 - s}}
\end{equation}
for appropriate constants $C_1$ and $C_2$, where
$s_2 \: = \: \sqrt{\tau_2}$.
Since the metrics $g(\tau)$
are uniformly comparable for $\tau\in [\tau_1,\tau_2]$, we conclude
(by a continuity argument
in $s$
) that the ${\mathcal L}$-geodesic $\gamma_v$
with $\frac{1}{2} \gamma_v^\prime(s_1)=v$ is defined on the whole
interval
$[s_1,s_2]$.
In particular, in terms of the original variable $\tau$,
for each
${\tau}\in [\tau_1,\tau_2]$
and each $p\in M$,
we get a globally defined and smooth ${\mathcal L}$-exponential map
${\mathcal L}\exp_\tau:T_pM\ra M$ which takes each $v\in T_pM$
to $\gamma({\tau})$, where $v \: = \: \lim_{\tau^\prime \rightarrow \tau_1}
\sqrt{\tau^\prime} \frac{d\gamma}{d\tau^\prime}$.
Note that unlike in the case of Riemannian geometry,
${\mathcal L}\exp_\tau(0)$ may
not be $p$, because of the $\nabla R$ term in
(\ref{Lgeod}).

We now fix $p\in M$, take $\tau_1=0$, and let $L(q,\bar\tau)$ be the
minimizer function as in Section \ref{I.7}.
We can imitate the traditional Riemannian geometry proof that geodesics
minimize for a short time.
Using the change of variable $s = \sqrt{\tau}$
and the implicit function theorem,
there is an $r=r(p)>0$ (which varies continuously with $p$) such that
for every $q\in M$ with $d(q,p)\leq 10r$ at $\tau=0$, and every
$0<\bar\tau\leq r^2$,
there is a unique ${\mathcal L}$-geodesic
$\gamma_{(q,\bar\tau)}:[0,\bar\tau]\ra M$,
starting at $p$ and ending at $q$, which remains within the ball $B(p,100r)$
(in the $\tau=0$ slice $(M,g(0))$),
and $\gamma_{(q,\bar\tau)}$ varies smoothly with $(q,\bar\tau)$.
Thus, the ${\mathcal L}$-length of $\gamma_{(q,\bar\tau)}$ varies smoothly
with $(q,\bar\tau)$, and defines a function $\hat L(q,\bar\tau)$ near $(p,0)$.
We claim that $\hat L = L$ near $(p, 0)$.
Suppose that $q \in B(p,r)$ and let $\alpha : [0, \bar \tau]
\rightarrow M$ be a smooth curve whose ${\mathcal L}$-length is
close to $L(q,\bar\tau)$. If $r$ is small, relative to the
assumed curvature bound, then $\alpha$ must stay within
$B(p, 10r)$. Equations
(\ref{grad1}) and (\ref{timed}) below imply that
\begin{equation}
\frac{d}{d\tau} \hat{L}(\alpha(\tau), \tau) \: = \:
\langle 2 \sqrt{\tau} X, \frac{d\alpha}{d\tau} \rangle \: + \:
\sqrt{\tau} (R - |X|^2) \: \le \:
\sqrt{\tau} \left( R + \Big| \frac{d\alpha}{d\tau} \Big|^2 \right) \: = \:
\frac{d}{d\tau} \: \left(
{\mathcal L} \text{length}(\alpha \big|_{[0, \tau]}) \right).
\end{equation}
Thus
$\gamma_{(q, \bar \tau)}$ minimizes when $(q, \bar \tau)$ is close to
$(p, 0)$.


We can now deduce that for all $(q,\bar\tau)$, there is an
${\mathcal L}$-geodesic
$\gamma:[0,\bar\tau]\ra M$ which has infimal ${\mathcal L}$-length
among all piecewise smooth curves starting at $p$ and ending at $q$
(with domain $[0,\bar\tau]$).  This can be done by  imitating the
usual broken geodesic argument, using the fact that
for $x,y$ in a given small ball of $M$ and
for sufficiently
small time intervals $[\tau^\prime, \tau^\prime+\epsilon] \subset
[0, \overline{\tau}]$, there is a unique
minimizer $\gamma$ for $\int_{\tau^\prime}^{\tau^\prime+\epsilon}
\sqrt{\tau} \left( R(\gamma(\tau)) \: +\: \left| \dot{\gamma}(\tau)
\right|^2 \right) \: d\tau$ with
$\gamma(\tau^\prime) \: = \: x$ and
$\gamma(\tau^\prime+\epsilon) \: = \: y$. Alternatively, using the
change of variable $s \: = \: \sqrt{\tau}$, one can
take a minimizer of ${\mathcal L}$ among $H^{1,2}$-regular curves.

Another technical issue is the justification of the change of variables
from $M$ to $T_pM$ in the proof of monotonicity of reduced volume.
Fix $p\in M$ and $\tau>0$, and let ${\mathcal L}\exp_\tau:T_pM\ra M$
be the map which takes $v\in T_pM$ to $\gamma_v(\tau)$, where
$\gamma_v:[0,\tau]\ra M$ is the unique ${\mathcal L}$-geodesic
with $\sqrt{\tau^\prime}\frac{d\gamma_v}{d\tau^\prime}\ra v$ as
$\tau^\prime\ra 0$.
Let ${\mathcal B}_\tau \subset M$ be the set of points which are
either endpoints of more than one minimizing ${\mathcal L}$-geodesic,
or which are the endpoint of a minimizing geodesic $\gamma_v:[0,\tau]\ra M$
where $v\in T_pM$ is a critical point of ${\mathcal L}\exp_\tau$.
We will call ${\mathcal B}_\tau$
the time-$\tau$ ${\mathcal L}$-cut locus of $p$. It is a closed
subset of $M$.
Let ${\mathcal G}_\tau \subset M$ be the complement of ${\mathcal B}_\tau$ and
let $\Omega_\tau \subset T_pM$ be the corresponding set of initial
conditions for minimizing ${\mathcal L}$-geodesics.
Then $\Omega_\tau$ is an open set, and ${\mathcal L}\exp_\tau$
maps it diffeomorphically onto ${\mathcal G}_\tau$.  We claim that
${\mathcal B}_\tau$
has measure zero.   By Sard's theorem, to prove this it suffices
to prove that the set ${\mathcal B}'_\tau$ of points
$q\in {\mathcal B}_\tau$ which are
regular values of ${\mathcal L}\exp_\tau$, has measure zero.
Pick $q\in {\mathcal B}'_\tau$, and distinct points $v_1,v_2\in T_pM$
such that $\gamma_{v_i}:[0,\tau]\ra M$ are both minimizing geodesics
ending at $q$.  Then ${\mathcal L}\exp_\tau$ is a local diffeomorphism
near each $v_i$. The first variation formula and the implicit function theorem
then show that there are neighborhoods
$U_i$ of $v_i$, and a smooth hypersurface $H$  passing through
$q$, such that if we have points $w_i\in U_i$ with
\begin{equation}
q' \: = \: {\mathcal L}\exp_\tau(w_1)={\mathcal L}\exp_\tau(w_2)\quad
\mbox{and}\quad{\mathcal L}\length(\gamma_{w_1})={\mathcal L}
\length(\gamma_{w_2}),
\end{equation}
then $q'$ lies on $H$.  Thus ${\mathcal B}'_\tau$
is contained in a countable union
of hypersurfaces, and hence has measure zero.

Therefore one may compute the integral of any integrable function on $M$
by pulling it back to $\Omega_\tau \subset T_pM$ and using the change
of variables formula.
Note that
if $\tau \: \le \: \tau^\prime$ then $\Omega_{\tau^\prime} \subset
\Omega_\tau$.


\section{I.(7.3)-(7.6). First derivatives of $L$}

In this section we do some preliminary calculations leading up to the
computation of the second variation of $L$.

A remark about the notation : $L$ is a function of
a point $q$ and a time $\overline{\tau}$. The notation
$L_{\overline{\tau}}$ refers to the partial derivative with
respect to $\overline{\tau}$, i.e. differentiation while keeping
$q$ fixed.  The notation $\frac{d}{d\tau}$ refers to differentiation
along an ${\mathcal L}$-geodesic, i.e. simultaneously
varying both the point and the time.

If $q$ is not in the time-$\overline{\tau}$
${\mathcal L}$-cut locus of $p$, let
$\gamma \: : \: [0, \overline{\tau}] \rightarrow M$ be the unique minimizing
${\mathcal L}$-geodesic from $p$ to $q$, with length $L(q, \overline{\tau})$.
If $c \: : \: (- \epsilon, \epsilon) \rightarrow M$ is a short curve
with $c(0) \: = \: q$, consider the $1$-parameter
family of minimizing ${\mathcal L}$-geodesics
$\tilde\gamma(s, \tau)$
with $\tilde\gamma(s, 0) \: = \: p$ and $\tilde\gamma(s, \overline{\tau})
\: =\: c(s)$. Putting $Y(\tau) \: =\:
\frac{\partial \tilde\gamma(s, \tau)}{\partial s} \Big|_{s=0}$,
equation (\ref{firstvariation}) gives
\begin{equation}
\langle \nabla L, c^\prime(0) \rangle \: = \:
\frac{dL(c(s), \overline{\tau})}{ds} \Big|_{s=0} \: = \:
2 \: \sqrt{\overline{\tau}} \:
\langle X(\overline{\tau}), Y(\overline{\tau}) \rangle.
\end{equation}
Hence
\begin{equation} \label{grad1}
(\nabla L)(q, \overline{\tau}) \: = \:
2 \: \sqrt{\overline{\tau}} \: X(\overline{\tau})
\end{equation}
and
\begin{equation} \label{gradsq}
|\nabla L|^2(q, \overline{\tau}) \: = \:
4 \: {\overline{\tau}} \: |X(\overline{\tau})|^2 \: = \:
- \: 4 \: {\overline{\tau}} \: R(q) \: + \:
4 \: {\overline{\tau}} \: \left( R(q) \: + \: |X(\overline{\tau})|^2 \right).
\end{equation}

If we simply extend the ${\mathcal L}$-geodesic $\gamma$
in $\overline{\tau}$, we
obtain
\begin{equation} \label{grad2}
\frac{dL(\gamma(\overline{\tau}),\overline{\tau})}{d\overline{\tau}} \: = \:
\sqrt{\overline{\tau}} \: \left( R(\gamma(\overline{\tau})) \: + \:
|X(\overline{\tau})|^2 \right).
\end{equation}
As
\begin{equation}
\frac{dL(\gamma(\overline{\tau}),\overline{\tau})}{d\overline{\tau}} \: = \:
L_{\overline{\tau}}(q, \overline{\tau}) \: + \:
\langle (\nabla L)(q, \overline{\tau}), X(\overline{\tau}) \rangle,
\end{equation}
equations (\ref{grad1}) and (\ref{grad2}) give
\begin{align} \label{timed}
L_{\overline{\tau}}(q, \overline{\tau}) \: & = \:
\sqrt{\overline{\tau}} \: \left( R(q) \: + \:
|X(\overline{\tau})|^2 \right) \: - \:
\langle (\nabla L)(q, \overline{\tau}), X(\overline{\tau}) \rangle \\
&=\:
2 \sqrt{\overline{\tau}} R(q) \:
- \: \sqrt{\overline{\tau}} \: \left( R(q) \: + \:
|X(\overline{\tau})|^2 \right). \notag
\end{align}

When computing $\frac{dl(\gamma(\tau), \tau)}{d\tau}$, it will be useful
to have a formula for $R(\gamma(\tau)) \: + \: \big| X(\tau) \big|^2$.
As in I.(7.3),
\begin{equation}
\frac{d}{d\tau} \left( R(\gamma(\tau)) \: + \: \big| X(\tau) \big|^2 \right)
\: = \:
R_\tau \: + \: \langle \nabla R, X \rangle \: + \: 2 \langle \nabla_X X, X
\rangle \: + \: 2 \Ric(X, X).
\end{equation}
Using the ${\mathcal L}$-geodesic equation (\ref{Lgeod}) gives
\begin{align} \label{ddt}
\frac{d}{d\tau} \left( R(\gamma(\tau)) \: + \: \big| X(\tau) \big|^2 \right)
\: & = \:
R_\tau \: + \: \frac{1}{\tau} \: R \: + \:
2 \langle \nabla R, X \rangle \:
- \: 2 \Ric(X, X) \: - \: \frac{1}{\tau} \: (R + |X|^2) \\
& = \: - \: H(X) \: - \: \frac{1}{\tau} \: (R + |X|^2), \notag
\end{align}
where
\begin{equation} \label{Hexp}
H(X) \: = \: - \: R_\tau \: - \: \frac{1}{\tau} \: R \: - \:
2 \langle \nabla R, X \rangle \:
+ \: 2 \Ric(X, X)
\end{equation}
is the expression of (\ref{Heqn})
after the change $\tau \: = \:- t$ and $X \rightarrow -X$.
Multiplying (\ref{ddt}) by $\tau^{\frac32}$ and integrating gives
\begin{equation} \label{parts}
\int_0^{\overline{\tau}} \tau^{\frac32} \:
\frac{d}{d\tau} \left( R(\gamma(\tau)) \: + \: \big| X(\tau) \big|^2 \right)
\: d\tau
\: = \: - \: K \: - \: L(q, \overline{\tau}),
\end{equation}
where
\begin{equation}
K \: = \: \int_0^{\overline{\tau}} \tau^{\frac32} \: H(X(\tau)) \: d\tau.
\end{equation}
Then integrating the left-hand side of (\ref{parts}) by parts gives
\begin{equation} \label{sub4}
\overline{\tau}^{\frac32} \:
\left( R(\gamma(\overline{\tau})) \: + \: \big| X(\overline{\tau})
\big|^2 \right)
\: = \: - \: K \: + \: \frac12 \: L(q, \overline{\tau}).
\end{equation}
Plugging this back into  (\ref{timed}) and (\ref{gradsq}) gives
\begin{equation} \label{touse1}
L_{\overline{\tau}}(q, \overline{\tau}) \: = \:
2 \sqrt{\overline{\tau}} R(q) \:
- \: \frac{1}{2\overline{\tau}} \: L(q,\overline{\tau}) \: + \:
\frac{1}{\overline{\tau}} \: K
\end{equation}
and
\begin{equation} \label{ttouse1}
|\nabla L|^2(q, \overline{\tau}) \: = \:
- \: 4 \: {\overline{\tau}} \: R(q) \: + \:
\frac{2}{\sqrt{\overline{\tau}}} \: L(q, \overline{\tau}) \: - \:
\frac{4}{\sqrt{\overline{\tau}}} \: K.
\end{equation}

\section{I.(7.7). Second variation of ${\mathcal L}$}


In this section we compute the second variation of ${\mathcal L}$.
We use it to compute the Hessian of $L$ on $M$.


To compute the second variation $\delta_Y^2 {\mathcal L}$, we start with
the first variation equation
\begin{equation}
\delta_Y {\mathcal L} \: = \:
\int_{0}^{\overline{\tau}} \sqrt{\tau} \: \left( \langle Y, \nabla R \rangle
\: + \: 2 \: \langle \nabla_Y X, X \rangle \right) \: d\tau.
\end{equation}
Recalling that $\delta_Y \gamma(\tau) \: =\: Y(\tau)$ and
$\delta_Y X(\tau) \: = \: (\nabla_Y X)(\tau)$,
the second variation is
\begin{align}
\delta_Y^2 {\mathcal L} \: & = \:
\int_{0}^{\overline{\tau}} \sqrt{\tau} \: \left( Y \cdot Y \cdot R
\: + \: 2 \: \langle \nabla_Y \nabla_Y X, X \rangle
\: + \: 2 \: \langle \nabla_Y X, \nabla_Y X \rangle \right) \: d\tau \\
& = \:
\int_{0}^{\overline{\tau}} \sqrt{\tau} \: \left( Y \cdot Y \cdot R
\: + \: 2 \: \langle \nabla_Y \nabla_X Y, X \rangle
\: + \: 2 \: \big| \nabla_X Y \big|^2 \right) \: d\tau \notag \\
& = \:
\int_{0}^{\overline{\tau}} \sqrt{\tau} \: \left( Y \cdot Y \cdot R
\: + \: 2 \: \langle \nabla_X \nabla_Y Y, X \rangle
\: + \: 2 \: \langle R(Y, X) Y, X \rangle
\: + \: 2 \: \big| \nabla_X Y \big|^2 \right) \: d\tau, \notag
\end{align}
where the notation $Z \cdot u$ refers to the directional derivative, i.e.
$Z\cdot u\;=\; i_Z\;du$.
In order to deal with the $\langle \nabla_X \nabla_Y Y, X \rangle$ term,
we have to compute $\frac{d}{d\tau} \langle \nabla_Y Y, X \rangle$.

From the general equation for the Levi-Civita
connection in terms of the metric
\cite[(1.29)]{Chavel}, if $g(\tau)$ is a $1$-parameter family of metrics,
with $\dot{g} \: = \: \frac{dg}{d\tau}$
and $\dot{\nabla} \: = \: \frac{d\nabla}{d\tau}$, then
\begin{equation}
2 \langle \dot{\nabla}_X Y, Z \rangle \: = \:
(\nabla_X \dot{g}) (Y, Z) \: + \:
(\nabla_Y \dot{g}) (Z, X) \: - \:
(\nabla_Z \dot{g}) (X, Y).
\end{equation}
In our case $\dot{g} \: = \: 2 \: \Ric$ and so
\begin{align} \label{xxx}
\frac{d}{d\tau} \langle \nabla_Y Y, X \rangle \: = \: &
\langle \nabla_X \nabla_Y Y, X \rangle \: + \:
\langle \nabla_Y Y, \nabla_X X \rangle \: + \:
2 \Ric ( \nabla_Y Y, X ) \: + \:
\langle \dot{\nabla}_Y Y, X \rangle \\
 = \: &
\langle \nabla_X \nabla_Y Y, X \rangle \: + \:
\langle \nabla_Y Y, \nabla_X X \rangle \: + \notag \\
& 2 \Ric ( \nabla_Y Y, X ) \: + \:
2 (\nabla_Y \Ric) (Y, X) \: - \:
(\nabla_X \Ric) (Y, Y). \notag
\end{align}

(Although we will not need it, we can write
\begin{align}
2 Y \cdot \Ric(Y, X) \: - \: X \cdot \Ric(Y, Y) \: = \: &
2 (\nabla_Y \Ric)(Y, X) \: + \:
2 \Ric(\nabla_Y Y, X) \: + 2 \Ric(Y, \nabla_Y X) \\
& - \:
(\nabla_X \Ric) (Y, Y) \: - \:
2 \Ric(\nabla_X Y, Y) \notag \\
 = \: &
2 \Ric ( \nabla_Y Y, X ) \: + \:
2 (\nabla_Y \Ric) (Y, X) \notag \\
& - \:
(\nabla_X \Ric) (Y, Y) \: - \:
2 \Ric([X, Y], Y).
\notag
\end{align}
We are assuming that the variation field $Y$ satisfies $[X, Y] \: = \: 0$
(this was used in deriving the ${\mathcal L}$-geodesic equation).
Hence one obtains the formula
\begin{equation}
\frac{d}{d\tau} \langle \nabla_Y Y, X \rangle \: = \:
\langle \nabla_X \nabla_Y Y, X \rangle \: + \:
\langle \nabla_Y Y, \nabla_X X \rangle \: +
2 Y \cdot \Ric(Y, X) \: - \: X \cdot \Ric(Y, Y)
\end{equation}
of I.7.)

Next, using (\ref{xxx}),
\begin{align}
2 \sqrt{\overline{\tau}} \langle \nabla_Y Y, X \rangle \: = \: &
2 \int_0^{\overline{\tau}} \frac{d}{d\tau}
\left( \sqrt{\tau} \langle \nabla_Y Y, X \rangle  \right) \: d\tau \\
 = \: &
\int_0^{\overline{\tau}}
\sqrt{\tau} \left[ \frac{1}{\tau} \:
\langle \nabla_Y Y, X \rangle \: + \:
2  \:  \frac{d}{d\tau} \: \langle \nabla_Y Y, X \rangle \right]
\: d\tau \notag \\
\: = \: &
\int_0^{\overline{\tau}}
\sqrt{\tau} \left[ \frac{1}{\tau} \:
\langle \nabla_Y Y, X \rangle \: + \:
2 \langle \nabla_X \nabla_Y Y, X \rangle \: + \:
2 \langle \nabla_Y Y, \nabla_X X \rangle \: + \: \right. \notag \\
& \left.
\: \: \: \: \:\: \: \: \: \: \: \: \: \: \: 4 \Ric ( \nabla_Y Y, X ) \: + \:
4 (\nabla_Y \Ric) (Y, X) \: - \: 2
(\nabla_X \Ric) (Y, Y) \right] \: d\tau \notag \\
\: = \: &
\int_0^{\overline{\tau}}
\sqrt{\tau} \left[
2 \langle \nabla_X \nabla_Y Y, X \rangle \: + \:
(\nabla_Y Y) R \: + \: \right. \notag \\
& \left.
\: \: \: \: \: \: \: \: \: \: \: \: \: \: \:
4 (\nabla_Y \Ric) (Y, X) \: - \: 2
(\nabla_X \Ric) (Y, Y) \right] \: d\tau \notag \\
& - \:
\int_0^{\overline{\tau}}
\sqrt{\tau} \left\langle \nabla_Y Y,
\nabla R
\: - \: 2 \nabla_X X  \: - \: 4 \: \Ric(X, \cdot) \: -\: \frac{1}{\tau}
\: X
\right\rangle \: d\tau. \notag
\end{align}
(Of course the last term vanishes if $\gamma$ is an ${\mathcal L}$-geodesic,
but we do not need to assume this here.)

The quadratic form $Q$ representing the Hessian of ${\mathcal L}$ on the
path space is given by
\begin{align}
Q(Y, Y) \:  = \: & \delta_Y^2 {\mathcal L} \: - \: \delta_{\nabla_Y Y}
{\mathcal L} \\
= \:
& \delta_Y^2 {\mathcal L} \: - \:
2 \sqrt{\overline{\tau}} \langle \nabla_Y Y, X \rangle \:
-  \notag \\
& \int_0^{\overline{\tau}}
\sqrt{\tau} \left\langle \nabla_Y Y,
\nabla R
\: - \: 2 \nabla_X X  \: - \: 4 \: \Ric(X, \cdot) \: -\: \frac{1}{\tau}
\: X
\right\rangle. \notag
\end{align}
It follows that
\begin{align} \label{Hess}
Q(Y, Y) \: = \: &
\int_0^{\overline{\tau}}
\sqrt{\tau} \left[ Y \cdot Y \cdot R \: - \: (\nabla_Y Y) R \: + \:
2 \langle R(Y, X) Y, X \rangle \: + \right. \\
& \left. 2 \: |\nabla_X Y|^2
\: - \: 4 (\nabla_Y \Ric) (Y, X) \: + \: 2
(\nabla_X \Ric) (Y, Y) \right] \: d\tau \notag \\
= \: & \int_0^{\overline{\tau}}
\sqrt{\tau} \left[\Hess_R (Y,Y) \: + \:
2 \langle R(Y, X) Y, X \rangle \: + \: 2 \: |\nabla_X Y|^2 \right. \notag \\
&
\: \: \: \: \: \: \: \: \: \: \: \: \: \: \:
 \left. - \: 4 (\nabla_Y \Ric) (Y, X) \: + \: 2
(\nabla_X \Ric) (Y, Y) \right] \: d\tau. \notag
\end{align}
There is an associated second-order differential
operator $T$ on vector fields $Y$ given by saying that
$2 \sqrt{\tau} \langle Y, TY\rangle$ equals the integrand of (\ref{Hess}) minus
$2 \frac{d}{d\tau} \left( \sqrt{\tau} \: \langle \nabla_X Y, Y \rangle
\right)$.
Explicitly,
\begin{equation}
TY \: = \: - \: \nabla_X \nabla_X Y \: - \: \frac{1}{2\tau} \: \nabla_X Y
\: + \: \frac12 \: \Hess_R(Y, \cdot) \: - \: 2 \: (\nabla_Y \Ric)(X, \cdot)
\: - \: 2 \: \Ric(\nabla_Y X, \cdot).
\end{equation}
Then
\begin{equation} \label{Qeqn}
Q(Y, Y) \: = \: 2 \int_0^{\overline{\tau}}
\sqrt{\tau} \langle Y, TY \rangle \: d\tau \: + \:
2 \: \sqrt{\overline{\tau}} \: \langle \nabla_X Y(\overline{\tau}),
Y(\overline{\tau}) \rangle.
\end{equation}
An ${\mathcal L}$-Jacobi field along an ${\mathcal L}$-geodesic
is a field $Y(\tau)$ that is annihilated by $T$.

The Hessian of the function
$L(\cdot, \overline{\tau})$ can be computed as follows.
Assume that $q \in M$ is outside of the time-$\overline{\tau}$
${\mathcal L}$-cut locus. Let $\gamma \: : \: [0, \overline{\tau}]
\rightarrow M$ be the minimizing ${\mathcal L}$-geodesic with
$\gamma(0) \: = \: p$ and $\gamma(\overline{\tau}) \: = \: q$.
Given $w \in T_qM$, take a short geodesic $c \: : \: (-\epsilon,
\epsilon) \rightarrow M$ with $c(0) \: = \: q$ and $c^\prime(0) \: = \: w$.
Form the $1$-parameter
family of ${\mathcal L}$-geodesics $\tilde\gamma(s, \tau)$ with
$\tilde\gamma(s, 0) \: = \: p$ and
$\tilde\gamma(s, \overline{\tau}) \: = \: c(s)$.
Then $Y(\tau) \: = \: \frac{\partial \tilde\gamma(s, \tau)}{\partial s}
\Big|_{s=0}$ is an ${\mathcal L}$-Jacobi field $Y$ along $\gamma$ with
$Y(0) \: = \: 0$ and $Y(\overline{\tau}) \: = \: w$.
We have
\begin{equation} \label{hesseq}
\Hess_L(w, w) \: = \: \frac{d^2 L(c(s),
\overline{\tau})}{ds^2} \Big|_{s=0} \: = \:
Q(Y, Y) \: = \:
2 \: \sqrt{\overline{\tau}} \: \langle \nabla_X Y(\overline{\tau}),
Y(\overline{\tau}) \rangle.
\end{equation}


From (\ref{Qeqn}),
a minimizer of
$Q(Y, Y)$, among fields $Y$ with given values at the endpoints,
is an ${\mathcal L}$-Jacobi field.
It follows that
$\Hess_L(w, w) \: \le \: Q(Y, Y)$ for {\em any} field $Y$ along $\gamma$
satisfying
$Y(0) \: = \: 0$ and $Y(\overline{\tau}) \: = \: w$.


\section{I.(7.8)-(7.9). Hessian bound for $L$} \label{7.9}

In this section we use a test variation field in order to estimate
the Hessian of $L$.

If $Y(\overline{\tau})$
is a unit vector at $\gamma(\overline{\tau})$, solve for
$\widetilde{Y}(\tau)$ in the equation
\begin{equation} \label{diffeq}
\nabla_X \widetilde{Y} \: = \: - \: \Ric(\widetilde{Y}, \cdot)
\: + \: \frac{1}{2\tau} \:\widetilde{Y},
\end{equation}
on the interval $0<\tau\leq \overline{\tau}$ with the endpoint condition
$\widetilde{Y}(\overline{\tau}) \: = \:
Y(\overline{\tau})$.
(For this section, we change notation from the
$Y$ in I.7 to $\widetilde{Y}$.)
Then
\begin{equation}
\frac{d}{d\tau} \langle \widetilde{Y}, \widetilde{Y} \rangle \: = \:
2 \Ric(\widetilde{Y}, \widetilde{Y})
\: + \: 2 \langle \nabla_X \widetilde{Y}, \widetilde{Y} \rangle \: = \:
\frac{1}{\tau} \langle \widetilde{Y}, \widetilde{Y} \rangle,
\end{equation}
so
$\langle \widetilde{Y}(\tau), \widetilde{Y}(\tau) \rangle \: = \:
\frac{\tau}{\overline{\tau}}$.
Thus we can extend $\widetilde{Y}$ continuously to the interval $[0,\overline{\tau}]$
by putting $\widetilde{Y}(0) \: = \: 0$.
Substituting into
(\ref{Hess}) gives
\begin{align} \label{sub1}
Q(\widetilde{Y}, \widetilde{Y}) \: = \: & \int_0^{\overline{\tau}}
\sqrt{\tau} \left[\Hess_R (\widetilde{Y},\widetilde{Y}) \: + \:
2 \langle R(\widetilde{Y}, X) \widetilde{Y}, X \rangle
\: + \: 2 \: \left| - \Ric(\widetilde{Y}, \cdot) \: + \:
\frac{1}{2\tau} \widetilde{Y} \right|^2 \right. \\
&
\: \: \: \: \: \: \: \: \: \: \: \: \: \: \:
 \left. - \: 4 (\nabla_{\widetilde{Y}} \Ric) (\widetilde{Y}, X) \: + \: 2
(\nabla_X \Ric) (\widetilde{Y}, \widetilde{Y}) \right] \: d\tau \notag \\
= \: & \int_0^{\overline{\tau}}
\sqrt{\tau} \left[\Hess_R (\widetilde{Y},\widetilde{Y}) \: + \:
2 \langle R(\widetilde{Y}, X) \widetilde{Y}, X \rangle  \: + \: 2
(\nabla_X \Ric) (\widetilde{Y}, \widetilde{Y}) \: \right. \notag \\
&
\: \: \: \: \: \: \: \: \: \: \: \: \: \: \:
\left. - \: 4 (\nabla_{\widetilde{Y}} \Ric) (\widetilde{Y}, X) \:
+ \: 2 \: |\Ric(\widetilde{Y}, \cdot)|^2 \: - \: \frac{2}{\tau}
\Ric(\widetilde{Y}, \widetilde{Y})
\: + \:
\frac{1}{2\tau \overline{\tau}}
\right] \: d\tau. \notag
\end{align}
From
\begin{align}
\frac{d}{d\tau} \Ric(\widetilde{Y}(\tau), \widetilde{Y}(\tau)) \: & = \:
\Ric_\tau(\widetilde{Y}, \widetilde{Y}) \: + \:
(\nabla_X \Ric)(\widetilde{Y}, \widetilde{Y}) \: + \: 2
\Ric(\nabla_X \widetilde{Y}, \widetilde{Y}) \\
& = \: \Ric_\tau(\widetilde{Y}, \widetilde{Y}) \: + \:
(\nabla_X \Ric)(\widetilde{Y}, \widetilde{Y}) \: + \:
\frac{1}{\tau} \: \Ric(\widetilde{Y}, \widetilde{Y}) \: - \:
2 |\Ric(\widetilde{Y}, \cdot)|^2, \notag
\end{align}
one obtains
\begin{align} \label{sub2}
& - \: 2 \sqrt{\overline{\tau}} \Ric({Y}({\overline{\tau}}),
{Y}({\overline{\tau}})) \: = \: - \: 2 \:
\int_0^{\overline{\tau}} \frac{d}{d\tau} \left(
\sqrt{\tau} \: \Ric(\widetilde{Y}, \widetilde{Y}) \right) \: d\tau \: =  \\
& - \: \int_0^{\overline{\tau}} \sqrt{\tau}
\left[ \frac{1}{\tau} \: \Ric(\widetilde{Y}, \widetilde{Y})
\: + \: 2 \Ric_\tau (\widetilde{Y}, \widetilde{Y})
\: + \: 2 (\nabla_X \Ric)(\widetilde{Y}, \widetilde{Y}) \: + \:
\frac{2}{\tau} \: \Ric(\widetilde{Y}, \widetilde{Y})
\: - \: 4 |\Ric(\widetilde{Y}, \cdot)|^2 \right]. \notag
\end{align}
Combining (\ref{sub1}) and (\ref{sub2}) gives
\begin{equation} \label{sub3}
\Hess_L(Y({\overline{\tau}}), Y({\overline{\tau}})) \: \le \:
Q(\widetilde{Y}, \widetilde{Y}) \: = \:
\frac{1}{\sqrt{\overline{\tau}}}
 - \: 2 \sqrt{\overline{\tau}} \Ric({Y}({\overline{\tau}}),
{Y}({\overline{\tau}})) \: - \: \int_0^{\overline{\tau}}
\sqrt{\tau} \: H(X, \widetilde{Y}) \: d\tau,
\end{equation}
where
\begin{align} \label{Htilde}
H(X, \widetilde{Y}) \: = \: & - \: \Hess_R
(\widetilde{Y},\widetilde{Y}) \: - \: 2 \langle
R(\widetilde{Y}, X)\widetilde{Y}, X \rangle
\: -\: 4 \left( \nabla_X \Ric(\widetilde{Y}, \widetilde{Y})
\: - \: \nabla_{\widetilde{Y}} \Ric(\widetilde{Y}, X) \right) \\
&  - \: 2 \Ric_\tau(\widetilde{Y}, \widetilde{Y}) \: + \:
2 \big| \Ric(\widetilde{Y}, \cdot) \big|^2
\: - \: \frac{1}{\tau} \: \Ric(\widetilde{Y}, \widetilde{Y}). \notag
\end{align}
is the expression appearing in
(\ref{Hexpression}), after the change $\tau = - t$ and $X \rightarrow -X$.
Note that $H(X, \widetilde{Y})$ is a quadratic form in
$\widetilde{Y}$. For its relation to the expression
$H(X)$ from (\ref{Hexp}), see Appendix \ref{appharnack}.


\section{I.(7.10). The Laplacian of $L$} \label{7.10}

In this section we estimate $\triangle L$.

Let $\{Y_i(\overline{\tau})\}_{i=1}^n$ be an orthonormal basis of
$T_{\gamma(\overline{\tau})}M$.
Solve for $\widetilde{Y}_i(\tau)$ from (\ref{diffeq}).
Putting
$\widetilde{Y}_i(\tau)
\: = \: \left( \frac{\tau}{\overline{\tau}} \right)^{1/2}
e_i(\tau)$,
the vectors $\{e_i(\tau)\}_{i=1}^n$ form an orthonormal basis of
$T_{\gamma(\tau)}M$.
Substituting into
(\ref{sub3}) and summing over $i$ gives
\begin{equation}
\triangle L \: \le \:
\frac{n}{\sqrt{\overline{\tau}}} \: - \: 2 \sqrt{\overline{\tau}} R
\: - \: \frac{1}{\overline{\tau}} \: \int_0^{\overline{\tau}}
\tau^{3/2} \: \sum_i H(X, e_i) \: d\tau.
\end{equation}
Then from (\ref{together}),
\begin{align} \label{touse2}
\triangle L \: \le \: &
\frac{n}{\sqrt{\overline{\tau}}} \: - \: 2 \sqrt{\overline{\tau}} R
\: - \: \frac{1}{\overline{\tau}} \: \int_0^{\overline{\tau}}
\tau^{3/2} \: H(X) \: d\tau \\
= \: & \frac{n}{\sqrt{\overline{\tau}}} \: - \: 2 \sqrt{\overline{\tau}} R
\: - \: \frac{1}{\overline{\tau}} \: K. \notag
\end{align}

\section{I.(7.11). Estimates on ${\mathcal L}$-Jacobi fields}

In this section we estimate the growth rate of an ${\mathcal L}$-Jacobi field.

Given an ${\mathcal L}$-Jacobi field $Y$, we have
\begin{equation}
\frac{d}{d\tau} |Y|^2 \: =\: 2 \Ric(Y, Y) \: + \: 2 \langle \nabla_X Y, Y
\rangle \: = \: 2 \Ric(Y, Y) \: + \: 2 \langle \nabla_Y X, Y \rangle.
\end{equation}
Thus
\begin{equation}
\frac{d |Y|^2}{d\tau} \: \Big|_{\tau = \overline{\tau}}
= \: 2 \Ric(Y(\overline{\tau}), Y(\overline{\tau}))
\: + \: \frac{1}{\sqrt{\overline{\tau}}} \:
\Hess_L(Y(\overline{\tau}), Y(\overline{\tau})),
\end{equation}
where we have used (\ref{hesseq}).

Let $\tilde{Y}$ be a field along $\gamma$ as in Section \ref{7.9},
satisfying (\ref{diffeq}) with $\tilde{Y}(\overline{\tau}) \: = \:
{Y}(\overline{\tau})$ and $|{Y}(\overline{\tau})| \: = \: 1$.
Then from (\ref{sub3}),
\begin{equation}
\frac{1}{\sqrt{\overline{\tau}}} \:
\Hess_L(Y(\overline{\tau}), Y(\overline{\tau})) \: \le \:
\frac{1}{\overline{\tau}} \: - \:
2 \Ric(Y(\overline{\tau}),Y(\overline{\tau})) \: - \:
\frac{1}{\sqrt{\overline{\tau}}} \int_0^{\overline{\tau}}
\sqrt{\tau} \: H(X, \tilde{Y}) \: d\tau.
\end{equation}
Thus
\begin{equation} \label{dY}
\frac{d |Y|^2}{d\tau} \: \Big|_{\tau = \overline{\tau}} \: \le \:
\frac{1}{\overline{\tau}} \: - \:
\frac{1}{\sqrt{\overline{\tau}}} \int_0^{\overline{\tau}}
\sqrt{\tau} \: H(X, \tilde{Y}) \: d\tau.
\end{equation}
The inequality is sharp if and only if the first inequality in
(\ref{sub3}) is an equality.
This is the case if and only if $\tilde{Y}$ is actually the
${\mathcal L}$-Jacobi field $Y$, in which case
\begin{equation} \label{sharp}
\frac{1}{\overline{\tau}} \: = \:
\frac{d |\tilde{Y}|^2}{d\tau} \: \Big|_{\tau = \overline{\tau}} \: = \:
\frac{d |Y|^2}{d\tau} \: \Big|_{\tau = \overline{\tau}}
\: = \:
2 \Ric(Y(\overline{\tau}),Y(\overline{\tau})) \: + \:
\frac{1}{\sqrt{\overline{\tau}}} \:
\Hess_L(Y(\overline{\tau}), Y(\overline{\tau})).
\end{equation}

\section{Monotonicity of
the reduced volume $\widetilde{V}$} \label{monoV}

In this section we show that the reduced volume $\widetilde{V}(\tau)$ is
monotonically nonincreasing in $\tau$.


Fix $p \in M$. Define $l(q, \tau)$
as in (\ref{reducedl}). In order to show that
$\tilde{V}(\tau)$ is well-defined in the noncompact case,
we will need a lower bound on $l(q, \tau)$. For later
use, we prove something slightly more general. Recall
that we are assuming that we have bounded curvature on
compact time intervals,
and that time slices are complete.

\begin{lemma}[estimates for reduced volume]
  \label{lem:kl-monotonicity-reduced-volume-estimates-reduced-volume}
  \dcref{Lemma 23.1}
  \group{Kleiner--Lott Reduced Geometry}
  \source{kleiner-lott}
 \label{lowerbound}
Given $0 < \overline{\tau}_1 \le \overline{\tau}_2$,
there constants $C_1, C_2 > 0$ so that for all $\tau \in
[\overline{\tau}_1, \overline{\tau}_2]$ and all
$q \in M$, we have
\begin{equation} \label{lbound}
l(q, \tau) \: \ge \:
C_1 \: d(p,q)^2 \:
- \: C_2.
\end{equation}
\end{lemma}
\begin{proof}
We write ${\mathcal L}$ in the form (\ref{changetos}).
Given an ${\mathcal L}$-geodesic $\gamma$ with
$\gamma(0) \: = \: p$ and $\gamma({\tau}) \: = \: q$, we obtain
\begin{equation}
{\mathcal L}(\gamma) \: \ge \:
\frac12 \int_0^{\sqrt{\tau}} \left| \frac{d\gamma}{ds} \right|^2 \: ds \:
- \: \const
\end{equation}
As the multiplicative change in the metric between times
$\overline{\tau}_1$ and $\overline{\tau}_2$ is bounded by a factor
$e^{\const \: (\overline{\tau}_2 - \overline{\tau}_1)}$, it follows that
$L(q, \tau) \: \ge \:
\const \: d(p,q)^2 \:
- \: \const$, where the distance is measured at time $\tau$.
The lemma follows.
\end{proof}

Define
$\tilde{V}(\tau)$
as in (\ref{reducedv}).
As the volume of time-$\tau$ balls in $M$ increases at most
exponentially fast in the radius, it follows that
$\tilde{V}(\tau)$ is well-defined.
From the discussion in Section \ref{Lremarks}, we can write
\begin{equation}
\tilde{V}(\tau) \: = \: \int_{T_pM} \tau^{-\frac{n}{2}} \:
e^{-l({\mathcal L}exp_\tau(v),\tau)} \: {\mathcal J}(v, \tau) \:
\chi_{\tau}(v) \:
dv,
\end{equation}
where ${\mathcal J}(v, \tau) \: = \:
\det \: d\left( {\mathcal L}exp_\tau \right)_v$ is the Jacobian factor
in the change of variable and
$\chi_\tau$ is the characteristic function of the
time-$\tau$ domain $\Omega_\tau$ of Section \ref{Lremarks}.

We first show
that for each $v$, the expression
$- \: \frac{n}{2} \ln(\tau) \:
- \: l({\mathcal L}exp_\tau(v), \tau) \: + \: \ln {\mathcal J}(v, \tau)$ is
nonincreasing in $\tau$. Let $\gamma$ be the ${\mathcal L}$-geodesic
with initial vector $v \in T_pM$.
From (\ref{grad2}) and (\ref{sub4}),
\begin{equation} \label{bbound1}
\frac{dl(\gamma(\tau), \tau)}{d\tau} \: \Big|_{\tau = \overline{\tau}} \: = \:
- \: \frac{1}{2\overline{\tau}} \: l(\gamma(\overline{\tau})) \: + \:
\frac12 \: \left( R(\gamma(\overline{\tau})) \: + \:
|X(\overline{\tau})|^2 \right) \: =\:
- \: \frac12 \: \overline{\tau}^{-\frac32} \: K.
\end{equation}

Next, let
$\{Y_i\}_{i=1}^n$ be a basis for the Jacobi fields along $\gamma$
that vanish at $\tau \: =\: 0$.
We can write
\begin{equation}
\ln {\mathcal J}(v, \tau)^2 \: = \:
\ln \det \: \left( \left( d\left( {\mathcal L}exp_\tau \right)_v \right)^*
d\left( {\mathcal L}exp_\tau \right)_v \right) \: = \: \ln \det(S(\tau))
\: + \: \const,
\end{equation}
where $S$ is the matrix
\begin{equation}
S_{ij}(\tau) \: = \: \langle Y_i(\tau), Y_j(\tau) \rangle.
\end{equation}
Then
\begin{equation}
\frac{d\ln {\mathcal J}(v, \tau)}{d\tau} \: = \: \frac12 \:
\Tr \left( S^{-1} \frac{dS}{d\tau} \right).
\end{equation}
To compute the derivative at $\tau \: = \: \overline{\tau}$, we can choose
a basis so that $S(\overline{\tau}) \: = \: I_n$, i.e.
$\langle Y_i(\overline{\tau}), Y_j(\overline{\tau}) \rangle
\: = \: \delta_{ij}$.
Then using (\ref{dY}) and computing as in Section \ref{7.10},
\begin{equation} \label{bbound2}
\frac{d\ln {\mathcal J}(v, \tau)}{d\tau}
\: \Big|_{\tau = \overline{\tau}}
\: = \: \frac12 \: \sum_{i=1}^n
\frac{d|Y_i|^2}{d\tau} \: \Big|_{\tau = \overline{\tau}}
\: \le \: \frac{n}{2\overline{\tau}} \: - \:
\frac{1}{2} \: \overline{\tau}^{-\frac32} K.
\end{equation}
If we have equality then (\ref{sharp}) holds for each
$Y_i$, i.e. $2 \Ric \: + \: \frac{1}{\sqrt{\overline{\tau}}} \:
\Hess_L \: =\: \frac{g}{\overline{\tau}}$ at $\gamma(\overline{\tau})$.

From (\ref{bbound1}) and (\ref{bbound2}), we deduce that
$\tau^{-\frac{n}{2}} \:
e^{-l({\mathcal L}exp_\tau(v),\tau)} \: {\mathcal J}(v, \tau)$ is
nonincreasing in $\tau$.
Finally, recall that if
$\tau \: \le \: \tau^\prime$ then $\Omega_{\tau^\prime} \subset
\Omega_\tau$, so $\chi_\tau(v)$ is nonincreasing in $\tau$. Hence
$\tilde{V}(\tau)$ is nonincreasing in $\tau$. If it is not
strictly decreasing then we must have
\begin{equation}
2 \Ric(\tau) \: + \: \frac{1}{\sqrt{\tau}} \: \Hess_{L(\tau)}
\: = \: \frac{g(\tau)}{\tau}.
\end{equation}
on all of $M$. Hence we have a gradient shrinking soliton solution.

\section{I.(7.15). A differential inequality for $L$} \label{I.(7.15)}

In this section we discuss an important differential inequality
concerning the reduced length $l$. We use the differential
inequality to estimate  $\min l(\cdot, \tau)$ from above.
We then give a lower bound on $l$.

With $\overline{L}(q, \tau) \: = \: 2 \sqrt{\tau} \: L(q, \tau)$,
equations (\ref{touse1}) and (\ref{touse2}) imply that
\begin{equation} \label{barrier}
\overline{L}_{\overline{\tau}} \: + \: \triangle \overline{L}
\: \le \: 2n
\end{equation}
away from the time-$\overline{\tau}$
${\mathcal L}$-cut locus of $p$.
We will eventually apply the maximum principle to this
differential inequality.  However to do so, we must
discuss several  senses in which
the inequality can be made global on $M$, i.e. how it can be interpreted
on the cut locus.

The first sense is that of a barrier differential inequality.
Given $f \in C(M)$ and a function $g$ on $M$, one says that
$\triangle f \: \le \: g$ {\em in the sense of barriers}
if for all $q \in M$ and $\epsilon > 0$,
there is a neighborhood $V$ of $q$ and some $u_\epsilon \in C^2(V)$ so that
 $u_\epsilon(q) = f(q)$, $u_\epsilon \ge f$ on $V$ and
$\triangle u_\epsilon \: \le \: g + \epsilon$ on $V$
\cite{Calabi}.
There is a similar spacetime definition for
$\triangle f \: -\: \frac{\partial f}{\partial t} \: \le \: g$
\cite{Dodziuk}. The point of barrier differential inequalities is
that they allow one to apply the maximum principle just as with
smooth solutions.

We illustrate this by constructing a barrier function for
$\overline{L}$ in (\ref{barrier}). Given the spacetime point
$(q, \overline{\tau})$, let $\gamma \: : \: [0, \overline{\tau}]
\rightarrow M$ be a minimizing ${\mathcal L}$-geodesic with
$\gamma(0) = p$ and $\gamma(\overline{\tau}) = q$. Given a small
$\epsilon > 0$, let
$u_\epsilon(q^\prime, \overline{\tau}^\prime)$ be the minimum of
\begin{equation}
\int_{\epsilon}^{\overline{\tau}^\prime} \sqrt{\tau} \: \left(
R(\gamma_2(\tau)) \: + \:
\big| \dot{\gamma_2}(\tau) \big|^2 \right) \: d\tau \: + \:
\int_{0}^{\epsilon} \sqrt{\tau} \: \left(
R(\gamma(\tau)) \: + \:
\big| \dot{\gamma}(\tau) \big|^2 \right) \: d\tau
\end{equation}
among curves $\gamma_2 \: : \: [\epsilon, \overline{\tau}^\prime]
\rightarrow M$
with $\gamma_2(\epsilon) = \gamma(\epsilon)$ and
$\gamma_2(\overline{\tau}^\prime) \: = \: q^\prime$.
Because the
new basepoint $(\gamma(\epsilon), \epsilon)$
is moved in along $\gamma$
from $p$, the minimizer $\ga_2$  will be unique and will vary smoothly with $q'$, when
$q'$ is close to $q$; otherwise a second minimizer or a ``conjugate
point'' would imply that $\ga$ was not minimizing. Thus the function $u_\epsilon$
is smooth in a spacetime
neighborhood $V$ of
$(q, \overline{\tau})$. Put
$U_\epsilon(q^\prime, \overline{\tau}^\prime) \: = \:
2 \sqrt{\overline{\tau}^\prime} \:
u_\epsilon(q^\prime, \overline{\tau}^\prime)$.
By construction, $U_\epsilon \: \ge \: \overline{L}$ in $V$ and
$U_\epsilon(q, \overline{\tau}) \: = \: \overline{L}(q, \overline{\tau})$.
For small $\epsilon$,
$(U_\epsilon)_{\overline{\tau}^\prime} \: + \: \triangle U_\epsilon$
will be bounded above
on $V$ by something close to $2n$. Hence
$\overline{L}$ satisfies (\ref{barrier}) globally on $M$ in the barrier sense.

As we are assuming bounded curvature on compact time intervals,
we can now apply the maximum principle of Appendix \ref{maxprin}
to conclude that the minimum of $\overline{L}(\cdot, \overline{\tau})
\: - \: 2n \overline{\tau}$ is nonincreasing in $\overline{\tau}$.
(Note that from Lemma \ref{lowerbound}, the minimum of
$\overline{L}(\cdot, \overline{\tau})
\: - \: 2n \overline{\tau}$ exists.)


\begin{lemma}[positivity for reduced l-geometry]
  \label{lem:kl-reduced-l-geometry-positivity-reduced-l-geometry}
  \dcref{Lemma 24.3}
  \group{Kleiner--Lott Reduced Geometry}
  \source{kleiner-lott}
 \label{neg}
For small positive $\overline{\tau}$, we have
$\min \overline{L}(\cdot, \overline{\tau})
\: - \: 2n \overline{\tau} < 0$.
\end{lemma}
\begin{proof}
Consider the
static curve at the point $p$. Then
for small $\overline{\tau}$, we have
$\overline{L}(\cdot, \overline{\tau}) \le \const \overline{\tau}^2$,
from which the claim follows.
\end{proof}

(Being a bit more careful with the
estimates in the proof of Lemma \ref{lowerbound}, one sees that
$\lim_{\overline{\tau} \rightarrow 0}
\min \overline{L}(\cdot, \overline{\tau}) \: = \: 0$.)
Then for $\overline{\tau} > 0$, we must have
$\min \overline{L}(\cdot, \overline{\tau}) \: \le \:
2n \overline{\tau}$, so $\min l(\cdot, \overline{\tau}) \: \le \:
\frac{n}{2}$.

The other sense of a differential inequality is the distributional
sense, i.e. $\triangle f \: \le \: g$ if for every nonnegative
compactly-supported smooth function $\phi$ on $M$,
\begin{equation}
\int_M (\triangle \phi) \: f \: dV \: \le \: \int_M \phi \: g \: dV.
\end{equation}
A general fact is that a barrier differential inequality implies
a distributional differential inequality
\cite{Ishii,Ye1}.

We illustrate this by giving an alternative proof that
$\tilde{V}(\tau)$ is nonincreasing in $\tau$.
From (\ref{touse1}), (\ref{ttouse1}) and
(\ref{touse2}),
one finds that in the barrier sense (and hence in the distributional sense
as well)
\begin{equation}
l_{\overline{\tau}} \: - \: \triangle l \: + \: |\nabla l|^2 \: - \: R
\: + \frac{n}{2\overline{\tau}} \: \ge \: 0
\end{equation}
or, equivalently, that
\begin{equation} \label{illustrate}
\left( \partial_{\overline{\tau}} \: -\: \triangle \right) \:
\left( \tau^{-\frac{n}{2}} \: e^{-l}  \: dV \right) \: \le \: 0.
\end{equation}
Then for all nonnegative $\phi \in C^\infty_c(M)$ and
$0 \: < \: \overline{\tau}_1 \: \le \overline{\tau}_2$, one obtains
\begin{align} \label{pphieqn}
& \int_M \phi \: \overline{\tau}_2^{-\frac{n}{2}} \:
e^{-l(\cdot,\overline{\tau}_2)} \: dV(\overline{\tau}_2) \: - \:
\int_M \phi \: \overline{\tau}_1^{-\frac{n}{2}} \:
e^{-l(\cdot,\overline{\tau}_1)} \: dV(\overline{\tau}_1) \:
\: = \: \\
& \int_{\overline{\tau}_1}^{\overline{\tau}_2} \int_M  \phi \:
\left( \partial_{\overline{\tau}} \: -\: \triangle \right) \: \left(
\tau^{-\frac{n}{2}} \:
e^{-l(\cdot,\tau)} \: dV(\tau) \right) \: d\tau \: + \:
\int_{\overline{\tau}_1}^{\overline{\tau}_2} \int_M (\triangle \phi) \:
\tau^{-\frac{n}{2}} \:
e^{-l(\cdot,\tau)} \: dV(\tau) \: d\tau \: \le \notag \\
&
\int_{\overline{\tau}_1}^{\overline{\tau}_2} \int_M (\triangle \phi) \:
\tau^{-\frac{n}{2}} \:
e^{-l(\cdot,\tau)} \: dV(\tau) \: d\tau.   \notag
\end{align}
We can find a sequence $\{\phi_i\}_{i=1}^\infty$ of such functions $\phi$
with range $[0,1]$ so
that $\phi_i$ is one on $B(p, i)$, vanishes outside of $B(p, i^2)$, and
$\sup_M |\triangle \phi_i| \: \le \: i^{-1}$,
uniformly in $\tau \in [\overline{\tau}_1, \overline{\tau}_2]$.
Then to finish the argument it suffices to have a good upper bound on
$e^{-l(\cdot,\tau)}$ in terms of $d(p, \cdot)$, uniformly in
$\tau \in [\overline{\tau}_1, \overline{\tau}_2]$.
This is given by Lemma \ref{lowerbound}.
The monotonicity of $\tilde{V}$ follows.

We also note the equation
\begin{equation} \label{illustrate2}
2 \triangle l \: - \: |\nabla l|^2  \: + \: R \:  + \:
\frac{l-n}{\overline{\tau}} \: \le \: 0,
\end{equation}
which follows from (\ref{ttouse1}) and
(\ref{touse2}).

Finally, suppose that the Ricci flow exists on a time interval
$\tau \in [0, \tau_0]$. From the maximum principle of
Appendix \ref{maxprin}, $R(\cdot, \tau) \: \ge \: - \: \frac{n}{2(\tau_0 - \tau)}$.
Then one obtains a lower bound on $l$ as before, using the
better lower bound for $R$.

\section{I.7.2. Estimates on the reduced length}

In this section we
suppose that our solution has nonnegative curvature operator.
We use this to derive estimates on the reduced length $l$.

We refer to Appendix \ref{appharnack} for
Hamilton's differential Harnack inequality.  We consider a Ricci flow defined on
a  time interval $t \in [0, \tau_0]$, with bounded nonnegative curvature operator,
and put $\tau \: = \: \tau_0 - t$. The differential Harnack inequality gives the
nonnegativity of the expression in (\ref{Hexpression}).
Comparing this with the formula for $H(X,Y)$ in (\ref{Htilde}), we can write
the nonnegativity as
\begin{equation}
\left( H(X, Y) \: + \: \frac{\Ric(Y,Y)}{\tau} \right) \: + \:
\frac{\Ric(Y,Y)}{\tau_0 - \tau} \: \ge \: 0,
\end{equation}
or
\begin{equation} \label{ineqtoshow}
H(X, Y) \: \ge \: - \: \Ric(Y, Y) \left( \frac{1}{\tau} \: + \:
\frac{1}{\tau_0 - \tau} \right).
\end{equation}
Then
\begin{equation}
H(X) \: \ge \: - \: R \: \left( \frac{1}{\tau} \: + \:
\frac{1}{\tau_0 - \tau} \right).
\end{equation}
As long as $\tau \: \le \: (1-c) \: \tau_0$, equations (\ref{gradsq}) and (\ref{sub4})
give
\begin{align}
4 \tau |\nabla l|^2 \: & = \: - \: 4 \tau R \: + \: 4l \: - \:
\frac{4}{\sqrt{\tau}} \int_0^\tau \tilde{\tau}^{3/2} H(X) \: d\tilde{\tau} \\
& \le \: - \: 4 \tau R \: + \: 4l \: + \:
\frac{4}{\sqrt{\tau}} \int_0^\tau \tilde{\tau}^{3/2} R \:
\left( \frac{1}{\tilde{\tau}} \: + \:
\frac{1}{\tau_0 - \tilde{\tau}} \right)
  \: d\tilde{\tau} \notag \\
& = \: - \: 4 \tau R \: + \: 4l \: + \:
\frac{4}{\sqrt{\tau}} \int_0^\tau \sqrt{\tilde\tau} R \:
\frac{\tau_0}{\tau_0 - \tilde{\tau}}
\: d\tilde{\tau} \notag \\
& \le \: - \: 4 \tau R \: + \: 4l \: + \:
\frac{4}{c \sqrt{\tau}} \int_0^\tau \sqrt{\tilde\tau} R
\: d\tilde{\tau} \notag \\
& \le \: - \: 4 \tau R \: + \: 4l \: + \:
\frac{8l}{c}, \notag
\end{align}
where the last line uses (\ref{defL}).
Thus
\begin{equation} \label{I.(7.16)}
|\nabla l|^2 \: + \: R \: \le \: \frac{Cl}{\tau}
\end{equation}
for a constant $C = C(c)$.
One shows similarly that
\begin{equation}
\frac{d}{d\tau} \ln |Y|^2 \: \le \: \frac{1}{\tau} \: (Cl+1)
\end{equation}
for an ${\mathcal L}$-Jacobi field $Y$, using (\ref{dY}).

\section{I.7.3. The no local collapsing theorem II} \label{I.7.3}

In this section we use the reduced volume to prove a no-local-collapsing
theorem for a Ricci flow on a finite time interval.

\begin{definition}[kappa noncollapsing]
  \label{def:kl-no-local-collapsing-theorem-ii-kappa-noncollapsing}
  \dcref{Definition 26.1}
  \group{Kleiner--Lott Reduced Geometry}
  \source{kleiner-lott}
 \label{noncollapsed2}
We now say that a Ricci flow solution $g(\cdot)$ defined on a time
interval $[0, T)$ is {\em $\kappa$-noncollapsed on the scale
$\rho$} if for each
$r < \rho$ and all $(x_0, t_0) \in M \times [0,T)$
with $t_0 \ge r^2$,
whenever it is true that
$|\Rm(x,t)| \: \le \: r^{-2}$ for every $x \in B_{t_0}(x_0, r)$
and $t \in [t_0 - r^2, t_0]$, then we also have
$\vol(B_{t_0}(x_0, r)) \: \ge \: \kappa r^n$.
\end{definition}

Definition \ref{noncollapsed2} differs from
Definition \ref{noncollapsed1} by the requirement that the
curvature bound holds in the entire parabolic region
$B_{t_0}(x_0, r) \times [t_0 - r^2, t_0]$ instead of just on
the ball $B_{t_0}(x_0, r)$ in the final time slice.  Therefore a Ricci
flow which is $\kappa$-noncollapsed in the sense of
Definition \ref{noncollapsed1} is also
$\kappa$-noncollapsed in the sense of
Definition \ref{noncollapsed2}.

\begin{theorem}[existence for kappa solutions]
  \label{thm:kl-no-local-collapsing-theorem-ii-existence-kappa-solutions}
  \dcref{Theorem 26.2}
  \group{Kleiner--Lott Reduced Geometry}
  \source{kleiner-lott}
  \uses{def:kl-no-local-collapsing-theorem-ii-kappa-noncollapsing}
 \label{nolocalcollapse}
Given numbers $n \in \Z^+$,  $T < \infty$
and $\rho, K, c > 0$, there is a number $\kappa = \kappa(n,K,c,\rho, T)>0$
with the following property.
Let $(M^n, g(\cdot))$ be a Ricci flow solution
defined on a
time interval $[0, T)$ with $T < \infty$, such that the curvature $|\Rm|$
is bounded on every compact subinterval $[0,T']\subset [0,T)$.
Suppose that $(M, g(0))$ is a complete
Riemannian manifold with $|\Rm| \: \le \: K$ and
$\inj(M, g(0)) \: \ge \: c \: > \: 0$. Then
the Ricci flow solution is $\kappa$-noncollapsed on the scale $\rho$,
in the sense of Definition \ref{noncollapsed2}.
Furthermore, with the other constants fixed, we can take
$\kappa$ to be nonincreasing in $T$.
\end{theorem}
\begin{proof}
We first observe that the existence of $\L$-geodesics and the
monotonicity of the reduced volume are valid in this setting; see Section \ref{Lremarks}.

Suppose that the theorem were false.  Then for given $T < \infty$
and $\rho, K, c > 0$, there are : \\
1. A sequence
$\{(M_k, g_k(\cdot))\}_{k=1}^\infty$ of Ricci flow solutions,
each defined on the
time interval $[0, T)$, with $|\Rm| \: \le \: K$ on $(M_k, g_k(0))$ and
$\inj(M_k, g_k(0)) \: \ge \: c$, \\
2. Spacetime points
$(p_k, t_k) \in M_k \times [0, T)$ and \\
3. Numbers $r_k \in (0, \rho)$ \\
having the following property :
$t_k \ge r_k^2$ and
if we put $B_k \: = \: B_{t_k}(p_k, r_k) \subset M_k$ then
$|\Rm|(x, t) \: \le \: r_k^{-2}$ whenever $x \in B_k$ and
$t \in [t_k - r_k^2, t_k]$, but
$\epsilon_k \: = \: r_k^{-1} \: \vol(B_k)^{\frac1n} \rightarrow 0$
as $k \rightarrow \infty$.
From short-time curvature estimates along with the assumed bounded geometry
at time zero, there is some $\overline{t} > 0$ so that
we have uniformly bounded geometry on the time interval
$[0, \overline{t}]$. In particular, we may assume that each
$t_k$ is greater than $\overline{t}$.

We define $\tilde{V}_k$ using curves
going backward in real time from the
basepoint $(p_k, t_k)$, i.e. forward in $\tau$-time from
$\tau = 0$.
The first step is to show that
$\tilde{V}_k(\epsilon_k r_k^2)$ is small. Note that
$\tau \:  = \: \epsilon_k r_k^2$ corresponds to a real time of
$t_k - \epsilon_k r_k^2$, which is very close to $t_k$.

Given an ${\mathcal L}$-geodesic $\gamma(\tau)$ with
$\gamma(0) \: = \: p_k$ and
velocity vector
$X(\tau) \: = \: \frac{d\gamma}{d\tau}$, its initial vector is
$v \: = \: \lim_{\tau \rightarrow 0} \sqrt{\tau} \: X(\tau)
\in T_{p_k} M_k$.
We first want
to show that if $|v| \: \le \: .1 \: \epsilon_k^{- 1/2}$ then
$\gamma$ does not escape from $B_k$ in time
$\epsilon_k r_k^2$.

We have
\begin{align}
\frac{d}{d\tau} \langle X(\tau), X(\tau) \rangle \: = \:
& 2 \: \Ric(X, X) \: + \: 2 \langle X, \nabla_X X \rangle \\
\: = \: &
2 \: \Ric(X, X) \: + \: \langle X,
\nabla R \: - \: \frac{1}{\tau} X \: - \: 4 \Ric(X, \cdot)
 \rangle \notag \\
\: = \: &
 - \: \frac{|X|^2}{\tau} \: - \: 2 \: \Ric(X, X) \: + \: \langle X,
\nabla R \rangle, \notag
\end{align}
so
\begin{equation}
\frac{d}{d\tau} \left( \tau |X|^2 \right) \: = \:
- \: 2 \: \tau \: \Ric(X, X) \: + \: \tau \: \langle X,
\nabla R \rangle.
\end{equation}
Letting $C$ denote a generic $n$-dependent constant,
for $x \in B(p_k, r_k/2)$ and $t \in [t_k \: - \: r_k^2/2,
t_k]$, the fact that $g_k$ satisfies the Ricci flow gives
an estimate $|\nabla R|(x, t) \: \le \: C r_k^{-3}$, as follows
from the case $l=0$, $m=1$ of
Appendix \ref{applocalder}.
Then in terms of dimensionless variables,
\begin{equation}
\Big| \frac{d}{d(\tau/r_k^2)} \left( \tau |X|^2 \right) \Big| \: \le \:
C \: \tau \: |X|^2 \: + \: C \: (\tau/r_k^2)^{1/2} \:
(\tau |X|^2)^{1/2}.
\end{equation}
Equivalently,
\begin{equation}
\Big| \frac{d}{d(\tau/r_k^2)} \left( \sqrt{\tau} \:
|X| \right) \Big| \: \le \:
C \: \sqrt{\tau} \: |X| \: + \: C \: \left(
\tau/r_k^2 \right)^{1/2}.
\end{equation}

Let us rewrite this as
\begin{equation} \label{diffmess}
\Big| \frac{d}{d( \frac{\tau}{\epsilon_k r_k^2})} \left(
\epsilon_k^{\frac12} \: \sqrt{\tau} \:
|X| \right) \Big| \: \le \:
C \: \epsilon_k \: \left( \epsilon_k^{\frac12} \: \sqrt{\tau} \: |X| \right)
\: + \: C \: \epsilon_k^2 \: \left(
\frac{\tau}{\epsilon_k r_k^2} \right)^{1/2}.
\end{equation}
We are interested in the time range when $\frac{\tau}{\epsilon_k r_k^2}
\in [0,1]$ and the initial condition satisfies
$\lim_{\tau \rightarrow 0} \epsilon_k^{\frac12} \: \sqrt{\tau} \:
|X|(\tau) \: \le \: .1$. Then because of the $\epsilon_k$-factors
on the right-hand side, it follows from (\ref{diffmess})
that for large $k$, we will have $\epsilon_k^{\frac12} \: \sqrt{\tau} \:
|X|(\tau) \: \le \: .11$ for all $\tau \in [0, \epsilon_k r_k^2]$.
Next,
\begin{equation}
\int_0^{\epsilon_k r_k^2} |X(\tau)| \: d\tau \: = \:
\epsilon_k^{- \: \frac12} \: \int_0^{\epsilon_k r_k^2}
\epsilon_k^{\frac12} \: \sqrt{\tau} \:
|X|(\tau) \: \frac{d\tau}{\sqrt{\tau}} \: \le \:
.11 \:
\epsilon_k^{- \: \frac12} \: \int_0^{\epsilon_k r_k^2}
\frac{d\tau}{\sqrt{\tau}} \: = \: .22 \: r_k.
\end{equation}

From the Ricci flow equation $g_\tau \: = \: 2 \: \Ric$, it follows
that the metrics $g(\tau)$ between $\tau \: = \: 0$ and
$\overline{\tau} \: = \: \epsilon_k r_k^2$
are $e^{C \epsilon_k}$-biLipschitz close to each other.  Then for
$\epsilon_k$ small,  the length
of $\gamma$, as measured with the metric at time $t_k$, will be at most
$.3 \: r_k$. This shows that $\gamma$ does not leave $B_k$ within
time $\epsilon_k r_k^2$.

Hence the contribution to $\tilde{V}_k(\epsilon_k r_k^2)$ coming from vectors
$v \in T_{p_k} M_k$
with $|v| \: \le \: .1 \: \epsilon_k^{- \: \frac12}$ is at most
$\int_{B_k} (\epsilon_k r_k^2)^{- \frac{n}{2}} \:
e^{- \: l(q, \epsilon_k r_k^2)} \: dq$.
We now want to give a lower
bound on $l(q, \epsilon_k r_k^2)$ for $q \in B_k$. Given the
${\mathcal L}$-geodesic $\gamma \: : \: [0, \epsilon_k r_k^2] \rightarrow M_k$
with $\gamma(0) \: = \: p_k$ and $\gamma(\epsilon_k r_k^2) \: = \: q$,
we have
\begin{equation}
{\mathcal L}(\gamma) \: \ge \: \int_0^{\epsilon_k r_k^2} \sqrt{\tau} \:
R(\gamma(\tau)) \: d\tau \: \ge \: - \:
\int_0^{\epsilon_k r_k^2} \sqrt{\tau} \:
n(n-1) \: r_k^{-2} \: d\tau \: = \: - \: \frac23 \: n(n-1) \:
\epsilon_k^{\frac32} \: r_k.
\end{equation}
Then
\begin{equation}
l(q, \epsilon_k r_k^2) \: \ge \: - \: \frac13 \: n(n-1) \: \epsilon_k.
\end{equation}
Thus the contribution to $\tilde{V}_k(\epsilon_k r_k^2)$ coming from vectors
$v \in T_{p_k}M_k$
with $|v| \: \le \: .1 \: \epsilon_k^{- \: \frac12}$ is at most
\begin{equation}
e^{\frac13 \: n(n-1) \: \epsilon_k} \:
(\epsilon_k r_k^2)^{- \frac{n}{2}} \: \vol_{t_k -
\epsilon_k r_k^2}(B_k) \: \le \:
e^{\frac13 \: n(n-1) \: \epsilon_k} \:
e^{\const \: \frac{1}{r_k^2} \: \epsilon_k r_k^2} \:
\epsilon_k^{\frac{n}{2}},
\end{equation}
which is less than $2 \epsilon_k^{\frac{n}{2}}$ for large $k$.

To estimate
the contribution to $\tilde{V}_k(\epsilon_k r_k^2)$ coming from vectors
$v \in T_{p_k}M_k$ with $|v| \: > \: .1 \: \epsilon_k^{- \: \frac12}$, we
can use the previously-shown monotonicity of the integrand in $\tau$.
As $\tau \rightarrow 0$, the Euclidean calculation of Section \ref{7Euclidean}
shows that
$\tau^{- n/2} \: e^{- \:
l({\mathcal L} exp_\tau(v), \tau)} \:
{\mathcal J}(v, \tau) \rightarrow 2^n \: e^{- \: |v|^2}$.
Then for all $\tau > 0$ and all $v \in \Omega_\tau$,
\begin{equation}
\tau^{- n/2} \: e^{- \:
l({\mathcal L}_\tau(v), \tau)} \:
{\mathcal J}(v, \tau) \: \le \: 2^n \: e^{- \: |v|^2},
\end{equation}
giving
\begin{equation}
\int_{T_{p_k}M_k - B(0, .1 \: \epsilon_k^{-1/2})}
\tau^{- n/2} \: e^{- \: l({\mathcal L}_\tau(v), \tau)} \:
J(v, \tau) \: \chi_\tau \: d^nv \: \le \:
2^n \int_{T_{p_k}M_k - B(0, .1 \: \epsilon_k^{-1/2})}
\: e^{- \: |v|^2} \: d^nv \: \le \:
 e^{- \:  \frac{1}{10 \epsilon_k}}
\end{equation}
for $k$ large.

The conclusion is that $\lim_{k \rightarrow \infty}
\tilde{V}_k(\epsilon_k r_k^2) \: = \: 0$. We now claim that there is a
uniform positive lower bound on $\tilde{V}_k(t_k)$.

To estimate $\tilde{V}_k(t_k)$ (where $\tau \: = \: t_k$ corresponds to
$t = 0$), we choose a point $q_k$ at time $t = \frac{\overline{t}}{2}$, i.e. at
$\tau \: = \: t_k - \frac{\overline{t}}{2}$, for which
$l(q_k, t_k - \overline{t}/2)
\: \le \: \frac{n}{2}$; see Section \ref{I.(7.15)}. Then we consider the concatenation of a
fixed curve $\gamma^{(k)}_1 \: : \:
[0,   t_k - \overline{t}/2] \rightarrow M_k$, having
$\gamma^{(k)}_1(0) \: = \: p_k$ and
$\gamma^{(k)}_1(t_k - \overline{t}/2) \: = \: q_k$,
with a fan of curves
$\gamma^{(k)}_2 \: : \: [t_k - \overline{t}/2, t_k] \rightarrow M_k$
having $\gamma^{(k)}_2(t_k - \overline{t}/2) \: = \: q_k$.
Because of the uniformly bounded geometry
in the spacetime region with $t \in [0, \overline{t}/2]$,
we can get an upper bound in
this way
for $l(\cdot, t_k)$ in a region around $q_k$. Integrating
$e^{-l(\cdot, t_k)}$, we get a positive lower bound on $\tilde{V}_k(t_k)$
that is uniform in $k$.

As
$\epsilon_k r_k^2 \rightarrow 0$, the monotonicity of $\tilde{V}$
implies that $\tilde{V}_k(t_k) \le \tilde{V}_k(\epsilon_k r_k^2)$
for large $k$, which is a contradiction.
\end{proof}

\section{I.8.3. Length distortion estimates}
\label{secI.8.3}

The distortion of distances under Ricci flow can be estimated in terms
of the Ricci tensor.
We first mention a crude estimate.

\begin{lemma}[length distortion estimates lemma]
  \label{lem:kl-length-distortion-estimates-length-distortion-estimates-lemma}
  \dcref{Lemma 27.1}
  \group{Kleiner--Lott Local Estimates}
  \source{kleiner-lott}
 \label{crudelemma}
If $\Ric \: \le \: (n-1) \: K$ then for $t_1 > t_0$,
\begin{equation}
\frac{\dist_{t_1}(x_0, x_1)}{\dist_{t_0}(x_0, x_1)} \: \ge \:
e^{-(n-1)K(t_1-t_0)}.
\end{equation}
\end{lemma}
\begin{proof}
For any curve $\gamma \: : \: [0,a] \rightarrow M$, we have
\begin{equation}
\frac{d}{dt} L(\gamma) \: =\:
\frac{d}{dt} \int_0^a
\sqrt{ \left\langle \frac{d\gamma}{ds}, \frac{d\gamma}{ds}
\right\rangle} \: ds \: = \:
- \: \int_0^a \Ric \left( \frac{d\gamma}{ds}, \frac{d\gamma}{ds}
\right) \: \frac{ds}{\left| \frac{d\gamma}{ds} \right|} \: \ge \:
- \: (n-1) \: K \: L(\gamma).
\end{equation}
Integrating gives
\begin{equation}
\frac{L(\gamma) \big|_{t_1}}{L(\gamma) \big|_{t_0}} \: \ge \:
e^{-(n-1)K(t_1-t_0)}.
\end{equation}
The lemma follows by taking
$\gamma$ to be a minimal geodesic at time $t_1$ between
$x_0$ and $x_1$.
\end{proof}

\begin{remark}[structural remark on length distortion estimates]
  \label{rem:kl-length-distortion-estimates-structural-remark-length-distortion-estimates}
  \dcref{Remark 27.5}
  \group{Kleiner--Lott Local Estimates}
  \source{kleiner-lott}
 \label{similar}
By a similar argument, if
$\Ric \: \ge \: - \: (n-1) \: K$ then for $t_1 > t_0$,
\begin{equation}
\frac{\dist_{t_1}(x_0, x_1)}{\dist_{t_0}(x_0, x_1)} \: \le \:
e^{(n-1)K(t_1-t_0)}.
\end{equation}
\end{remark}

We can write the conclusion of Lemma \ref{crudelemma} as
\begin{equation} \label{conclusion}
\frac{d}{dt} \dist_t(x_0, x_1) \: \ge \: - \: (n-1) K \:
\dist_t(x_0, x_1),
\end{equation}
where the derivative is interpreted in the sense of forward difference quotients.

The estimate in Lemma \ref{crudelemma} is multiplicative.
We now give an estimate that is additive in the distance.

\begin{lemma}[distance distortion lemma]
  \label{lem:kl-length-distortion-estimates-distance-distortion-lemma}
  \dcref{Lemma 27.8}
  \group{Kleiner--Lott Local Estimates}
  \source{kleiner-lott}
 \label{distancedist} (cf. Lemma I.8.3(b))
Suppose $\dist_{t_0}(x_0, x_1) \: \ge \: 2r_0$, and
$\Ric(x, t_0) \: \le \: (n-1) \: K$ for all $x\in B_{t_0}(x_0, r_0) \cup
B_{t_0}(x_1, r_0)$. Then
\begin{equation}
\frac{d}{dt} \dist_t(x_0, x_1) \: \ge \:
- \: 2 \: (n-1) \: \left( \frac23 \: K \: r_0 \: + \: r_0^{-1} \right)
\end{equation}
at time $t = t_0$.
\end{lemma}
\begin{proof}
If $\gamma$ is a normalized minimal geodesic from $x_0$ to $x_1$ with
velocity field $X(s) = \frac{d\gamma}{ds}$ then
for any piecewise-smooth
normal
vector field $V$ along $\gamma$ that vanishes at the
endpoints, the second variation formula gives
\begin{equation}
\int_0^{d(x_0, x_1)} \left( \Big| \nabla_X V \Big|^2 \: + \:
\langle R(V, X)V, X \rangle \right) \: ds \: \ge \: 0.
\end{equation}
Let $\{e_i(s)\}_{i=1}^{n-1}$ be a parallel orthonormal frame along
$\gamma$ that is perpendicular to $X$.
Put $V_i(s) \: = \: f(s) \: e_i(s)$, where
\begin{equation}
f(s) \: = \:
\begin{cases}
\frac{s}{r_0} & \text{ if } 0 \le s \le r_0, \\
1 & \text{ if } r_0 \le s \le d(x_0, x_1)-r_0, \\
\frac{d(x_0, x_1)-s}{r_0} & \text{ if } d(x_0, x_1)-r_0 \le s \le d(x_0, x_1).
\end{cases}
\end{equation}
Then $\Big| \nabla_X V_i \Big| \: = \: |f^\prime(s)|$ and
\begin{equation}
\int_0^{d(x_0, x_1)}  \Big| \nabla_X V_i \Big|^2 \: ds \: = \:
2 \int_0^{r_0} \frac{1}{r_0^2} \: ds \: = \: \frac{2}{r_0}.
\end{equation}
Next,
\begin{align}
\int_0^{d(x_0, x_1)}
\langle R(V_i, X)V_i, X \rangle \: ds \: = &
\int_0^{r_0} \frac{s^2}{r_0^2} \langle R(e_i, X)e_i, X \rangle \: ds \: +  \\
& \int_{r_0}^{d(x_0,x_1)-r_0} \langle R(e_i, X)e_i, X \rangle \: ds \: + \notag \\
& \int_{d(x_0,x_1)-r_0}^{d(x_0,x_1)} \frac{(d(x_0,x_1) - s)^2}{r_0^2} \langle R(e_i, X)e_i, X \rangle \: ds.
\notag
\end{align}
Then
\begin{align}
0 \: & \le \: \sum_{i=1}^{n-1} \int_0^{d(x_0, x_1)} \left( \Big| \nabla_X V_i \Big|^2 \: + \:
\langle R(V_i, X)V_i, X \rangle \right) \: ds \\
& = \:
\frac{2(n-1)}{r_0} \: - \: \int_0^{d(x_0, x_1)} \Ric(X,X) \: ds \: + \:
\int_0^{r_0} \left( 1 - \frac{s^2}{r_0^2} \right) \: \Ric(X, X) \: ds \: + \notag \\
& \: \: \: \: \: \:
\int_{d(x_0,x_1)-r_0}^{d(x_0,x_1)}  \left( 1 - \frac{(d(x_0,x_1) - s)^2}{r_0^2} \right) \: \Ric(X, X) \: ds. \notag
\end{align}
This gives
\begin{align}
\frac{d}{dt} \: \dist_t(x_0,x_1) \: & = \: - \int_0^{d(x_0,x_1)} \Ric(X, X) \: ds \\
& \ge \:
- \: \frac{2(n-1)}{r_0} \: - \: \int_0^{r_0} \left( 1 - \frac{s^2}{r_0^2} \right) \: \Ric(X, X) \: ds \notag \\
& \: \: \: \: \: \: - \:
\int_{d(x_0,x_1)-r_0}^{d(x_0,x_1)}  \left( 1 - \frac{(d(x_0,x_1) - s)^2}{r_0^2} \right) \: \Ric(X, X) \: ds
\notag \\
& \ge \: - \: \frac{2(n-1)}{r_0} \: - \: 2(n-1) K \cdot \frac23 \: r_0, \notag
\end{align}
which proves the lemma.
\end{proof}

We now give an additive version of Lemma \ref{crudelemma}.
\begin{corollary}[estimates for length distortion estimates]
  \label{cor:kl-length-distortion-estimates-estimates-length-distortion-estimates}
  \dcref{Corollary 27.16}
  \group{Kleiner--Lott Local Estimates}
  \source{kleiner-lott}
 \label{bound1}
\cite[Theorem 17.2]{Hamilton}
If $\Ric \: \le \: K$ with $K > 0$
then for all $x_0, x_1 \in M$,
\begin{equation}
\frac{d}{dt} \: \dist_t(x_0, x_1) \: \ge \: - \: \const(n) \: K^{1/2}.
\end{equation}
\end{corollary}
\begin{proof}
  \uses{lem:kl-length-distortion-estimates-distance-distortion-lemma}
Put $r_0 \: = \: K^{-1/2}$.
If $\dist_t(x_0, x_1) \: \le 2r_0$ then the corollary follows from
(\ref{conclusion}).
If $\dist_t(x_0, x_1) \: > 2r_0$ then
it follows from Lemma \ref{distancedist}.
\end{proof}

The proof of the next lemma is similar to that of Lemma
\ref{distancedist} and is given in I.8.

\begin{lemma}[estimates for distance distortion]
  \label{lem:kl-length-distortion-estimates-estimates-distance-distortion}
  \dcref{Lemma 27.18}
  \group{Kleiner--Lott Local Estimates}
  \source{kleiner-lott}
 (cf. Lemma I.8.3(a)) \label{I.8.3a}
Suppose that
$\Ric(x, t_0) \: \le \: (n-1) \: K$ on $B_{t_0}(x_0, r_0)$. Then
the distance function $d(x,t) = \dist_t(x,x_0)$ satisfies
\begin{equation}
d_t - \triangle d \: \ge \:
- \: (n-1) \: \left( \frac23 \: K \: r_0 \: + \: r_0^{-1} \right)
\end{equation}
at time $t = t_0$, outside of $B_{t_0}(x_0, r_0)$. The inequality must
be understood in the barrier sense (see Section \ref{I.(7.15)}) if
necessary.
\end{lemma}

\section{I.8.2. No local collapsing propagates forward in time and
to larger scales}

This section is concerned with a localized version of the
no-local-collapsing theorem.  The main result, Theorem \ref{8.2thm},
says that noncollapsing propagates forward in time and to a
larger distance scale.

We first give a local version of Definition \ref{noncollapsed2}.

\begin{definition}[kappa solutions definition]
  \label{def:kl-no-local-collapsing-propagates-forward-time-larger-scales-kappa-solutions-definition}
  \dcref{Definition 28.1}
  \group{Kleiner--Lott Local Estimates}
  \source{kleiner-lott}
 (cf. Definition of I.8.1)
A Ricci flow solution is said to be {\em $\kappa$-collapsed at
$(x_0, t_0)$, on the scale $r>0$}, if $|\Rm|(x,t) \: \le \: r^{-2}$
for all $(x,t)\in B_{t_0}(x_0,r)\times [t_0-r^2,t_0]$,
but $\vol(B_{t_0}(x_0,r^2)) \: \le \:
\kappa r^n$.
\end{definition}

\begin{theorem}[existence for kappa solutions]
  \label{thm:kl-no-local-collapsing-propagates-forward-time-larger-scales-existence-kappa-solutions}
  \dcref{Theorem 28.2}
  \group{Kleiner--Lott Local Estimates}
  \source{kleiner-lott}
 (cf. Theorem I.8.2) \label{8.2thm}
For any $0<A<\infty$, there is some $\kappa = \kappa(A)>0$ with
the following property.   Let
$g(\cdot)$ be a Ricci flow solution defined for
$t \in [0, r_0^2]$, having complete time slices and uniformly bounded
sectional curvature.
Suppose that $\vol(B_0(x_0, r_0)) \geq A^{-1} r_0^n$
and that $|\Rm|(x,t) \le \frac{1}{n r_0^2}$ for all
$(x,t)\in B_0(x_0,r_0)\times [0,r_0^2]$. Then the
solution cannot
be $\kappa$-collapsed on a scale less than $r_0$ at any point
$(x, r_0^2)$ with $x \in B_{r_0^2}(x_0,Ar_0)$.
\end{theorem}

\begin{remark}[estimates for no local collapsing propagates forward in time and to larger scales]
  \label{rem:kl-no-local-collapsing-propagates-forward-time-larger-scales-estimates-no-local-collapsing-propagates-forward-time-larger-scales}
  \dcref{Remark 28.3}
  \group{Kleiner--Lott Local Estimates}
  \source{kleiner-lott}
 \label{8.2fix}
In \cite[Theorem I.8.2]{Perelman} the assumption is that
$|\Rm|(x,t) \le r_0^{-2}$.
We make the slightly stronger assumption
that $|\Rm|(x,t) \le \frac{1}{n r_0^2}$.
The extra factor of $n$
is needed in order to assert in the proof that the region
$\{(y, t) \: : \:
\dist_{\frac12}(y,x_0)
\le \frac{1}{10}, \: t \in [0, \frac12 ]\}$ has
bounded geometry; see below.  Clearly this change of
hypothesis does not make any substantial difference in the
sequel.
\end{remark}

\begin{proof}
  \uses{rem:kl-length-distortion-estimates-structural-remark-length-distortion-estimates,thm:kl-no-local-collapsing-theorem-ii-existence-kappa-solutions}
We follow the lines of the proof of Theorem \ref{nolocalcollapse}.
By scaling, we can take $r_0 = 1$. Choose $x \in M$ with
$\dist_{1}(x,x_0) < A$. Define the reduced volume
$\tilde{V}(\tau)$ by
means of curves starting at $(x, 1)$. An effective lower bound
on $\tilde{V}(1)$ would imply that the solution is not $\kappa$-collapsed
at $(x, 1)$, on a scale less than $1$, for an appropriate $\kappa > 0$.

We first note that the geometry of the region $\{(y, t) \: : \:
\dist_{\frac12}(y,x_0)
\le \frac{1}{10}, \: t \in [0, \frac12 ]\}$ is uniformly bounded.
To see this, the upper sectional curvature bound implies that
$\Ric \: \le \: 1$, so the
distance distortion estimate of
Section \ref{secI.8.3} implies that
$B_{\frac12}(x_0, \frac{1}{10}) \subset B_0(x_0, 1)$.
In particular, $|\Rm|(y,t) \: \le \: \frac{1}{n}$
on the region. By Remark \ref{similar},
if $\dist_{\frac12}(y,x_0)
\le \frac{1}{10}$ and $t \in [0, \frac12 ]$ then
$B_0(y, \frac{1}{1000}) \subset B_t(y, \frac{1}{100})$.
The Bishop-Gromov inequality gives a lower bound for the
time-zero volume of $B_0(y, \frac{1}{1000})$, of the form
$\vol_0(B_0(y, \frac{1}{1000})) \: \ge \: C_1(n,A)$.
The Ricci flow equation then gives a lower bound for
the time-$t$ volume of $B_0(y, \frac{1}{1000})$, of the form
$\vol_t(B_0(y, \frac{1}{1000})) \: \ge \: C_2(n,A)$.
Thus the time-$t$ volume of $B_t(y, \frac{1}{100})$ satisfies
$\vol(B_t(y, \frac{1}{100})) \: \ge \: C_2(n,A)$.
This, along with the uniform sectional curvature bound,
implies that the region has uniformly bounded geometry.

If we have an effective upper bound on $\min_y l(y, \frac12)$,
where $y$ ranges over points that satisfy $\dist_{\frac12}(y,x_0)
\le \frac{1}{10}$, then we obtain a lower bound on
$\tilde{V} \left( 1 \right)$.
Thus it suffices to obtain an effective upper bound on $\min_y
l(y, \frac{1}{2})$ or,
equivalently, on $\min_y \overline{L}(y, \frac{1}{2})$
(as defined using ${\mathcal L}$-geodesics
from $(x,1)$) for $y$ satisfying $\dist_{\frac{1}{2}}(y, x_0) \: \le \:
\frac{1}{10}$.
Applying the maximum principle to (\ref{barrier}) gave an upper bound
on $\inf_M \overline{L}$. The idea is to spatially
localize this estimate near $x_0$,
by means of a radial function $\phi$.

Let $\phi = \phi(u)$ be a smooth function that equals $1$ on
$(- \infty, \frac{1}{20})$, equals infinity on
$(\frac{1}{10}, \infty)$ and is increasing on
$(\frac{1}{20}, \frac{1}{10})$, with
\begin{equation} \label{phieqn}
2 (\phi^\prime)^2/\phi - \phi^{\prime \prime} \: \ge \:
(2A+100n) \phi^\prime - C(A) \phi
\end{equation}
for some constant $C(A) < \infty$.
To satisfy (\ref{phieqn}),
it suffices to take $\phi(u) \: = \: \frac{1}{e^{(2A+100n)
(\frac{1}{10} - u)} - 1}$ for $u$ near $\frac{1}{10}$.

We claim that $\overline{L} \: + \: 2n \: + \: 1 \: \ge 1$ for
$t \ge \frac12$. To see this, from the end of Section \ref{I.(7.15)},
\begin{equation}
R(\cdot, \tau) \: \ge \: - \: \frac{n}{2(1-\tau)}.
\end{equation}
Then for $\overline{\tau} \in [0, \frac12]$,
\begin{equation}
L(q, \overline{\tau}) \: \ge \: - \:
\int_0^{\overline{\tau}} \sqrt{\tau} \: \frac{n}{2(1-\tau)}
d\tau \: \ge \: - \: n \:
\int_0^{\overline{\tau}} \sqrt{\tau} \:
d\tau \: = \: - \: \frac{2n}{3} \:
{\overline{\tau}}^{3/2}.
\end{equation}
Hence
\begin{equation}
\overline{L}(q, \overline{\tau}) \: = \: 2 \: \sqrt{\overline{\tau}} \:
L(q, \overline{\tau}) \: \ge \: - \frac{4n}{3} \:
\overline{\tau}^2 \: \ge \:
- \: \frac{n}{3},
\end{equation}
which proves the claim.

Now put
\begin{equation}
h(y,t) \: = \: \phi(d(y,t) - A(2t-1)) \:
(\overline{L}(y,1-t) +2n+1),
\end{equation}
where $d(y,t) = \dist_t(y,x_0)$.
It follows from the above claim that $h(y,t) \: \ge \: 0$
if $t \ge \frac12$.
Also,
\begin{equation}
\min_y h(y,1) \: \le \: h(x,1) \: = \:
\phi(\dist_1(x,x_0) - A) \cdot (2n+1) \: = \: 2n+1.
\end{equation}
As $\phi$ is infinite on $(\frac{1}{10}, \infty)$ and
$\overline{L}(\cdot, \frac12) +2n+1 \ge 1$, the
minimum of $h(\cdot, \frac12)$ is achieved at some
$y$ satisfying $d(y,\frac12) \: \le \: \frac{1}{10}$.

The calculations in I.8 give
\begin{equation}
\square h \: \ge \: -(2n+C(A)) h
\end{equation}
at a minimum point of $h$,
where $\square \: = \: \partial_t - \triangle$.
Then $\frac{d}{dt} h_{min}(t) \: \ge \: -(2n+C(A)) \: h_{min}(t)$,
so
\begin{equation}
h_{min} \left( \frac12 \right)
\: \le \: e^{n+\frac{C(A)}{2}} \: h_{min}(1)
\: \le \: (2n+1) \: e^{n+\frac{C(A)}{2}}.
\end{equation}
It follows that
\begin{equation}
\min_{y \: : \: d(y, \frac12) \le \frac{1}{10}}
\overline{L}(y, \frac12) +2n+1 \: \le
\: (2n+1) \: e^{n+\frac{C(A)}{2}}.
\end{equation}
This implies the theorem.
\end{proof}

\section{I.9. Perelman's differential Harnack inequality}
\label{I.9}

This section is concerned with a localized version of the ${\mathcal W}$-functional.
It is mainly used in I.10.

Let $g(\cdot)$ be a Ricci flow solution on a manifold $M$, defined for
$t \in (a,b)$. Put
$\square \: = \: \partial_t - \triangle$. For $f_1, f_2 \in C^\infty_c((a,b) \times M)$, we have
\begin{align}
0 \: & = \: \int_a^b \frac{d}{dt} \int_M f_1(t,x) f_2(t,x) \: dV \: dt \\
&  = \:
\int_a^b \int_M ((\partial_t - \triangle) f_1) \: f_2 \: dV \: + \:
\int_a^b \int_M f_1 \: (\partial_t + \triangle \: - \: R) f_2 \: dV \notag \\
&  = \:
\int_a^b \int_M (\square f_1) \: f_2 \: dV \: - \:
\int_a^b \int_M f_1 \:  \square^* f_2 \: dV, \notag
\end{align}
where $\square^* \: = \: - \: \partial_t - \triangle \: + \: R$. In this sense,
$\square^*$ is the formal adjoint to $\square$.

Now suppose that the Ricci flow is defined for $t \in [0,T)$.
Suppose that
\begin{equation}
u \: = \: (4\pi(T-t))^{- \: \frac{n}{2}} \: e^{-f}
\end{equation}
satisfies $\square^* u \: = \: 0$. Put
\begin{equation} \label{veqn}
v \: = \: [(T-t)(2\triangle f - |\nabla f|^2 + R) + f - n ] \: u.
\end{equation}
If $M$ is compact then using (\ref{id}),
\begin{equation}
{\mathcal W}(g_{ij},f,T-t) \: = \: \int_M v \: dV.
\end{equation}

\begin{proposition}[Harnack inequalities proposition]
  \label{prop:kl-perelman-s-differential-harnack-inequality-harnack-inequalities-proposition}
  \dcref{Proposition 29.5}
  \group{Kleiner--Lott Differential Harnack}
  \source{kleiner-lott}
 (cf. Proposition I.9.1) \label{localized}
\begin{equation} \label{I.(9.1)}
\square^*v \: = \: - \: 2(T-t) \: \Big| R_{ij} + \nabla_i \nabla_j f - \frac{g_{ij}}{2(T-t)}
\Big|^2 \: u.
\end{equation}
\end{proposition}
\par\medskip\noindent\textit{Source argument (not certified as complete).}\ \begingroup\itshape
We note that the right-hand side of I.(9.1) should be multiplied by $u$.

To prove the proposition, we first claim that
\begin{equation} \label{der}
\frac{d\triangle}{dt} \: = \: 2 \: R_{ij} \: \nabla_i \nabla_j.
\end{equation}
To see this, for $f_1, f_2 \in C^\infty_c(M)$, we have
\begin{equation}
\int_M f_1 \: \triangle f_2 \: dV \: = \: - \:
\int_M \langle df_1, df_2 \rangle \: dV.
\end{equation}
Differentiating with respect to $t$ gives
\begin{equation}
\int_M f_1 \: \frac{d\triangle}{dt} f_2 \: dV \: - \:
\int_M f_1 \: \triangle f_2 \: R \: dV
\: = \: - \: 2 \:
\int_M \Ric(df_1, df_2) \: dV
\: + \: \int_M \langle df_1, df_2 \rangle \: R \: dV,
\end{equation}
so
\begin{equation}
\frac{d\triangle}{dt} f_2 \: - \:
R \: \triangle f_2
\: = \: 2 \: \nabla_i (R_{ij} \: \nabla_j f_2)
\: - \: \nabla_i (R \:  \nabla_i f_2).
\end{equation}
Then (\ref{der}) follows from the traced second Bianchi identity.

Next, one can check that $\square^* u \: = \: 0$ is equivalent to
\begin{equation}
(\partial_t \: + \: \triangle) \: f \: = \: \frac{n}{2} \: \frac{1}{T-t}
\: + \: |\nabla f|^2 \: - \: R.
\end{equation}
Then one obtains
\begin{align}
u^{-1} \: \square^* v \: = \: &
- \: (\partial_t \: + \: \triangle) \: \left[ (T - t) \:
(2 \triangle f \: - \: |\nabla f|^2 \: + \: R) \: + \: f \right]
\: - \\
& 2 \langle \nabla \left[ (T - t) \:
(2 \triangle f \: - \: |\nabla f|^2 \: + \: R) \: + \: f \right],
u^{-1} \nabla u \rangle \notag \\
\: = \: & 2 \triangle f \: - \: |\nabla f|^2 \: + \: R \: - \:
(T-t) \: (\partial_t \: + \: \triangle) \:
(2 \triangle f \: - \: |\nabla f|^2 \: + \: R) \notag \\
&- \:
(\partial_t \: + \: \triangle) \: f  \: + \:
2 \: (T-t) \: \langle \nabla
(2 \triangle f \: - \: |\nabla f|^2 \: + \: R),
\nabla f \rangle \: + \: 2 \: |\nabla f|^2. \notag
\end{align}
Now
\begin{align}
(\partial_t \: + \: \triangle) \:
(2 \triangle f \: - \: |\nabla f|^2 \: + \: R) \: = \:
& 2 (\partial_t \triangle) f \: + \: 2 \triangle \:
(\partial_t \: + \: \triangle) \: f \\
& - \: (\partial_t \: + \: \triangle) \: |\nabla f|^2 \: + \:
(\partial_t \: + \: \triangle) \: R \notag \\
 \: = \: & 4 \: R_{ij} \: \nabla_i \nabla_j f \: \: + \:
2 \: \triangle \: (|\nabla f|^2 \: - \: R) \: - \: 2 \:
\Ric(df, df) \notag \\
& - \: 2 \langle \nabla f_t, \nabla f \rangle \: - \:
\triangle |\nabla f|^2
\: + \: \triangle R \: + \: 2 \: |\Ric|^2 \: + \: \triangle R \notag \\
 \: = \: & 4 \: R_{ij} \: \nabla_i \nabla_j f \: \: + \:
2 \: \triangle \: |\nabla f|^2 \: - \: 2 \:
\Ric(df, df) \notag \\
& - \: 2 \langle \nabla
(- \: \triangle f \: + \: |\nabla f|^2 \: - \: R),
\nabla f \rangle \: - \:
\triangle |\nabla f|^2  \: + \: 2 \: |\Ric|^2. \notag
\end{align}
Hence the term in $u^{-1} \: \square^* v$ proportionate to
$(T-t)^{-1}$ is
\begin{equation}
- \: \frac{n}{2} \: \frac{1}{T-t}.
\end{equation}
The term proportionate to $(T-t)^0$ is
\begin{equation}
2 \triangle f \: - \: |\nabla f|^2 \: + \: R \: - \: |\nabla f|^2
\: + \: R \: + \: 2 |\nabla f|^2 \: = \: 2 (\triangle f \: + \: R).
\end{equation}
The term proportionate to $(T-t)$ is $(T-t)$ times
\begin{align}
& - \: 4 \: R_{ij} \: \nabla_i \nabla_j f \: \: - \:
2 \: \triangle \: |\nabla f|^2 \: + \: 2 \:
\Ric(df, df) \: + \\
&  2 \langle \nabla
(- \: \triangle f \: + \: |\nabla f|^2 \: - \: R),
\nabla f \rangle \: + \:
\triangle |\nabla f|^2  \: - \: 2 \: |\Ric|^2 \: + \notag \\
& 2 \: \langle \nabla
(2 \triangle f \: - \: |\nabla f|^2 \: + \: R),
\nabla f \rangle \: = \notag \\
&- \: 4 \: R_{ij} \: \nabla_i \nabla_j f \: \: - \:
\triangle \: |\nabla f|^2 \: + \: 2 \:
\Ric(df, df) \: + \:  2 \langle \nabla \triangle f,
\nabla f \rangle \: - \: 2 \: |\Ric|^2 \: = \notag \\
&- \: 4 \: R_{ij} \: \nabla_i \nabla_j f \: \: - \:
2 \: |\Hess(f)|^2 \: - \: 2 \: |\Ric|^2. \notag
\end{align}
Putting this together gives
\begin{equation}
\square^* v \: = \: - \: 2 \: (T-t) \:
\Big|R_{ij} \: + \: \nabla_i \nabla_j f \: - \: \frac{1}{2(T-t)} \: g_{ij}
\Big|^2 \: u.
\end{equation}
This proves the proposition.
\endgroup\par\medskip

As a consequence of Proposition \ref{localized},
\begin{align}
\frac{d}{dt} {\mathcal W}(g_{ij}, f, T-t) \: & = \:
\frac{d}{dt} \int_M v \: dV \: = \:
\int_M (\partial_t \: + \: \triangle \: - \: R)v \: dV \\
& = \: 2 \: (T-t) \: \int_M
\Big|R_{ij} \: + \: \nabla_i \nabla_j f \: - \: \frac{1}{2(T-t)} \: g_{ij}
\Big|^2 \: u \: dV. \notag
\end{align}
In this sense, Proposition \ref{localized} is a local version of
the monotonicity of ${\mathcal W}$.

\begin{corollary}[monotonicity for Harnack inequalities]
  \label{cor:kl-perelman-s-differential-harnack-inequality-monotonicity-harnack-inequalities}
  \dcref{Corollary 29.19}
  \group{Kleiner--Lott Differential Harnack}
  \source{kleiner-lott}
 (cf. Corollary I.9.2) \label{whenever}
If $M$ is closed, or whenever the maximum principle holds, then
$\max v/u$ is nondecreasing in $t$.
\end{corollary}
\begin{proof}
We note that the statement of Corollary I.9.2
should have $\max v/u$ instead of
$\min v/u$.

To prove the corollary, we have
\begin{equation}
(\partial_t \: + \: \triangle) \: \frac{v}{u} \: = \:
\frac{v \square^* u \: - \: u \square^* v}{u^2} \: - \:
\frac{2}{u} \: \left\langle \nabla u, \nabla \frac{v}{u}
\right\rangle.
\end{equation}
As $\square^* u = 0$ and $\square^* v \le 0$,
the corollary now follows from the maximum principle.
\end{proof}


We now assume that the Ricci flow solution is defined on
the closed interval $[0, T]$.


\begin{corollary}[Harnack inequalities corollary]
  \label{cor:kl-perelman-s-differential-harnack-inequality-harnack-inequalities-corollary}
  \dcref{Corollary 29.21}
  \group{Kleiner--Lott Differential Harnack}
  \source{kleiner-lott}
 (cf. Corollary I.9.3) \label{corI.9.3}
Under the same assumptions, if the solution is defined for
$t \in [0,T]$
and $u$ tends to a $\delta$-function as
$t \rightarrow T$ then $v \le 0$ for all $t < T$.
\end{corollary}
\begin{proof}
Suppose that $h$ is a positive solution of $\square h = 0$. Then
\begin{equation}
\frac{d}{dt} \: \int_M h v \: dV \: = \:
\int_M \left( (\square h) \: v \: - \:
h \square^* v \right) \:  dV \: = \:
- \: \int_M h \: \square^* v \: dV \: \ge \: 0.
\end{equation}
As $t \rightarrow T$, the computation of $\int_M hv$ approaches
the flat-space calculation, which one finds to be zero; see
\cite{Ni} for details.
(Strictly speaking, the paper \cite{Ni} deals with the case when
$M$ is closed. It is indicated that the proof should extend to the
noncompact setting.)
Thus
$\int_M h(t_0) v(t_0) \: dV$ is nonpositive for all $t_0<T$. As $h(t_0)$ can be
taken to be an arbitrary positive function, and then flowed forward to
a positive solution of $\square h = 0$, it follows that
$v(t_0) \le 0$ for all $t_0 < T$.
\end{proof}

The next result compares the function $f$ used in the ${\mathcal  W}$-functional
and the function $l$ used in the reduced volume.

\begin{corollary}[convergence for Harnack inequalities]
  \label{cor:kl-perelman-s-differential-harnack-inequality-convergence-harnack-inequalities}
  \dcref{Corollary 29.23}
  \group{Kleiner--Lott Differential Harnack}
  \source{kleiner-lott}
 (cf. Corollary I.9.5) \label{compare}
Under the assumptions of the previous corollary, let $p \in M$ be the
point where the limit $\delta$-function is concentrated. Then
$f(q,t) \le l(q, T-t)$, where $l$ is the reduced distance defined using
curves starting from $(p, T)$.
\end{corollary}
\begin{proof}
  \uses{cor:kl-perelman-s-differential-harnack-inequality-monotonicity-harnack-inequalities}
Equation (\ref{illustrate}) implies that
$\square^* \left( (4 \pi \tau)^{-n/2} \: e^{-l} \right) \: \le \: 0$.
(This corrects the statement at the top of page 23 of I.)
From this and the fact that
$\square^* \left( (4 \pi \tau)^{-n/2} \: e^{-f} \right) \: = \: 0$,
the argument of the proof of Corollary \ref{whenever}  gives that
$\max e^{f-l}$ is nondecreasing in $t$, so
$\max (f-l)$ is nondecreasing in $t$. As $t \rightarrow T$ one
obtains the flat-space result, namely that $f \: - \: l$
vanishes.  Thus $f(t) \: \le \: l(T - t)$ for all $t \in [0, T)$.
\end{proof}

\begin{remark}[structural remark on Harnack inequalities]
  \label{rem:kl-perelman-s-differential-harnack-inequality-structural-remark-harnack-inequalities}
  \dcref{Remark 29.24}
  \group{Kleiner--Lott Differential Harnack}
  \source{kleiner-lott}
  \uses{cor:kl-perelman-s-differential-harnack-inequality-convergence-harnack-inequalities}

To give an alternative proof of Corollary \ref{compare}, putting $\tau \: = \: T \: - \: t$,
Corollary I.9.4 of \cite{Perelman} says that for any smooth curve $\gamma$,
\begin{equation}
\frac{d}{d\tau} \: f(\gamma(\tau), \tau) \: \le \:
\frac{1}{2} \left( R(\gamma(\tau), \tau) \: + \: \big| \dot{\gamma}(t)
\big|^2 \right) \: - \: \frac{1}{2\tau} \: f(\gamma(\tau), \tau),
\end{equation}
or
\begin{equation}
\frac{d}{d\tau} \: \left( \tau^{1/2} f(\gamma(\tau), \tau) \right) \: \le \:
\frac{1}{2} \: \tau^{1/2} \:
\left( R(\gamma(\tau), \tau) \: + \: \big| \dot{\gamma}(t)
\big|^2 \right).
\end{equation}
Take $\gamma$ to be a curve emanating from
$(p, T)$.
For small $\tau$,
\begin{equation}
f(\gamma(\tau), \tau) \sim d(p, \gamma(\tau))^2/4\tau =
O(\tau^0).
\end{equation}
Then
integration gives $\tau^{1/2} f \: \le \: \frac12 \: L$, or $f \: \le \: l$.
\end{remark}

\section{The statement of the pseudolocality theorem}
\label{statement of I.10.1}

The next theorem says that, in a localized sense, if the initial
data of a Ricci flow solution
has a lower bound on the scalar curvature and satisfies
an isoperimetric inequality close to that of Euclidean space
then there is a sectional curvature bound in a forward region.
The result is not used in the sequel.

\begin{theorem}[pseudolocality theorem]
  \label{thm:kl-statement-pseudolocality-theorem-pseudolocality-theorem}
  \dcref{Theorem 30.1}
  \group{Kleiner--Lott Pseudolocality}
  \source{kleiner-lott}
 (cf. Theorem I.10.1) \label{thmI.10.1}
For every $\alpha > 0$ there exist $\delta,\epsilon > 0$ with the
following property.  Suppose that we have a
smooth
pointed Ricci flow solution
$(M, (x_0, 0), g(\cdot))$
defined for $t \in [0, (\epsilon r_0)^2]$,
such that each time slice
is complete.
Suppose that
for any $x \in B_0(x_0, r_0)$ and $\Omega \subset B_0(x_0, r_0)$,
we have $R(x,0) \ge - r_0^{-2}$ and $\vol(\partial \Omega)^n
\ge (1-\delta) \: c_n \: \vol(\Omega)^{n-1}$, where $c_n$ is the
Euclidean isoperimetric constant.  Then
$|\Rm|(x,t) <
\alpha t^{-1} \: + \: (\epsilon r_0)^{-2}$
 whenever
$0 < t \le (\epsilon r_0)^2$ and
$d(x,t) = \dist_t(x, x_0) \le
\epsilon r_0$.
\end{theorem}

The sectional curvature bound $|\Rm|(x,t) <
\alpha t^{-1} \: + \: (\epsilon r_0)^{-2}$ necessarily blows up
as $t \rightarrow 0$, as nothing was assumed about
the sectional curvature at $t = 0$.

We first sketch the idea of the proof of Theorem
\ref{thmI.10.1}.
It is an argument by contradiction.
One takes a Ricci flow solution that satisfies the assumptions and
picks a point $(\overline{x}, \overline{t})$ where the
desired curvature bound does not hold.  One can assume, roughly speaking,
that $(\overline{x}, \overline{t})$ is the first point in the given
solution where the bound does not hold. (This will give the curvature
bound needed for taking
a limit in a sequence of counterexamples.)  One now
considers the solution $u$ to the conjugate heat equation, starting as
a $\delta$-function at $(\overline{x}, \overline{t})$, and the
corresponding function $v$. We know that $v \: \le \: 0$. The first
goal is to get a negative upper bound for the integral of $v$ over
an appropriate ball $B$ at
a time $\widetilde{t}$ near $\overline{t}$; see Section \ref{Claim3}.
The argument to get such
a bound is by contradiction.  If there were not such a bound then one
could consider a rescaled sequence of counterexamples with
$\int_B v \: dV \: \rightarrow 0$, and try to take a
limit. If one has the injectivity radius bounds needed to take a limit
then one obtains a limit solution with
$\int_B v \: dV \: = \: 0$, which implies that the limit solution is a
gradient shrinking soliton, which violates curvature assumptions. If one
doesn't have the injectivity radius bounds
then one can do a further rescaling to see that
in fact $\int_B v \: dV \: \rightarrow -\infty$ for some subsequence,
which is a contradiction.

If $M$ is compact then
$\int_M v \: dV$ is monotonically nondecreasing in $t$.
As (\ref{I.(9.1)}) is a localized version of this statement,
whether $M$ is compact or noncompact
we can use a cutoff function $h$ and equation (\ref{I.(9.1)})
to get a negative upper bound on $\int_M h v \: dV$ at time $t = 0$.
Finally, $\int_M v \: dV$ is the expression that
appears in the
logarithmic Sobolev inequality. If the isoperimetric constant
is sufficiently close to the Euclidean value $c_n$ then one concludes
that $\int_M hv \: dV$ must be bounded below by a constant close to zero,
which contradicts the negative upper bound on $\int_M hv \: dV$.

\section{Claim 1 of I.10.1. A
point selection
argument}

In Theorem \ref{thmI.10.1},
we can assume that $r_0 = 1$ and $\alpha < \frac{1}{100n}$. Fix $\alpha$
and put $M_\alpha = \{(x,t) \: : \: |\Rm(x,t)| \: \ge \: \alpha t^{-1} \}$.

The next lemma says that if we have a point $(x,t)$ where the conclusion of
Theorem \ref{thmI.10.1} does not hold then there is another point
$(\overline{x}, \overline{t})$ with $|\Rm(\overline{x}, \overline{t})|$
large (relative to $\overline{t}^{-1}$) so that any other such point
$(x^\prime, t^\prime)$ either has
$t^\prime > \overline{t}$
or is much farther from $x_0$ than $\overline{x}$ is.

\begin{lemma}[a point selection argument lemma]
  \label{lem:kl-point-selection-argument-point-selection-argument-lemma}
  \dcref{Claim 1}
  \group{Kleiner--Lott Pseudolocality}
  \source{kleiner-lott}
 (cf. Claim 1 of I.10.1) \label{Claim 1}
For any $A > 0$, if $g(\cdot)$ is a Ricci flow solution for
$t \in [0, \epsilon^2]$, with $A \epsilon < \frac{1}{100n}$, and
$|\Rm|(x,t) \ge \alpha t^{-1} + \epsilon^{-2}$
for some
$(x,t)$ satisfying
$t \in (0, \epsilon^2]$
and $d(x,t) \le \epsilon$, then
one can find $(\overline{x}, \overline{t}) \in M_\alpha$ with
$\overline{t} \in (0, \epsilon^2]$ and $d(\overline{x}, \overline{t}) <
(2A+1) \epsilon$, such that
\begin{equation} \label{I.(10.1)}
|\Rm(x^\prime,t^\prime)| \: \le \: 4 \: |\Rm(\overline{x}, \overline{t})|
\end{equation}
whenever
\begin{equation} \label{I.(10.2)}
(x^\prime, t^\prime) \in M_\alpha, \: \: \: \: t^\prime \in (0, \overline{t}],
 \: \: \: \: d(x^\prime, t^\prime) \le d(\overline{x}, \overline{t}) +
A |\Rm|^{- \: \frac12}(\overline{x}, \overline{t}).
\end{equation}
\end{lemma}
\begin{proof}
The proof is by a point selection argument as in
Appendix \ref{pointpicking}. By assumption, there is a point $(x,t)$
satisfying $t \in (0, \epsilon^2]$,
$d(x, t) \le \epsilon$ and
$|\Rm(x, t)| \ge \alpha t^{-1} + \epsilon^{-2}$.
Clearly $(x,t) \in M_\alpha$. Define points
$(x_k, t_k)$ inductively as follows. First, $(x_1, t_1) = (x, t)$. Next,
suppose that $(x_k, t_k)$ is constructed but cannot be taken for
$(\overline{x}, \overline{t})$. Then there is some point
$(x_{k+1}, t_{k+1}) \in M_\alpha$ such that $0 < t_{k+1} \le t_k$,
$d(x_{k+1}, t_{k+1}) \le d(x_k, t_k) + A |\Rm|^{- \: \frac12}(x_k,t_k)$ and
$|\Rm|(x_{k+1},t_{k+1}) > 4|\Rm|(x_k, t_k)$. Continuing in this way,
the point $(x_k, t_k)$ constructed has
$|\Rm|(x_k, t_k) \ge 4^{k-1} |\Rm|(x_1, t_1) \ge 4^{k-1} \epsilon^{-2}$.
Then
\begin{align}
d(x_k, t_k) & \le d(x_1, t_1) + A |\Rm|^{- \: \frac12}(x_1, t_1) + \ldots
+ A |\Rm|^{- \: \frac12}(x_{k-1}, t_{k-1}) \\
& \le \epsilon + 2A
|\Rm|^{- \: \frac12}(x_1, t_1) \le (2A+1) \epsilon. \notag
\end{align}
As the solution is smooth, the induction process must terminate after a
finite number of steps and the last value $(x_k, t_k)$ can be taken
for $(\overline{x}, \overline{t})$.
\end{proof}

\section{Claim 2 of I.10.1. Getting parabolic regions}

In Lemma \ref{Claim 1}, we know that (\ref{I.(10.1)}) is satisfied under the
condition (\ref{I.(10.2)}). The spacetime region described in (\ref{I.(10.2)})
is
not a product region, due to the fact that $d(x, t)$ is time-dependent.
The next goal is to obtain the estimate (\ref{I.(10.1)}) on a product region in
spacetime; this will be necessary when taking limits of Ricci flow solutions.
To get the estimate on a product region, one needs to bound
how fast distances are changing with respect to $t$.

\begin{lemma}[getting parabolic regions lemma]
  \label{lem:kl-getting-parabolic-regions-getting-parabolic-regions-lemma}
  \dcref{Claim 2}
  \group{Kleiner--Lott Pseudolocality}
  \source{kleiner-lott}
  \uses{lem:kl-point-selection-argument-point-selection-argument-lemma}
 (cf. Claim 2 of I.10.1) \label{Claim 2}
For the point $(\overline{x}, \overline{t})$ constructed in Lemma \ref{Claim 1},
\begin{equation} \label{I.(10.1a)}
|\Rm(x^\prime,t^\prime)| \: \le \: 4 \: |\Rm(\overline{x}, \overline{t})|
\end{equation}
holds whenever
\begin{equation} \label{I.(10.3)}
\overline{t} \: - \: \frac12 \: \alpha Q^{-1} \: \le \:
t^\prime \: \le \: \overline{t},
\: \: \: \:
\dist_{\overline{t}}(x^\prime, \overline{x}) \: \le \: \frac{1}{10} \: A Q^{- \: \frac12},
\end{equation}
where $Q = |\Rm(\overline{x}, \overline{t})|$.
\end{lemma}
\begin{proof}
  \uses{lem:kl-length-distortion-estimates-distance-distortion-lemma,lem:kl-point-selection-argument-point-selection-argument-lemma}
We first claim that if $(x^\prime, t^\prime)$ satisfies
$\overline{t} \: - \: \frac12 \: \alpha Q^{-1} \: \le \: t^\prime \: \le \:
\overline{t}$ and $d(x^\prime, t^\prime) \: \le \: d(\overline{x}, \overline{t})
\: + \: A Q^{-1/2}$ then $|\Rm|(x^\prime, t^\prime) \: \le \: 4Q$. To see this,
if $(x^\prime, t^\prime) \in M_\alpha$ then it is true by Lemma \ref{Claim 1}.  If
$(x^\prime, t^\prime) \notin M_\alpha$ then
$|\Rm|(x^\prime, t^\prime) \: < \: \alpha (t^\prime)^{-1}$.
As $(\overline{x}, \overline{t}) \in M_\alpha$, we know that
$Q \: \ge \: \alpha \overline{t}^{-1}$. Then $t^\prime \: \ge \:
\overline{t} \: - \: \frac12 \: \alpha Q^{-1} \: \ge \:
\frac12 \: \overline{t}$ and so $|\Rm|(x^\prime, t^\prime) \: < \: 2 \:
\alpha \overline{t}^{-1} \: \le \: 2Q$.

Thus we have a uniform curvature bound on the time-$t^\prime$ distance ball
$B(x_0, d(\overline{x}, \overline{t})
\: + \: A Q^{-1/2})$, provided that
$\overline{t} \: - \: \frac12 \: \alpha Q^{-1} \: \le \: t^\prime \: \le \:
\overline{t}$. We now claim that the time-$\overline{t}$ ball
$B(x_0, d(\overline{x}, \overline{t})
\: + \: \frac{1}{10} \: A Q^{-1/2})$ lies in the time-$t^\prime$ distance ball
$B(x_0, d(\overline{x}, \overline{t})
\: + \: A Q^{-1/2})$. To see this,
applying Lemma \ref{distancedist} with $r_0 \: = \:
\frac{1}{100} A Q^{-1/2}$ and the above curvature bound, if
$x^\prime$ is in the time-$\overline{t}$ ball
$B(x_0, d(\overline{x}, \overline{t})
\: + \: \frac{1}{10} \: A Q^{-1/2})$ then
\begin{equation} \label{32.4}
\dist_{t^\prime}(x_0, x^\prime) \: - \:
\dist_{\overline{t}}(x_0, x^\prime) \: \le \:
\frac12 \: \alpha Q^{-1} \: \cdot \: 2(n-1) \:
\left( \frac23 \cdot 4Q ( \frac{1}{100} A Q^{-1/2}) \: + \:
100 \: A^{-1} Q^{1/2} \right).
\end{equation}
Assuming that $A$ is sufficiently large
(we'll take $A \rightarrow \infty$ later)
and using the fact that $\alpha \: < \: \frac{1}{100n}$, it follows that
$d(x^\prime, t^\prime) \: \le \:
d(x^\prime, \overline{t}) \: + \:
\frac{1}{2} A Q^{-1/2} \: \le \:
d(\overline{x},\overline{t}) + A Q^{- \: \frac12}$,
which is what we want to show.  We note that
the argument also shows that is indeed self-consistent to use the curvature
bounds in the application of Lemma \ref{distancedist}.

Now suppose that  $(x^\prime, t^\prime)$ satisfies
(\ref{I.(10.3)}). By the triangle inequality,
$x^\prime$ lies in the time-$\overline{t}$ distance ball
$B(x_0, d(\overline{x}, \overline{t}) \: + \: \frac{1}{10} \: A Q^{-1/2})$.
Then $x^\prime$ is in the time-$t^\prime$ distance ball
$B(x_0, d(\overline{x}, \overline{t})
\: + \: A Q^{-1/2})$ and so $|\Rm(x^\prime, t^\prime)| \: \le \: 4Q$,
which proves the lemma.
\end{proof}

\section{Claim 3 of I.10.1. An upper bound on the integral of $v$} \label{Claim3}

We first make some remarks about the fundamental solution to the
backward heat equation. Let $(M, (\overline{x}, b), g(\cdot))$ be a smooth
one-parameter family of complete pointed Riemannian manifolds,
parametrized by $t \in (a,b]$. The fundamental solution
$u$ of the backward heat equation is a positive solution of
$\square^* u = 0$ on $M \times  (a,b)$
such that
$ u(\cdot, t)$ converges  to $\delta_{\overline{x}}$
in the distributional sense,  as $t\ra b^-$.
It is constructed as follows (cf. \cite[Section 3]{Dodziuk}).
Let $\{ D_i \}_{i=1}^\infty$ be an exhaustion of $M$ by
an increasing sequence of smooth compact codimension-zero
submanifolds-with-boundary containing $\overline{x}$ in
the interior. Let $u^{(i)}$ be the unique solution of
$\square^* u^{(i)} = 0$ on $D_i \times (a,b)$ with
$\lim_{t \rightarrow b^-} u^{(i)}(x, t) = \delta_{\overline{x}}(x)$,
as constructed using Dirichlet boundary conditions on $D_i$.
If $D_i \subset D_j$ then $u^{(i)} \le u^{(j)}$ on $D_i$,
using the maximum principle as in \cite[Lemma 3.1]{Dodziuk}.
Then the fundamental solution  is defined to be the limit
$u=\lim_{i \rightarrow \infty} u^{(i)}$, with smooth  convergence
 on compact subsets of $M \times  (a,b)$.
 The function $u$ is
independent of the choice of exhaustion sequence $\{D_i\}_{i=1}^\infty$.
For any $t \in (a,b)$, we have $\int_M u(x,t) \: dV(x) \: \le \: 1$.
If $\int_M u(x,t) \: dV(x) \: = \: 1$ for all $t$ then
we say that $(M, (\overline{x}, b), g(\cdot))$ is stochastically complete
for $\square^*$.
This will be the case if one has
bounded curvature on compact time intervals, but need not be the case in
general.

\begin{lemma}[integral bound in pseudolocality]
  \label{lem:kl-upper-bound-integral-integral-bound-pseudolocality}
  \dcref{Lemma 33.1}
  \group{Kleiner--Lott Pseudolocality}
  \source{kleiner-lott}
 \label{stoccomplete}
Let $\{(M_k, (\overline{x}_k, b), g_k(\cdot)) \}_{k=1}^\infty$ be a
sequence of
manifolds as above,
each defined on the time interval
$(a,b]$.
Suppose that $\lim_{k \rightarrow \infty}
(M_k, (\overline{x}_k, b), g_k(\cdot)) \: = \:
(M_\infty, (\overline{x}_\infty, b), g_\infty(\cdot))$
in the pointed smooth topology, and that
$(M_\infty, (\overline{x}_\infty, b), g_\infty(\cdot))$
is stochastically complete for $\square^*$.
Then after passing to a subsequence,
the fundamental solutions $\{ u_k \}_{k=1}^\infty$ converge smoothly
on compact subsets of $M_\infty \times (a,b)$ to the fundamental
solution $u_\infty$.
(Of course, we use the
pointed diffeomorphisms inherent in the statement of pointed convergence
in order to compare the $u_k$'s with $u_\infty$.)
\end{lemma}
\begin{proof}
From the uniform upper $L^1$-bound on $\{ u_k(\cdot, t) \}_{k=1}^\infty$
and parabolic
regularity, after passing to a subsequence we can assume that
$\{ u_k\}_{k=1}^\infty$ converges smoothly on compact subsets of
$M_\infty \times (a,b)$ to some function $U$. From the construction of
$u_\infty$, it follows easily that $u_\infty \le U$. For any
$t \in (a,b)$, we have
\begin{align}
\int_{M_\infty} ( U(x,t) \: - \: u_\infty(x,t) ) \: \dvol(x) \: & = \:
\int_{M_\infty} \liminf_k ( u_k(x, t) \: - \: u_\infty(x,t) ) \: \dvol(x) \\
& \le \: \liminf_k \int_{M_\infty} ( u_k(x, t) \: - \: u_\infty(x,t) ) \:
\dvol(x) \: \le \: 0, \notag
\end{align}
so $U \: = \: u_\infty$.
\end{proof}


Starting the proof of Theorem \ref{thmI.10.1}, we suppose that the
theorem is not true. Then there are
sequences $\epsilon_k \rightarrow 0$ and
$\delta_k \rightarrow 0$, and
pointed Ricci flow solutions $(M_k, (x_{0,k}, 0), g_k(\cdot))$
which
satisfy the hypotheses of the theorem but for which there is a point
$(x_k, t_k)$ with $0 < t_k \le \epsilon_k^2$,
$d(x_k, t_k) \le \epsilon_k$
and
$|\Rm|(x_k, t_k) \geq \alpha t_k^{-1} + \epsilon_k^{-2}$.
Given the flow $(M_k, g_k(\cdot))$, we reduce $\epsilon_k$ as much as
possible so that there is still such a point $(x_k, t_k)$. Then
\begin{equation} \label{largebound}
|\Rm|(x, t) < \alpha t_k^{-1} + 2\epsilon_k^{-2}
\end{equation}
whenever
$0 < t \le \epsilon_k^2$ and
$d(x, t) \le \epsilon_k$.
Put $A_k \: = \: \frac{1}{100n\epsilon_k}$. Construct points
$(\overline{x}_k, \overline{t}_k)$ as in Lemma \ref{Claim 1}.
Consider fundamental
solutions $u_k = (4\pi(\overline{t}_k - t))^{- \: \frac{n}{2}} \:
e^{-f_k}$ of $\square^* u_k = 0$ satisfying
$\lim_{t \rightarrow \overline{t}_k^-} u(x, t) = \delta_{\overline{x}_k}(x)$.
Construct the corresponding functions
$v_k$ from (\ref{veqn}).

\begin{lemma}[integral bound in pseudolocality]
  \label{lem:kl-upper-bound-integral-integral-bound-pseudolocality-2}
  \dcref{Claim 3}
  \group{Kleiner--Lott Pseudolocality}
  \source{kleiner-lott}
 (cf. Claim 3 of I.10.1) \label{Claim 3}
There is some $\beta > 0$ so that for
all sufficiently large $k$,
there is some
$\widetilde{t}_k \in [\overline{t}_k \: - \: \frac12 \: \alpha Q_k^{-1}, \overline{t}_k]$ with
$\int_{B_k} v_k \: dV_k \: \le \: - \beta$, where $Q_k \: = \: |\Rm|(\overline{x}_k,
\overline{t}_k)$ and $B_k$ is the time-$\widetilde{t}_k$ ball of radius
$\sqrt{\overline{t}_k - \widetilde{t}_k}$ centered at $\overline{x}_k$.
\end{lemma}
\begin{proof}
  \uses{cor:kl-perelman-s-differential-harnack-inequality-harnack-inequalities-corollary,lem:kl-getting-parabolic-regions-getting-parabolic-regions-lemma,lem:kl-upper-bound-integral-integral-bound-pseudolocality}
Suppose that the claim is not true.  After passing to a subsequence, we
can assume that for any choice of $\widetilde{t}_k$,
$\liminf_{k \rightarrow \infty}
\int_{B_k} v_k \: dV_k \: \ge \: 0$.


Consider the pointed solution $(M_k, (\overline{x}_k, \overline{t}_k), g_k(\cdot))$
parabolically rescaled by $Q_k$.
Suppose first that there is a subsequence so that
the injectivity radii of the scaled metrics at
$(\overline{x}_k, \overline{t}_k)$ are bounded away from zero.
Since $A_k \rightarrow \infty$, we can use
Lemma \ref{Claim 2} and
Appendix \ref{subsequence} to
take a subsequence that converges to a complete Ricci flow solution
$(M_\infty, (\overline{x}_\infty, \overline{t}_\infty), g_\infty(\cdot))$ on a
time interval $(\overline{t}_\infty - \frac12 \alpha, \overline{t}_\infty]$, with
$|\Rm| \le 4$ and $|\Rm|(\overline{x}_\infty, \overline{t}_\infty) = 1$.
Consider the fundamental solution $u_\infty$ of $\square^*$ on $M_\infty$
with $\lim_{t \rightarrow \overline{t}_\infty^-} u_\infty(x_\infty, t) \: = \:
\delta_{\overline{x}_\infty}(x_\infty)$.
As before,
let $u_k$ be the fundamental solution
of $\square^*$ on $M_k$
with $\lim_{t \rightarrow \overline{t}_k^-} u_k(x_k, t) \: = \:
\delta_{\overline{x}_k}(x_k)$.
In view of the pointed convergence of the rescalings of
$(M_k, (\overline{x}_k, \overline{t}_k), g_k(\cdot))$ to
$(M_\infty, (\overline{x}_\infty, \overline{t}_\infty),
g_\infty(\cdot))$,
Lemma \ref{stoccomplete} implies that
after passing to a further subsequence we can ensure that
$\lim_{k \rightarrow \infty} u_k \: = \: u_\infty$, with
smooth convergence on
compact subsets of $M_\infty \times (\overline{t}_\infty
 - \frac12 \alpha, \overline{t}_\infty)$.
(The curvature bounds on
$(M_\infty, (\overline{x}_\infty, \overline{t}_\infty),
g_\infty(\cdot))$ ensure that it is stochastically complete
for $\square^*$.)
From Corollary \ref{corI.9.3}, $v_\infty \: \le \: 0$.
Note that we are applying Corollary \ref{corI.9.3} on
$M_\infty \times (\overline{t}_\infty
- \frac12 \alpha, \overline{t}_\infty)$,
where we have the curvature bounds needed to use the maximum
principle.


Given $\widetilde{t}_\infty \in (\overline{t}_\infty -
\frac12 \alpha, \overline{t}_\infty)$,
let $B_\infty$ be the time-$\widetilde{t}_\infty$ ball of radius
$\sqrt{\overline{t}_\infty - \widetilde{t}_\infty}$
centered at $\overline{x}_\infty$.
In view of the smooth convergence
$\lim_{k \rightarrow \infty} u_k \: = \: u_\infty$ on
compact subsets of
$M_\infty \times (\overline{t}_\infty - \frac12 \alpha,
\overline{t}_\infty)$,
it follows that
$\int_{B_\infty} v_\infty \: dV_\infty \: = \: 0$ at time $\widetilde{t}_\infty$, so
$v_\infty$ vanishes on $B_\infty$ at time $\widetilde{t}_\infty$. Let $h$ be a solution
to $\square h = 0$ on $M_\infty \times [\widetilde{t}_\infty, \overline{t}_\infty)$ with
$h(\cdot, \widetilde{t}_\infty)$ a nonnegative nonzero function supported in
$B_\infty$. As in the proof of Corollary \ref{corI.9.3},
$\int_{M_\infty} h v_\infty \: dV_\infty$ is nondecreasing in $t$ and vanishes for
$t \: = \: \widetilde{t}_\infty$ and  $t \: \rightarrow \: \overline{t}_\infty$.
Thus $\int_{M_\infty} h v_\infty \: dV_\infty$ vanishes for all
$t \in [\widetilde{t}_\infty, \overline{t}_\infty)$. However, for $t \in
(\widetilde{t}_\infty, \overline{t}_\infty)$, $h$ is strictly positive and
$v_\infty$ is nonpositive.  Thus $v_\infty$ vanishes on $M_\infty$ for all
$t \in (\widetilde{t}_\infty, \overline{t}_\infty)$, and so
\begin{equation}
\Ric(g_\infty) \: + \: \Hess f_\infty \: - \:
\frac{1}{2(\overline{t} - t)} \: g_{\infty} \: = \: 0.
\end{equation}
on this interval.
We know that $|\Rm| \le 4$ on
$M_\infty \times (\overline{t}_\infty -
\frac12 \alpha, \overline{t}_\infty]$.
From the evolution equation,
\begin{equation}
\frac{dg_{\infty}}{dt} \: = \: - \: 2 \Ric(g_\infty) \: = \:
2 \Hess f_\infty \: - \:
\frac{1}{\overline{t} - t} \: g_{\infty}.
\end{equation}
It follows that the supremal and infimal
sectional curvatures of $g_\infty(\cdot,t)$ go like
$(\overline{t}_\infty \: - \: t)^{-1}$.
Hence $g_\infty$ is flat,
which contradicts the fact that
$|\Rm|(\overline{x}_\infty, \overline{t}_\infty) \: = \: 1$.


Suppose now that there is a subsequence so that
the injectivity radii of the scaled metrics
at $(\overline{x}_k, \overline{t}_k)$ tend to zero.
Parabolically rescale $(M_k,(\overline{x}_k, \overline{t}_k), g_k(\cdot))$
further so that the injectivity radius becomes one. After passing to a
subsequence we will have convergence to a flat Ricci flow solution
$(-\infty, 0] \times L$. The complete
flat manifold $L$ can be described as the total space of a flat
orthogonal
$\R^m$-bundle over a flat compact manifold $C$. After separating
variables, the
fundamental solution $u_\infty$ on $L$ will be Gaussian in the fiber directions and will
decay exponentially fast to a constant in the base directions, i.e.
$u(x, \tau) \: \sim \: (4\pi \tau)^{-m/2} \: e^{- \frac{|x|^2}{4\tau}} \:
\frac{1}{\vol(C)}$, where $|x|$ is the fiber norm. With this for $u_\infty$,
one finds that
$v_\infty \: = \: (m - n) \: \left( 1 + \frac12 \ln(4 \pi \tau) \right) \: u_\infty$.
With $B_\tau$ the ball around a basepoint
of radius $\sqrt{\tau}$, the integral of $u$ over
$B_\tau$ has a positive limit as $\tau \rightarrow \infty$, and so
$\lim_{\tau \rightarrow \infty} \int_{B_\tau} v_\infty \: dV_\infty \: = \: - \: \infty$.
Then there are times $\widetilde{t}_k
\in [\overline{t}_k \: - \: \frac12 \: \alpha Q_k^{-1}, \overline{t}_k]$ so that
$\lim_{k \rightarrow \infty}
\int_{B_k} v_k \: dV_k \: = \: - \infty$, which is a contradiction.
\end{proof}

\section{Theorem I.10.1. Proof of the pseudolocality theorem} \label{pfI.10.1}

Continuing with the proof of Theorem \ref{thmI.10.1}, we now use
Lemma \ref{Claim 3} to get a contradiction to a log Sobolev inequality.
For simplicity of notation, we drop the subscript $k$ and deal with a
particular $(M_k, (\overline{x}_k, \overline{t}_k), g_k(\cdot))$ for
$k$ large.
Define a smooth function $\phi$ on $\R$ which is
one on $(- \infty, 1]$, decreasing on $[1,2]$ and zero on $[2, \infty]$,
with $\phi^{\prime \prime} \ge  - 10 \phi^\prime$ and
$(\phi^\prime)^2 \le 10 \phi$. To construct $\phi$ we can
take the function which is $1$ on $(-\infty, 1]$,
$1 \: - \: 2(x-1)^2$ on $[1, 3/2]$, $2(x-2)^2$ on $[3/2, 2]$ and
$0$ on $[2, \infty)$, and
smooth it slightly.

Put $\widetilde{d}(y,t) = d(y, t) + 200 n \sqrt{t}$. We claim that if
$10A \epsilon \le \widetilde{d}(y,t) \le 20A\epsilon$ then
$d_t(y,t) \: - \: \triangle d(y,t) \: + \: \frac{100n}{\sqrt{t}} \: \ge \: 0$.
To see this, recalling that $t \in [0, \epsilon^2]$,
if $10A \epsilon \le \widetilde{d}(y,t) \le 20A\epsilon$ and $A$ is
sufficiently large then
$9A\epsilon \: \le \: d(y, t) \: \le \:
21A\epsilon$. We apply
Lemma \ref{I.8.3a} with the parameter $r_0$ of
Lemma \ref{I.8.3a} equal to $\sqrt{t}$. As
$r_0 \le \epsilon$, we have $y \notin B(x_0, r_0)$.
From (\ref{largebound}), on $B(x_0, r_0)$ we have
$|\Rm|(\cdot,t) \: \le \: \alpha \: t^{-1} \: + \: 2 \: \epsilon^{-2}$.
Then from Lemma \ref{I.8.3a}, at $(y, t)$ we have
\begin{align}
d_t \: - \: \triangle d \: & \ge \: - \: (n-1) \:
\left( \frac23 (\alpha \: t^{-1} \: + \: 2 \: \epsilon^{-2} ) t^{1/2}
\: + \: t^{-1/2} \right) \\
& = \: - \: (n-1) \:
\left( 1 + \frac23 \alpha + \frac43 \:
\epsilon^{-2} t \right) \: t^{-1/2}. \notag
\end{align}
It follows that
$d_t \: - \: \triangle d \: + \: \frac{100n}{\sqrt{t}} \: \ge \: 0$.

Now put $h(y,t) = \phi \left( \frac{\widetilde{d}(y,t)}{10A\epsilon} \right)$.
Then $\square h = \frac{1}{10A\epsilon} \left(
d_t \: - \: \triangle d \: + \: \frac{100n}{\sqrt{t}} \right) \phi^\prime \: - \:
\frac{1}{(10A\epsilon)^2} \phi^{\prime \prime}$, where the arguments of
$\phi^\prime$ and $\phi^{\prime \prime}$ are
$\frac{\widetilde{d}(y,t)}{10A\epsilon}$. Where
$\phi^\prime \neq 0$,
we have $d_t \: - \: \triangle d \: + \: \frac{100n}{\sqrt{t}} \: \ge \: 0$.
The fundamental solution $u(x,t) = (4\pi(\overline{t}-t))^{- \: \frac{n}{2}} \:
e^{-f(x,t)}$ of $\square^*$ is positive for $t \in [0, \overline{t})$
and we have
$\int_M u \: dV \: \le \: 1$ for all $t$.
(Recall that we are not assuming stochastic completeness.)
Then
\begin{align}
\left( \int_M hu \: dV \right)_t \: & = \: \int_M ((\square h) u \: - \: h
\square^* u) \: dV \:  = \: \int_M (\square h) u \: dV \: \le \: - \:
\frac{1}{(10A\epsilon)^2} \int_M \phi^{\prime \prime} u \: dV \\
& \le \frac{10}{(10A\epsilon)^2} \int_M \phi u \: dV \:
 \le \:
\frac{10}{(10A\epsilon)^2} \int_M u \: dV \: \le \: \frac{10}{(10A\epsilon)^2}. \notag
\end{align}
Hence
\begin{equation} \label{ut}
\int_M hu \: dV \Big|_{t=0} \ge \int_M hu \: dV \Big|_{t=\overline{t}} -
\frac{\overline{t}}{(A\epsilon)^2} \ge 1 - A^{-2}.
\end{equation}

Similarly, using Proposition \ref{localized} and Corollary \ref{corI.9.3},
\begin{align}
\left(- \int_M hv \: dV \right)_t \: & = \: - \: \int_M
((\square h) v - h \square^* v) \: dV \: \le \: - \int_M (\square h) v
\: dV \\
& \le \:
\frac{1}{(10A\epsilon)^2} \: \int_M \phi^{\prime \prime} v \: dV \: \le \: - \:
\frac{10}{(10A\epsilon)^2} \: \int_M \phi v \: dV \: = \: - \:
\frac{10}{(10A\epsilon)^2} \: \int_M h v \: dV. \notag
\end{align}
Consider the time $\widetilde{t}$ of Lemma \ref{Claim 3}. As
$(\overline{x}, \overline{t}) \in M_\alpha$,
$\widetilde{t} \in
[\overline{t}/2, \overline{t}]$. Then $\sqrt{\overline{t} -
\widetilde{t}} \: \le \: 2^{-1/2} \: \epsilon$ and so for large $A$,
$h$ will be one on the ball $B$
at time $\widetilde{t}$ of radius $\sqrt{\overline{t} -
\widetilde{t}}$ centered at $\overline{x}$, using
(\ref{32.4}). Then at time $\widetilde{t}$,
\begin{equation}
- \int_M hv \: dV \: \ge \: - \int_B v \: dV \: \ge \: \beta.
\end{equation}
Thus
\begin{equation}
- \int_M hv \: dV \Big|_{t=0} \: \ge \: \beta \: e^{-
\frac{\widetilde{t}}{(A\epsilon)^2}} \: \ge \:
\beta \: e^{- \frac{\overline{t}}{(A\epsilon)^2}} \: \ge \: \beta
\left( 1 - \frac{\overline{t}}{(A \epsilon)^2} \right)
\: \ge \: \beta (1 - A^{-2}).
\end{equation}

Working at time $0$, put $\widetilde{u} = hu$ and
$\widetilde{f} = f - \log h$.
In what follows we implicitly integrate over $\supp(h)$.
We have
\begin{equation}
\beta (1 - A^{-2}) \le - \int_M hv \: dV \: = \:
\int_M [(-2\triangle f + |\nabla f|^2 - R) \overline{t}
- f + n]hu \: dV.
\end{equation}

We claim that
\begin{equation}
\int_M \left( - 2 \triangle f \: + \: |\nabla f|^2 \right) \: h e^{-f} \:
dV \: = \: \int_M \left( - \: |\nabla \widetilde{f}|^2 \: + \:
\frac{|\nabla h|^2}{h^2} \right) \: he^{-f} \: dV.
\end{equation}
This follows from
\begin{align}
\int_M \left( - 2 \triangle f \: + \: |\nabla f|^2 \right) \: h e^{-f} \:
dV \: & = \:
\int_M \left( 2 \langle \nabla f, \nabla(he^{-f}) \rangle \: + \:
|\nabla f|^2 \: h e^{-f} \right) \:
dV \\
& = \: \int_M \left( 2 \langle \nabla f, \frac{\nabla h}{h} -
\nabla f \rangle \: + \:
|\nabla f|^2  \right) \: h e^{-f} \:
dV \notag \\
& = \: \int_M \langle \nabla f, 2 \frac{\nabla h}{h} -
\nabla f \rangle \: h e^{-f} \: dV \notag \\
& = \:
\int_M \langle \nabla \widetilde{f} + \frac{\nabla h}{h}, \frac{\nabla h}{h} -
\nabla \widetilde{f} \rangle \: h e^{-f} \: dV \notag \\
& = \:
\int_M \left( - \: |\nabla \widetilde{f}|^2 \: + \:
\frac{|\nabla h|^2}{h^2} \right) \: he^{-f} \: dV. \notag
\end{align}
Then
\begin{align}
& \int_M [(-2\triangle f + |\nabla f|^2 - R) \overline{t}
- f + n]hu \: dV \: = \\
& \int_M [- \overline{t} |\nabla \widetilde{f}|^2 - \widetilde{f} + n]
\widetilde{u} \: dV \: + \:
\int_M [\overline{t}(|\nabla h|^2/h - Rh) - h \log h ] u \: dV. \notag
\end{align}

Next, $\frac{|\nabla h|^2}{h} \: \le \: \frac{10}{(10A\epsilon)^2}$ and
$- \: Rh \: \le \: 1$ (from the assumed lower bound on $R$ at time zero). Then
\begin{equation}
\int_M \overline{t} \: \left( \frac{|\nabla h|^2}{h} \: - \: R h \right)
\: u \: dV \: \le \: \epsilon^2 \: \left( \frac{10}{(10A\epsilon)^2} \: + \:
1 \right) \: \le \: A^{-2} \: + \: \epsilon^2.
\end{equation}
Also,
\begin{align}
- \int_M uh \log{h} \: dV \: & = \: -
\int_{B(x_0, 20 A \epsilon) - B(x_0, 10 A \epsilon)} uh \log{h} \: dV \:
\le \: \int_{M - B(x_0, 10 A \epsilon)} u \: dV \\
&\le \: 1 \: - \: \int_{B(x_0, 10 A \epsilon)} u \: dV. \notag
\end{align}
Putting $\overline{h}(y) = \phi \left( \frac{d(y)}{5A\epsilon} \right)$, a result
similar to (\ref{ut}) shows that
\begin{equation}
\int_{B(x_0, 10 A \epsilon)} u \: dV \: \ge \:
\int_{M} \overline{h} u \: dV \: \ge \:
 1 \: - \: c A^{-2}
 \end{equation}
for an appropriate constant $c$.
Putting this together gives
\begin{equation} \label{somme}
\beta(1 - A^{-2}) \: \le \: \int_M \left( - \overline{t} |\nabla \widetilde{f}|^2 - \widetilde{f} + n
\right)
\widetilde{u} \: dV \: + \: (1+c) A^{-2} + \epsilon^2.
\end{equation}

Put $\widehat{g} = \frac{1}{2\overline{t}} g$,
$\widehat{u} = (2 \overline{t})^{\frac{n}{2}} \widetilde{u}$ and define
$\widehat{f}$ by $\widehat{u} \: = \: (2\pi)^{-\frac{n}{2}} \: e^{-\widehat{f}}$.
From (\ref{ut}) and (\ref{somme}), if we restore the subscript $k$ then
$\lim_{k \rightarrow \infty} \int_M \widehat{u}_k \: d\widehat{V}_k \: = \: 1$ and for large $k$,
\begin{equation}
\frac12 \: \beta \: \le \: \int_{M_k}
\left( - \: \frac12 \: |\nabla \widehat{f}_k|^2 - \widehat{f}_k + n
\right)
\widehat{u}_k \: d\widehat{V}_k.
\end{equation}
If we normalize $\widehat{u}_k$ by putting
$U_k \: = \: \frac{\widehat{u}_k}{\int_{M_k} \widehat{u}_k \: d\widehat{V}_k}$,
and
define $F_k$ by $U_k \: = \: (2\pi)^{- \: \frac{n}{2}} \: e^{- \: F_k}$, then
for large $k$, we also have
\begin{equation} \label{contra1}
\frac12 \: \beta \: \le \: \int_{M_k}
\left( - \frac12 |\nabla {F}_k|^2 - {F}_k + n
\right)
{U}_k \: d\widehat{V}_k.
\end{equation}

On the other hand, the
logarithmic Sobolev inequality for $\R^n$ \cite[I.(8)]{Beckner} says that
\begin{equation} \label{Beckner}
\int_{\R^n} \left( - \: \frac12 \: |\nabla F|^2 \: - \: F \: + \: n \right) U \: dV
\: \le \: 0,
\end{equation}
provided that the compactly-supported function
$U \: = \: (2 \pi)^{-n/2} \: e^{-F}$ satisfies
$\int_{\R^n} U \: dV \: = \: 1$.
As was mentioned to us by Peter Topping, one can get a sharper inequality
by applying (\ref{Beckner}) to the rescaled function
$U_c(x) \: = \: c^n \: U(cx)$ and optimizing with respect to $c$.
The result is
\begin{equation}
\int_{\R^n} |\nabla F|^2 \: U \: dV \: \ge \: n \: e^{1 \: - \: \frac{2}{n} \:
\int_{\R^n} F \: U \: dV}.
\end{equation}
Given this inequality on $\R^n$,
one can use a symmetrization argument to prove the same
inequality for a compactly-supported function on any complete Riemannian
manifold, provided that the Euclidean isoperimetric
inequality holds for domains in the support of $U$.
See, for example, \cite[Proposition 4.1]{Ni3} which gives
the symmetrization argument for
(\ref{Beckner}), attributing it to Perelman.
Again using the inequality for $\R^n$, if instead we have
$\vol(\partial \Omega)^n
\ge (1-\delta_k) \: c_n \: \vol(\Omega)^{n-1}$ for domains
$\Omega \subset \supp(U_k)$ then the symmetrization argument gives
\begin{equation} \label{contra2}
\int_{M_k} |\nabla {F}_k|^2 \: {U}_k \: d\widehat{V}_k \: \ge \:
(1 - \delta_k)^{\frac{2}{n}} \:
n \: e^{1 \: - \: \frac{2}{n} \:
\int_{M_k} {F}_k \: U_k \: d\widehat{V}_k}.
\end{equation}

Equations (\ref{contra1}) and (\ref{contra2}) imply that
\begin{equation}
\frac{n}{2} \: \left( (1 - \delta_k)^{\frac{2}{n}} \: e^{1 - \frac{2}{n} \int_{M_k}
F_k \: U_k \: d\widehat{V}_k} \: - \: 1 \: - \:
\left( 1 - \frac{2}{n} \int_{M_k}
F_k \: U_k \: d\widehat{V}_k \right) \right) \: \le \: - \: \frac{\beta}{2}.
\end{equation}
However,
\begin{equation}
\lim_{k \rightarrow \infty} \inf_{x \in \R}
\left(
(1 - \delta_k)^{\frac{2}{n}} \: e^{x} \: - \: 1 \: - \: x \right) \: = \: 0.
\end{equation}
This is a contradiction.


\section{I.10.2. The volumes of future balls}

The next result gives a lower bound on the volumes of future
balls.

\begin{corollary}[volume corollary]
  \label{cor:kl-volumes-future-balls-volume-corollary}
  \dcref{Corollary 35.1}
  \group{Kleiner--Lott Pseudolocality}
  \source{kleiner-lott}
  \uses{thm:kl-statement-pseudolocality-theorem-pseudolocality-theorem}
 (cf. Corollary I.10.2) \label{corI.10.2}
Under the assumptions of Theorem \ref{thmI.10.1}, for
$0 < t \le (\epsilon r_0)^2$ we have
$\vol(B_t(x, \sqrt{t})) \: \ge \: c t^{\frac{n}{2}}$ for $x \in B_0(x_0, \epsilon r_0)$,
where $c = c(n)$ is a universal constant.
\end{corollary}
\par\medskip\noindent\textit{Source argument (not certified as complete).}\ \begingroup\itshape  (Sketch)
If the corollary were not true then taking a sequence of counterexamples,
we can center ourselves around the
collapsing balls $B(x, \sqrt{t})$ to obtain functions $f$ as in Section \ref{pfI.10.1}.
As in the proof of Theorem \ref{noncollthm}, the volume condition
along with the fact that $\int_M (2 \pi)^{-n/2} \: e^{-f} \: dV \rightarrow 1$
means that $f \rightarrow - \infty$, which implies that
$\int_M \left( - \: \frac12 \: |\nabla f|^2 \: - \: f \: + \: n \right) u \: dV
\rightarrow \infty$. This contradicts the logarithmic Sobolev inequality.
\endgroup\par\medskip

\section{I.10.4. $\kappa$-noncollapsing at future times}

The next result gives $\kappa$-noncollapsing at future times.

\begin{corollary}[existence for kappa solutions]
  \label{cor:kl-noncollapsing-at-future-times-existence-kappa-solutions}
  \dcref{Corollary 36.1}
  \group{Kleiner--Lott Pseudolocality}
  \source{kleiner-lott}
 (cf. Corollary I.10.4)
There are $\delta,\epsilon > 0$ such that for any $A > 0$ there exists
$\kappa = \kappa(A) > 0$ with the following property.
Suppose that we have a Ricci flow solution
$g(\cdot)$ defined for $t \in [0, (\epsilon r_0)^2]$ which
has bounded $|\Rm|$ and complete time slices. Suppose that
for any $x \in B(x_0, r_0)$ and $\Omega \subset B(x_0, r_0)$,
we have $R(x,0) \ge - r_0^{-2}$ and $\vol(\partial \Omega)^n
\ge (1-\delta) \: c_n \: \vol(\Omega)^{n-1}$, where $c_n$ is the
Euclidean isoperimetric constant. If $(x,t)$ satisfies
$A^{-1} (\epsilon r_0)^2 \le t \le (\epsilon r_0)^2$ and
$\dist_t(x, x_0) \le A r_0$ then $g(\cdot)$ is not
$\kappa$-collapsed at $(x,t)$ on scales less than $\sqrt{t}$.
\end{corollary}
\begin{proof}
  \uses{cor:kl-volumes-future-balls-volume-corollary,thm:kl-no-local-collapsing-propagates-forward-time-larger-scales-existence-kappa-solutions,thm:kl-statement-pseudolocality-theorem-pseudolocality-theorem}
Using Theorem \ref{thmI.10.1} and
Corollary \ref{corI.10.2}, we can apply
Theorem \ref{8.2thm} starting at time $A^{-1} (\epsilon r_0)^2$.
\end{proof}

\section{I.10.5. Diffeomorphism finiteness}
\label{I.10.5}

In this section we prove the diffeomorphism finiteness of Riemannian
manifolds with local isoperimetric inequalities, a lower bound on
scalar curvature and an upper bound on volume.

\begin{theorem}[existence for Riemannian metrics]
  \label{thm:kl-diffeomorphism-finiteness-existence-riemannian-metrics}
  \dcref{Theorem 37.1}
  \group{Kleiner--Lott Pseudolocality}
  \source{kleiner-lott}

Given $n \in \Z^+$, there is a $\delta > 0$ with the following property.  For any
$r_0, V > 0$, there are finitely many diffeomorphism types
of compact $n$-dimensional Riemannian manifolds $(M, g_0)$ satisfying \\
1. $R \: \ge \: - r_0^{-2}$. \\
2. $\vol(M, g_0) \: \le \: V$. \\
3. Any domain $\Omega \subset M$ contained in a metric
$r_0$-ball satisfies $\vol(\partial \Omega)^n
\ge (1-\delta) \: c_n \: \vol(\Omega)^{n-1}$, where $c_n$ is the
Euclidean isoperimetric constant.
\end{theorem}
\begin{proof}
  \uses{cor:kl-volumes-future-balls-volume-corollary,thm:kl-statement-pseudolocality-theorem-pseudolocality-theorem}
Choose $\alpha > 0$. Let $\delta$ and $\epsilon$ be the parameters
of Theorem \ref{thmI.10.1}.
Consider Ricci flow $g(\cdot)$ starting from $(M, g_0)$.
Let $T > 0$ be the maximal number so that a smooth flow exists
for $t \in [0, T)$.  If $T < \infty$ then
$\lim_{t \rightarrow T^-} \sup_{x \in M} |\Rm(x, t)| \: = \: \infty$. It
follows from Theorem \ref{thmI.10.1} that $T > (\epsilon r_0)^2$.
Put $\widehat{g} = g((\epsilon r_0)^2)$. Theorem \ref{thmI.10.1} gives a
uniform double-sided sectional curvature bound on $(M, \widehat{g})$.
Corollary \ref{corI.10.2} gives a uniform lower bound on the volumes
of $(\epsilon r_0)$-balls in $(M, \widehat{g})$. Let $\{x_i\}_{i=1}^N$ be a
maximal $(2\epsilon r_0)$-separated net in $(M, \widehat{g})$.

From the lower bound $R \: \ge \: - r_0^{-2}$ on $(M, g_0)$ and the
maximum principle, we have
$R(x,t)  \: \ge \: - r_0^{-2}$ for $t \in [0, (\epsilon r_0)^2]$.
Then the Ricci flow equation gives
a uniform upper bound on $\vol(M, \widehat{g})$. This implies a uniform
upper bound on $N$ or, equivalently, a uniform upper bound on
$\diam(M, \widehat{g})$. The theorem now follows from the diffeomorphism
finiteness of $n$-dimensional Riemannian
manifolds with double-sided sectional curvature bounds,
upper bounds on diameter and lower bounds on volume.
\end{proof}




\section{I.11.1. $\kappa$-solutions}
\label{secI.11.1}

\begin{definition}[kappa noncollapsing]
  \label{def:kl-solutions-kappa-noncollapsing}
  \dcref{Definition 38.1}
  \group{Kleiner--Lott Kappa Solutions}
  \source{kleiner-lott}

Given $\kappa > 0$, a {\em $\kappa$-solution} is a Ricci flow
solution $(M,g(\cdot))$ that is defined on a time interval of the
form $(-\infty,C)$
(or $(-\infty,C]$)
such that

$\bullet$ The curvature $|\Rm|$ is bounded on each compact time interval
$[t_1,t_2]\subset (-\infty,C)$
(or $(-\infty,C]$), and each time slice $(M,g(t))$ is complete.

$\bullet$ The curvature operator is nonnegative and the scalar curvature
is everywhere positive.

$\bullet$ The Ricci flow is $\kappa$-noncollapsed at all scales.
\end{definition}


By abuse of terminology, we may sometimes write that
``$(M, g(\cdot))$ is a $\kappa$-solution'' if it is a $\kappa$-solution
for some $\kappa > 0$.


From Appendix \ref{appharnack},
$R_t  \:  \ge \: 0$ for an ancient solution.
This implies the essential equivalence of the notions of
$\kappa$-noncollapsing in
Definitions \ref{noncollapsed1} and \ref{noncollapsed2} , when
restricted to ancient solutions.
Namely, if a solution is $\kappa$-collapsed in the sense of
Definition \ref{noncollapsed2} then it is
automatically $\kappa$-collapsed in the sense of
Definition \ref{noncollapsed1}.  Conversely,
if a time-$t_0$ slice of an ancient solution is collapsed in the sense
of Definition \ref{noncollapsed1}
then the fact that $R_t  \:  \ge \: 0$, together with bounds on
distance distortion, implies that it is collapsed in the sense of
Definition \ref{noncollapsed2} (possibly for a different value of $\kappa$).

The relevance of $\kappa$-solutions is that a blowup limit of a
finite-time singularity on a compact manifold will be a $\kappa$-solution.

For examples of $\kappa$-solutions, if $n \ge 3$ then
there is a $\kappa$-solution
on the cylinder $\R \times S^{n-1}(r)$,
where the radius satisfies $r^2(t) \: = \: r_0^2 \: - \:
2(n-2)t$. There is also a $\kappa$-solution on the
$\Z_2$-quotient $\R \times_{\Z_2} S^{n-1}(r)$, where
the generator of $\Z_2$ acts by reflection on $\R$ and by
the antipodal map on $S^{n-1}$.
On the other hand,  the quotient solution on
$S^1 \times S^{n-1}(r)$ is not $\kappa$-noncollapsed for any
$\kappa > 0$,
as can be seen by looking at large negative time.

Bryant's gradient steady soliton is a
three-dimensional $\kappa$-solution given by
$g(t) \: = \: \phi_t^* g_0$, where $g_0 \: = \: dr^2 \: + \: \mu(r) \:
d\Theta^2$ is a certain rotationally symmetric metric on $\R^3$. It has
sectional curvatures that go like $r^{-1}$, and $\mu(r) \sim r$.
The gradient function $f$ satisfies $R_{ij} \: + \: \nabla_i \nabla_j f
\: = \: 0$, with $f(r) \: \sim \: -2r$. Then for $r$ and $r \: - \: 2t$
large, $\phi_t(r, \Theta) \: \sim \: (r-2t, \Theta)$. In particular,
if $R_0 \in C^\infty(\R^3)$ is the scalar curvature function of
$g_0$ then $R(t, r, \Theta) \sim R_0(r-2t, \Theta)$.

To check the conclusion of Corollary \ref{corI.11.8}
in this case,
given a point $(r_0, \Theta) \in \R^3$ at time $0$, the scalar
curvature goes like $r_0^{-1}$. Multiplying the soliton metric by
$r_0^{-1}$ and sending $t \rightarrow r_0 t$ gives the asymptotic metric
\begin{equation}
d(r/\sqrt{r_0})^2 \:
+ \: \frac{r-2r_0t}{r_0} \: d\Theta^2.
\end{equation}
Putting $u \: = \: (r - r_0)/\sqrt{r_0}$,
the rescaled metric is approximately
\begin{equation}
du^2 \: + \:
\left( 1 \: + \: \frac{u}{\sqrt{r_0}} \: - \: {2t} \right)
\: d\Theta^2.
\end{equation}
Given $\epsilon \: > 0$, this will be $\epsilon$-biLipschitz close to
the evolving cylinder
$du^2 \: + \:
\left( 1 \: - \: {2t} \right)
\: d\Theta^2$ provided that $|u| \: \le \: \epsilon \sqrt{r_0}$,
i.e. $|r-r_0| \: \le \: \epsilon r_0$. To have an $\epsilon$-neck,
we want this to hold whenever $|r \: - \: r_0|^2 \: \le \:
(\epsilon r_0^{-1})^{-1}$. This will be the case if $r_0 \: \ge \:
\epsilon^{-3}$. Thus $M_\epsilon$ is approximately
\begin{equation}
\{ (r, \Theta) \in \R^3 \: : \: r \: \le \: \epsilon^{-3} \}
\end{equation}
and $Q \: = \: R(x_0,0) \sim \epsilon^3$.
Then $\diam(M_\epsilon) \sim \epsilon^{-3}$
and at the origin $0 \in M_\epsilon$,
$R(0, 0) \sim \epsilon^0$. It follows that for the value of
$\kappa$ corresponding to this solution, $C(\epsilon, \kappa)$ must
grow at least as fast as $\epsilon^{-3}$ as $\epsilon \rightarrow 0$.

\section{I.11.2. Asymptotic solitons} \label{I.11.2}

This section shows that every $\kappa$-solution has a gradient shrinking
soliton buried inside of it, in an asymptotic sense as $t \rightarrow -\infty$.
Such a soliton will be called an {\em asymptotic soliton}.

Heuristically, the existence of an asymptotic soliton is
a consequence of the compactness results and
the monotonicity of the reduced volume.  Taking an appropriate
sequence of
spacetime points going backward in time, one constructs a
limiting rescaled solution. As the limit reduced volume
is constant in time, the monotonicity formula implies that
this limit solution is a gradient shrinking soliton. This is
the basic idea but
the rigorous argument is a bit more subtle.

Pick an arbitrary point $(p, t_0)$ in the $\kappa$-solution $(M, g(\cdot))$. Define
the reduced volume $\widetilde{V}(\tau)$ and the reduced length
$l(q, \tau)$ as in Section \ref{I.7}, by means of curves starting from
$(p, t_0)$, with $\tau = t_0 - t$. From
Section \ref{I.(7.15)}, for each $\tau> 0$ there is some
$q(\tau)\in M$ such that $l(q(\tau), \tau) \le \frac{n}{2}$. (Note that
$l \ge 0$ from the curvature assumption.)

\begin{proposition}[existence for Ricci solitons]
  \label{prop:kl-asymptotic-solitons-existence-ricci-solitons}
  \dcref{Proposition 39.1}
  \group{Kleiner--Lott Kappa Solutions}
  \source{kleiner-lott}
 (cf. Proposition I.11.2) \label{asympsoliton}
There is a sequence $\overline{\tau}_i \rightarrow \infty$ so that if we
consider the solution $g(\cdot)$ on the time interval
$[t_0 - \overline{\tau}_i, t_0 - \frac12 \overline{\tau}_i]$  and
parabolically rescale it
at the point $(q(\overline{\tau}_i), t_0 -
\overline{\tau}_i)$ by the factor $\overline{\tau_i}^{-1}$ then as
$i \rightarrow \infty$, the
rescaled solutions
converge to a nonflat gradient shrinking soliton (restricted to $[-1, - \frac12]$).
\end{proposition}
\begin{proof}
  \uses{prop:kl-perelman-s-differential-harnack-inequality-harnack-inequalities-proposition}

Equation (\ref{I.(7.16)}) implies that
$|\nabla l^{1/2}|^2 \: \le \: \frac{C}{4\tau}$, and so
\begin{equation}
|l^{1/2}(q,\tau) \: - \: l^{1/2}(q(\tau),\tau)| \:
\le \: \sqrt{\frac{C}{4\tau}}
\: \dist_{t_0-\tau}(q, q(\tau)).
\end{equation}
We apply this estimate initially at some fixed time $\tau = \overline{\tau}$, to
obtain
\begin{equation}
l(q, \overline{\tau}) \: \le \: \left( \sqrt{\frac{C}{4\overline{\tau}}}
\: \dist_{t_0- \overline{\tau}}(q, q(\overline{\tau})) \: + \:
\sqrt{\frac{n}{2}} \right)^2.
\end{equation}
From (\ref{touse1}), (\ref{ttouse1}) and (\ref{I.(7.16)}),
\begin{equation}
\partial_\tau l \: = \: \frac{R}{2} \: - \: \frac{|\nabla l|^2}{2}
\: - \: \frac{l}{2\tau} \: \ge \: - \: \frac{(1+C)l}{2\tau}.
\end{equation}
This implies that for $\tau \in \left[ \frac12 \overline{\tau},
\overline{\tau} \right]$,
\begin{equation} \label{upperl}
l(q, {\tau}) \: \le \: \left( \frac{\overline{\tau}}{\tau}
\right)^{\frac{1+C}{2}} \:
\left( \sqrt{\frac{C}{4\overline{\tau}}}
\: \dist_{t_0- \overline{\tau}}(q, q(\overline{\tau})) \: + \:
\sqrt{\frac{n}{2}} \right)^2.
\end{equation}
Also from (\ref{I.(7.16)}), we have
$\tau R \: \le \: Cl$.
Then we can plug in the previous bound on
$l$ to get an upper bound on $\tau R$ for
$\tau \in \left[ \frac12 \overline{\tau}, \overline{\tau} \right]$.
The upshot is that for any $\epsilon > 0$, one can find $\delta > 0$
so that both $l(q,\tau)$ and $\tau R(q, t_0-\tau)$ do not exceed
$\delta^{-1}$ whenever $\frac12 \overline{\tau} \le \tau \le \overline{\tau}$
and $\dist^2_{t_0-\overline{\tau}} (q, q(\overline{\tau})) \: \le \:
\epsilon^{-1} \overline{\tau}$.

Varying $\overline{\tau}$, as the rescaled solutions
(with basepoints at $(q(\overline{\tau}), t_0 - \overline{\tau})$)
are uniformly noncollapsing and have
uniform curvature bounds on balls, Appendix \ref{subsequence}
implies that
we can take
a sequence $\overline{\tau}_i \rightarrow \infty$
to get a pointed limit
$(\overline{M}, \overline{q}, \overline{g}(\cdot))$ that is
a complete Ricci flow solution (in the backward parameter $\tau$) for
$\frac12 \: < \: \tau \: < \: 1$.
We may assume that we have locally Lipschitz convergence of $l$ to
a limit function $\overline{l}$.

We define the reduced volume $\tilde{V}(\tau)$ for the limit solution
using the limit function $\overline{l}$.
We claim that for any
$\tau \in \left( \frac12, 1 \right)$, if we put $\tau_i = \tau \overline{\tau}_i$
then the number $\tilde{V}(\tau)$ for
the limit solution is the limit of numbers $\tilde{V}(\tau_i)$ for the
original solution. One wishes to apply dominated convergence to the
integrals $\int_M e^{-l(q, \tau_i)} \:
\tau_i^{-n/2} \: \dvol(q,t_0 - \tau_i)$.
(Note that $\tau_i^{-n/2} \: \dvol(q,t_0 - \tau_i) \: = \:
\tau^{-n/2} \: \overline{\tau}_i^{-n/2} \: \dvol(q,t_0 - \tau_i)$ and
$\overline{\tau}_i^{-n/2} \: \dvol(q,t_0 - \tau_i)$ is the volume form for the
rescaled metric $\overline{\tau}_i^{-1} \: g(t_0 - \tau_i)$.)
However, to do so one needs uniform lower bounds on
$l(q, \tau^\prime)$ for the original solution in terms of
$d_{t_0-\tau^\prime}(q, q(\tau^\prime))$, for $\tau^\prime \in (-\infty, 0)$.
By an argument of Perelman,
written in detail in \cite{Ye1}, one does indeed have a lower bound of the
form
\begin{equation} \label{Perelmanbound}
l(q, \tau^\prime) \: \ge \: - \: l(q(\tau^\prime)) \: - \: 1 \: + \: C(n) \:
\frac{d_{t_0-\tau^\prime}(q, q(\tau^\prime))^2}{\tau^\prime}.
\end{equation}
The nonnegative curvature gives polynomial volume growth for
distance balls,
so using (\ref{Perelmanbound})
one can apply dominated convergence to the integrals
$\int_M e^{-l(q, \tau_i)} \:
\tau_i^{-n/2} \: \dvol(q,t_0 - \tau_i)$.
Thus $\lim_{i \rightarrow \infty} \tilde{V}(\tau_i) \: = \:
\tilde{V}(\tau)$.

As (\ref{upperl}) gives
a uniform upper bound on $l$ on an appropriate
ball around $q(\tau_i)$, and there is a lower
volume bound on the ball, it follows that as $i \rightarrow \infty$,
$\tilde{V}(\tau_i)$ is uniformly bounded away from zero. From this
argument and the monotonicity of $\tilde{V}$,
$\tilde{V}(\tau)$ is a positive constant $c$ as a function of
$\tau$, namely the limit of the reduced volume of the original
solution as real time goes to $- \infty$. As the original solution is
nonflat, the constant $c$ is strictly less than the limit of the
reduced volume of the original solution
as real time goes to zero, which is $(4 \pi)^{\frac{n}{2}}$.

Next, we will apply (\ref{illustrate}) and (\ref{illustrate2}).
As (\ref{illustrate}) holds distributionally
for each rescaled solution, it follows that
it holds distributionally for $\overline{l}$. In particular, the
nonpositivity implies that the left-hand side of (\ref{illustrate}), when
computed for the limit solution, is actually a nonpositive measure.
If the left-hand side of
(\ref{illustrate}) (for the limit solution)
were not strictly zero then using (\ref{pphieqn})  we would conclude
that $\frac{d\tilde{V}}{d\tau}$
is somewhere negative, which is a contradiction.
(We use the fact that (\ref{Perelmanbound}) passes to the limit
to give a similar lower bound on $\overline{l}$.)
Thus we must have equality in (\ref{illustrate}) for the limit solution.
This implies equality in (\ref{touse2}),
which implies equality in  (\ref{illustrate2}).
Writing
 (\ref{illustrate2}) as
\begin{equation}
(4 \triangle \: - \: R) \: e^{-\frac{l}{2}} \: = \:
\frac{l-n}{\tau} \: e^{-\frac{l}{2}},
\end{equation}
elliptic theory gives smoothness of $\overline{l}$.

In I.11.2 it is said that equality in  (\ref{illustrate2}) implies equality
in (\ref{bbound2}), which implies that one has a gradient shrinking soliton.
There is a problem with this argument, as the use of (\ref{bbound2})
implicitly
assumes that the solution is defined for all $\tau \ge 0$, which we do not
know. (The function $\overline{l}$ is only defined by a limiting
procedure, and not in terms of ${\mathcal L}$-geodesics on some
Ricci flow solution.) However, one can instead use Proposition \ref{localized}, with
$f = \overline{l}$.
Equality in (\ref{illustrate2}) implies that $v = 0$, so (\ref{I.(9.1)})
directly gives the gradient shrinking soliton equation.
(The problem with the argument using (\ref{bbound2}), and its resolution using
(\ref{I.(9.1)}), were pointed out by the UCSB group.)

If the gradient shrinking soliton $\overline{g}(\cdot)$ is
flat then, as it will be $\kappa$-noncollapsed
at all scales, it must be $\R^n$. From the soliton equation,
$\partial_i \partial_j \overline{l} \: = \:
\frac{\overline{g}_{ij}}{2\tau}$ and
$\triangle \overline{l} \: = \: \frac{n}{2\tau}$. Putting this
into the equality (\ref{illustrate2}) gives
$|\nabla \overline{l}|^2 \: = \: \frac{\overline{l}}{\tau}$.
It follows that the level sets of $\overline{l}$ are distance spheres.
Then (\ref{illustrate}) implies
that with an appropriate choice of
origin,
$\overline{l} \: = \: \frac{|x|^2}{4\tau}$.  The reduced volume
$\tilde{V}(\tau)$ for the limit solution is now computed to be
$(4 \pi)^{\frac{n}{2}}$, which is a contradiction.  Therefore the
gradient shrinking soliton is not flat.
\end{proof}

We remark that the gradient soliton constructed here does
not, {\it a priori}, have bounded curvature on
compact time intervals,
i.e.
it may not be a $\kappa$-solution.  In the $2$ and
$3$-dimensional cases
one can prove this using additional
reasoning.  See Section \ref{altpf11.3} where
it is shown that $2$-dimensional $\kappa$-solutions are
round spheres, and Section \ref{pf11.7} where the $3$-dimensional
case is discussed.


\section{I.11.3. Two dimensional $\kappa$-solutions} \label{I.11.3}

The next result is a classification of two-dimensional
$\kappa$-solutions.  It is important when doing dimensional
reduction.

\begin{corollary}[kappa solutions corollary]
  \label{cor:kl-two-dimensional-solutions-kappa-solutions-corollary}
  \dcref{Corollary 40.1}
  \group{Kleiner--Lott Kappa Solutions}
  \source{kleiner-lott}
 (cf. Corollary I.11.3) \label{2Dsoliton}
The only oriented two-dimensional $\kappa$-solution is
the shrinking round $2$-sphere.
\end{corollary}
\begin{proof}
  \uses{prop:kl-asymptotic-solitons-existence-ricci-solitons}
First, the
only nonflat oriented
nonnegatively curved gradient shrinking 2-D soliton is the
round $S^2$. The reference \cite{Hamsurface} given in I.11.3 for this
fact does not actually cover it, as the reference
only deals with compact solitons. A proof using
Proposition \ref{asympsoliton} to rule
out the noncompact case appears in \cite{Ye2}.

Given this, the limit solution in Proposition \ref{asympsoliton}
is a shrinking round $2$-sphere.  Thus the rescalings
$\overline{\tau}_i^{-1} g(t_0 - \overline{\tau}_i)$ converge to
a round $2$-sphere as $i \rightarrow \infty$.  However, by \cite{Hamsurface}
the Ricci flow makes an
almost-round $2$-sphere become more round. Thus any given time slice
of the original $\kappa$-solution must be a round $2$-sphere.
\end{proof}

\begin{remark}[structural remark on Ricci solitons]
  \label{rem:kl-two-dimensional-solutions-structural-remark-ricci-solitons}
  \dcref{Remark 40.2}
  \group{Kleiner--Lott Kappa Solutions}
  \source{kleiner-lott}
  \uses{cor:kl-two-dimensional-solutions-kappa-solutions-corollary}

One can employ a somewhat different
line of reasoning to prove Corollary \ref{2Dsoliton}; see Section
\ref{altpf11.3}.
\end{remark}

\section{I.11.4. Asymptotic scalar curvature and asymptotic volume ratio}
\label{11.4proof}

In this section we first show that the asymptotic scalar curvature ratio of
a $\kappa$-solution is infinite.  We then show that the asymptotic
volume ratio vanishes. The proofs are somewhat rearranged from
those in I.11.4.  They are logically independent of
Section \ref{I.11.3}, i.e. also cover the case $n=2$.
We will use results from
Appendices \ref{appharnack} and \ref{Alexandrov},
in particular (\ref{harnackinequality}).

\begin{definition}[asymptotic scalar curvature ratio definition]
  \label{def:kl-asymptotic-scalar-curvature-asymptotic-volume-ratio-asymptotic-scalar-curvature-ratio-definition}
  \dcref{Definition 41.1}
  \group{Kleiner--Lott Kappa Solutions}
  \source{kleiner-lott}

If $M$ is a complete connected Riemannian manifold then its {\em asymptotic
scalar curvature ratio} is ${\mathcal R} \: = \:
\limsup_{x \rightarrow \infty} \: R(x) \: d(x, p)^2$. It is independent
of the choice of basepoint $p$.
\end{definition}

\begin{theorem}[compactness for kappa solutions]
  \label{thm:kl-asymptotic-scalar-curvature-asymptotic-volume-ratio-compactness-kappa-solutions}
  \dcref{Theorem 41.2}
  \group{Kleiner--Lott Kappa Solutions}
  \source{kleiner-lott}
 \label{asympscalar}
Let $(M,g(\cdot))$ be a noncompact $\kappa$-solution.
Then the asymptotic scalar curvature ratio ${\mathcal R}$
is infinite for each time slice.
\end{theorem}
\begin{proof}
Suppose that $M$ is $n$-dimensional,
with $n\geq 2$.   Pick $p\in M$ and consider a time-$t_0$
slice $(M,g(t_0))$.
We deal with the cases ${\mathcal R} \in (0, \infty)$ and
${\mathcal R} = 0$ separately and show that they lead to contradictions. \\ \\
{\bf Case 1: $0<{\mathcal R}<\infty$.}  We choose a sequence
$x_k\in M$ such that $d_{t_0}(x_k,p)\ra \infty$ and
$R(x_k,t_0)d^2(x_k,p)\ra{\mathcal R}$.
Consider the rescaled pointed solution
$(M,x_k,g_k(t))$ with
$g_k(t) =R(x_k,t_0)g(t_0 + \frac{t}{R(x_k,t_0)})$ and $t\in (-\infty,0]$.
We have $R_k(x_k,0)=1$, and for all $b>0$, for sufficiently large $k$,
we have $R_k(x,t)\leq R_k(x,0)\leq \frac{2{\mathcal R}}{d_k^2(x,p)}$
for all $x$ such that $d_k(x,p)>b$.
Fix numbers $b, B > 0$ so that  $b<\sqrt{{\mathcal R}}<B$.
The $\kappa$-noncollapsing assumption gives a uniform positive
lower bound on the
injectivity radius of $g_k(0)$ at $x_k$,
and so by Appendix \ref{subsequence}
we may extract a pointed limit solution
$(M_\infty,x_\infty,g_\infty(\cdot))$, defined on a time interval
$(-\infty,0]$ from the sequence
$(M_k,x_k,g_k(\cdot))$ where $M_k =\{x\in M\mid b<d_k(x,p)<B\}$.
Note that $g_\infty$ has nonnegative curvature operator and
the time slice $(M_\infty,g_\infty(0))$ is locally
isometric to an annular
portion of a nonflat metric cone, since $(M_k,p,g_k(0))$
Gromov-Hausdorff converges to the Tits cone $\ctits(M,g(t_0))$.
(We use the word ``locally'' because the annulus in $\ctits(M,g(t_0))$
need not be geodesically convex in $\ctits(M,g(t_0))$, so we are only
saying that the distance functions in small balls match up.)
When $n=2$ this contradicts the fact that $R_\infty(x_\infty,0)=1$.
When $n\geq 3$, we will derive a contradiction from
Hamilton's curvature evolution equation
\begin{equation}
\label{curvevol}
\Rm_t=\Delta \Rm+Q(\Rm).
\end{equation}

Let $d_v : \ctits(M,g(t_0)) \rightarrow \R$ be the distance function
from the vertex and let
$\rho:M_\infty\ra \R$ be the pullback of $d_v$ under the inclusion of
the annulus $M_\infty$ in $\ctits(M,g(t_0))$.

\end{proof}
\begin{lemma}[scalar curvature lemma]
  \label{lem:kl-asymptotic-scalar-curvature-asymptotic-volume-ratio-scalar-curvature-lemma}
  \dcref{Lemma 41.4}
  \group{Kleiner--Lott Kappa Solutions}
  \source{kleiner-lott}

The metric cone structure on $(M_\infty,g_\infty(0))$  is
smooth, i.e. $\rho$ is a smooth function.
\end{lemma}
\begin{proof}
Consider a unit speed geodesic segment $\ga$ in the Tits cone
$\ctits(M,g(t_0))$,
such that $\ga$ is disjoint from the vertex $v\in \ctits(M,g(t_0))$.
Note that since $\ctits(M,g(t_0))$ is a Euclidean cone over
the Tits boundary $\tits (M,g(t_0))$, the
geodesic $\ga$ lies in the cone over a geodesic segment
$\hat\ga\subset\tits (M,g(t_0))$.
Thus $\ga$ lies in a
$2$-dimensional
locally convex flat subspace of $\ctits(M,g(t_0))$.
Also, as in a flat $2$-dimensional cone,
the second derivative of the composite function
$d_v^2\circ\ga$ is identically $2$.

Since $\rho$ is obtained from $d_v$ by composition with a locally isometric
embedding $(M_\infty,g_\infty(0))\ra \ctits(M,g(t_0))$, the composition
of $\rho^2$ with any unit speed geodesic segment in $M_\infty$
also has second derivative identically equal to $2$.


Since $\rho^2$ is Lipschitz, Rademacher's theorem implies that $\rho^2$
is differentiable almost everywhere.
Let $y\in M_\infty$ be a point of differentiability of $\rho^2$.
 If the injectivity radius of $g_\infty(0)$ at
$y$ is $>a$, then the function $h_y$ given by  the composition
$$
T_yM_\infty\supset B(0,a)\stackrel{\exp_y}{\lra}\;M_\infty
\stackrel{\rho^2}{\lra}\R
$$
has the property that its second radial derivative is identically $2$,
and it is differentiable at the origin $0\in T_yM_\infty$.
Therefore
$h_y$ is a second order polynomial with Hessian identically $2$,
and is smooth.  As the injectivity radius is a continuous function, this
implies that $\rho^2$ is smooth everywhere in $M_\infty$.

Since $\rho^2$ is strictly positive on $M_\infty$,
it follows that $\rho=\sqrt{\rho^2}$
is smooth as well.
\end{proof}
\begin{proof}
  \proves{thm:kl-asymptotic-scalar-curvature-asymptotic-volume-ratio-compactness-kappa-solutions}

By the lemma, we may choose a smooth local orthonormal
frame $e_1,\ldots,e_n$ for $(M_\infty,g_\infty(0))$
near $x_\infty$ such that $e_1$ points radially outward (with respect to the
cone structure), and $e_2,e_3$ span a $2$-plane at $x_\infty$ with strictly
positive curvature;
such a $2$-plane exists because $R_\infty(x_\infty, 1) = 1$.
Put $P = e_1 \wedge e_2$.
In terms of the curvature operator, the fact that
$\Rm_\infty(e_1,e_2,e_2,e_1)=0$ is equivalent to
$\langle P, \Rm_\infty P \rangle \: = \: 0$.
As the curvature operator is nonnegative, it follows that
$\Rm_\infty P \: = \: 0$. (In fact, this is true for any metric cone.)
Differentiating gives
\begin{equation} \label{diff1}
\left( \nabla_{e_i} \Rm_\infty \right) P \: + \: \Rm_\infty (\nabla_{e_i} P) \: = \: 0
\end{equation}
and
\begin{equation} \label{diff2}
\left( \triangle \Rm_\infty \right) P \: + \: 2 \sum_i \left( \nabla_{e_i} \Rm_\infty \right)
\nabla_{e_i} P \: + \: \Rm_\infty \left( \triangle P \right) \: = \: 0.
\end{equation}
Taking the inner product of (\ref{diff2}) with $P$ gives
\begin{align}
0 \: & = \: \langle P, \left( \triangle \Rm_\infty \right) P \rangle  \: + \: 2 \sum_i
\langle P, \left( \nabla_{e_i} \Rm_\infty \right)
\nabla_{e_i} P \rangle \\
& = \: \langle P, \left( \triangle \Rm_\infty \right) P \rangle  \: + \: 2 \sum_i
\langle \nabla_{e_i} P, \left( \nabla_{e_i} \Rm_\infty \right)
P \rangle. \notag
\end{align}
Then  (\ref{diff1}) gives
\begin{equation}
\langle P, (\triangle \Rm_\infty) P \rangle \: = \:  2 \:
\sum_i \langle \nabla_{e_i} P,  \Rm_\infty \left( \nabla_{e_i} P \right)  \rangle.
\end{equation}

As the sphere of distance $r$ from the vertex in a metric cone has
principal curvatures $\frac{1}{r}$, we have
$\nabla_{e_3} e_1 \: = \: - \: \frac{1}{r} \: e_3$.
Then
\begin{equation}
\nabla_{e_3} (e_1 \wedge e_2) \: = \:
\left( \nabla_{e_3} e_1 \right) \wedge e_2 \: + \:
e_1 \wedge \nabla_{e_3} e_2 \: = \: \frac{1}{r} \left( e_2 \wedge e_3 \right) \: + \:
e_1 \wedge \nabla_{e_3} e_2.
\end{equation}
This shows that $\nabla_{e_3} P$ has a nonradial component
$\frac{1}{r} \: e_2 \wedge e_3$. Thus
$(\Delta \Rm_\infty)(e_1,e_2,e_2,e_1)>0$.
The zeroth order quadratic term $Q(\Rm)$ appearing in (\ref{curvevol})
is nonnegative when $\Rm$ is nonnegative, so we conclude that
$\partial_t\Rm_\infty(e_1,e_2,e_2,e_1)>0$ at $t=0$.  This means that
$\Rm_\infty(-\eps)(e_1,e_2,e_2,e_1)<0$ for $\eps>0$ sufficiently small,
which is impossible. \\ \\
{\bf Case 2: ${\mathcal R}=0$.}
Let us take any sequence $x_k\in M$ with
$d_{t_0}(x_k,p)\ra \infty$.  Set $r_k = d_{t_0}(x_k,p)$, put
\begin{equation}
g_k(t) =  r_k^{-2}g(t_0 + r_k^2 t)
\end{equation}
for $t\in (-\infty,0]$, and let $d_k(\cdot, \cdot)$ be the distance function
associated to $g_k(0)$.
For any $0<b<B$, put
\begin{equation}
M_k(b,B) \: = \: \{x\in M\mid 0<b<d_k(x,p)<B\}.
\end{equation}
Since ${\mathcal R}=0$, we get that $\sup_{x\in M_k(b,B)}|\Rm_k(x,0)|\ra 0$
as $k\ra\infty$.  Invoking the $\kappa$-noncollapsed assumption as in
the previous case, we
may assume that $(M,p,g_k(0))$ Gromov-Hausdorff converges to a metric
cone $(M_\infty,p_\infty,g_\infty)$ (the Tits cone $\ctits(M,g(t_0))$)
which is flat and smooth away from the
vertex $p_\infty$, and the convergence is smooth away from $p_\infty$.

The ``unit sphere'' in $\ctits(M,g(t_0))$
defines a compact smooth hypersurface $S_\infty$
in $(M_\infty - \{p_\infty\},g_\infty(0))$ whose principal curvatures
are identically $1$. If $n \ge 3$ then $S_\infty$ must be a quotient of
the standard $(n-1)$-sphere by the free action of a finite group of
isometries. We have a sequence $S_k\subset M_k$
of approximating smooth hypersurfaces whose principal curvatures
(with respect to $g_k(0)$) go to $1$ as $k\ra\infty$.  In view of the
convergence to $(M_\infty,p_\infty,g_\infty)$, for
sufficiently large $k$,
the inward principal curvatures of $S_k$ with respect to $g_k(0)$
are close to $1$. As $M$ has nonnegative curvature,
$S_k$ is diffeomorphic to a sphere \cite[Theorem A]{Eschenburg}.
Thus $S_\infty$ is isometric to the standard $(n-1)$-sphere,
and so
$\ctits(M,g(t_0))$ is isometric to $n$-dimensional Euclidean space.
Then  $(M,g(t_0))$ is isometric to $\R^n$, which
contradicts the definition of a $\kappa$-solution.

In the case $n=2$ we know that $S_\infty$ is diffeomorphic to a circle but
we do not know {\it a priori} that it has length $2\pi$.
To handle the case $n = 2$, we use the fact that $g_k(t)$ is
a Ricci flow solution, to
extract a limiting smooth incomplete
time-independent Ricci flow solution
$(M_\infty \backslash p_\infty, g_\infty(t))$ for $t \in [-1, 0]$.
Note that this solution is unpointed.
In view of the
convergence to the limiting solution,
for sufficiently large $k$,
the inward principal curvatures of $S_k$ with respect to $g_k(t)$
are close to $1$ for all $t\in [-1,0]$.
This implies that
$S_k$ bounds a domain $B_k\subset M$ whose diameter with
respect to $g_k(t)$ is uniformly bounded above, say by $10$
(see Appendix \ref{Alexandrov}).

Applying the Harnack
inequality (\ref{harnackinequality})
with $y_k\in S_k$ (at time $0$) and $x\in B_k$
(at time $-1$), we see that
$\sup_{x\in B_k}|\Rm_k(x,-1)|\ra 0$ as $k\ra \infty$.
Thus $(B_k,p,g_k(-1))$ Gromov-Hausdorff converges to a flat manifold
$(B_\infty,\bar p_\infty,g_\infty(-1))$
with convex boundary. As all of the principal curvatures of
$\partial B_\infty$ are $1$,
$B_\infty$ must be isometric to a Euclidean unit ball.
This implies that $S_\infty$ is isometric to the standard $S^1$ of
length $2\pi$,
and we obtain a contradiction as before.
\end{proof}

\begin{definition}[asymptotic volume ratio definition]
  \label{def:kl-asymptotic-scalar-curvature-asymptotic-volume-ratio-asymptotic-volume-ratio-definition}
  \dcref{Definition 41.12}
  \group{Kleiner--Lott Kappa Solutions}
  \source{kleiner-lott}

If $M$ is a complete $n$-dimensional
Riemannian manifold with nonnegative Ricci curvature
then its {\em asymptotic
volume ratio} is ${\mathcal V} \: = \:
\lim_{r \rightarrow \infty} \: r^{-n} \: \vol(B(p, r))$. It is independent
of the choice of basepoint $p$.
\end{definition}

\begin{proposition}[existence for kappa solutions]
  \label{prop:kl-asymptotic-scalar-curvature-asymptotic-volume-ratio-existence-kappa-solutions}
  \dcref{Proposition 41.13}
  \group{Kleiner--Lott Kappa Solutions}
  \source{kleiner-lott}
 (cf. Proposition I.11.4) \label{propI.11.4}
Let $(M,g(\cdot))$ be a noncompact $\kappa$-solution.
Then the asymptotic volume ratio ${\mathcal V}$
vanishes for each time slice $(M,g(t_0))$.  Moreover, there is a
sequence of points $x_k\in M$ going to infinity such that the
pointed sequence $\{ (M,(x_k,t_0),g(\cdot)) \}_{k=1}^\infty$ converges,
modulo rescaling
by $R(x_k,t_0)$, to a $\kappa$-solution which isometrically splits off
an $\R$-factor.
\end{proposition}
\begin{proof}
Consider the time-$t_0$ slice.
Suppose that ${\mathcal V}>0$.
As ${\mathcal R}=\infty$,  there
are sequences $x_k\in M$ and  $s_k>0$
such that $d_{t_0}(x_k,p)\ra \infty$, $\frac{s_k}{d_{t_0}(x_k,p)}\ra 0$,
$R(x_k, t_0 )s_k^2\ra\infty$, and $R(x, t_0)\leq 2R(x_k, t_0)$ for all
$x\in B_{t_0}(x_k,s_k)$
\cite[Lemma 22.2]{Hamilton}.
Consider the rescaled pointed solution
$(M,x_k,g_k(t))$ with
$g_k(t) =R(x_k,t_0) \: g(t_0 + \frac{t}{R(x_k,t_0)})$ and $t\in (-\infty,0]$.
As $R_t \ge 0$, we have
$R_k(x,t)\leq 2$ whenever $t\leq 0$ and
$d_k(x,x_k)\leq R(x_k, t_0)^{1/2} s_k$,
where $d_k$ is the distance function  for $g_k(0)$
and $R_k(\cdot,\cdot)$ is the scalar curvature of $g_k(\cdot)$.
The $\kappa$-noncollapsing assumption gives a uniform positive
lower bound on the
injectivity radius of $g_k(0)$ at $x_k$,
so by Appendix \ref{subsequence}
we may extract a complete pointed limit solution
$(M_\infty,x_\infty,g_\infty(t))$, $t\in (-\infty,0]$, of a subsequence of
the sequence of pointed Ricci flows.
By relative volume comparison, $(M_\infty,x_\infty,g_\infty(0))$
has positive asymptotic volume ratio.  By Appendix \ref{Alexandrov},
the Riemannian manifold
$(M_\infty,x_\infty,g_\infty(0))$ is isometric to an Alexandrov
space which splits off a line, which means that it is a Riemannian
product $\R\times N$. This implies a product structure for
earlier times;
see Appendix \ref{maxprin}.
Now when $n=2$, we have a contradiction,
since $R(x_\infty,0)=1$ but $(M_\infty,g_\infty(0))$ is a product
surface, and must therefore be flat.  When $n>2$ we obtain
a $\kappa$-solution on an $(n-1)$-manifold with positive asymptotic
volume ratio at time zero, and by induction this is impossible.
\end{proof}




\section{In a $\kappa$-solution, the curvature and the normalized
volume control each other}

In this section we show that, roughly speaking, in a
$\kappa$-solution
the curvature and the normalized volume control each other.


\begin{corollary}[existence for kappa solutions]
  \label{cor:kl-solution-curvature-normalized-volume-control-each-other-existence-kappa-solutions}
  \dcref{Corollary 42.1}
  \group{Kleiner--Lott Kappa Solutions}
  \source{kleiner-lott}
 \label{11.4cors}
1.  If $B(x_0,r_0)$ is a ball in a time slice of a $\kappa$-solution,
then the normalized volume $r_0^{-n}\vol(B(x_0,r_0))$ is controlled
(i.e. bounded away from zero) $\Longleftrightarrow$
the normalized scalar curvature $r_0^2R(x_0)$ is controlled
(i.e. bounded above).

2. If $B(x_0,r_0)$ is a ball in a time slice of a $\kappa$-solution,
then the normalized volume $r_0^{-n}\vol(B(x_0,r_0))$ is almost
maximal $\Longleftrightarrow$
the normalized scalar curvature $r_0^2R(x_0)$ is almost zero.

3. (Precompactness)
If $(M_k,(x_k,t_k),g_k(\cdot))$ is a sequence of pointed $\kappa$-solutions
(without the assumption that $R(x_k, t_k) = 1$)
and for some $r>0$, the $r$-balls $B(x_k,r)\subset (M_k,g_k(t_k))$
have controlled normalized volume, then a subsequence converges
to an ancient solution $(M_\infty,(x_\infty,0),g_\infty(\cdot))$
which has nonnegative curvature operator, and is $\kappa$-noncollapsed
(though {\it a priori} the curvature may be unbounded on a given time slice).

4.
There is a constant
$\eta=\eta(n,\kappa)$ such that for every $n$-dimensional
$\kappa$-solution $(M,g(\cdot))$, and all $x\in M$, we have
$|\nabla R|(x,t)\leq \eta R^{\frac{3}{2}}(x,t)$ and
$|R_t|(x,t) \leq \eta R^2(x,t)$.
More generally, there are scale invariant bounds on
all derivatives of the curvature tensor, that only depend
on
$n$ and
$\kappa$. That is, for
each $\rho,k,l<\infty$ there is a constant
$C=C(n,\rho,k,l,\kappa) < \infty$
such that
$\left|\frac{\D^k}{\D t^k}\;
\nabla^l\Rm\right|(y,t) \: \le \: C\;R(x,t)^{(k+\frac{l}{2}+1)}$ for any
$y \in B_t(x,\rho R(x,t)^{-\frac12})$.


5.  There is a function $\al:[0,\infty)\ra[0,\infty)$ depending only
on $\kappa$ such that $\lim_{s\ra\infty}\al(s)=\infty$,
and  for every  $\kappa$-solution $(M,g(\cdot))$ and
$x,y\in M$, we have $R(y)d^2(x,y)\geq \al(R(x)d^2(x,y))$.
\end{corollary}
\begin{proof}
  \uses{cor:kl-curvature-bounds-ricci-flow-solutions-nonnegative-curvature-operator-ass-existence-curvature-operators,cor:kl-length-distortion-estimates-estimates-length-distortion-estimates,cor:kl-necklike-behavior-at-infinity-three-dimensional-solution-strong-version-existence-kappa-solutions,cor:kl-necklike-behavior-at-infinity-three-dimensional-solution-weak-version-existence-kappa-solutions,cor:kl-two-dimensional-solutions-kappa-solutions-corollary,cor:kl-volume-bound-existence-curvature-operators,def:kl-notation-terminology-smooth-closeness-epsilon-necks,lem:kl-alternate-proof-corollary-2dsoliton-using-proposition-propi-11-4-corolla-existence-kappa-solutions,lem:kl-getting-parabolic-regions-getting-parabolic-regions-lemma,lem:kl-length-distortion-estimates-distance-distortion-lemma,lem:kl-necklike-behavior-at-infinity-three-dimensional-solution-strong-version-existence-kappa-solutions,lem:kl-point-selection-argument-point-selection-argument-lemma,prop:kl-asymptotic-scalar-curvature-asymptotic-volume-ratio-existence-kappa-solutions,thm:kl-compactness-space-three-dimensional-solutions-compactness-kappa-solutions,thm:kl-statement-pseudolocality-theorem-pseudolocality-theorem}
Assertion 1, $\Longrightarrow$.
Suppose we have a sequence of $\kappa$-solutions
$(M_k,g_k(\cdot))$, and sequences
$t_k\in (-\infty, 0]$, $x_k\in M_k$, $r_k>0$, such that at time $t_k$,
the normalized volume of $B(x_k,r_k)$ is
$\geq c>0$, and $R(x_k,t_k)r_k^2\ra \infty$.
By Appendix \ref{pointpicking},
for each $k$, we can find $y_k\in B(x_k,5r_k)$,
$\bar r_k\leq r_k$, such that $R(y_k,t_k)\bar r_k^2\geq R(x_k,t_k)r_k^2$,
and $R(z,t_k)\leq 2R(y_k,t_k)$ for all $z\in B(y_k,\bar r_k)$.
Note that by relative volume comparison,
whenever $\widetilde{r}_k \le \overline{r}_k$ we have
\begin{equation}
\label{volbig}
\frac{\vol(B(y_k,{\widetilde r}_k))}{\widetilde{r}_k^n}\geq
\frac{\vol(B(y_k,\bar r_k))}{\bar r_k^n}\geq
\frac{\vol(B(y_k,10r_k))}{(10r_k)^n}
\geq \frac{\vol(B(x_k,r_k))}{(10r_k)^n}\geq\frac{c}{10^n}.
\end{equation}
Rescaling the sequence of pointed solutions $(M_k,(y_k,t_k),g_k(\cdot))$
by $R(y_k,t_k)$, we get a sequence satisfying the hypotheses
of Appendix \ref{subsequence} (we use here the fact that $R_t \ge 0$
for an ancient solution), so it accumulates on a limit
flow $(M_\infty,(y_\infty,0),g_\infty(\cdot))$ which is
a $\kappa$-solution.
By (\ref{volbig}), the asymptotic volume ratio of $(M,g_\infty(0))$ is
$\geq \frac{c}{10^n}>0$.  This contradicts Proposition \ref{propI.11.4}.


\bigskip
Assertion 3.  By relative volume comparison, it follows that  every
$r$-ball in $(M_k,g_k(t_k))$ has normalized volume bounded below
by a ($k$-independent) function of its distance to $x_k$.
By 1, this implies that  the curvature of $(M_k,g_k(t_k))$ is bounded
by a $k$-independent function of the distance to $x_k$,
and hence we can apply Appendix \ref{subsequence} to extract
a smoothly converging subsequence.

\bigskip
Assertion 1, $\Longleftarrow$.  Suppose we have a sequence
$(M_k,g_k(\cdot))$ of $\kappa$-solutions, and sequences
$x_k\in M_k$, $r_k>0$, such that $R(x_k,t_k)r_k^2<c$ for all
$k$, but $r_k^{-n}\vol(B(x_k,r_k))\ra 0$.
For large $k$, we can choose $\bar r_k \in (0, r_k)$ such that
$\bar r_k^{-n}\vol(B(x_k,\bar r_k))=\frac{1}{2}c_n$
where $c_n$ is the volume of the unit Euclidean $n$-ball.
By relative volume comparison, $\frac{\bar r_k}{r_k}\ra 0$.
Applying 3, we see that the pointed sequence $(M_k,(x_k,t_k),g_k(\cdot))$,
rescaled by the factor $\bar r_k^{-2}$, accumulates
on a pointed ancient solution $(M_\infty,(x_\infty,0),g_\infty(\cdot))$,
such that the ball $B(x_\infty,1)\subset (M_\infty,g_\infty)$ has
normalized volume $\frac{1}{2}c_n$ at $t=0$.

Suppose the ball $B(x_\infty,1)\subset (M_\infty,g_\infty(0))$ were flat.
Then by the Harnack inequality (\ref{harnackinequality}) (applied to the approximators)
we would have $R_\infty(x,t)=0$ for all $x\in M_\infty$, $t\leq 0$,
i.e. $(M_\infty,g_\infty(t))$ would be a time-independent flat manifold.
It cannot be $\R^n$ since $\vol(B(x_\infty,1)) \: = \: \frac12 \: c_n$.
But flat manifolds other than Euclidean space have zero asymptotic volume ratio
(as follows from the Bieberbach theorem that if $N=\R^n/\Gamma$ is
a flat manifold and $\Ga$ is nontrivial then there is a $\Gamma$-invariant
affine subspace $A\subset \R^n$ of dimension at least $1$
on which $\Gamma$ acts cocompactly).
This contradicts the assumption that the sequence $(M_k,g_k(\cdot))$
is $\kappa$-noncollapsed.  Thus
$B(x_\infty,1)\subset (M_\infty,g_\infty(0))$ is not flat, which means,
by the Harnack inequality, that the
scalar curvature of $g_\infty(0)$ is strictly positive
everywhere.  Therefore, with respect to $g_k$, we have
\begin{equation}
\liminf_{k\ra \infty}R(x_k,t_k)r_k^2=
\liminf_{k\ra\infty}(R(x_k,t_k))\bar
r_k^2)\left(\frac{r_k}{\bar r_k}\right)^2
\geq \const
\liminf_{k\ra\infty}\left(\frac{r_k}{\bar r_k}\right)^2=\infty,
\end{equation}
which is a contradiction.

\bigskip
Assertion 2, $\Longrightarrow$.  Apply 1,
the precompactness criterion, and the fact
that a nonnegatively-curved manifold whose balls
have normalized volume $c_n$ must be flat.

\bigskip
Assertion 2, $\Longleftarrow$. Apply 1, the precompactness criterion,
and the Harnack inequality (\ref{harnackinequality}) (to the approximators).

\bigskip
Assertion 4.  This follows by rescaling $g$ so that $R(x,t)=1$, and
applying 1 and 3.

\bigskip
Assertion 5.  The quantity $R(z)d^2(u,v)$ is scale invariant.
If the assertion failed then  we would have sequences $(M_k,g_k(\cdot))$,
$x_k,y_k\in M_k$, such that $R(y_k)=1$ and
$d(x_k,y_k)$ remains bounded, but the curvature at $x_k$ blows up.
This contradicts $1$ and $3$.

\end{proof}

\section{An alternate proof of Corollary \ref{2Dsoliton}
using Proposition \ref{propI.11.4} and Corollary \ref{11.4cors}}
\label{altpf11.3}

In this section we give an alternate proof of Corollary \ref{2Dsoliton}.
It uses Proposition \ref{propI.11.4} and Corollary \ref{11.4cors}
To clarify the chain of logical dependence, we remark that
this section is concerned with $2$-dimensional $\kappa$-solutions,
and does not use anything from Sections \ref{I.11.2} or \ref{I.11.3}.   It does use
Proposition \ref{propI.11.4}.  However, we avoid circularity here because the proof
of Proposition \ref{propI.11.4} given in Section \ref{11.4proof}, unlike the
proof in \cite{Perelman}, does not use Corollary \ref{2Dsoliton}.
\begin{lemma}[existence for kappa solutions]
  \label{lem:kl-alternate-proof-corollary-2dsoliton-using-proposition-propi-11-4-corolla-existence-kappa-solutions}
  \dcref{Lemma 43.1}
  \group{Kleiner--Lott Kappa Solutions}
  \source{kleiner-lott}

\label{2dvol}
There is a constant $v=v(\kappa)>0$ such that if
$(M,g(\cdot))$ is a 2-dimensional $\kappa$-solution
({\it a priori} either compact or noncompact), $x,y\in M$ and
$r = d(x,y)$ then
\begin{equation}
\label{volnotsmall}
\vol(B_t(x,r))\geq vr^2.
\end{equation}
\end{lemma}
\begin{proof}
  \uses{cor:kl-solution-curvature-normalized-volume-control-each-other-existence-kappa-solutions}
If the lemma were not true then there would be a sequence $(M_k,g_k(\cdot))$
of 2-dimensional $\kappa$-solutions, and sequences $x_k,y_k\in M_k$, $t_k\in\R$
such that $r_k^{-2}\vol(B_{t_k}(x_k,r_k))\ra 0$, where
$r_k = d(x_k,y_k)$.  Let $z_k$ be the midpoint of a shortest segment
from $x_k$ to $y_k$ in the $t_k$-time slice $(M_k,g_k(t_k))$. For large $k$,
choose
$\bar r_k \in (0, r_k/2)$ such that
\begin{equation}
\label{halfvol}
\bar r_k^{-2}\vol(B_{t_k}(z_k,\bar r_k))=\frac{\pi}{2},
\end{equation}
 i.e.
half the area of the  unit disk in $\R^2$.
As
\begin{equation}
\frac{\pi}{2} = \bar r_k^{-2}\vol(B_{t_k}(z_k,\bar r_k)) \le
\bar r_k^{-2}\vol(B_{t_k}(x_k, r_k)) =
(\bar r_k / r_k)^{-2} \:  r_k^{-2} \vol(B_{t_k}(x_k, r_k)),
\end{equation}
it follows that
$\lim_{k\ra\infty}\frac{\bar r_k}{ r_k}=0$.
Then by part 3 of Corollary \ref{11.4cors},
the sequence of pointed Ricci flows $(M_k,(z_k,t_k),g_k(\cdot))$,
when rescaled by $\bar r_k^{-2}$, accumulates on a complete Ricci flow
$(M_\infty,(z_\infty,0),g_\infty(\cdot))$.  The segments from $z_k$
to $x_k$ and $y_k$ accumulate on a line in $(M_\infty,g_\infty(0))$,
and hence $(M_\infty,g_\infty(0))$ splits off a line.  By (\ref{halfvol}),
$(M_\infty,g_\infty(0))$ cannot be isometric to $\R^2$, and hence
must be a cylinder.  Considering the approximating
Ricci flows, we get a contradiction to the $\kappa$-noncollapsing assumption.
\end{proof}
\begin{proof}
  \proves{cor:kl-two-dimensional-solutions-kappa-solutions-corollary}

\bigskip
Lemma \ref{2dvol} implies that the asymptotic volume ratio of any
noncompact 2-dimensional $\kappa$-solution is at least $v>0$.
By Proposition \ref{propI.11.4}
we therefore conclude that every 2-dimensional $\kappa$-solution
is compact.  (This was implicitly assumed in the
proof of Corollary I.11.3 in \cite{Perelman}, as its reference
\cite{Hamsurface} is about compact surfaces.)

Consider the family $\f$
of 2-dimensional $\kappa$-solutions $(M,(x,0),g(\cdot))$
with $\diam(M,g(0))=1$.  By Lemma \ref{2dvol}, there is uniform
lower bound on the volume of the $t=0$ time slices of $\kappa$-solutions in
$\f$. Thus $\f$  is compact in the smooth topology by part 3 of
Corollary \ref{11.4cors}
(the precompactness leads to compactness in view of the diameter bound).
This implies
(recall that $R > 0$) that there is a constant $K\geq 1$ such that
every time slice of every 2-dimensional $\kappa$-solution has
$K$-pinched curvature.

Hamilton has shown that volume-normalized Ricci
flow on compact surfaces with positively pinched initial data
converges exponentially fast to a constant curvature metric
\cite{Hamsurface}.
His argument shows that there is a small $\eps>0$, depending
continuously on the initial data, so that when the volume of the
(unnormalized) solution
has  gone down by a factor of
at least $\eps^{-1}$,
the pinching is at most the square
root of the initial pinching.  By the compactness of the family $\f$,
this $\eps$ can be chosen
uniformly when we take the initial data to be the $t=0$ time slice
of a  $\kappa$-solution in $\f$.

Now let $K_0$ be the worst pinching of
a 2-dimensional $\kappa$-solution, and let
 $(M,g(\cdot))$ be a $\kappa$-solution where the curvature pinching of
$(M,g(0))$ is $K_0$.  Choosing $t<0$ such that
$\eps\vol(M,g(t))=\vol(M,g(0))$,
the previous paragraph implies the curvature pinching of $(M,g(t))$ is at
least $K_0^2$. This would
 contradict the fact that $K_0$ is the upper bound on the
pinching for all $\kappa$-solutions, unless $K_0=1$.



\section{I.11.5. A volume bound}

In this section we give a consequence of Proposition \ref{propI.11.4}
concerning the volumes of metric balls in Ricci flow solutions
with nonnegative curvature operator.

\end{proof}
\begin{corollary}[existence for curvature operators]
  \label{cor:kl-volume-bound-existence-curvature-operators}
  \dcref{Corollary 44.1}
  \group{Kleiner--Lott Kappa Solutions}
  \source{kleiner-lott}
 (cf. Corollary I.11.5) \label{corI.11.5}
For every $\epsilon > 0$, there is an $A < \infty$ with the following
property.  Suppose that we have a sequence of (not necessarily
complete) Ricci flow solutions $g_k(\cdot)$ with nonnegative
curvature operator, defined on $M_k \times [t_k, 0]$, such that \\
1.  For each $k$, the time-zero ball $B(x_k, r_k)$ has compact closure
in $M_k$. \\
2. For all $(x, t) \in B(x_k, r_k) \times
[t_k, 0]$, $\frac12 R(x,t) \le R(x_k, 0) = Q_k$. \\
3. $\lim_{k \rightarrow \infty} t_k Q_k = - \infty$. \\
4. $\lim_{k \rightarrow \infty} r_k^2 Q_k = \infty$. \\
Then for large $k$, $\vol(B(x_k, A Q_k^{- \: \frac12})) \le
\epsilon (A Q_k^{- \: \frac12})^n$ at time zero.
\end{corollary}
\begin{proof}
  \uses{cor:kl-length-distortion-estimates-estimates-length-distortion-estimates,prop:kl-asymptotic-scalar-curvature-asymptotic-volume-ratio-existence-kappa-solutions}
Given $\epsilon > 0$, suppose that the corollary is not true.
Then there is a sequence of such Ricci flow solutions with
$\vol(B(x_k, A_k Q_k^{- \: \frac12})) \: > \: \epsilon
(A_k Q_k^{- \: \frac12})^n$ at time zero, where
 $A_k \rightarrow \infty$.
 By Bishop-Gromov, $\vol(B(x_k, Q_k^{- \: \frac12})) \: > \: \epsilon
Q_k^{- \: \frac{n}{2}}$ at time zero, so we can parabolically rescale by
$Q_k$ and take a convergent subsequence. The limit $(M_\infty,
g_\infty(\cdot))$ will be a nonflat complete ancient solution with
nonnegative curvature operator, bounded curvature and
 ${\mathcal V}(0) \: > \: 0$.
 By Proposition \ref{propI.11.4}, it cannot be $\kappa$-noncollapsed
 for any $\kappa$.
Thus for each $\kappa \: > \: 0$,
there are a point $(x_\kappa, t_\kappa) \in M_\infty \times (-\infty,0]$ and a radius
$r_\kappa$ so that $|\Rm(x_\kappa, t_\kappa)| \: \le \:
r_\kappa^{-2}$ on the time-$t_\kappa$ ball $B(x_\kappa, r_\kappa)$,
but $\vol(B(x_\kappa, r_\kappa)) \: < \: \kappa \:
r_\kappa^n$. From the Bishop-Gromov inequality,
${\mathcal V}(t_\kappa) \: < \: \kappa$ for the limit solution.

We claim that
${\mathcal V}(t)$ is nonincreasing in $t$. To see this, we have
$\frac{\dvol(U)}{dt} \: = \: \int_U R \: dV  \ge 0$ for any domain $U
\subset M_\infty$.
Also, as $R \le 2$ on $M_\infty \times (-\infty, 0]$,
Corollary \ref{bound1}
gives that distances on $M_\infty$
decrease at most linearly in $t$,
which implies the claim.

Thus ${\mathcal V}(0) = 0$ for the limit solution, which is a contradiction.
\end{proof}

\section{I.11.6. Curvature bounds for Ricci flow solutions with
nonnegative curvature operator, assuming a lower volume bound}

In this section we show that for a Ricci flow solution with nonnegative
curvature operator, a lower bound on the volume of a ball
implies an earlier upper curvature bound on a slightly smaller ball.
This will be used in Section \ref{I.12.3}.

\begin{corollary}[existence for curvature operators]
  \label{cor:kl-curvature-bounds-ricci-flow-solutions-nonnegative-curvature-operator-ass-existence-curvature-operators}
  \dcref{Corollary 45.1}
  \group{Kleiner--Lott Kappa Solutions}
  \source{kleiner-lott}
 (cf. Corollary I.11.6) \label{corI.11.6}
For every $w > 0$, there are $B = B(w) < \infty$, $C = C(w) < \infty$ and
$\tau_0 = \tau_0(w) > 0$ with the following properties. \\
(a) Take  $t_0 \in  [- r_0^2, 0)$.
Suppose that we have a (not necessarily complete) Ricci flow
solution $(M, g(\cdot))$, defined for $t \in [t_0, 0]$, so that at time zero
the metric ball $B(x_0, r_0)$ has compact closure. Suppose that
for each $t \in [t_0, 0]$,
$g(t)$ has nonnegative curvature operator and $\vol(B_t(x_0, r_0))
\ge w r_0^n$. Then
\begin{equation} \label{Rbound3}
R(x,t) \le C r_0^{-2} + B (t - t_0)^{-1}
\end{equation}
whenever
$\dist_t(x, x_0) \le \frac14 r_0$. \\
(b)
Suppose that we have a (not necessarily complete) Ricci flow
solution $(M, g(\cdot))$, defined for $t \in [- \tau_0 r_0^2, 0]$,
so that at time zero
the metric ball $B(x_0, r_0)$ has compact closure. Suppose that
for each $t \in [- \tau_0 r_0^2, 0]$,
$g(t)$ has nonnegative curvature operator.
If we assume a time-zero volume bound $\vol(B_0(x_0, r_0))
\ge w r_0^n$ then
\begin{equation}
R(x,t) \le C r_0^{-2} + B (t + \tau_0 r_0^2)^{-1}
\end{equation}
whenever $t \in [- \tau_0 r_0^2,0]$ and $\dist_t(x, x_0) \le \frac14 r_0$.
\end{corollary}
\begin{remark}[structural remark on curvature operators]
  \label{rem:kl-curvature-bounds-ricci-flow-solutions-nonnegative-curvature-operator-ass-structural-remark-curvature-operators}
  \dcref{Remark 45.4}
  \group{Kleiner--Lott Kappa Solutions}
  \source{kleiner-lott}
  \uses{cor:kl-curvature-bounds-ricci-flow-solutions-nonnegative-curvature-operator-ass-existence-curvature-operators}

The statement in \cite[Corollary 11.6(a)]{Perelman}
does not have any constraint on $t_0$.
In our proof we seem to need that $-t_0 \le c r_0^2$ for some arbitrary
but fixed
constant $c < \infty$.
(The statement $R(\overline{x}, \overline{t}) >
C + B (\overline{t} - t_0)^{-1}$ in
\cite[Proof of Corollary 11.6(a)]{Perelman} is the issue.)
For simplicity we take
$- t_0 \le r_0^2$. This point does not affect the proof of
Corollary \ref{corI.11.6}(b), which is what ends up getting used.
\end{remark}
\begin{proof}
  \proves{cor:kl-curvature-bounds-ricci-flow-solutions-nonnegative-curvature-operator-ass-existence-curvature-operators}
  \uses{cor:kl-volume-bound-existence-curvature-operators,lem:kl-getting-parabolic-regions-getting-parabolic-regions-lemma,lem:kl-length-distortion-estimates-distance-distortion-lemma,lem:kl-point-selection-argument-point-selection-argument-lemma,thm:kl-statement-pseudolocality-theorem-pseudolocality-theorem}
For part (a), we can assume that $r_0 \: = \: 1$. Given $B,C > 0$,
suppose that $g(\cdot)$ is a Ricci flow solution for
$t \in [t_0, 0]$ that satisfies the hypotheses of the corollary,
with $R(x, t) \: > \: C \: + \: B(t - t_0)^{-1}$
for some $(x, t)$ satisfying $\dist_t(x, x_0) \: \le \: \frac{1}{4}$.
Following the notation of the proof of Theorem \ref{thmI.10.1},
except changing the $A$ of Theorem \ref{thmI.10.1} to
$\widehat{A}$,
put $\widehat{A} \: = \: \lambda C^{\frac12}$ and $\alpha \: = \:
\min(\lambda^2 C^{\frac12}, B)$, where we will take $\lambda$ to be a
sufficiently small number that only depends on $n$. Put
\begin{equation}
M_\alpha \: = \: \{(x^\prime, t^\prime) \: : \: R(x^\prime, t^\prime) \: \ge \:
\alpha (t^\prime-t_0)^{-1}\}.
\end{equation}
Clearly $(x, t) \in M_\alpha$.

We first go through the analog of the proof of Lemma \ref{Claim 1}.
We claim that there is some $(\overline{x}, \overline{t}) \in
M_\alpha$, with $\overline{t} \in (t_0, 0]$ and
$\dist_{\overline{t}}(\overline{x}, x_0) \: \le \: \frac13$, such that
$R(x^\prime, t^\prime) \: \le \: 2 Q \: =
\: 2R(\overline{x}, \overline{t})$ whenever
$(x^\prime, t^\prime) \in M_\alpha$, $t^\prime \in (t_0, \overline{t}]$
and $\dist_{t^\prime}(x^\prime, x_0) \le
\dist_{\overline{t}}(\overline{x}, x_0) \: + \: \widehat{A}Q^{- \: \frac12}$.
Put $(x_1, t_1) = (x, t)$. Inductively, if we cannot take
$(x_k, t_k)$ for $(\overline{x}, \overline{t})$ then there is some
$(x_{k+1}, t_{k+1}) \in M_\alpha$ with
$t_{k+1} \in (t_0, t_k]$,
$R(x_{k+1}, t_{k+1}) > 2
R(x_k, t_k)$ and $\dist_{t_{k+1}}(x_{k+1}, x_0) \le
\dist_{t_k}(x_k, x_0) \: + \: \widehat{A}R(x_k, t_k)^{- \: \frac12}$. As the
process must terminate, we end up with $(\overline{x}, \overline{t})$
satisfying
\begin{equation}
\dist_{\overline{t}}(\overline{x}, x_0) \: \le \: \frac14
\: + \: \frac{1}{1-\sqrt{1/2}} \: \widehat{A} \: R(x,t)^{- \: \frac12}
\: \le \: \frac13
\end{equation}
if $\lambda$ is sufficiently small.

Next, we go through the analog of the proof of Lemma \ref{Claim 2}.
As in the proof of Lemma \ref{Claim 2}, $R(x^\prime, t^\prime) \le 2
R(\overline{x}, \overline{t})$ whenever
$\overline{t}-\frac12 \alpha Q^{-1} \le t^\prime \le \overline{t}$
and $\dist_{t^\prime}(x^\prime, x_0) \le
\dist_{\overline{t}}(\overline{x}, x_0) + \widehat{A} Q^{- \: \frac12}$. We
claim that the time-$\overline{t}$ ball
$B(x_0, \dist_{\overline{t}}(\overline{x}, x_0) + \frac{1}{10} \widehat{A}
Q^{-1/2})$ is contained in the time-$t^\prime$ ball
$B(x_0, \dist_{\overline{t}}(\overline{x}, x_0) + \widehat{A} Q^{- \: \frac12})$.
To see this, we apply Lemma \ref{distancedist} with
$r_0 \: = \: \frac{1}{2} Q^{-1/2}$ to give
\begin{equation}
\dist_t(x_0, \overline{x}) \: - \:
\dist_{\overline{t}}(x_0, \overline{x}) \: \le \: \const(n) \:
\alpha Q^{-1/2} \: \le \: \lambda \const(n) \widehat{A} Q^{-1/2}.
\end{equation}
If $\lambda$ is sufficiently small then the claim follows. The argument
also shows that it is consistent to use the curvature bound when
applying Lemma \ref{distancedist}.

Hence $R(x^\prime, t^\prime) \le 2
R(\overline{x}, \overline{t})$ whenever
$\overline{t} \: - \: \frac{1}{2}  \alpha
Q^{-1} \: \le \: t^\prime \: \le \: \overline{t}$ and
$\dist_{\overline{t}}(x^\prime, x_0) \: \le \:
\dist_{\overline{t}}(\overline{x}, x_0) + \frac{1}{10} \widehat{A}
Q^{-1/2}$. It follows that
$R(x^\prime, t^\prime) \le 2
R(\overline{x}, \overline{t})$ whenever
$\overline{t} \: - \: \frac{1}{2}  \alpha
Q^{-1} \: \le \: t^\prime \: \le \: \overline{t}$ and
$\dist_{\overline{t}}(x^\prime, \overline{x}) \: \le \: \frac{1}{10}
\widehat{A} Q^{-1/2}$. This shows that there is an $A^\prime = A^\prime
(B, C)$, which goes
to infinity as $B, C \rightarrow \infty$, so that
$R(x^\prime, t^\prime) \le 2
R(\overline{x}, \overline{t})$ whenever
$\overline{t} \: - \: A^\prime
Q^{-1} \: \le \: t^\prime \: \le \: \overline{t}$ and
$\dist_{\overline{t}}(x^\prime, \overline{x}) \: \le \: A^\prime
Q^{-1/2}$.

Now suppose that Corollary \ref{corI.11.6}(a) is not true.
Fixing $w > 0$,
for any sequences $\{B_k\}_{k=1}^\infty$ and $\{C_k\}_{k=1}^\infty$
going to infinity and
for each $k$, there is a Ricci flow solution $g_k(\cdot)$
which satisfies the hypotheses of
the corollary but for which $R(x_k,t_k) \ge C_k \: + \: B_k (t_k - t_{0,k})^{-1}$
for some point $(x_k, t_k)$ satisfying $\dist_{t_k}(x_k, x_{0,k}) \le \frac14$.
We can assume that $\lambda^2 \: C_k^{\frac12} \: \ge \: B_k$.
From the preceding discussion, there is a sequence $A^\prime_k \rightarrow
\infty$ and points $(\overline{x}_k, \overline{t}_k)$
with $\dist_{\overline{t}_k}(\overline{x}_k, {x}_{0,k}) \le \frac13$ so that
$R(x_k^\prime, t_k^\prime) \le 2
R(\overline{x}_k, \overline{t}_k)$ whenever
$\overline{t}_k \: - \: A_k^\prime
Q_k^{-1} \: \le \: t_k^\prime \: \le \: \overline{t}_k$ and
$\dist_{\overline{t}_k}(x_k^\prime, \overline{x}_k) \: \le \: A_k^\prime
Q_k^{-1/2}$, where
\begin{equation}
Q_k \: = \: R(\overline{x}_k, \overline{t}_k) \: \ge \:
B_k \: (\overline{t}_k - t_{0,k})^{-1} \: \ge \: B_k.
\end{equation}
By Corollary \ref{corI.11.5}, for any $\epsilon > 0$ there is some
$A = A(\epsilon)< \infty$
so that for large $k$,
\begin{equation}
\vol(B(\overline{x}_k, A/\sqrt{Q_k})) \: \le \: \epsilon
(A/\sqrt{Q_k})^n
\end{equation}
at time zero. By the Bishop-Gromov inequality,
$\vol(B(\overline{x}_k, 1)) \: \le \: \epsilon$ for large $k$, since
$Q_k \rightarrow \infty$.
If we took $\epsilon$ sufficiently small from the beginning then we
would get a contradiction to the fact that
\begin{equation}
\vol(B(\overline{x}_k, 1)) \: \ge \: \vol \left( B \left( x_{0,k}, \frac23 \right) \right) \: \ge
\left( \frac23 \right)^n \vol(B(x_{0,k}, 1)) \: \ge \left( \frac23 \right)^n w.
\end{equation}


For part (b), the idea is to choose the parameter
$\tau_0$ sufficiently small so that
we will still have the estimate
$\vol(B(x_0, r_0)) \: \ge \: 5^{-n} \: w \: r_0^n$ for the time-$t$ ball
$B(x_0, r_0)$ when $t \in [-\tau_0 r_0^2, 0]$,
and so we can apply part (a) with $w$ replaced by $\frac{w}{5}$.
The value of $\tau_0$ will emerge from the proof.
More precisely, putting $r_0 \: = \: 1$ and with a given
$\tau_0$, let $\tau$ be the
largest number in $[0, \tau_0]$ so that the time-$t$ ball
$B(x_0, 1)$ satisfies
$\vol(B(x_0, 1)) \: \ge \: 5^{-n} \: w$ whenever $t \in [-\tau, 0]$.
If $\tau \: < \: \tau_0$ then at time $- \tau$, we have
$\vol(B(x_0, 1)) \: = \: 5^{-n} \: w$.
The conclusion of part (a) holds in the sense that
\begin{equation} \label{cc}
R(x, t) \: \le \: C(5^{-n} w) \: + \: B(5^{-n} w) ( t + \tau)^{-1}
\end{equation}
whenever $t \in [- \tau, 0]$ and $\dist_t(x, x_0) \: \le \: \frac14$.
Lemma \ref{distancedist}, along with (\ref{cc}), implies that the
time-$(- \tau)$ ball $B(x_0, \frac14)$ contains the
time-$0$ ball $B(x_0, \frac14 \: - \: 10(n-1)(\tau \sqrt{C} \: + \:
2 \sqrt{B \tau}))$. From the nonnegative curvature, the
time-$(- \tau)$ volume of the first ball is at least as large as the
time-$0$ volume of the second ball.  Then
\begin{align} \label{lll}
5^{-n} \: w \: & = \: \vol(B(x_0, 1))
\: \ge \: \vol(B(x_0, \frac14)) \\
& \ge \:
\vol(B(x_0, \frac14 \: - \: 10(n-1)(\tau \sqrt{C} \: + \:
2 \sqrt{B \tau}))) \notag \\
& \ge
(\frac14 \: - \: 10(n-1)(\tau \sqrt{C} \: + \:
2 \sqrt{B \tau}))^n \: \vol(B(x_0, 1)) \notag \\
& \ge \:
(\frac14 \: - \: 10(n-1)(\tau \sqrt{C} \: + \:
2 \sqrt{B \tau}))^n \: w, \notag
\end{align}
where the balls on the top line of (\ref{lll}) are at time-($- \tau$),
and the other balls are at time-$0$.
Thus $\frac14 \: - \: 10(n-1)(\tau \sqrt{C} \: + \:
2 \sqrt{B \tau}) \: \le \: \frac15$.
This contradicts our assumption
that $\tau \: < \: \tau_0$ provided that
$\frac14 \: - \: 10(n-1)(\tau_0 \sqrt{C} \: + \:
2 \sqrt{B \tau_0}) \: = \: \frac15$.
\end{proof}

Finally, we give a version of Corollary \ref{corI.11.6}(b)  where instead of assuming a
nonnegative curvature operator, we assume that the
curvature operator
in the time-dependent ball of radius $r_0$ around $x_0$ is bounded below by
$- \: r_0^{-2}$.
\begin{corollary}[existence for curvature operators]
  \label{cor:kl-curvature-bounds-ricci-flow-solutions-nonnegative-curvature-operator-ass-existence-curvature-operators-3}
  \dcref{Corollary 45.13}
  \group{Kleiner--Lott Kappa Solutions}
  \source{kleiner-lott}
 (cf. end of Section I.11.6)\label{I.11.6end}
For every $w > 0$, there are $B = B(w) < \infty$, $C = C(w) < \infty$ and
$\tau_0 = \tau_0(w) > 0$ with the following property.
Suppose that we have a (not necessarily complete) Ricci flow
solution $(M, g(\cdot))$, defined for $t \in [- \tau_0 r_0^2, 0]$, so that at time zero
the metric ball $B(x_0, r_0)$ has compact closure. Suppose that
for each $t \in [- \tau_0 r_0^2, 0]$, the
curvature operator in the
time-$t$ ball $B(x_0, r_0)$ is bounded below by $- \: r_0^{-2}$.
If we assume a time-zero volume bound $\vol(B_0(x_0, r_0))
\ge w r_0^n$ then
\begin{equation}
R(x,t) \le C r_0^{-2} + B (t + \tau_0 r_0^2)^{-1}
\end{equation}
whenever $t \in [- \tau_0 r_0^2,0]$ and $\dist_t(x, x_0) \le \frac14 r_0$.
\end{corollary}
\begin{proof}
The blowup argument goes through as before.
The only real difference is that the volume of the
time-$(- \tau)$ ball $B(x_0, \frac14)$ will be at least
$e^{-\const \tau r_0^{-2}}$ times the volume of the
time-$0$ ball $B(x_0, \frac14 \: - \: 10(n-1)(\tau \sqrt{C} \: + \:
2 \sqrt{B \tau}))$.
\end{proof}

\section{I.11.7. Compactness of the space of three-dimensional
$\kappa$-solutions}
\label{pf11.7}

In this section we prove a compactness result for the space of three-dimensional
$\kappa$-solutions.  The three-dimensionality assumption is used to
show that
the limit solution has bounded curvature.

If a three-dimensional $\kappa$-solution $M$ is compact then it is
diffeomorphic to a quotient of $S^3$ or $\R \times S^2$, as it has
nonnegative curvature and is nonflat.
If its asymptotic
soliton (see Section \ref{I.11.2}) is also closed then
$M$ is a quotient of the round $S^3$ or $\R \times S^2$.
There are $\kappa$-solutions on $S^3$ and $\R P^3$ with noncompact
asymptotic soliton; see \cite[Section 1.4]{Perelman2}.
They are not isometric to the round metric; this corrects the statement in
the first paragraph of \cite[Section 11.7]{Perelman}.
\begin{theorem}[compactness for kappa solutions]
  \label{thm:kl-compactness-space-three-dimensional-solutions-compactness-kappa-solutions}
  \dcref{Theorem 46.1}
  \group{Kleiner--Lott Kappa Solutions}
  \source{kleiner-lott}
 (cf. Theorem I.11.7) \label{thmI.11.7}
Given $\kappa > 0$, the set of oriented
three-dimensional $\kappa$-solutions is compact modulo scaling.  That is, from any
sequence of such solutions and points $(x_k, 0)$,
after appropriate dilations
we can extract a
smoothly converging subsequence that satisfies the same conditions.
\end{theorem}
\begin{proof}
  \uses{cor:kl-solution-curvature-normalized-volume-control-each-other-existence-kappa-solutions,cor:kl-two-dimensional-solutions-kappa-solutions-corollary,prop:kl-asymptotic-scalar-curvature-asymptotic-volume-ratio-existence-kappa-solutions}
If $(M_k, (x_k, 0), g_k(\cdot))$ is a sequence of such $\kappa$-solutions with
$R(x_k, 0) = 1$ then
parts 1 and 3 of Corollary \ref{11.4cors} imply there is a subsequence that
converges to an ancient solution $(M_\infty, (x_\infty, 0), g_\infty(\cdot))$
which has nonnegative curvature operator and is $\kappa$-noncollapsed.
The remaining issue is to show that it has bounded curvature.
Note that $R_t\geq 0$ since $g_\infty(\cdot)$  is  a limit of a sequence
of Ricci flows satisfying $R_t\geq 0$.
Hence it is
enough to show that $(M_\infty, g(0))$ has bounded scalar curvature.

If not, there is a sequence of points $y_i$ going to infinity in $M_\infty$
such that $R(y_i, 0) \rightarrow \infty$ and
$R(y,0) \le 2 R(y_i,0)$ for $y \in B(y_i, A_i R(y_i,0)^{- \: \frac12})$, where
$A_i \rightarrow \infty$; compare \cite[Lemma 22.2]{Hamilton}.
Using the $\kappa$-noncollapsing, a subsequence of the rescalings
$(M_\infty, y_i, R(y_i, 0) g_\infty)$ will converge to a limit manifold
$N_\infty$.
As in the proof of Proposition \ref{propI.11.4}
from Appendix \ref{Alexandrov},
$N_\infty$ will split off a line.
By Corollary \ref{2Dsoliton} or Section \ref{altpf11.3},
$N_\infty$ must be the standard solution
on $\R \times S^2$. Thus $(M_\infty, g(0))$ contains a sequence
$D_i$ of neck regions, with their
cross-sectional radii
tending to zero as
$i \rightarrow \infty$.

Note that $M_\infty$
has to be 1-ended.  Otherwise, it would contain a line,
and would therefore have to split off a line isometrically
\cite[Theorem 8.17]{Cheeger-Ebin}.
But then $M_\infty$, the product of a line and a surface,
could not have neck regions with cross-sections tending to zero.

From the theory of nonnegatively curved manifolds
\cite[Chapter 8.5]{Cheeger-Ebin},  there is an
exhaustion $M_\infty \: = \: \bigcup_{t\ge 0}C_t$
by nonempty totally convex compact sets $C_t$ so that
$(t_1 \le t_2) \Rightarrow (C_{t_1} \subset C_{t_2})$, and
\begin{equation} \label{C}
C_{t_1} \: = \: \{ q \in C_{t_2} \: : \: \dist(q, \partial C_{t_2})
\: \ge \: t_2 - t_1\}.
\end{equation}
Now consider a
neck region $D$ which is  close to a cylinder.  Note by
triangle comparison -- or simply because the distance function
in $D$ is close to that of a product metric --  any
minimizing geodesic  segment $\ga\subset D$
of length large compared to cross-sectional radius of $D$ must
be nearly orthogonal to the cross-section.
It follows from this and (\ref{C}) that if $t>0$ and
$\partial C_t$ contains a point $p\in D$ such that $d(p,\D D)$
is large compared to the cross-section of $D$, then
 $\partial C_t\cap D$ is an approximate $2$-sphere
cross-section of $D$.
Fix such a neck region $D_0$ and let $C_{t_0}$ be the
corresponding convex set.
As $M_\infty$
has one end, $\partial C_{t_0}$ has only one connected
component,
namely the approximate $2$-sphere cross-section.

For all $t \: > \: t_0$, there is a distance-nonincreasing retraction
$r \: : \: C_t \rightarrow C_{t_0}$ which maps $C_t \: - \: C_{t_0}$ onto
$\partial C_{t_0}$ \cite{Sharafutdinov}. Let $D$ be a neck region
with a very small cross-section and let $C_t$ be a convex set so that
$\partial C_t$ intersects $D$ in an approximate $2$-sphere cross-section.
Then $\partial C_t$ consists entirely of this approximate cross-section.
The restriction of $r$ to $\partial C_t$ is distance-nonincreasing, but
will map the $2$-sphere $\partial C_t$ onto
the $2$-sphere $\partial C_{t_0}$. This is a contradiction.
\end{proof}

\begin{remark}[compactness for kappa solutions]
  \label{rem:kl-compactness-space-three-dimensional-solutions-compactness-kappa-solutions}
  \dcref{Remark 46.3}
  \group{Kleiner--Lott Kappa Solutions}
  \source{kleiner-lott}
 The statement of \cite[Theorem 11.7]{Perelman} is about
noncompact $\kappa$-solutions but the proof works whether the
solutions are compact or noncompact.
\end{remark}


\begin{remark}[compactness for kappa solutions]
  \label{rem:kl-compactness-space-three-dimensional-solutions-compactness-kappa-solutions-3}
  \dcref{Remark 46.4}
  \group{Kleiner--Lott Kappa Solutions}
  \source{kleiner-lott}
  \uses{cor:kl-two-dimensional-solutions-kappa-solutions-corollary}

 One may wonder where we have used the fact that we have a
Ricci flow solution, i.e. whether the curvature is bounded for
any $\kappa$-noncollapsed Riemannian $3$-manifold with nonnegative
sectional curvature.  Following the above argument, we could again split off
a line in a rescaling around high-curvature points.  However, we would
not necessarily know that the ensuing nonnegatively-curved
surface is compact. ({\it A priori}, it could be a smoothed-out cone, for
example.)  In the case of a Ricci flow, the compactness comes from
Corollary \ref{2Dsoliton} or Section \ref{altpf11.3}.
\end{remark}



\begin{corollary}[kappa solutions corollary]
  \label{cor:kl-compactness-space-three-dimensional-solutions-kappa-solutions-corollary}
  \dcref{Corollary 46.5}
  \group{Kleiner--Lott Kappa Solutions}
  \source{kleiner-lott}

Let $(M,g(\cdot))$ be a $3$-dimensional $\kappa$-solution.
Then any asymptotic soliton constructed as in Section \ref{I.11.2} is
also a $\kappa$-solution.
\end{corollary}


\section{I.11.8. Necklike behavior at infinity of a three-dimensional
$\kappa$-solution - weak version}

The next corollary says that outside of a compact region,  any oriented noncompact
three-dimensional $\kappa$-solution looks necklike (after rescaling).
In this section we give a simple argument to prove the corollary, except for
a diameter bound on the compact region.  In the next section we give an
argument that also proves the diameter bound.

More information on three-dimensional $\kappa$-solutions is in
Section \ref{II.1}.

\begin{definition}[kappa solution definition]
  \label{def:kl-necklike-behavior-at-infinity-three-dimensional-solution-weak-version-kappa-solution-definition}
  \dcref{Definition 47.1}
  \group{Kleiner--Lott Kappa Solutions}
  \source{kleiner-lott}
  \uses{def:kl-notation-terminology-smooth-closeness-epsilon-necks}

Fix $\epsilon > 0$. Let $(M, g(\cdot))$ be an oriented three-dimensional $\kappa$-solution.
We say that a point $x_0 \in M$ is the {\em center of an $\epsilon$-neck} if the
solution $g(\cdot)$ in the set
$\{(x,t) \: : \: - (\epsilon Q)^{-1} < t \le 0, \dist_0(x, x_0)^2 < (\epsilon Q)^{-1}\}$,
where $Q = R(x_0, 0)$, is, after scaling with the factor $Q$, $\epsilon$-close
in some fixed smooth topology to the corresponding subset of the
evolving round cylinder (having scalar curvature one at time zero).
(See Definition \ref{closenessdef} below for a more precise statement.)

We let $M_\epsilon$ denote the points in $M$ that are not centers of
$\epsilon$-necks.
\end{definition}
\begin{corollary}[existence for kappa solutions]
  \label{cor:kl-necklike-behavior-at-infinity-three-dimensional-solution-weak-version-existence-kappa-solutions}
  \dcref{Corollary 47.2}
  \group{Kleiner--Lott Kappa Solutions}
  \source{kleiner-lott}
 (cf. Corollary I.11.8) \label{corI.11.8}
For any $\epsilon > 0$, there exists $C = C(\epsilon, \kappa) > 0$ such that
if $(M, g(\cdot))$ is an oriented noncompact three-dimensional $\kappa$-solution then \\
1.  $M_\epsilon$ is compact with $\diam(M_\epsilon) \le C Q^{- \: \frac12}$ and \\
2. $C^{-1} Q \le R(x,0) \le CQ$ whenever $x \in M_\epsilon$, \\
where $Q = R(x_0, 0)$ for some $x_0 \in \partial M_\epsilon$.
\end{corollary}
\begin{proof}
  \uses{cor:kl-solution-curvature-normalized-volume-control-each-other-existence-kappa-solutions,prop:kl-asymptotic-scalar-curvature-asymptotic-volume-ratio-existence-kappa-solutions,thm:kl-compactness-space-three-dimensional-solutions-compactness-kappa-solutions}
We  prove here the claims of
Corollary \ref{corI.11.8}, except for the
diameter bound. In the next section we give another argument which
also proves the diameter bound.

We claim first that $M_\epsilon$ is compact.  Suppose not.
Then there is a sequence of points $x_k \in M_\epsilon$ going to
infinity. Fix a basepoint $x_0 \in M$. Then
$R(x_0) \: \dist_0^2(x_0, x_k) \rightarrow \infty$. By
part 5 of Corollary \ref{11.4cors},
$R(x_k) \: \dist_0^2(x_0, x_k) \rightarrow \infty$.  Rescaling
around $(x_k, 0)$ to make its scalar curvature one, we can use
Theorem \ref{thmI.11.7} to extract a convergent subsequence
$(M_\infty, x_\infty)$.
As in the proof of Proposition \ref{propI.11.4}, we can say that
$(M_\infty, x_\infty)$ splits off a line.  Hence for large $k$,
$x_k$ is the center of an $\epsilon$-neck, which is a contradiction.

Next we claim that for any $\epsilon$, there exists
$C = C(\epsilon, \kappa) > 0$ such that if $g_{ij}(t)$ is a $\kappa$-solution
then for
any point $x \in M_\epsilon$, there is a point $x_0 \in
\partial M_\epsilon$ such that $\dist_0(x, x_0) \: \le \: CQ^{-1/2}$
and $C^{-1} Q \: \le \: R(x, 0) \: \le \: CQ$, where
$Q \: = \: R(x_0, 0)$.

If not then there is a sequence $\{M_i\}_{i=1}^\infty$
of $\kappa$-solutions along with points
$x_i \in M_{i,\epsilon}$ such that for each $y_i \in
\partial M_{i,\epsilon}$, we have \\
1. $\dist_0^2(x_i, y_i) \: R(y_i, 0) \: \ge \: i$ or \\
2. $R(y_i, 0) \: \ge \: i \: R(x_i, 0)$ or \\
3. $R(x_i, 0) \: \ge \: i \: R(y_i, 0)$.

Rescale the metric on $M_i$ so that $R(x_i, 0) \: = \: 1$. From
Theorem \ref{thmI.11.7}, a subsequence of the pointed spaces
$(M_i, x_i)$ will converge smoothly to a $\kappa$-solution
$(M_\infty, x_\infty)$. Also, $x_\infty \in M_{\infty, \epsilon}$.

Taking a subsequence, we
can assume that 1. occurs for each $i$, or 2. occurs for each $i$,
or 3. occurs for each $i$. If $M_\infty \neq M_{\infty, \epsilon}$,
choose $y_\infty \in
\partial M_{\infty,\epsilon}$. Then $y_\infty$ is the limit of
a subsequence of points $y_i \in \partial M_{i,\epsilon}$.

If 1. occurs for each $i$ then $\dist_0^2(x_\infty, y_\infty) \:
R(y_\infty, 0) \: = \: \infty$, which is impossible.
If 2. occurs for each $i$ then $R(y_\infty, 0) \: = \: \infty$, which
is impossible.
If 3. occurs for each $i$ then $R(y_\infty, 0) \: = \: 0$.
It follows from (\ref{harnackinequality}) that $M_\infty$ is flat,
which is impossible, as $R(x_\infty, 0) \: = \: 1$.

Hence $M_\infty \: = \: M_{\infty,\epsilon}$, i.e. no point in the
noncompact ancient solution $M_\infty$ is the center of an $\epsilon$-neck.
This contradicts the previous conclusion that
$M_{\infty,\epsilon}$ is compact.
\end{proof}

\section{I.11.8. Necklike behavior at infinity of a three-dimensional
$\kappa$-solution - strong version}
\label{strongversion}

The following corollary is an application of the compactness result
Theorem \ref{thmI.11.7}.
it is a refinement of
\cite[Cor. I.11.8]{Perelman}.

\begin{corollary}[existence for kappa solutions]
  \label{cor:kl-necklike-behavior-at-infinity-three-dimensional-solution-strong-version-existence-kappa-solutions}
  \dcref{Corollary 48.1}
  \group{Kleiner--Lott Kappa Solutions}
  \source{kleiner-lott}

\label{neckstructure}
For all $\kappa>0$, there exists an $\eps_0>0$ such that
for all $0<\eps<\eps_0$ there exists an $\al=\al(\eps,\kappa)$
with the property that for  any $\kappa$-solution $(M,g(\cdot))$,
and at any time $t$, precisely one of the following holds
($M_\eps$ denotes the set of points which are not centers of $\eps$-necks
at time $t$):

A.  $(M,g(\cdot))$ is round cylindrical flow, and
so every point at every time is the center of an $\eps$-neck for all
$\eps>0$.

B.  $M$ is noncompact, $M_\eps\neq\emptyset$, and for
all $x,y\in M_\eps$, we have $R(x)d^2(x,y)<\al$.

C.  $M$ is compact, and there is a pair of points
$x,y\in M_\eps$ such that $R(x)d^2(x,y)>\al$,
\begin{equation}
M_\eps\subset B(x,\al R(x)^{-\frac{1}{2}})\cup
B(y,\al R(y)^{-\frac{1}{2}}),
\end{equation}
and there is a minimizing geodesic $\ol{xy}$ such that
every $z\in M - M_\eps$ satisfies $R(z)d^2(z,\ol{xy})<\al$.

D. $M$ is compact and there exists a point $x\in M_\eps$
such that $R(x)d^2(x,z)<\al$ for all $z\in M$.
\end{corollary}
\begin{lemma}[existence for kappa solutions]
  \label{lem:kl-necklike-behavior-at-infinity-three-dimensional-solution-strong-version-existence-kappa-solutions}
  \dcref{Lemma 48.3}
  \group{Kleiner--Lott Kappa Solutions}
  \source{kleiner-lott}

\label{annoy}
For all $\eps>0$, $\kappa>0$,  there exists
$\al=\al(\eps,\kappa)$  with the following
property.  Suppose $(M,g(\cdot))$ is any
$\kappa$-solution, $x,y,z\in M$, and
at time $t$ we have $x, y \in M_\epsilon$
and $R(x)d^2(x,y)>\al$.  Then at time $t$ either
$R(x)d^2(z,x)<\al$ or $R(y)d^2(z,y)<\al$ or ($R(z)d^2(z,\ol{xy})<\al$ and
$z \notin M_\epsilon$).
\end{lemma}
\begin{proof}
  \uses{cor:kl-solution-curvature-normalized-volume-control-each-other-existence-kappa-solutions,thm:kl-compactness-space-three-dimensional-solutions-compactness-kappa-solutions}
Pick $\eps>0,\,\kappa>0$, and suppose no such $\al$ exists.  Then
there is a sequence $\al_k\ra\infty$, a
sequence of $\kappa$-solutions $(M_k,g_k(\cdot))$,
and sequences $x_k,y_k,z_k\in M_k$, $t_k\in \R$ violating the $\al_k$-version
of the statement
for all $k$.  In particular, $x_k, y_k \in (M_k)_\epsilon$ and
\begin{equation}
R(x_k,t_k)d^2_{t_k}(x_k,y_k)\ra\infty,\;
R(x_k,t_k)d^2_{t_k}(z_k,x_k)\ra\infty,\;
\mbox{and}\,R(y_k,t_k)d^2_{t_k}(z_k,y_k)\ra\infty.
\end{equation}

 Let $z_k'\in\ol{x_ky_k}$ be a point in $\ol{x_ky_k}$
nearest $z_k$ in $(M_k,g_k(t_k))$.



We first show that $R(x_k,t_k)d_{t_k}^2(z_k',x_k)\ra\infty$.
If not, we may pass to a subsequence on which $R(x_k,t_k)d_{t_k}^2(z_k',x_k)$
remains bounded.   Applying Theorem \ref{thmI.11.7}, we may
pass to a subsequence and
rescale by $R(x_k,t_k)$,
to make the sequence $(M_k,(x_k,t_k),g_k(\cdot))$
converge to a $\kappa$-solution $(M_\infty,(x_\infty,0),g_\infty(\cdot))$,
the segments $\ol{x_ky_k}\subset (M_k,g_k(t_k))$
converge to a ray $\ol{x_\infty\xi}\subset (M_\infty,g_\infty(0))$,
and the segments $\ol{z_k'z_k}$ converge to
a ray $\ol{z_\infty'\eta}$.  Recall that the comparison angle
$\cangle_{z^\prime_\infty}(u,v)$ tends to the Tits angle
$\tangle(\xi,\eta)$ as $u\in\ol{z_\infty'\xi}$, $v\in \ol{z_\infty'\eta}$
tend to infinity.  Since $d(z_k,z_k')=d(z_k,\ol{x_ky_k})$ we must
have $\tangle(\xi,\eta)\geq\frac{\pi}{2}$.  Now consider a sequence
$u_k\in \ol{z_\infty'\xi}$ tending to infinity.  By
Theorem \ref{thmI.11.7}, part 5 of Corollary \ref{11.4cors},
and the remarks about Alexandrov spaces in Appendix \ref{Alexandrov},
if we rescale $(M_\infty,(u_k,0),g_\infty(\cdot))$ by $R(u_k,0)$,
we get round cylindrical flow as a limit.   When $k$ is sufficiently
large,
we may find an almost product region $D\subset (M_\infty,g_\infty(\cdot))$
containing $u_k$ which is disjoint from $\ol{z_\infty'\eta}$,
and whose cross-section $\Si\times\{0\}\subset\Si\times (-1,1)\simeq D$
intersects the ray $\ol{z_\infty'\xi}$ transversely at a single point.
This implies that $\Si\times \{0\}$ separates the two ends
of $\ol{z_\infty'\xi}\cup\ol{z_\infty'\eta}$ from each other; hence
$M_\infty$ is two-ended, and $(M_\infty,g_\infty(\cdot))$ is round
cylindrical flow.   This contradicts the assumption that
$x_k$ is not the center of
an $\eps$-neck.  Hence $R(x_k,t_k)d_{t_k}^2(z_k',x_k)\ra\infty$,
and similar reasoning shows that $R(y_k,t_k)d_{t_k}^2(z_k',y_k)\ra\infty$.

By part 5 of Corollary \ref{11.4cors}, we therefore have
$R(z_k',t_k)d_{t_k}^2(z_k',x_k)\ra\infty$ and
$R(z_k',t_k)d_{t_k}^2(z_k',y_k)\ra\infty$.  Rescaling the sequence
$(M_k,(z_k',t_k),g_k(\cdot))$ by
$R(z_k',t_k)$, we get convergence to round cylindrical flow
(since any limit flow contains a line), and $\ol{z_k'z_k}$
subconverges to a segment orthogonal to the $\R$-factor, which
implies that $R(z_k',t_k)d_{t_k}^2(z_k,z_k')$ is bounded
and $z_k$ is the center of an $\eps$-neck for large $k$.  This contradicts
our assumption that the $\al_k$-version of the lemma is
violated for each $k$.
\end{proof}
\begin{proof}
  \proves{cor:kl-necklike-behavior-at-infinity-three-dimensional-solution-strong-version-existence-kappa-solutions}

\bigskip
\noindent
{\em Proof of Corollary \ref{neckstructure}.}
Let $(M,g(\cdot))$
be a $\kappa$-solution, and $\eps>0$.

\medskip
{\em Case 1: Every $x\in (M,g(t))$
is the center of an $\eps$-neck.}  In this case, if $\eps>0$ is sufficiently
small, $M$ fibers over a
$1$-manifold with fiber $S^2$.  If the
$1$-manifold is homeomorphic to $\R$, then
$M$ has two ends, which implies that the flow
$(M,g(\cdot))$ is an evolving round cylinder.  If the
base of the fibration were a circle, then the universal
cover $(\tilde M,\tilde g(t))$ would split off
a line, which would imply that the universal covering
flow would be a round cylindrical flow; but this
would violate the $\kappa$-noncollapsed assumption
at very negative times.  Thus A holds in this
case.


{\em Case 2: There exist $x,y\in M_\eps$ such that
$R(x)d^2(x,y)>\al$.}  By Lemma \ref{annoy} and Corollary \ref{11.4cors}
part 5, for all $z \in M - \left( B(x,\al R(x)^{-\frac{1}{2}})\cup
B(y,\al R(y)^{-\frac{1}{2}}) \right)
$, we have $R(z)d^2(z,\ol{xy})<\alpha$ and $z \notin M_\epsilon$.
This implies (again by
Corollary \ref{11.4cors} part 5) that there exists a
$\gamma=\gamma(\eps,\kappa)$ such that for every $z\in M$
there is a $z'\in\ol{xy}$ for which  $R(z')d^2(z',z)<\gamma$,
which means that $M$ must be compact, and C holds.

{\em Case 3: $M_\eps\neq \emptyset$, and for all $x,y\in M_\eps$,
we have $R(x)d^2(x,y)<\al$.}
If $M$ is noncompact then
we are in case B and are done, so assume that $M$ is compact.
Pick $x\in M_\eps$, and suppose  $z\in M$ maximizes
$R(x)d^2(\cdot,x)$.
If $R(x)d^2(z,x)\geq\al$,  then $z$ is the center of an $\eps$-neck, and we
may look at the cross-section $\Si$ of the neck region.
If $\Si$ separates $M$, then when $\eps>0$ is sufficiently
small, we  get a contradiction to the assumption
that $z$ maximizes $R(x)d^2(x,\cdot)$.  Hence
$\Si$ cannot separate $M$, and there is a loop
passing through $x$ which intersects $\Si$ transversely
at one point.  It follows that the universal covering flow
$(\tilde M,\tilde g(\cdot))$ is cylindrical flow, a contradiction.
Hence $R(x)d^2(x,z)<\al$ for all $z\in M$, so D holds.
\end{proof}



\section{More properties of $\kappa$-solutions}

In  this section we prove some additional properties of
$\kappa$-solutions. In particular, Corollary \ref{longsegment} implies
that if
$z$
lies in a geodesic segment $\eta$ in a $\kappa$-solution $M$
and if the endpoints of $\eta$ are sufficiently far from
$z$ (relative to $R(z)^{- \: \frac12}$) then $z \notin M_\epsilon$.
The results of this section will be used in the proof of Theorem
\ref{thmI.12.1}.

\begin{proposition}[existence for kappa solutions]
  \label{prop:kl-more-properties-solutions-existence-kappa-solutions}
  \dcref{Proposition 49.1}
  \group{Kleiner--Lott Kappa Solutions}
  \source{kleiner-lott}

\label{neckcrit}
For all $\kappa>0$, $\alpha>0$, $\theta>0$,
there exists a $\beta(\kappa,\alpha,\theta)<\infty$
such that if $(M,g(t))$ is a time slice of a $\kappa$-solution,
$x,y_1,y_2\in M$, $R(x)d^2(x,y_i)>\beta$ for $i=1,2$, and $\cangle_x(y_1,y_2)\geq \theta$,
then (a) $x$ is the center of
 an $\alpha$-neck, and (b) $\cangle_x(y_1,y_2)\geq \pi-
\alpha$.
\end{proposition}
\begin{proof}
  \uses{lem:kl-necklike-behavior-at-infinity-three-dimensional-solution-strong-version-existence-kappa-solutions}
The proof of this is similar to the first part of the proof of
Lemma \ref{annoy}.  Note that when $\alpha$ is small, then after enlarging
$\beta$ if necessary, the neck region around $x$ will separate
$y_1$ from $y_2$; this implies (b).
\end{proof}

\begin{corollary}[existence for kappa solutions]
  \label{cor:kl-more-properties-solutions-existence-kappa-solutions}
  \dcref{Corollary 49.2}
  \group{Kleiner--Lott Kappa Solutions}
  \source{kleiner-lott}
 \label{longsegment}
For all $\kappa>0$, $\eps>0$, there exists a $\rho=\rho(\kappa,\eps)$ such that
if $(M,g(t))$ is a time slice of a $\kappa$-solution, $\eta\subset (M,g(t))$
is a minimizing geodesic segment with endpoints $y_1,y_2$, $z\in M$, $z'\in \eta$
is a point in $\eta$ nearest $z$, and $R(z')d^2(z',y_i)>\rho$
for $i=1,2$, then
$z,z'$ are centers of
$\eps$-necks, and $\max(R(z)d^2(z,z'),R(z')d^2(z,z'))<4\pi^2$.
\end{corollary}
\begin{proof}
Pick $\eps'>0$.  Under the assumptions, if
\begin{equation}
\min(R(z')d^2(z',y_1),R(z')d^2(z',y_2))
\end{equation}
is sufficiently large, we can apply the preceding proposition
to the triple $z',y_1,y_2$, to conclude that $z'$ is the center of an
$\eps'$-neck.  Since the shortest segment from $z$ to $z'$
is orthogonal to $\eta$, when $\eps'$ is small enough
the segment $\ol{zz'}$ will lie close to an $S^2$ cross-section
in the approximating round cylinder, which
gives $R(z')d^2(z,z')\lesssim 2\pi^2$.
\end{proof}

\section{I.11.9. Getting a uniform value of $\kappa$}
\label{secI.11.9}


\begin{proposition}[existence for kappa solutions]
  \label{prop:kl-getting-uniform-value-existence-kappa-solutions}
  \dcref{Proposition 50.1}
  \group{Kleiner--Lott Kappa Solutions}
  \source{kleiner-lott}
 \label{propI.11.9}
There is a $\kappa_0 > 0$ so that if $(M, g(\cdot))$ is an oriented
three-dimensional $\kappa$-solution, for some
$\kappa > 0$, then it is a $\kappa_0$-solution or
it is a quotient of the round shrinking $S^3$.
\end{proposition}
\begin{proof}
  \uses{cor:kl-three-dimensional-noncompact-noncollapsed-gradient-shrinkers-standard-compactness-kappa-solutions,prop:kl-asymptotic-solitons-existence-ricci-solitons,thm:kl-no-local-collapsing-theorem-ii-existence-kappa-solutions}
Let $(M, g(\cdot))$ be a $\kappa$-solution.
Suppose that for some $\kappa^\prime>0$,
the solution is $\kappa^\prime$-collapsed at some scale.
After
rescaling, we can assume that there is a point $(x_0, 0)$ so that
$|\Rm(x, t)| \: \le \: 1$ for all $(x, t)$ satisfying
$\dist_0(x, x_0) \: < \: 1$ and $t \in [-1, 0]$, with
$\vol(B_0(x_0, 1)) \: < \: \kappa^\prime$.  Let $\widetilde{V}(t)$ denote
the reduced volume as a function of $t \in (- \infty, 0]$, as
defined using curves from $(x_0, 0)$. It is nondecreasing in $t$.
As in the proof of Theorem \ref{nolocalcollapse}, there is an estimate
$\widetilde{V}(-\kappa^\prime) \: \le \: 3 (\kappa^\prime)^{3/2}$. Take a sequence
of times $t_i \rightarrow - \infty$. For each $t_i$, choose $q_i \in M$
so that $l(q_i, t_i) \: \le \: \frac{3}{2}$. From the proof of Proposition
\ref{asympsoliton}, for all $\epsilon > 0$
there is a $\delta > 0$ such that $l(q, t)$
does not exceed $\delta^{-1}$ whenever $t \in [t_i, t_i/2]$ and
$\dist_{t_i}^2(q, q_i) \: \le \: \epsilon^{-1} t_i$. Given the
monotonicity of $\widetilde{V}$ and the upper bound on $l(q,t)$,
we obtain an upper bound on the volume of the time-$t_i$ ball
$B(q_i, \sqrt{t_i/\epsilon})$ of the form $\const \: t_i^{3/2} \:
e^{\delta^{-1}} \: (\kappa^\prime)^{3/2}$.

On the other hand, from
Proposition \ref{asympsoliton}, a subsequence of the rescalings
of the ancient solution around $(q_i, t_i)$ converges to a nonflat
gradient shrinking soliton.  If the gradient shrinking soliton is compact
then it must be a quotient of the round shrinking $S^3$
\cite{Hamiltonnnnn}. Otherwise,
Corollary \ref{cornosoliton} says that if the gradient shrinking soliton is
noncompact then it must be an evolving cylinder or its
$\Z_2$-quotient.  Fixing $\epsilon$, this gives a lower bound on
$\vol(B_{t_i} (q_i, \sqrt{t_i/\epsilon}))$ in terms of the noncollapsing
constants of the evolving cylinder and its $\Z_2$-quotient. Hence there
is a universal constant $\kappa_0$ so that if
$\kappa^\prime<\kappa_0$ then we obtain a contradiction to the
assumption of $\kappa^\prime$-collapsing.
\end{proof}
\begin{remark}[existence for Ricci solitons]
  \label{rem:kl-getting-uniform-value-existence-ricci-solitons}
  \dcref{Remark 50.2}
  \group{Kleiner--Lott Kappa Solutions}
  \source{kleiner-lott}
  \uses{cor:kl-three-dimensional-noncompact-noncollapsed-gradient-shrinkers-standard-compactness-kappa-solutions,prop:kl-asymptotic-solitons-existence-ricci-solitons,thm:kl-compactness-space-three-dimensional-solutions-compactness-kappa-solutions}
 The hypotheses of Corollary \ref{cornosoliton}
assume a global upper bound on the sectional curvature of any time slice,
which in the $n$-dimensional case
is not {\it a priori} true for the asymptotic soliton of
Proposition \ref{asympsoliton}.  However, in our $3$-dimensional case,
the argument of Theorem \ref{thmI.11.7} shows that there is such an upper bound.
\end{remark}

\section{II.1.2. Three-dimensional noncompact $\kappa$-noncollapsed
gradient shrinkers are standard}
\label{II.1.2}

In this section we show that any complete oriented
$3$-dimensional noncompact
$\kappa$-noncollapsed
gradient shrinking soliton with
bounded nonnegative curvature
is either the evolving round cylinder $\R \times S^2$ or its
$\Z_2$-quotient.

The basic example of a gradient shrinking soliton is the
metric on $\R \times S^2$ which gives the $2$-sphere
a radius of $\sqrt{-2t}$ at time $t \in (-\infty, 0)$. With
coordinates $(s, \theta)$ on $\R \times S^2$, the function $f$ is given by
$f(t, s, \theta) \: = \: - \: \frac{s^2}{4t}$.

\begin{lemma}[existence for kappa solutions]
  \label{lem:kl-three-dimensional-noncompact-noncollapsed-gradient-shrinkers-standard-existence-kappa-solutions}
  \dcref{Lemma 51.1}
  \group{Kleiner--Lott Kappa Solutions}
  \source{kleiner-lott}
 \label{nosoliton} (cf. Lemma of II.1.2)
There is no complete oriented $3$-dimensional noncompact
$\kappa$-noncollapsed gradient shrinking soliton with
bounded positive sectional curvature.
\end{lemma}
\par\medskip\noindent\textit{Source argument (not certified as complete).}\ \begingroup\itshape
The idea of the proof is to show that the soliton has the qualitative
features of a shrinking cylinder, and then to get a contradiction to
the assumption of positive sectional curvature.

Applying $\nabla_i$ to the gradient shrinker equation
\begin{equation} \label{IIeqn}
\nabla_i \nabla_j f \: + \: R_{ij} \: + \: \frac{1}{2t} \: g_{ij}
\: = \: 0
\end{equation}
gives
\begin{equation}
\triangle \nabla_j f \: + \: \nabla_i R_{ij} \: = \: 0.
\end{equation}
As $\nabla_i R_{ij} \: = \: \frac{1}{2} \: \nabla_j R$ and
$\triangle \nabla_j f \: = \: \nabla_j \triangle f \: + \:
R_{jk} \nabla_k f \: = \: \nabla_j \left( -R-\frac{n}{2t} \right) \: + \:
R_{jk} \nabla_k f$,
we obtain
\begin{equation} \label{gradR}
\nabla_i R \: = \: 2 \: R_{ij} \: \nabla_j f.
\end{equation}

Fix a basepoint $x_0 \in M$ and consider a normalized minimal geodesic
$\gamma \: : \: [0, \overline{s}] \rightarrow M$ in the time $-1$ slice
with $\gamma(0) \: = \: x_0$. Put $X(s) \: = \: \frac{d\gamma}{ds}$.
As in the proof of Lemma \ref{distancedist},
$\int_0^{\overline{s}} \Ric(X, X) \: ds \: \le \:
\const$ for some constant independent of $\overline{s}$.
If $\{Y_i\}_{i=1}^3$ are orthonormal parallel vector fields along $\gamma$
then
\begin{equation}
\left( \int_0^{\overline{s}} |\Ric(X, Y_1)| \: ds \right)^2 \: \le \:
\overline{s} \: \int_0^{\overline{s}} |\Ric(X, Y_1)|^2 \: ds \: \le \:
\overline{s} \: \sum_{i=1}^3 \int_0^{\overline{s}} |\Ric(X, Y_i)|^2 \: ds.
\end{equation}
Thinking of $\Ric$ as a self-adjoint linear operator on $TM$,
$\sum_{i=1}^3 |\Ric(X, Y_i)|^2 \: = \: \langle X, \Ric^2 X \rangle$.
In terms of a pointwise orthonormal frame  $\{e_i\}$
of eigenvectors of $\Ric$,
with eigenvalues $\lambda_i$, write $X \: = \: \sum_{i=1}^3 X_i e_i$.
Then
\begin{equation}
\langle X, \Ric^2 X \rangle \: = \: \sum_{i=1}^3 \lambda_i^2 \: X_i^2
\: \le \: (\sum_{i=1}^3 \lambda_i) \: (\sum_{i=1}^3 \lambda_i \: X_i^2)
\: = \: R \cdot \Ric(X, X).
\end{equation}
Hence
\begin{equation}
\left( \int_0^{\overline{s}} |\Ric(X, Y_1)| \: ds \right)^2 \: \le \:
(\sup_M R) \: \overline{s} \:
\int_0^{\overline{s}} \Ric(X, X) \: ds \: \le \:
\const \: \overline{s}.
\end{equation}

Multiplying (\ref{IIeqn}) by $X^i X^j$ and summing gives
$\frac{d^2f(\gamma(s))}{ds^2} \: + \: \Ric(X, X) \: - \: \frac{1}{2}
\: =\: 0$. Then
\begin{equation}
\frac{d f(\gamma(s))}{ds} \Big|_{s = \overline{s}} \: = \:
\frac{d f(\gamma(s))}{ds} \Big|_{s = 0} \: + \:
\frac{1}{2} \: \overline{s} \: - \: \int_0^{\overline{s}} \Ric(X, X) \: ds
\: \ge \: \frac{1}{2} \: \overline{s} \: - \: \const
\end{equation}
This implies that there is a compact subset of $M$ outside of which
$f$ has no critical points.

If $Y$ is a unit vector field perpendicular to $X$ then
multiplying (\ref{IIeqn}) by $X^i Y^j$ and summing gives
$\frac{d}{ds} (Y \cdot f)(\gamma(s)) \: + \: \Ric(X, Y)
\: =\: 0$. Then
\begin{equation}
(Y \cdot f)(\gamma(\overline{s})) \: = \:
(Y \cdot f)(\gamma(0)) \: - \:
\int_0^{\overline{s}} \: \Ric(X, Y) \: ds
\end{equation}
and
\begin{equation}
|(Y \cdot f)(\gamma(\overline{s}))| \: \le \:
\const (\sqrt{\overline{s}} \: +\: 1).
\end{equation}
For large $\overline{s}$, $|(Y \cdot f)(\gamma(\overline{s}))|$
is small compared to $(X \cdot f)(\gamma(\overline{s}))$. This
means that as one approaches infinity, the gradient of $f$ becomes
more and more parallel to the gradient of the distance function
from $x_0$, where by the latter we mean the vectors $X$ that are
tangent to minimal geodesics.

The gradient flow of $f$ is given by the equation
\begin{equation}
\frac{dx}{du} \: =\: (\nabla f)(x).
\end{equation}
Then along a flowline, equation (\ref{gradR}) implies that
\begin{equation}
\frac{dR(x)}{du} \: = \: \left\langle \nabla R, \frac{dx}{du} \right\rangle
\: = \: 2 \Ric(\nabla f, \nabla f).
\end{equation}
In particular, outside of a compact set, $R$ is strictly increasing
along the flowlines. Put $\overline{R} \: = \: \limsup_{x \rightarrow \infty}
R$.
Take points $x_\alpha$ tending toward infinity, with
$R(x_\alpha) \rightarrow \overline{R}$. Putting
${r}_\alpha \: = \: \sqrt{\dist_{-1}(x_0, x_\alpha)}$, we have
$\frac{{r}_\alpha}{\dist_{-1}(x_0, x_\alpha)} \rightarrow 0$ and
$R(x_\alpha) \: {r}_\alpha^2 \rightarrow \infty$.
Then the argument of the proof of Proposition \ref{propI.11.4}
shows that any convergent
subsequence of the rescalings around $(x_\alpha, -1)$ splits off
a line.  Hence the limit is a shrinking round cylinder with
scalar curvature $\overline{R}$ at time $-1$. Because our original
solution exists up to time zero, we must have
$\overline{R} \le 1$.
Now equation
(\ref{shrinkingflow}) says that the Ricci flow
is given by $g(-t) \: = \: - \: t \: \phi_t^* g(-1)$, where
$\phi_t$ is the flow generated by $\nabla f$.
It follows that $\inf_{x \in M} R(x,t) \: = \: Ct^{-1}$ for some $C > 0$.
That is, the curvature blows up uniformly as $t \rightarrow 0$.
Comparing this with the
singularity time of the shrinking round cylinder implies that
$\overline{R} = 1$. Performing a similar argument with any
sequence of $x_\alpha$'s tending toward infinity, with the property
that $R(x_\alpha)$ has a limit, shows that
$\lim_{x \rightarrow \infty} R(x) \: = 1$.

Let $N$ denote a (connected component of a)
level surface of $f$. At a point of $N$,
choose an orthonormal frame
$\{e_1, e_2, e_3\}$ with $e_3 \: = \: X$ normal to $N$.
From the Gauss-Codazzi equation,
\begin{equation}
R^N \: = \: 2 \: K^N(e_1, e_2)
\: = \: 2 (K^M(e_1, e_2) \: + \: \det(S)),
\end{equation}
where $S$ is the shape operator. As
$R \: = \: 2 (K^M(e_1, e_2) \: + \: K^M(e_1, e_3) \: + \:
K^M(e_2, e_3))$ and
$\Ric(X, X) \: = \: K^M(e_1, e_3) \: + \:
K^M(e_2, e_3)$, we obtain
\begin{equation}
R^N \: = \: R \: - \: 2 \Ric(X, X) \: + \: 2 \det(S).
\end{equation}
The shape operator is given by
$S \: = \: \frac{Hess f |_{TN}}{|\nabla f|}$.
From (\ref{IIeqn}),
$\Hess f \: = \: \frac{1}{2} \: - \: \Ric$. We can
diagonalize $\Ric \Big|_{TN}$ to write
$\Ric \: = \:
\begin{pmatrix}
r_1 & 0 & c_1 \\
0 & r_2 & c_2 \\
c_1 & c_2 & r_3
\end{pmatrix}$, where $r_3 \: = \: \Ric(X, X)$. Then
\begin{align}
\det \left( \Hess f \Big|_{TN} \right) \: & = \:
\left( \frac{1}{2} \: - \: r_1 \right) \left( \frac{1}{2} \: - \: r_2 \right)
\: = \:
\frac{1}{4} \: \left( (1 - r_1 - r_2)^2 \: - \: (r_1 - r_2)^2 \right)
\\
& \le \: \frac{1}{4} \: (1 - r_1 - r_2)^2
\:  = \:
\frac{1}{4} \: (1 - R + \Ric(X, X))^2. \notag
\end{align}
This shows that the scalar curvature of $N$ is bounded above by
\begin{equation}
R \: - \: 2 \Ric(X, X) \: + \: \frac{(1-R+\Ric(X, X))^2}{2|\nabla f|^2}.
\end{equation}

If $|\nabla f|$ is large then
$1 \: - \: R \: + \: \Ric(X, X) \: < \: 2 \: |\nabla f|^2$.
As $1 \: - \: R \: + \: \Ric(X, X)$ is positive when the
distance from $x$ to $x_0$ is large enough,
\begin{align}
\left( 1 \: - \: R \: + \: \Ric(X, X) \right)^2 \: & < \:
2 (1 \: - \: R \: + \: \Ric(X, X)) |\nabla f|^2 \\
& \le \:
2 (1 \: - \: R \: + \: \Ric(X, X)) |\nabla f|^2 \: + \:
2 |\nabla f|^2 \Ric(X, X) \notag
\end{align}
and so
\begin{equation}
\frac{\left( 1 \: - \: R \: + \: \Ric(X, X) \right)^2}{2|\nabla f|^2}
 \: < \:
1 \: - \: R \: + \: 2 \Ric(X, X).
\end{equation}
Hence
\begin{equation}
R \: - \: 2 \Ric(X, X) \: + \: \frac{(1-R+\Ric(X, X))^2}{2|\nabla f|^2}
\: < \: 1.
\end{equation}
This shows that $R^N < 1$ if $N$ is sufficiently far from $x_0$.

If $Y$ is a unit vector that is tangential to $N$ then from (\ref{IIeqn}),
\begin{equation}
\nabla_Y \nabla_Y f \: = \: \frac12 \: - \: \Ric(Y,Y).
\end{equation}
If $\{Y, Z, W\}$ is an orthonormal basis then
\begin{equation}
\Ric(Y, Y) \: = \: K^M(Y,Z) \: + \: K^M(Y, W) \: \le \:
K^M(Y,Z) \: + \: K^M(Y, W) \: + \: K^M(Z, W) \: = \: \frac12 \: R.
\end{equation}
Hence $\nabla_Y \nabla_Y f \: \ge \: \frac12 \: (1-R)$, which is
positive if $N$ is sufficiently far from $x_0$. Thus $N$ is convex
and so
the area of the level set increases as the level increases.
On the other hand,
we can take points $x_\alpha$ on the level sets going to
infinity, apply the previous splitting argument
and use the fact that $\grad f$ becomes almost parallel to
$\grad d(\cdot, x_0)$. Within one of the approximate cylinders coming from the
splitting argument, there is a projection $\pi$ to its base $S^2$. As
a tangent plane of $N$ is almost perpendicular to $\Ker(d\pi)$,
the restriction of $\pi$ to $N$ is an almost-isometry from
$N$ to $S^2$.
By the monotonicity of $\area(N)$, we conclude that
$\area(N) \: \le \: 8\pi$ if $N$ is sufficiently far from $x_0$.
However as $N$ is a topologically a $2$-sphere,
the Gauss-Bonnet theorem says that
$\int_N R^N \: dA \: = \: 8\pi$. This contradicts the facts that
$R^N \: < 1$ and $\area(N) \: \le \: 8\pi$.
\endgroup\par\medskip

\begin{corollary}[compactness for kappa solutions]
  \label{cor:kl-three-dimensional-noncompact-noncollapsed-gradient-shrinkers-standard-compactness-kappa-solutions}
  \dcref{Corollary 51.22}
  \group{Kleiner--Lott Kappa Solutions}
  \source{kleiner-lott}
 \label{cornosoliton}
The only complete oriented $3$-dimensional noncompact
$\kappa$-noncollapsed gradient shrinking solitons with
bounded nonnegative sectional curvature are the round
evolving $\R \times S^2$ and its $\Z_2$-quotient
$\R \times_{\Z_2} S^2$.
\end{corollary}
\begin{proof}
  \uses{cor:kl-two-dimensional-solutions-kappa-solutions-corollary,lem:kl-three-dimensional-noncompact-noncollapsed-gradient-shrinkers-standard-existence-kappa-solutions,thm:kl-maximum-principles-existence-curvature-operators}
Let $(M, g(\cdot))$ be a complete oriented $3$-dimensional noncompact
$\kappa$-noncollapsed gradient shrinking soliton with
bounded nonnegative sectional curvature. By Lemma \ref{nosoliton}, $M$ cannot
have positive sectional curvature.  From Theorem \ref{splitting},
$M$ must locally split off an $\R$-factor.  Then the universal cover splits
off an $\R$-factor and so, by Corollary \ref{2Dsoliton} or Section
\ref{altpf11.3}, must be the standard $\R \times S^2$. From the
$\kappa$-noncollapsing, $M$ must be $\R \times S^2$ or $\R \times_{\Z_2} S^2$.
\end{proof}

\section{I.12.1. Canonical neighborhood theorem}
\label{I.12.1}

In this section we show that a high-curvature region of a three-dimensional
Ricci flow is modeled by part of a $\kappa$-solution.


We first define the notion of $\Phi$-almost nonnegative curvature.

\begin{definition}[almost nonnegative curvature]
  \label{def:kl-canonical-neighborhood-theorem-almost-nonnegative-curvature}
  \dcref{Definition 52.1}
  \group{Kleiner--Lott Canonical Neighborhoods}
  \source{kleiner-lott}
 (cf. I.12) \label{pinchingdef}
Let $\Phi \in C^\infty(\R)$ be a positive nondecreasing function
such that
for positive $s$, $\frac{\Phi(s)}{s}$ is a decreasing function which
tends to zero as $s \rightarrow \infty$.
A Ricci flow solution is said to have {\em $\Phi$-almost nonnegative
curvature} if for all $(x,t)$, we have
\begin{equation} \label{phieqnn}
\Rm(x,t) \: \ge \: - \: \Phi(R(x,t)).
\end{equation}
\end{definition}

\begin{remark}[existence for canonical neighborhoods]
  \label{rem:kl-canonical-neighborhood-theorem-existence-canonical-neighborhoods}
  \dcref{Remark 52.3}
  \group{Kleiner--Lott Canonical Neighborhoods}
  \source{kleiner-lott}
  \uses{def:kl-canonical-neighborhood-theorem-almost-nonnegative-curvature}

Note that $\Phi$-almost nonnegative curvature implies that the scalar curvature is uniformly
bounded below by $- \: 6 \: \Phi(0)$.
The formulation of the pinching condition in \cite[Section 12]{Perelman}
is that there is a decreasing function $\phi$, tending to zero at infinity,
so that $\Rm(x,t) \: \ge \: - \: \phi(R(x,t)) \: R(x,t)$ for each
$(x,t)$. This formulation has a problem when $R(x,t) < 0$, if one
takes $\phi$ to be defined on all of $\R$. The condition in
Definition \ref{pinchingdef} is what comes out of the
three-dimensional
Hamilton-Ivey pinching result (applied to the rescaled metric
$\tilde{g}(t) = \frac{g(t)}{t}$) if we assume normalized initial
conditions; see Appendix \ref{phiappendix}.
\end{remark}

We note that since the sectional curvatures have to add up to
$R$, the lower bound (\ref{phieqnn}) implies a double-sided bound on the
sectional curvatures.  Namely,
\begin{equation} \label{double}
- \Phi(R) \: \le \:
\Rm \: \le \: \frac{R}{2} \: + \: \left( \frac{n(n-1)}{2} - 1 \right)
\: \Phi(R).
\end{equation}

The main use of the pinching condition is to show that
blowup limits have nonnegative sectional curvature.

\begin{lemma}[convergence for canonical neighborhoods]
  \label{lem:kl-canonical-neighborhood-theorem-convergence-canonical-neighborhoods}
  \dcref{Lemma 52.5}
  \group{Kleiner--Lott Canonical Neighborhoods}
  \source{kleiner-lott}
 \label{almost nonneg}
Let $\{(M_k, p_k, {g}_k)\}_{k=1}^\infty$ be a sequence of
complete pointed Riemannian manifolds with $\Phi$-almost
nonnegative curvature.
Given a sequence $Q_k \rightarrow \infty$, put
$\overline{g}_k = Q_k g_k$ and
suppose that there
is a pointed smooth limit
$(M_\infty, p_\infty, \overline{g}_\infty) = \lim_{k \rightarrow \infty}
(M_k, p_k, \overline{g}_k)$. Then
$M_\infty$ has nonnegative sectional curvature.
\end{lemma}
\begin{proof}
  \uses{def:kl-canonical-neighborhood-theorem-almost-nonnegative-curvature}
First, the $\Phi$-almost
nonnegative curvature condition implies that the scalar
curvature of $M_k$ is bounded below uniformly in $k$.
For $m \in M_\infty$, let $m_k \in (M_k,\bar g_k)$ be
a sequence of approximants to $m$.  Then
$\lim_{k \rightarrow \infty} \overline{\Rm}(m_k) \: = \:
\overline{\Rm}_\infty(m)$, where $\overline{\Rm}(m_k) \: = \:
Q_k^{-1} \: {\Rm}(m_k)$.
There are two possibilities : either the numbers $R(m_k)$ are
uniformly bounded above or they are not.  If they are uniformly
bounded above then (\ref{double}) implies that
$\Rm(m_k)$ is uniformly bounded above and below,
so $\overline{\Rm}_\infty(m)
= 0$. Suppose on the other hand that a subsequence of the numbers
$R(m_k)$ tends to infinity.  We pass to this subsequence. Now
$\overline{R}_\infty(m) = \lim_{k \rightarrow \infty} \overline{R}(m_k)$
exists by assumption and is nonnegative.
Applying (\ref{phieqnn})
gives that $\overline{\Rm}_\infty(m)$, the limit of
\begin{equation}
Q_k^{-1} \: {\Rm}(m_k) \: = \:
\overline{R}(m_k) \: \frac{{\Rm}(m_k)}{R(m_k)},
\end{equation}
is nonnegative.
\end{proof}

We now prove the first version of the ``canonical neighborhood''
theorem.

\begin{theorem}[compactness for canonical neighborhoods]
  \label{thm:kl-canonical-neighborhood-theorem-compactness-canonical-neighborhoods}
  \dcref{Theorem 52.7}
  \group{Kleiner--Lott Canonical Neighborhoods}
  \source{kleiner-lott}
 (cf. Theorem I.12.1) \label{thmI.12.1}
Given $\epsilon, \kappa, \sigma > 0$ and a function $\Phi$ as above, one can
find $r_0 > 0$ with the following property.  Let $g(\cdot)$ be a Ricci flow
solution on a three-manifold $M$, defined for $0 \le t \le T$ with
$T \ge 1$. We suppose that for each $t$, $g(t)$ is complete,
and the sectional curvature is bounded on compact time intervals.
Suppose that the Ricci flow has
$\Phi$-almost nonnegative curvature and is $\kappa$-noncollapsed
on scales less than $\sigma$. Then for any point $(x_0, t_0)$ with $t_0 \ge 1$
and $Q = R(x_0, t_0) \ge r_0^{-2}$, the solution in
$\{(x,t) \: : \: \dist_{t_0}^2(x, x_0) < (\epsilon Q)^{-1}, t_0 - (\epsilon Q)^{-1}
\le t \le t_0\}$ is, after scaling by the factor $Q$, $\epsilon$-close to the
corresponding subset of a $\kappa$-solution.
\end{theorem}

\begin{remark}[existence for canonical neighborhoods]
  \label{rem:kl-canonical-neighborhood-theorem-existence-canonical-neighborhoods-5}
  \dcref{Remark 52.8}
  \group{Kleiner--Lott Canonical Neighborhoods}
  \source{kleiner-lott}
  \uses{rem:kl-canonical-neighborhood-theorem-structural-remark-canonical-neighborhoods,thm:kl-canonical-neighborhood-theorem-compactness-canonical-neighborhoods}
 \label{wangg}
Our statement of Theorem \ref{thmI.12.1} differs slightly from that
in \cite[Theorem 12.1]{Perelman}.
First, we allow $M$ to be noncompact, provided that
there is bounded sectional curvature on compact time intervals.
This generalization will be useful for later work.
More importantly, the statement in \cite[Theorem 12.1]{Perelman}
has noncollapsing at scales less than $r_0$, whereas we require
noncollapsing at scales less than $\sigma$. See Remark \ref{wang} for further
comment.

In the phrase ``$t_0 \ge 1$'' there is an implied scale which comes
from the $\Phi$-almost nonnegativity assumption, and similarly for the
statement  ``scales less than $\sigma$''.
\end{remark}
\par\medskip\noindent\textit{Source argument (not certified as complete).}\ \begingroup\itshape
We give a proof which differs in some points from the proof in
\cite{Perelman} but which has the same ingredients. We first outline
the argument.

Suppose that the theorem is false.  Then for some $\epsilon, \kappa,\sigma>0$,
we have a sequence of such $\Phi$-nonnegatively curved
$3$-dimensional Ricci flows $(M_k,g_k(\cdot))$
defined on intervals $[0,T_k]$,
and sequences $r_k \rightarrow 0$, $\hat{x}_k\in M_k$, $\hat{t}_k\geq 1$ such
that $M_k$ is $\kappa$-noncollapsed on scales $< \sigma$ and
$Q_k = R(\hat{x}_k,\hat{t}_k) \ge r_k^{-2}$, but the $Q_k$-rescaled solution
in $B_{(\epsilon Q_k)^{-1/2}}(\hat{x}_k) \times
[\hat{t}_k - (\epsilon Q_k)^{-1}, \hat{t}_k]$ is not $\epsilon$-close to the
corresponding subset of any $\kappa$-solution.

We note that if the statement were false for $\epsilon$ then it would
also be false for any smaller $\epsilon$. Because of
this, somewhat paradoxically, we will begin the argument with a
given $\epsilon$ but will allow ourselves to make $\epsilon$ small
enough later so that the argument works.
To be clear, we will eventually get a contradiction using a fixed
(small) value of $\epsilon$, but as the proof goes along we will
impose some upper bounds on this value in order for the proof to work.
(If we tried to list all of
the constraints at the beginning of the argument then they would
look unmotivated.)

The goal is to get a contradiction based on the ``bad'' points
$(\widehat{x}_k, \widehat{t}_k)$.
In a sense, the method of proof of Theorem \ref{thmI.12.1} is an induction on
the curvature scale. For example, if we were to make the additional
assumption in the theorem
 that $R(x, t) \le R(x_0, t_0)$ for all $x \in M$ and $t \le t_0$
then the theorem would be very easy to prove. We would just take
a convergent subsequence of the rescaled solutions, based at
$(\hat{x}_k, \hat{t}_k)$, to get a $\kappa$-solution; this would
give a contradiction. This simple argument
can be considered to be the first step
in a proof by induction on curvature scale. In the proof of
Theorem \ref{thmI.12.1} one effectively proves the result at a
given curvature scale inductively by assuming that the result
is true at higher curvature scales.

The actual proof consists of four steps.
Step 1 consists of replacing the sequence $(\widehat{x}_k, \widehat{t}_k)$
by another sequence of ``bad'' points $(x_k, t_k)$ which have the
property that points near
$(x_k, t_k)$ with distinctly higher scalar curvature are ``good'' points.
It then suffices to get a contradiction based on the existence of the
sequence $(x_k, t_k)$.

In steps 2-4 one uses the points $(x_k, t_k)$
to build up a $\kappa$-solution, whose existence then contradicts the
``badness'' of the points $(x_k, t_k)$.
More precisely, let
$(M_k,(x_k,t_k),\bar g_k(\cdot))$ be the result of rescaling $g_k(\cdot)$
by $R(x_k,t_k)$.
We will show that the sequence of pointed
flows $(M_k,(x_k,t_k),\bar g_k(\cdot))$ accumulates on a $\kappa$-solution
$(M_\infty, (x_\infty, t_0), \overline{g}_\infty(\cdot))$,
thereby obtaining a contradiction.

In step 2 one takes a pointed limit of the manifolds
$(M_k, x_k,\bar g_k(t_k))$ in order to construct what
will become the final time slice of the $\kappa$-solution,
$(M_\infty, x_\infty, \overline{g}_\infty(t_0))$. In order to take this
limit, it is necessary to show that the manifolds
$(M_k, x_k,\bar g_k(t_k))$ have uniformly bounded curvature on distance
balls of a fixed radius.  If this were not true then for some radius,
a subsequence of the
manifolds $(M_k, x_k,\bar g_k(t_k))$ would have curvatures that
asymptotically blowup on the ball of that radius.  One shows that
geometrically, the curvature blowup is due to the asymptotic formation of
a cone-like point at the blowup radius.  Doing a further
rescaling at this cone-like point, one obtains a Ricci flow solution
that ends on a part of a nonflat metric cone. This gives a contradiction
as in the case $0 < {\mathcal R} < \infty$ of Theorem \ref{asympscalar}.

Thus one can construct the pointed limit
$(M_\infty, x_\infty, \overline{g}_\infty(t_0))$. The goal now is
to show that $(M_\infty, x_\infty, \overline{g}_\infty(t_0))$ is the
final time slice of a $\kappa$-solution
$(M_\infty, (x_\infty, t_0), \overline{g}_\infty(\cdot))$. In step 3
one shows that $(M_\infty, x_\infty, \overline{g}_\infty(t_0))$ extends
backward
to a Ricci flow solution on some time interval $[t_0 - \Delta, t_0]$,
and that
the time slices have bounded nonnegative curvature.  In step 4 one
shows that the Ricci flow solution can be extended all the way to time
$(-\infty, t_0]$, thereby constructing a $\kappa$-solution

\bigskip
{\em Step 1: Adjusting the choice of basepoints.}

We first modify the points
$(\widehat{x}_k, \widehat{t}_k)$ slightly in time
to points $(x_k, t_k)$
so that in the given Ricci flow solution, there are no other ``bad'' points with
much larger scalar curvature in
a earlier time interval whose length is large compared to
$R(x_k, t_k)^{-1}$.
(The phrase ``nearly the smallest curvature $Q$'' in
\cite[Proof of Theorem 12.1]{Perelman} should read
``nearly the largest curvature $Q$''. This is clear from
the sentence in parentheses that follows.) The proof of
the next lemma is by pointpicking, as in Appendix \ref{pointpicking}.

\begin{lemma}[canonical neighborhoods lemma]
  \label{lem:kl-canonical-neighborhood-theorem-canonical-neighborhoods-lemma}
  \dcref{Lemma 52.9}
  \group{Kleiner--Lott Canonical Neighborhoods}
  \source{kleiner-lott}
  \uses{thm:kl-canonical-neighborhood-theorem-compactness-canonical-neighborhoods}
  \label{ptpicking}
We can find $H_k\ra\infty$, $x_k\in M_k$, and $t_k\geq\frac{1}{2}$
such that $Q_k =  R(x_k,t_k)\ra\infty$, and for all $k$
the conclusion of Theorem \ref{thmI.12.1} fails at $(x_k,t_k)$,
but holds for
any $(y,t)\in M_k\times[t_k-H_kQ_k^{-1},t_k]$
for which $R(y,t)\geq 2R(x_k,t_k)$.
\end{lemma}
\begin{proof}
  \uses{thm:kl-canonical-neighborhood-theorem-compactness-canonical-neighborhoods}
Choose $H_k\ra\infty$ such that
$H_k(R(\widehat{x}_k,\widehat{t}_k))^{-1}\leq \frac{1}{10}$
for all $k$.  For each $k$, initially set $(x_k,t_k) = (\hat x_k,\hat t_k)$.
Put $Q_k = R(x_k, t_k)$ and look for a point in
$M_k\times[t_k-H_kQ_k^{-1},t_k]$ at which
Theorem \ref{thmI.12.1} fails,
and the scalar curvature is at least $2R(x_k,t_k)$.
If such a point
exists, replace $(x_k,t_k)$  by this point; otherwise do nothing.
Repeat this until the second alternative occurs.
This process must terminate with a new choice of $(x_k,t_k)$
satisfying the lemma.
\end{proof}

Hereafter we use this modified sequence $(x_k, t_k)$.
Let $(M_k,(x_k,t_k),\bar g_k(\cdot))$ be the result of rescaling $g_k(\cdot)$
by $Q_k=R(x_k,t_k)$. We use $\bar R_k$ to denote its scalar curvature;
in particular, $\bar R_k (x_k, t_k) = 1$.
Note that the rescaled time interval of Lemma \ref{ptpicking}
has duration $H_k \rightarrow \infty$;
this is what we want in order to try to extract an ancient solution.

\bigskip
{\em Step 2: For every $\rho<\infty$, the scalar curvature $\bar R_k$
is uniformly bounded
on the $\rho$-balls $B(x_k,\rho)\subset (M_k,\bar g_k(t_k))$
(the argument for this is essentially equivalent to
\cite[Pf. of Claim 2 of Theorem I.12.1]{Perelman}).}
Before proceeding, we need some bounds which come from our choice
of basepoints, and the derivative bounds inherited (by approximation)
from $\kappa$-solutions.

\begin{lemma}[existence for canonical neighborhoods]
  \label{lem:kl-canonical-neighborhood-theorem-existence-canonical-neighborhoods}
  \dcref{Lemma 52.10}
  \group{Kleiner--Lott Canonical Neighborhoods}
  \source{kleiner-lott}
 \label{grad}
There is a constant $C = C(\kappa)$ so that for any $(x, t)$ in a
Ricci flow solution, if $R(x, t) > 0$ and the solution in
$B_t(x, (\epsilon R(x, t))^{-1/2}) \times [t-(\epsilon R(x, t))^{-1}, t]$
is $\epsilon$-close to a corresponding subset of a $\kappa$-solution
then $|\nabla R^{-1/2}|(x, t) \: \le \: C$ and
$|\partial_t R^{-1}|(x,t) \: \le \: C$.
\end{lemma}
\begin{proof}
  \uses{thm:kl-compactness-space-three-dimensional-solutions-compactness-kappa-solutions}
This follows from the compactness in Theorem \ref{thmI.11.7}.
\end{proof}

Note that the same value of $C$ in Lemma \ref{grad}
also works for smaller $\epsilon$.

\begin{lemma}[canonical neighborhoods lemma]
  \label{lem:kl-canonical-neighborhood-theorem-canonical-neighborhoods-lemma-8}
  \dcref{Lemma 52.11}
  \group{Kleiner--Lott Canonical Neighborhoods}
  \source{kleiner-lott}
 (cf. Claim 1 of I.12.1)
\label{claim1} For each $(\overline{x}, \overline{t})$ with
$t_k \: - \: \frac12 \: H_k \: Q_k^{-1} \le \overline{t} \le t_k$, we have
$R_k(x, t) \le 4 \overline{Q}_k$ whenever
$\overline{t} \: - \: c \: \overline{Q}_k^{-1} \le t \le \overline{t}$ and
$\dist_{\overline{t}}(x, \overline{x}) \: \le \: c \: \overline{Q}_k^{-1/2}$,
where $\overline{Q}_k = Q_k + |R_k(\overline{x}, \overline{t})|$ and
$c = c(\kappa) > 0$ is a small constant.
\end{lemma}
\begin{proof}
  \uses{lem:kl-canonical-neighborhood-theorem-canonical-neighborhoods-lemma,lem:kl-canonical-neighborhood-theorem-existence-canonical-neighborhoods}
If $R_k(x, t) \le 2Q_k$ then there is nothing to show.  If
$R_k(x, t) > 2Q_k$, consider a spacetime curve $\gamma$ that goes linearly from
$(x,t)$ to $(x, \overline{t})$, and then goes from
$(x, \overline{t})$ to $(\overline{x}, \overline{t})$ along a
minimizing geodesic.  If there is a point on $\gamma$ with curvature
$2Q_k$, let $p$ be the nearest such point to $(x, t)$. If not, put
$p = (\overline{x}, \overline{t})$. From the conclusion of Lemma
\ref{ptpicking}, we can apply
Lemma \ref{grad} along $\gamma$ from $(x, t)$ to $p$.
The claim follows from integrating the ensuing
derivative bounds along $\gamma$.
\end{proof}

\begin{lemma}[canonical neighborhoods lemma]
  \label{lem:kl-canonical-neighborhood-theorem-canonical-neighborhoods-lemma-9}
  \dcref{Lemma 52.12}
  \group{Kleiner--Lott Canonical Neighborhoods}
  \source{kleiner-lott}
 \label{rescaled}
In terms of the rescaled solution $\overline{g}_k(\cdot)$,
for each $(\overline{x}, \overline{t})$ with
$t_k \: - \: \frac12 \: H_k \le \overline{t} \le t_k$, we have
$\overline{R}_k(x, t) \le 4 \widetilde{Q}_k$ whenever
$\overline{t} \: - \: c \: \widetilde{Q}_k^{-1} \le t \le \overline{t}$ and
$\dist_{\overline{t}}(x, \overline{x}) \: \le \: c \: \widetilde{Q}_k^{-1/2}$,
where $\widetilde{Q}_k = 1 + |\overline{R}_k(\overline{x}, \overline{t})|$.
\end{lemma}
\begin{proof}
  \uses{lem:kl-canonical-neighborhood-theorem-canonical-neighborhoods-lemma-8}
This is just the rescaled version of Lemma \ref{claim1}.
\end{proof}

For all $\rho\geq 0$, put
\begin{equation}
D(\rho) = \sup\{\bar R_k(x,t_k)\mid k\geq 1,\;x\in B(x_k,\rho)\subset
(M_k,\bar g_k(t_k))\},
\end{equation}
and let $\rho_0$ be the supremum of the $\rho$'s for which
$D(\rho)<\infty$.  Note that $\rho_0>0$, in view of Lemma \ref{rescaled}
(taking $(\bar x,\bar t)=(x_k,t_k)$).
Suppose that $\rho_0<\infty$.
After passing to a subsequence if necessary, we can find
a sequence $y_k\in M_k$ with $\dist_{t_k}(x_k,y_k) \rightarrow \rho_0$
and $\bar R(y_k,t_k)\ra\infty$.    Let $\eta_k\subset (M_k,\bar g_k(t_k))$
be a minimizing geodesic segment from $x_k$ to $y_k$.
Let $z_k\in \eta_k$ be the point on $\eta_k$ closest to $y_k$
at which $\bar R(z_k,t_k)=2$, and let $\ga_k$
be the subsegment of $\eta_k$ running from $y_k$ to $z_k$.  By
Lemma \ref{rescaled} the length of $\ga_k$ is bounded away from zero
independent of $k$.  Due to the $\Phi$-pinching
(see (\ref{double})),
for all $\rho<\rho_0$, we have a uniform bound on $|\Rm|$
on the balls $B(x_k,\rho)\subset (M_k,\bar g_k(t_k))$.  The
injectivity radius is also controlled in $B(x_k,\rho)$,
in view of the curvature bounds and the $\kappa$-noncollapsing.
Therefore after passing
to a subsequence, we can assume that the pointed sequence
$(B(x_k,\rho_0),\bar g_k(t_k),x_k)$ converges in the pointed
Gromov-Hausdorff topology (i.e. for all $\rho<\rho_0$ we have the
usual Gromov-Hausdorff convergence) to a pointed $C^1$-Riemannian manifold
$(Z,\bar g_\infty,x_\infty)$,  the segments $\eta_k$ converge to a  segment (missing
an endpoint)
$\eta_\infty\subset Z$ emanating from $x_\infty$, and $\gamma_k$
converges to  $\ga_\infty\subset\eta_\infty$.  Let $\bar Z$
denote the completion of $(Z,\bar g_\infty)$, and $y_\infty\in \bar Z$
the limit point of $\eta_\infty$.    Note that by
Lemma \ref{ptpicking} and
part 4 of Corollary \ref{11.4cors}, the Riemannian structure near
$\ga_\infty$ may be chosen to be many times differentiable.
(Alternatively, this follows from Lemma \ref{rescaled} and the Shi
estimates of Appendix \ref{applocalder}.)
In particular
the scalar curvature $\bar R_\infty$ is defined, differentiable,
and satisfies the bound in Lemma \ref{grad} near $\ga_\infty$.

\begin{lemma}[existence for canonical neighborhoods]
  \label{lem:kl-canonical-neighborhood-theorem-existence-canonical-neighborhoods-10}
  \dcref{Lemma 52.14}
  \group{Kleiner--Lott Canonical Neighborhoods}
  \source{kleiner-lott}

\label{nearlyround}
1. There is a function $c:(0,\infty)\ra \R$ depending only
on $\kappa$, with
$\lim_{t\ra 0}c(t)=\infty$, such that if $w\in\ga_\infty$ then
$\bar R_\infty(w) \: d(y_\infty,w)^2 \: > \: c(\eps)$.

2. There is a function $\eps':(0,\infty)\ra \R\cup \{\infty\}$ depending only
on $\kappa$, with $\lim_{t\ra 0}\eps'(t)=0$, such that if
$w\in \ga_\infty$ and $d(y_\infty,w)$ is sufficiently small
then the pointed manifold $(Z,w, \bar R_\infty(w)\bar g_\infty)$ is
$2\eps^\prime(\epsilon)$-close to a round cylinder in the $C^2$ topology.
\end{lemma}
\begin{proof}
  \uses{lem:kl-canonical-neighborhood-theorem-canonical-neighborhoods-lemma,lem:kl-canonical-neighborhood-theorem-existence-canonical-neighborhoods,prop:kl-more-properties-solutions-existence-kappa-solutions}
It follows from Lemma \ref{ptpicking} that
for all $w\in \ga_\infty$, the pointed Riemannian
manifold
$(Z,w,\bar R_\infty(w)\bar g_\infty)$ is $2\eps$-close
to (a time slice of) a pointed $\kappa$-solution.
From the definition of pointed closeness, there is an
embedded region around $w$, large on the scale defined by
$\bar R_\infty(w)$, which is close to the corresponding subset of
a pointed $\kappa$-solution.  This gives a lower bound on
the distance $\rho_0-d(w,x_\infty)$ to the point of curvature
blowup, thereby proving part 1 of the lemma.

We know that $(Z,w, \bar R_\infty(w) \bar g_\infty)$ is $2\eps$-close to
a pointed $\kappa$-solution $(N,\star,h(t))$ in the pointed
$C^2$-topology.   By Lemma \ref{grad},
we know  $\bar R_\infty(w)$ tends to $\infty$
as $d(w,x_\infty)\ra\rho_0$, for $w\in\ga_\infty$.
In particular, $\bar R_\infty(w)d^2(w,x_\infty)\ra \infty$.
From part 1 above, we can choose $\epsilon$ small enough in order to
make
$\bar R_\infty(w)d^2(w,y_\infty)$ large enough to apply Proposition
\ref{neckcrit}. Hence the pointed
manifold $(N,\star,R(\star)h(t))$ is $\eps^{\prime \prime}(\eps)$-close to
a round cylinder, where $\eps^{\prime \prime}:(0,\infty)\ra\R$ is a function with
$\lim_{t\ra 0}\eps^{\prime \prime}(t)=0$. The lemma follows.
\end{proof}


Note that the function $\epsilon^\prime$ in Lemma
\ref{nearlyround} is independent of the particular manifold
$Z$ that arises in our proof.

From part 2 of Lemma \ref{nearlyround}, if
$\eps$ is small and $ w \in \ga_\infty$ is
sufficiently close to $y_\infty$ then
$(Z, w, \overline{R}_\infty(w) \overline{g}_\infty)$ is
$C^2$-close to a round cylinder.  The cross-section
of the cylinder has diameter approximately
$\pi \: (\overline{R}_\infty(w)/2)^{- \: \frac12}$.
If we form the
union of the balls $B(w,2\pi(R_\infty(w)/2)^{- \: \frac{1}{2}})$, as $w$
ranges over such points in $\ga_\infty$, then we obtain a
connected Riemannian manifold $W$.
By adding in the point $y_\infty$, we get a
metric space $\bar W$ which is locally complete, and geodesic
near $y_\infty$. As the original manifolds $M_k$ had
$\Phi$-almost nonnegative curvature, it follows
from Lemma \ref{almost nonneg}
that $W$ is
nonnegatively curved.  Furthermore,
$y_\infty$ cannot be an interior point of any geodesic segment in
$\bar W$, since such a geodesic would have to pass through a cylindrical
region near $y_\infty$
twice.  The usual
proof of the Toponogov triangle comparison inequality now applies near
$y_\infty$ since
minimizers remain in the smooth nonnegatively curved part of $\bar W$.
Then $\overline{W}$ has nonnegative curvature in the Alexandrov sense.


This implies that blowups of $(\bar W,y_\infty)$ converge
to the tangent cone $C_{y_\infty}\bar W$.
As $\overline{W}$ is three-dimensional, so is
$C_{y_\infty}\bar W$
\cite[Corollary 7.11]{Burago-Gromov-Perelman}.
It will be
$C^2$-smooth away from the vertex and nowhere flat, by
part 2 of Lemma \ref{nearlyround}.
Pick $z\in C_{y_\infty}\bar W $ such that $d(z,y_\infty)=1$.
Then
for any $\delta < \frac12$,
the ball $B(z,\delta)\subset C_{y_\infty}\bar W $
is the Gromov-Hausdorff limit of a sequence of
rescaled balls
$B(\widehat{z}_k,\widehat{r}_k)\subset (M_k,\bar g_k(t_k))$
where $\widehat{r}_k\ra 0$, whose center points $(\widehat{z}_k,t_k)$
satisfy the conclusions of Theorem
\ref{thmI.12.1}.  Applying Lemma \ref{rescaled}
and Appendix \ref{phiappendix},
if $\delta = \delta(\kappa)$ is taken small enough then
we get the curvature bounds needed to extract a limiting
Ricci flow solution whose time zero slice is isometric
to $B(z,\delta)$.
Now we can apply
the reasoning from the $0<\r<\infty$ case of Theorem
\ref{asympscalar} to get a contradiction.  This completes step 2.

\bigskip
{\em Step 3: The  sequence of pointed
flows $(M_k,\bar g_k(\cdot),(x_k,t_k))$ accumulates on a
pointed Ricci flow $(M_\infty,\bar g_\infty(\cdot),(x_\infty,t_0))$
which is defined on a time interval $[t',t_0]$ with $t'< t_0$.}
By step 2, we know that the scalar curvature   of $(M_k,\bar g_k(t_k))$
at $y\in M_k$ is bounded by a function of  the distance from
$y$ to $x_k$.
Lemma \ref{rescaled} extends this curvature control to a
backward parabolic neighborhood centered at $y$ whose radius depends only
on $d(y,x_k)$.  Thus we can conclude,
using $\Phi$-pinching (\ref{double}) and
Shi's estimates (Appendix \ref{applocalder}),
that all derivatives of the
curvature $(M_k,\bar g_k(t_k))$ are controlled as a function
of the distance from $x_k$, which means that the
sequence of pointed manifolds $(M_k,\bar g_k(t_k),x_k)$ accumulates
to a smooth manifold $(M_\infty,\bar g_\infty)$.

From Lemma \ref{almost nonneg},
$M_\infty$ has nonnegative sectional curvature.
We claim that $M_\infty$ has bounded curvature.
If not then there is a sequence of points
$q_k \in M_\infty$ so that
$\lim_{k \rightarrow \infty} \overline{R}(q_k) \: = \: \infty$
and $\overline{R}(q) \le 2 \overline{R}(q_k)$ for
$q \in B(y_k, A_k \overline{R}(q_k)^{- \frac12})$, where
$A_k \rightarrow \infty$; compare \cite[Lemma 22.2]{Hamilton}.
Lemma \ref{ptpicking} implies that for large $k$, a rescaled neighborhood of
$(M_\infty, q_k)$ is $\epsilon$-close to the corresponding subset
of a time slice of a $\kappa$-solution. As in the proof of
Theorem \ref{thmI.11.7}, we obtain a sequence of neck-like regions
in $M_\infty$ with smaller
and smaller cross-sections, which contradicts the
existence of the Sharafutdinov retraction.

By Lemma \ref{rescaled} again,
we now get curvature control
 on $(M_k,\bar g_k(\cdot))$ for a time
interval $[t_k-\De,t_k]$ for some $\De>0$, and hence
we can extract a  subsequence which converges to a pointed Ricci flow
$(M_\infty,(x_\infty,t_0),\bar g_\infty(\cdot))$ defined
for $t\in[t_0 - \De,t_0]$, which has nonnegative curvature and bounded
curvature on
compact time intervals.


\bigskip
{\em Step 4: Getting an ancient solution.}
Let $(t',t_0]$ be the maximal time interval on which
we can extract a limiting solution $(M_\infty,\bar g_\infty(\cdot))$
with bounded curvature on compact time intervals.
Suppose that $t' > -\infty$.
By Lemma \ref{rescaled} the maximum of the scalar curvature
on the time slice $(M_\infty,\bar g_\infty(t))$ must tend
to infinity as $t\ra t'$.  From the trace Harnack inequality,
$R_t \: + \: \frac{R}{t - t^\prime} \: \ge \: 0$, and so
\begin{equation}
\bar R_\infty(x,t)\leq Q \frac{t_0 - t^\prime}{t-t'},
\end{equation}
where $Q$ is the maximum of the scalar curvature
on $(M_\infty,\bar g_\infty(t_0))$.  Combining this
with Corollary \ref{bound1},
we get
\begin{equation}
\frac{d}{dt}d_t(x,y)\geq \const \: \sqrt{Q \frac{t_0 - t^\prime}{t-t'}}.
\end{equation}
Since the right hand side is integrable on $(t',t_0]$, and using
the fact that distances are nonincreasing in time (since $\Rm \geq 0$),
it follows that there is a constant $C$ such that
\begin{equation}
\label{addbd}
|d_t(x,y)-d_{t_0}(x,y)|<C
\end{equation}
for all $x,y\in M_\infty$, $t\in(t',t_0]$.

If $M_\infty$ is compact then by (\ref{addbd}) the diameter of
$(M_\infty,\bar g_\infty(t))$ is bounded independent of
$t\in (t',t_0]$.  Since the minimum of the scalar
curvature is increasing in
time, it is also bounded independent of $t$.  Now the argument
in Step 2 shows that the curvature is bounded everywhere
independent of $t$.  (We can apply the argument of Step 2 to the
time-$t$ slice because the main ingredient was
Lemma \ref{ptpicking}, which holds for rescaled time $t$.)

We may therefore assume $M_\infty$ is noncompact.
To be consistent with the notation of I.12.1, we now relabel
the basepoint $(x_\infty, t_0)$ as $(x_0, t_0)$.  Since
nonnegatively curved manifolds are asymptotically conical
(see Appendix \ref{Alexandrov}),
there is a constant $D$ such that if $y\in M_\infty$,
and $d_{t_0}(y,x_0)>D$, then there is a point $x\in M_\infty$
such that
\begin{equation}
\label{dists}
d_{t_0}(x,y)=d_{t_0}(y,x_0)\quad \mbox{and}
\quad d_{t_0}(x,x_0)\geq \frac{3}{2}d_{t_0}(y,x_0);
\end{equation}
by (\ref{addbd})
the same conditions hold at all times $t\in(t',t_0]$, up to error
$C$.   If for some such $y$, and some $t\in (t',t_0]$ the scalar
curvature were large, then $(M_\infty,(y,t),\bar R_\infty(y,t)\bar
g_\infty(t))$
would be $2\eps$-close to a $\kappa$-solution $(N,h(\cdot),(z,t_0))$.  When
$\eps$ is small we could  use Proposition \ref{neckcrit} to
see that $y$ lies in a neck region $U$ in $(M_\infty,\bar g_\infty(t))$
of diameter $\approx R(y,t)^{-\frac{1}{2}}\ll 1$.

We claim that $U$ separates
$x_0$ from $x$ in the sense that $x_0$ and $x$ belong
to disjoint components of $M_\infty - U$, where $x_0, y$, and $x$
satisfy (\ref{dists}).
To see this, we
recall that if $M_\infty$ has more than one end then it splits
isometrically, in which case the claim is clear. If $M_\infty$ has
one end then we consider an exhaustion of $M_\infty$ by totally
convex compact sets $C_t$ as in Section \ref{pf11.7}. One of the
sets $C_t$ will have boundary consisting of an approximate $2$-sphere
cross-section in the neck region $U$, giving the separation.

(For another argument, suppose that $U$ does not separate $x_0$ from
$x$. Since $\cangle_y(x_0, x) > \frac{\pi}{2}$ (from Proposition \ref{neckcrit}),
the segments $\overline{yx_0}$ and $\overline{yx}$ must exit $U$ by
different ends.  If $x_0$ and $x$ can be joined by a curve avoiding $U$
then there is a nonzero element of
$\HH^1(M_\infty, \Z)$. The corresponding infinite cyclic cover of $M_\infty$ will then
isometrically split off a line and its quotient $M_\infty$ will be
compact, which is a contradiction.  We thank one of the referees for
this argument.)

Obviously, at time $t_0$ the set
$U$ still separates $x_0$ from $x$.
Since $\bar g_\infty$ has nonnegative curvature,
we have $\diam_{t_0}(U)\leq \diam_t(U)\ll 1$.
Since $(M_\infty,\bar g_\infty({t_0}))$ has bounded geometry,
there cannot be topologically separating subsets of arbitrarily small
diameter.
Thus there must be a uniform upper bound on $R(y, t)$
and the curvature of $(M_\infty,\bar g_\infty)$ is uniformly bounded
(in space and time)
outside a set of uniformly bounded diameter.  Repeating the reasoning
from Step 2, we get uniform bounds everywhere.  This contradicts
our assumption that the curvature blows up as $t\ra t'$.

It remains to show that the ancient solution is a $\kappa$-solution.
The only remaining point is to show that it is $\kappa$-noncollapsed
at all scales. This follows from the fact that the original Ricci
flow solutions $(M_k, g_k(\cdot))$ were $\kappa$-noncollapsed on
scales less than the fixed number $\sigma$.
\endgroup\par\medskip

\begin{remark}[structural remark on canonical neighborhoods]
  \label{rem:kl-canonical-neighborhood-theorem-structural-remark-canonical-neighborhoods}
  \dcref{Remark 52.19}
  \group{Kleiner--Lott Canonical Neighborhoods}
  \source{kleiner-lott}
  \uses{rem:kl-canonical-neighborhood-theorem-existence-canonical-neighborhoods-5,thm:kl-canonical-neighborhood-theorem-compactness-canonical-neighborhoods,thm:kl-no-local-collapsing-theorem-ii-existence-kappa-solutions}
 \label{wang}
As mentioned in Remark \ref{wangg}, the statement of
\cite[Theorem 12.1]{Perelman} instead assumes noncollapsing at all
scales less than $r_0$. Bing Wang pointed out that with this
assumption, after constructing the ancient solution in Step 4
of the proof, one only gets that it is $\kappa$-collapsed at
all scales less than one. Hence it may not be a $\kappa$-solution.
The literal statement of \cite[Theorem 12.1]{Perelman} is not
used in the rest of \cite{Perelman,Perelman2}, but rather its
method of proof.  Because of this, the change of hypotheses
does not seem to lead to any problems. The method of proof
of Theorem \ref{thmI.12.1} is used in two different ways.
The first way is to construct the Ricci flow with surgery on a
fixed finite time interval, as in Section \ref{II.5}.
In this case the
noncollapsing at a given scale $\sigma$ comes from
Theorem \ref{nolocalcollapse}, and its extension when surgeries
are allowed.  The second way is to analyze the large-time
behavior of the Ricci flow, as in the next few sections.
\end{remark}

\section{I.12.2. Later scalar curvature bounds on bigger balls
from curvature and volume bounds}
\label{I.12.2}

The next theorem roughly says that if one has a sectional curvature bound on a
ball, for a certain time interval, and
a lower bound on the volume of the ball at the initial time, then one obtains
an upper scalar curvature bound on a larger ball at the final time.

We first write out the corrected version of the theorem
(see II.6.2).
\begin{theorem}[positivity for scalar curvature]
  \label{thm:kl-later-scalar-curvature-bounds-bigger-balls-from-curvature-volume-bounds-positivity-scalar-curvature}
  \dcref{Theorem 53.1}
  \group{Kleiner--Lott Canonical Neighborhoods}
  \source{kleiner-lott}
 \label{thmI.12.2}
For any $A < \infty$, there exist $K = K(A) < \infty$ and
$\rho = \rho(A) >  0$ with the following property. Suppose in
dimension three we have a Ricci flow solution with
$\Phi$-almost nonnegative curvature.
Given $x_0 \in M$ and $r_0 > 0$, suppose that
$r_0^2 \: \Phi(r_0^{-2}) \: < \: \rho$,
the solution is defined for
$0 \le t \le r_0^2$ and it has
$|\Rm|(x, t) \le \frac{1}{3r_0^{2}}$
for all $(x, t)$ satisfying $\dist_0(x, x_0) < r_0$. Suppose in
addition that the volume of the metric ball $B(x_0, r_0)$ at
time zero is at least $A^{-1} r_0^3$.
Then $R(x, r_0^2) \le K r_0^{-2}$ whenever $\dist_{r_0^2}(x, x_0) <
A r_0$.
\end{theorem}

\begin{remark}[existence for noncollapsing]
  \label{rem:kl-later-scalar-curvature-bounds-bigger-balls-from-curvature-volume-bounds-existence-noncollapsing}
  \dcref{Remark 53.2}
  \group{Kleiner--Lott Canonical Neighborhoods}
  \source{kleiner-lott}
  \uses{prop:kl-proof-double-sided-curvature-bound-thick-part-modulo-two-propositions-existence-canonical-neighborhoods,rem:kl-no-local-collapsing-propagates-forward-time-larger-scales-estimates-no-local-collapsing-propagates-forward-time-larger-scales,thm:kl-no-local-collapsing-propagates-forward-time-larger-scales-existence-kappa-solutions}
 \label{rmkI.12.2}
The added restriction that $r_0^2 \: \Phi(r_0^{-2}) \: < \: \rho$ (see II.6.2)
imposes an upper bound on $r_0$.  This is necessary, as
otherwise the conclusion would imply that neck pinches cannot occur.

There is an apparent gap in the proof of \cite[Theorem 12.2]{Perelman},
in the sentence
(There is a little subtlety...).  We instead follow the proof
of II.6.3(b,c) (see Proposition \ref{propII.6.3}(b,c)),
which proves the same statement in the
presence of surgeries.


The volume assumption in the theorem is used to guarantee
noncollapsing, by means of Theorem \ref{8.2thm}. The reason
for the ``$3$'' in the hypothesis $|\Rm|(x, t) \le \frac{1}{3r_0^{2}}$
comes from Remark \ref{8.2fix}.
\end{remark}
\begin{proof}
  \proves{thm:kl-later-scalar-curvature-bounds-bigger-balls-from-curvature-volume-bounds-positivity-scalar-curvature}
  \uses{lem:kl-canonical-neighborhood-theorem-canonical-neighborhoods-lemma-9,lem:kl-length-distortion-estimates-distance-distortion-lemma,thm:kl-canonical-neighborhood-theorem-compactness-canonical-neighborhoods,thm:kl-no-local-collapsing-propagates-forward-time-larger-scales-existence-kappa-solutions}
The proof is in two steps.  In the first step one shows that
if $R(x, r_0^2)$ is large then a parabolic neighborhood of
$(x, r_0^2)$ is close to the corresponding subset of a
$\kappa$-solution.  In the second part one uses this to
prove the theorem.

The first step is the following lemma.

\end{proof}

\begin{lemma}[existence for kappa solutions]
  \label{lem:kl-later-scalar-curvature-bounds-bigger-balls-from-curvature-volume-bounds-existence-kappa-solutions}
  \dcref{Lemma 53.3}
  \group{Kleiner--Lott Canonical Neighborhoods}
  \source{kleiner-lott}
 \label{12.2lemma}
For any $\epsilon > 0$ there exists
$K = K(A, \epsilon) < \infty$
so that for any $r_0$,
whenever we have a solution as in the statement of the theorem
and $\dist_{r_0^2}(x, x_0) <
A r_0$ then  \\
(a) $R(x, r_0^2) < K r_0^{-2}$ or \\
(b) The solution in
$\{(x^\prime, t^\prime) \: : \: \dist_{t}(x^\prime, x) <
(\epsilon Q)^{-1}, t - (\epsilon Q)^{-1} \le t^\prime \le t \}$ is,
after scaling by the factor $Q$, $\epsilon$-close to the
corresponding subset of a $\kappa$-solution.

Here $t = r_0^2$ and
$Q = R(x,t)$.
\end{lemma}
\begin{remark}[structural remark on thick--thin geometry]
  \label{rem:kl-later-scalar-curvature-bounds-bigger-balls-from-curvature-volume-bounds-structural-remark-thick-thin-geometry}
  \dcref{Remark 53.4}
  \group{Kleiner--Lott Canonical Neighborhoods}
  \source{kleiner-lott}
  \uses{thm:kl-canonical-neighborhood-theorem-compactness-canonical-neighborhoods}

One can think of this lemma as a
localized analog of Theorem \ref{thmI.12.1}, where ``localized'' refers to
the fact that both the hypotheses and the conclusion involve the
point $x_0$.
\end{remark}

\begin{proof}
  \proves{lem:kl-later-scalar-curvature-bounds-bigger-balls-from-curvature-volume-bounds-existence-kappa-solutions}
  \uses{lem:kl-canonical-neighborhood-theorem-canonical-neighborhoods-lemma-9,lem:kl-length-distortion-estimates-distance-distortion-lemma,thm:kl-canonical-neighborhood-theorem-compactness-canonical-neighborhoods,thm:kl-no-local-collapsing-propagates-forward-time-larger-scales-existence-kappa-solutions}
To prove the lemma,
suppose that there is
a sequence of such pointed solutions
$(M_k, x_{0,k}, g_k(\cdot))$, along with points
$\hat{x}_k \in M_k$, so that $\dist_{r_0^2}(\hat{x}_k, x_{0,k}) <
A r_0$ and $r_0^2 \: R(\hat{x}_k, r_0^2) \rightarrow \infty$, but
$(\hat{x}_k, r_0^2)$ does not satisfy conclusion (b) of the
lemma. As in the proof of Theorem \ref{thmI.12.1}, we will allow ourselves to make
$\epsilon$ smaller during the course of the proof.

We first show that there is a sequence $D_k \rightarrow \infty$
and modified points $(\overline{x}_k, \overline{t}_k)$
with $\frac{3}{4} r_0^{2} \le \overline{t}_k \le r_0^{2}$,
$\dist_{\overline{t}_k}(\overline{x}_k, x_{0,k}) < (A+1) r_0$ and
$Q_k = R(\overline{x}_k, \overline{t}_k) \rightarrow \infty$,
so that any point
$(x^\prime_k, t^\prime_k)$ with $R(x^\prime_k, t^\prime_k) > 2Q_k$,
$\overline{t}_k - D_k^2 Q_k^{-1} \le
t^\prime_k \le \overline{t}_k$ and
$\dist_{t^\prime_k}(x^\prime_k, {x}_{0,k})
< \dist_{\overline{t}_k}(\overline{x}_k,
x_{0,k}) +
D_k Q_k^{-1/2}$ satisfies conclusion (b) of the lemma,
but $(\overline{x}_k, \overline{t}_k)$ does not satisfy conclusion (b) of
the lemma. (Of course, in saying ``$(x^\prime_k, t^\prime_k)$
satisfies conclusion (b)'' or
``$(\overline{x}_k, \overline{t}_k)$ does not satisfy conclusion (b)'', we mean that
the $(x,t)$ in conclusion (b) is replaced by $(x^\prime_k, t^\prime_k)$ or
$(\overline{x}_k, \overline{t}_k)$, respectively.)

The construction of $(\overline{x}_k, \overline{t}_k)$ is by a pointpicking
argument.
Put $D_k = \frac{r_0 R(\hat{x}_k, r_0^2)^{1/2}}{10}$.
Start with $(x_k, t_k) = (\hat{x}_k, r_0^2)$
and look if there is a point
$(x^\prime_k, t^\prime_k)$ with $R(x^\prime_k, t^\prime_k) > 2
R(x_k, t_k)$,
$\overline{t}_k - D_k^2 R(x_k, t_k)^{-1} \le
t^\prime_k \le \overline{t}_k$ and
$\dist_{t^\prime_k}(x^\prime_k, {x}_{0,k})
< \dist_{{t}_k}({x}_k,
x_{0,k}) +
D_k R(x_k, t_k)^{-1/2}$, but
which does not have a neighborhood that is
$\epsilon$-close to the corresponding subset of a $\kappa$-solution.
If there is such a point, we replace $(x_k, t_k)$ by $(x^\prime_k,
t^\prime_k)$ and repeat the process.  The process must terminate after
a finite number of steps to give a point
$(\overline{x}_k, \overline{t}_k)$ with the desired property.

(Note that the condition
$\dist_{t^\prime_k}(x^\prime_k, {x}_{0,k})
< \dist_{\overline{t}_k}(\overline{x}_k,
x_{0,k}) +
D_k Q_k^{-1/2}$ involves the metric at time $t^\prime_k$. In order
to construct an ancient solution, one of the issues will be to
replace this by a condition that only involves the metric at time
$\overline{t}_k$, i.e. that involves a parabolic neighborhood
around $(\overline{x}_k, \overline{t}_k)$.)

Let $\overline{g}_k(\cdot)$ denote
the rescaling of the solution $g_k(\cdot)$ by $Q_k$. We
normalize the time interval of the rescaled solution
by fixing a number ${t}_\infty$
and saying that for all $k$,
the time-$\overline{t}_k$ slice of $(M_k, g_k)$ corresponds to the
time-${t}_\infty$ slice of $(M_k, \overline{g}_k)$.
Then the scalar
curvature $\overline{R}_k$ of $\overline{g}_k$ satisfies
$\overline{R}_k(\overline{x}_k, {t}_\infty) \: = \: 1$.

By the argument of Step 2 of the proof of Theorem \ref{thmI.12.1},
a subsequence of the
pointed spaces $(M_k, \overline{x}_k, \overline{g}_k({t}_\infty))$
will smoothly converge to a nonnegatively-curved pointed space $(M_\infty,
\overline{x}_\infty,\overline{g}_\infty)$.
By the pointpicking, if $m \in M_\infty$ has $\overline{R}(m) \ge 3$
then a parabolic neighborhood of $m$ is $\epsilon$-close
to the corresponding region in a $\kappa$-solution.
It follows,
as in Step 3 of the proof of Theorem
\ref{thmI.12.1}, that the sectional curvature of $M_\infty$ will be bounded above by
some $C < \infty$. Using Lemma \ref{rescaled}, the metric on
$M_\infty$ is the time-$t_\infty$ slice of
a nonnegatively-curved Ricci flow solution defined on some time interval
$[t_\infty - c , t_\infty]$, with $c > 0$,
and one has convergence of a subsequence $\overline{g}_k(t)
\rightarrow \overline{g}_\infty(t)$ for $t \in [t_\infty - c , t_\infty]$.
As $\overline{R}_t \ge 0$, the scalar curvature on this time interval
will be uniformly bounded above by $6C$ and so from  the
$\Phi$-almost nonnegative curvature (see (\ref{double})), the
sectional curvature will be uniformly bounded above on the time interval.
Hence we can apply Lemma \ref{distancedist} to get a uniform additive bound on the
length
distortion between times $t_\infty - c$ and $t_\infty$ (see Step 4 of the
proof of Theorem \ref{thmI.12.1}). More precisely, in applying Lemma
\ref{distancedist}, we
use the curvature bound coming from the hypotheses of the theorem near
$x_0$, and the just-derived upper curvature bound near $\overline{x}_k$.

It follows that for a given $A^\prime > 0$, for large
$k$, if $t_k^\prime \in [\overline{t}_k - c Q_k^{-1}/2, \overline{t}_k]$ and
$\dist_{\overline{t}_k}(x_k^\prime, \overline{x}_k) < A^\prime Q_k^{-1/2}$
then $\dist_{t^\prime_k}(x^\prime_k, {x}_{0,k})
< \dist_{\overline{t}_k}(\overline{x}_k,
x_{0,k}) +
D_k Q_k^{-1/2}$. In particular, if a point
$(x_k^\prime, t_k^\prime)$ lies in the
parabolic neighborhood given by
$t_k^\prime \in [\overline{t}_k - c Q_k^{-1}/2, \overline{t}_k]$ and
$\dist_{\overline{t}_k}(x_k^\prime, \overline{x}_k) < A^\prime Q_k^{-1/2}$,
and  has
$R(x_k^\prime, t_k^\prime) > 2Q_k$, then it has a neighborhood that
is $\epsilon$-close to the corresponding subset of a $\kappa$-solution.

As in Step 4 of the proof of Theorem \ref{thmI.12.1}, we now extend
$(M_\infty,
\overline{g}_\infty, \overline{x}_\infty)$ backward to an ancient solution
$\overline{g}_\infty(\cdot)$, defined for $t \in (-\infty, t_\infty]$.
To do so, we use the fact that if the solution
is defined backward to a time-$t$ slice then
the length distortion bound, along with the pointpicking,
implies that a point $m$ in a time-$t$ slice with $\overline{R}_\infty(m)
> 3$ has a neighborhood that is $\epsilon$-close to the corresponding
subset of a $\kappa$-solution. The ancient solution is $\kappa$-noncollapsed
at all scales since the original solution was $\kappa$-noncollapsed at
some scale, by Theorem \ref{8.2thm}.
Then we obtain smooth convergence
of parabolic regions of the points $(\overline{x}_k, \overline{t}_k)$ to
the $\kappa$-solution, which is a contradiction to the choice of the
$(\overline{x}_k, \overline{t}_k)$'s.
\end{proof}
\begin{proof}
  \proves{thm:kl-later-scalar-curvature-bounds-bigger-balls-from-curvature-volume-bounds-positivity-scalar-curvature}

We now know that regions of high scalar curvature are modeled by
corresponding
regions in $\kappa$-solutions. To continue with the proof of the theorem,
fix $A$ large.
Suppose that the theorem is not true.  Then there are \\
1. Numbers $\rho_k \rightarrow 0$, \\
2. Numbers $r_{0,k}$ with
$r_{0,k}^{2} \: \Phi(r_{0,k}^{-2}) \: \le \: \rho_k$,
\\
3. Solutions
$(M_k, g_k(\cdot))$ defined for $0 \le t \le r_{0,k}^2$, \\
4. Points $x_{0,k} \in M_k$ and \\
5. Points $x_k \in M_k$\\
so that \\
a. $|\Rm|(x, t) \le r_{0,k}^{-2}$ for
all $(x, t) \in M_k \times [0, r_{0,k}^2]$
satisfying $\dist_0(x, x_{0,k}) < r_{0,k}$, \\
b. The volume of the metric ball $B(x_{0,k}, r_{0,k})$ at
time zero is at least $A^{-1} r_{0,k}^3$ and \\
c.  $\dist_{r_{0,k}^2}(x_k, x_{0,k}) <
A r_{0,k}$, but \\
d. $r_{0,k}^2 \: R(x_k, r_{0,k}^2) \rightarrow \infty$.

We now apply Step 2 of the proof of Theorem \ref{thmI.12.1} to obtain a contradiction.
That is, we take a subsequence of
$\{(M_k, x_{0,k}, r_{0,k}^{-2} \: g_k(r_{0,k}^2))\}_{k=1}^\infty$ that converges on a maximal ball.
The only difference is that in Theorem \ref{thmI.12.1}, the nonnegative curvature
of $W$ came from the $\Phi$-almost nonnegative curvature assumption on the
original manifolds $M_k$ along with the fact (with the notation of the proof of
Theorem \ref{thmI.12.1})
that the numbers $Q_k = R(x_k, t_k)$, which
we used to rescale, go to infinity.
In the present case
the rescaled scalar curvatures $r_{0,k}^2 \: R(x_{0,k}, r_{0,k}^2)$ at the
basepoints $x_{0,k}$
stay bounded.
However, if a point $y \in W$ is a limit of points
$\widetilde{x}_k \in M_k$ then the equations
\begin{equation}
\Rm(\widetilde{x}_k, r_{0,k}^2) \: \ge \: - \: \Phi(R(\widetilde{x}_k, r_{0,k}^2))\end{equation} in the form
\begin{equation}
r_{0,k}^2 \: \Rm(\widetilde{x}_k, r_{0,k}^2) \: \ge \: - \:
\frac{\Phi \left(  r_{0,k}^{2}
R(\widetilde{x}_k, r_{0,k}^2) \cdot r_{0,k}^{-2}
\right)}{ R(\widetilde{x}_k, r_{0,k}^2)} \:
r_{0,k}^2 \: R(\widetilde{x}_k, r_{0,k}^2)
\end{equation}
pass to the limit to give $\Rm(y) \: \ge 0$
(using that $y \in W$, so $r_{0,k}^2 \: \Rm(\widetilde{x}_k, r_{0,k}^2) \rightarrow
R(y) > 0$).
This is enough to carry out the argument.
\end{proof}

\section{I.12.3. Earlier scalar curvature bounds on smaller balls from
lower curvature bounds and volume bounds}
\label{I.12.3}

The main result of this section
says that if one has a lower bound on volume and sectional curvature
on a ball at a certain time then one obtains an upper scalar curvature bound on
a smaller ball at an earlier time.

We first prove a result in Riemannian geometry saying that under
certain hypotheses, metric balls have subballs of a controlled size
with almost-Euclidean volume.

\begin{lemma}[existence for scalar curvature]
  \label{lem:kl-earlier-scalar-curvature-bounds-smaller-balls-from-lower-curvature-bound-existence-scalar-curvature}
  \dcref{Lemma 54.1}
  \group{Kleiner--Lott Canonical Neighborhoods}
  \source{kleiner-lott}
 \label{Alexlemma}
Given $w^\prime > 0$ and $n \in \Z^+$, there is a number $c = c(w^\prime,n) > 0$ with the
following property.
Let $B$ be a radius-$r$ ball with compact closure in an $n$-dimensional Riemannian
manifold. Suppose that the sectional curvatures of $B$ are bounded below by
$- \: r^{-2}$.
Suppose that $\vol(B) \ge w^\prime r^n$. Then there is a subball $B^\prime \subset B$ of
radius $r^\prime \ge cr$ so that $\vol(B^\prime) \: \ge \: \frac12 \: \omega_n \:
(r^\prime)^n$, where $\omega_n$ is the volume of the unit ball in $\R^n$.
\end{lemma}
\begin{proof}
Suppose that the lemma is not true. Rescale
so that $r \: = \: 1$. Then there is a sequence of Riemannian
manifolds $\{M_i\}_{i=1}^\infty$ with balls $B(x_i, 1) \subset M_i$
having compact closure
so that $\Rm \Big|_{B(x_i, 1)} \: \ge \: - \: 1$ and
$\vol(B(x_i, 1)) \: \ge \: w^\prime$, but with the property that all balls
$B(x_i^\prime, r^\prime) \subset B(x_i, 1)$ with $r^\prime \: \ge \:
i^{-1}$ satisfy
$\vol(B(x_i^\prime, r^\prime)) \: < \: \frac{1}{2} \: \omega_n
(r^\prime)^n$. After taking a subsequence, we can assume that
$\lim_{i \rightarrow \infty} (B(x_i, 1), x_i) \: = \: (X, x_\infty)$ in the
pointed Gromov-Hausdorff topology. From
\cite[Theorem 10.8]{Burago-Gromov-Perelman}, the Riemannian volume
forms $\dvol_{M_i}$ converge weakly to the three-dimensional Hausdorff
measure $\mu$ of $X$. From \cite[Corollary 6.7 and
Section 9]{Burago-Gromov-Perelman}, for any $\epsilon > 0$,
there are small balls $B(x_\infty^\prime,
r^\prime) \subset X$ with compact closure in $X$ such that $\mu(B(x_\infty^\prime,
r^\prime)) \: \ge \: (1-\epsilon) \: \omega_n (r^\prime)^n$.
This gives a contradiction.
\end{proof}

\begin{theorem}[estimates for scalar curvature]
  \label{thm:kl-earlier-scalar-curvature-bounds-smaller-balls-from-lower-curvature-bound-estimates-scalar-curvature}
  \dcref{Theorem 54.2}
  \group{Kleiner--Lott Canonical Neighborhoods}
  \source{kleiner-lott}
 (cf. Theorem I.12.3) \label{thmI.12.3}
For any $w > 0$ there exist $\tau = \tau(w) > 0$, $K = K(w) < \infty$ and
$\rho = \rho(w) > 0$ with the following property.  Suppose that
$g(\cdot)$ is a Ricci flow on a closed three-manifold $M$, defined for
$t \in [0, T)$, with $\Phi$-almost nonnegative curvature. Let $(x_0, t_0)$
be a spacetime point and let $r_0 > 0$ be a radius with $t_0 \ge 4 \tau r_0^2$
and
$r_0^2 \: \Phi(r_0^{-2}) \: < \: \rho$.
Suppose that $\vol_{t_0}(B(x_0,r_0)) \ge w r_0^3$ and
the time-$t_0$ sectional curvatures on $B(x_0, r_0)$ are bounded below by
$- \: r_0^{-2}$. Then $R(x,t) \le K r_0^{-2}$ whenever
$t \in [t_0 - \tau r_0^2, t_0]$ and
$\dist_t(x, x_0) \: \le \: \frac14 \: r_0$.
\end{theorem}
\begin{proof}
  \uses{cor:kl-curvature-bounds-ricci-flow-solutions-nonnegative-curvature-operator-ass-existence-curvature-operators-3,lem:kl-earlier-scalar-curvature-bounds-smaller-balls-from-lower-curvature-bound-existence-scalar-curvature,rem:kl-later-scalar-curvature-bounds-bigger-balls-from-curvature-volume-bounds-existence-noncollapsing,thm:kl-later-scalar-curvature-bounds-bigger-balls-from-curvature-volume-bounds-positivity-scalar-curvature}
Let $\tau_0(w)$, $B(w)$ and $C(w)$ be the constants of Corollary
\ref{I.11.6end}. Put $\tau(w) \: = \: \frac12 \: \tau_0(w)$ and
$K(w) = C(w) + 2 \frac{B(w)}{\tau_0(w)}$. The function $\rho(w)$
will be specified in the course of the proof.

Suppose that the theorem is not true. Take a counterexample with
a point $(x_0, t_0)$ and a radius $r_0 > 0$ such that the
time-$t_0$ ball $B(x_0, r_0)$ satisfies
the assumptions of the theorem, but the conclusion of the theorem
fails.  We claim there is a counterexample coming from a point
$(\widehat{x}_0, \widehat{t}_0)$
and a radius $\widehat{r}_0 > 0$, with the additional property that for any
$(x^\prime_0, t^\prime_0)$ and $r^\prime_0$ having
$t^\prime_0 \in [\widehat{t}_0 - 2 \tau \widehat{r}_0^2, \widehat{t}_0]$ and $r^\prime_0 \le
\frac{1}{2} \widehat{r}_0$, if $\vol_{t_0^\prime}(B(x_0^\prime,r_0^\prime)) \ge w
(r_0^\prime)^3$ and
the time-$t_0^\prime$ sectional curvatures on $B(x_0^\prime, r_0^\prime)$
are bounded below by
$- \: (r_0^\prime)^{-2}$ then $R(x,t) \le K (r_0^\prime)^{-2}$ whenever
$t \in [t_0^\prime - \tau (r_0^\prime)^2, t_0^\prime]$ and
$\dist_t(x, x_0^\prime) \: \le \: \frac14 \: r_0^\prime$.
This follows from a
pointpicking argument - suppose that it is not
true for the original $x_0, t_0, r_0$.  Then there are
$(x^\prime_0, t^\prime_0)$ and $r^\prime_0$ with
$t^\prime_0 \in [t_0 - 2 \tau r_0^2, t_0]$ and $r^\prime_0 \le
\frac{1}{2} r_0$, for which
the assumptions of the theorem hold but the conclusion
does not. If the triple
$(x^\prime_0, t^\prime_0, r^\prime_0)$ satisfies the claim then we
stop, and otherwise we iterate the procedure.  The iteration must
terminate, which provides the desired triple
$(\widehat{x}_0, \widehat{t}_0, \widehat{r}_0)$. Note that
$\widehat{t}_0 \: > \: t_0 - 4 \tau r_0^2 \: \ge \: 0$.

We relabel $(\widehat{x}_0, \widehat{t}_0, \widehat{r}_0)$ as
$({x}_0, {t}_0, {r}_0)$.
For simplicity, let us assume that
the time-$t_0$ sectional curvatures on $B(x_0, r_0)$ are strictly
greater than
$- \: r_0^{-2}$; the general case will follow from continuity.
Let $\tau^\prime > 0$ be the largest number such that
$\Rm(x,t) \: \ge \: -\: r_0^{-2}$ whenever
$t \in [t_0 - \tau^\prime r_0^2, t_0]$ and $\dist_t(x, x_0) \le r_0$.
If $\tau^\prime \ge 2 \tau = \tau_0(w)$ then Corollary
\ref{I.11.6end} implies that
$R(x,t) \le C r_0^{-2} + B (t-t_0+ 2 \tau r_0^2)^{-1}$ whenever
$t \in [t_0 - 2 \tau r_0^2, t_0]$ and
$\dist_t(x, x_0) \: \le \: \frac14 \: r_0$.
In particular, $R(x,t) \le K$ whenever
$t \in [t_0 - \tau r_0^2, t_0]$ and
$\dist_t(x, x_0) \: \le \: \frac14 \: r_0$, which contradicts our
assumption that the conclusion of the theorem fails.

Now suppose that $\tau^\prime < 2 \tau$. Put
$t^\prime = t_0 - \tau^\prime r_0^2$.
From estimates on
the length and volume distortion under the Ricci flow, we know that
there are numbers $\alpha = \alpha(w) > 0$ and
$w^\prime = w^\prime(w) > 0$ so that the time-$t^\prime$ ball
$B(x_0, \alpha r_0)$ has volume at least $w^\prime (\alpha r_0)^3$.
From Lemma \ref{Alexlemma}, there is a subball
$B(x^\prime, r^\prime) \subset B(x_0, \alpha r_0)$ with
$r^\prime \ge c \alpha r_0$ and $\vol(B(x^\prime, r^\prime)) \: \ge \:
\frac12 \: \omega_3 \: (r^\prime)^3$. From the preceding
pointpicking argument, we have the estimate $R(x,t) \: \le \: K
(r^\prime)^{-2}$ whenever $t \in [t^\prime - \tau (r^\prime)^2,
t^\prime]$ and $\dist_t(x, x^\prime) \: \le \: \frac14 r^\prime$.
From the $\Phi$-almost nonnegative curvature,
we have a bound
$|\Rm|(x,t) \: \le \:
\const \: K \: (r^\prime)^{-2} \: + \: \const \: \Phi( K (r^\prime)^{-2})$
at
such a point (see (\ref{double})).  If $\rho(w)$ is taken sufficiently
small then we can ensure that $r_0$ is small enough, and
hence $r^\prime$ is small enough, to make
$K \: (r^\prime)^{-2} \: + \: \Phi( K (r^\prime)^{-2}) \: \le \:
2 \: K \: (r^\prime)^{-2}$.
Then we can apply
Theorem \ref{thmI.12.2} to a time interval ending at time
$t^\prime$, after a redefinition of its
constants, to obtain a bound of the form
$R(x, t^\prime) \: \le \: K^\prime (r^\prime)^{-2}$
whenever $\dist_{t}(x, x^\prime) \: \le \: 10 r_0$, where
$K^\prime$ is related to the constant
$K$ of Theorem \ref{thmI.12.2}. (We also obtain a similar estimate at
times slightly less than $t^\prime$.)  Thus at such a point,
$\Rm(x, t^\prime) \: \ge \: - \:
\Phi(K^\prime (r^\prime)^{-2})$.  If we choose $\rho(w)$ to be
sufficiently small to force $r^\prime$ to be sufficiently small to force
$- \:
\Phi(K^\prime (r^\prime)^{-2})\: > \: - \: r_0^{-2}$
then we have
$\Rm > - r_0^{-2}$ on $\overline{B_{t^\prime}(x_0, r_0)} \subset
B_{t^\prime}(x^\prime, 10r_0)$, which
contradicts the assumed maximality of $\tau^\prime$.

We note that in the application of Theorem \ref{thmI.12.2}
at the end of the proof, we must take into account the
extra hypothesis, in the notation of
Theorem \ref{thmI.12.2}, that
$r_0^2 \: \Phi(r_0^{-2}) \: < \: \rho$
(see Remark \ref{rmkI.12.2}).  This
will be satisfied if the $r_0$ in Theorem \ref{thmI.12.2} is small
enough, which is ensured
by taking the $\rho$ of Theorem \ref{thmI.12.3} small enough.
\end{proof}

\section{I.12.4. Small balls with strongly negative curvature are
volume-collapsed}
\label{I.12.4}

In this section we show that under certain hypotheses, if the infimal sectional
curvature on an $r$-ball is exactly $- \: r^{-2}$ then the volume of the ball
is small compared to $r^3$.

\begin{corollary}[positivity for curvature]
  \label{cor:kl-small-balls-strongly-negative-curvature-volume-collapsed-positivity-curvature}
  \dcref{Corollary 55.1}
  \group{Kleiner--Lott Canonical Neighborhoods}
  \source{kleiner-lott}
 (cf. Corollary I.12.4) \label{corI.12.4}
For any $w > 0$, one can find $\rho > 0$ with the following property.
Suppose that $g(\cdot)$ is a $\Phi$-almost nonnegatively curved
Ricci flow solution on a closed three-manifold $M$, defined for
$t \in [0, T)$ with $T \ge 1$. If $B(x_0, r_0)$ is a metric ball
at time $t_0 \ge 1$ with $r_0 < \rho$ and if
$\inf_{x \in B(x_0,r_0)}  \Rm(x, t_0)  = - r_0^{-2}$ then
$\vol(B(x_0, r_0)) \le w r_0^3$.
\end{corollary}
\begin{proof}
  \uses{thm:kl-earlier-scalar-curvature-bounds-smaller-balls-from-lower-curvature-bound-estimates-scalar-curvature,thm:kl-later-scalar-curvature-bounds-bigger-balls-from-curvature-volume-bounds-positivity-scalar-curvature}
Fix $w > 0$.
The number $\rho$ will be specified in the course of the proof.
Suppose that the corollary is not true, i.e.
there is a Ricci flow solution as in the statement of the corollary along
with a metric ball $B(x_0, r_0)$ at a time $t_0 \ge 1$
so that  $\inf_{x \in B(x_0,r_0)}  \Rm(x, t_0)  = - r_0^{-2}$
and $\vol(B(x_0, r_0)) \: > \: w r_0^n$.
The idea is to use
Theorem \ref{thmI.12.3}, along with the $\Phi$-almost nonnegative
curvature, to get a double-sided sectional
curvature bound on a smaller
ball at an earlier time. Then one goes forward in time using
Theorem \ref{thmI.12.2}, along with the $\Phi$-almost nonnegative
curvature, to get a lower sectional curvature bound on the
original ball, thereby obtaining a contradiction.

Looking at the hypotheses of Theorem \ref{thmI.12.3}, if
we require $r_0 < (4\tau)^{- \: \frac12}$ then
$4 \tau r_0^2 < 1 \le t_0$.
From Theorem \ref{thmI.12.3}, $R(x, t) \: \le \: K r_0^{-2}$
whenever  $t \in [t_0 - \tau r_0^2, t_0]$
and $\dist_t(x, x_0) \: \le \: \frac{1}{4} r_0$, provided that $r_0$ is
small enough that
$r_0^2 \: \Phi(r_0^{-2})$
is less than the $\rho$ of Theorem \ref{thmI.12.3}.
If in addition $r_0$ is sufficiently small then it follows that
$|\Rm(x,t)| \: \le \: \const \: \Phi(K r_0^{-2}) \: \le \: r_0^{-2}$.

From the Bishop-Gromov inequality and the bounds on length and
volume distortion under Ricci flow, there is a small number $c$ so
that we are ensured that $|\Rm(x,t)| \le (cr_0)^{-2}$ for all $(x,t)$
satisfying $\dist_{t_0 - (cr_0)^2}(x, x_0) < cr_0$ and
$t \in [t_0 - (cr_0)^2, t_0]$, and in addition the volume of
$B(x_0, cr_0)$ at time $t_0 - (cr_0)^2$ is at least
$c(cr_0)^3$. Choosing the constant $A$ of Theorem \ref{thmI.12.2}
appropriately in terms of $c$, we can apply Theorem \ref{thmI.12.2}
to the ball $B(x_0, cr_0)$ and the time interval $[t_0 - (cr_0)^2, t_0]$
to conclude that at time $t_0$,
$R(\cdot, t_0) \Big|_{B(x_0, r_0)} \: \le \: K(A) \:
(cr_0)^{-2}$, where $K(A)$ is as in the statement of
Theorem \ref{thmI.12.2}.
From the $\Phi$-almost nonnegative curvature condition,
\begin{equation}
\Rm \Big|_{B(x_0, r_0)} \: \ge \: - \:
\Phi \left(K(A) \:
(cr_0)^{-2} \right).
\end{equation}
If $r_0$ is sufficiently small then we contradict the assumption that
$\inf \Rm(x, t_0)
\Big|_{B(x_0,r_0)} \: = \: - \:
\frac{1}{r_0^2}$.
\end{proof}

\section{I.13.1. Thick-thin decomposition for nonsingular flows} \label{I.13.1}

The main result of this section says that if $g(\cdot)$ is a Ricci flow
solution on a closed oriented three-dimensional manifold $M$ that
exists for $t \in [0, \infty)$ then for large $t$, $(M, g(t))$ has a
thick-thin decomposition.
A fuller description is in Sections \ref{II.7.1}-\ref{II.7.4}.

We assume that at time zero, the sectional curvatures are bounded below
by $-1$. This can always be achieved by rescaling the initial
metric.  Then we have the $\Phi$-almost nonnegative curvature result of
(\ref{phi}).

If the metric $g(t)$ has nonnegative sectional curvature then it must be flat, as we are
assuming that the Ricci flow exists for all time. Let us assume that $g(t)$ is not flat, so
it has some negative sectional curvature.
Given $x \in M$, consider the time-$t$ ball $B_t(x,r)$. Clearly if $r$ is
sufficiently small then
$\Rm \Big|_{B_t(x,r)} \: > \: - \: r^{-2}$, while if $r$ is sufficiently large
(maybe greater than the diameter of $M$) then it is not true that
$\Rm \Big|_{B_t(x,r)} \: > \: - \: r^{-2}$.
Let $\widehat{r}(x,t) > 0$ be the unique number such that
$\inf \Rm \Big|_{B_t(x,\widehat{r})} \: = \: - \: \widehat{r}^{-2}$.
Let $M_{thin}(w,t)$ be the set of points $x \in M$ for which
\begin{equation} \label{thinpart}
\vol(B_t(x, \widehat{r}(x,t))) \: < \: w \: \widehat{r}(x,t)^3.
\end{equation}
Put $M_{thick}(w,t) = M - M_{thin}(w,t)$.

As the statement of (\ref{phi}) is invariant under parabolic rescaling
(although we must take $t \ge t_0$ for (\ref{phi}) to apply), if $t \ge t_0$
and we are interested in the Ricci flow at time $t$ then we can apply
Theorem \ref{thmI.12.2}, Theorem \ref{thmI.12.3} and Corollary \ref{corI.12.4}
to the rescaled flow $g(t^\prime) \: = \: t^{-1} g(tt^\prime)$.
From Corollary \ref{corI.12.4},
for any $w > 0$ we can find $\widehat{\rho} = \widehat{\rho}(w) > 0$
so that if $\widehat{r}(x,t) < \widehat{\rho} \sqrt{t}$ then $x \in M_{thin}(w,t)$,
provided that $t$ is
sufficiently large (depending on $w$). Equivalently, if
$t$ is sufficiently large (depending on $w$) and $x \in M_{thick}(w,t)$ then
$\widehat{r}(x,t) \ge \widehat{\rho} \sqrt{t}$.

\begin{theorem}[existence for thick--thin geometry]
  \label{thm:kl-thick-thin-decomposition-nonsingular-flows-existence-thick-thin-geometry}
  \dcref{Theorem 56.2}
  \group{Kleiner--Lott Thick--Thin Decomposition}
  \source{kleiner-lott}
 (cf. I.13.1) \label{thmI.13.1}
There are numbers $T = T(w) > 0$, $\overline{\rho} = \overline{\rho}(w) > 0$
and $K = K(w) < \infty$ so that if $t \ge T$ and
$x \in M_{thick}(w,t)$ then $|\Rm| \le K t^{-1}$ on $B_t(x, \overline{\rho} \sqrt{t})$,
and $\vol(B_t(x, \overline{\rho} \sqrt{t})) \: \ge \: \frac{1}{10} \:
w \left( \overline{\rho} \sqrt{t} \right)^3$.
\end{theorem}
\begin{proof}
  \uses{cor:kl-small-balls-strongly-negative-curvature-volume-collapsed-positivity-curvature,thm:kl-earlier-scalar-curvature-bounds-smaller-balls-from-lower-curvature-bound-estimates-scalar-curvature,thm:kl-later-scalar-curvature-bounds-bigger-balls-from-curvature-volume-bounds-positivity-scalar-curvature}
The method of proof is the same as in Corollary \ref{corI.12.4}.  By assumption,
$\Rm \Big|_{B_t(x, \widehat{r}(x,t))} \ge - \: \widehat{r}(x,t)^{-2}$ and
$\vol(B_t(x, \widehat{r}(x,t))) \: \ge \: w \: \widehat{r}(x,t)^3$.   As
$\widehat{r}(x,t) \ge \widehat{\rho} \sqrt{t}$, for any $c \in (0,1)$ we have
$\Rm \Big|_{B_t(x, c\widehat{\rho} \sqrt{t})} \ge - \: (c\widehat{\rho})^{-2} t^{-1}$. By
the Bishop-Gromov inequality,
\begin{align}
\vol(B_t(x, c\widehat{\rho} \sqrt{t})) \: & \ge \:
\frac{
\int_0^{ \frac{c \widehat{\rho} \sqrt{t}}{\widehat{r}(x,t)}} \sinh^2(u) \: du}{\int_0^1 \sinh^2(u) \: du}
\: w \: \widehat{r}(x,t)^3
\: \ge \: \frac{1}{3\int_0^1 \sinh^2(u) \: du} \: w \:
(c\widehat{\rho})^3 \: t^{\frac32} \\
& \ge \:  \frac{1}{10} \: w \:
(c\widehat{\rho})^3 \: t^{\frac32}. \notag
\end{align}
Considering Theorem \ref{thmI.12.3} with its $w$ replaced by $\frac{w}{10}$,
if $c = c(w)$ is taken sufficiently small
(to ensure $t \ge 4 \tau (c \widehat{\rho} \sqrt{t})^2$)
and $t$ is larger than a certain
$w$-dependent constant
(to ensure
$\frac{\Phi((c \widehat{\rho} \sqrt{t})^{-2})}{(c
\widehat{\rho} \sqrt{t})^{-2}} <\rho$)
then we can apply Theorem \ref{thmI.12.3} with
$r_0 =  c \widehat{\rho} \sqrt{t}$ to obtain
$R(x^\prime,t^\prime) \le K^\prime(w) c^{-2} \widehat{\rho}^{-2} t^{-1}$  whenever
$t^\prime \in [t - \tau c^2 \widehat{\rho}^2 t, t]$ and
$\dist_{t^\prime}(x^\prime, x) \: \le \: \frac14 \: c \widehat{\rho} \sqrt{t}$. From the
$\Phi$-almost nonnegative curvature (see (\ref{double})),
\begin{equation}
|\Rm|(x^\prime, t^\prime) \: \le \: \const \:
K^\prime c^{-2} \widehat{\rho}^{-2} t^{-1} \: + \: \const \:
\Phi(K^\prime c^{-2} \widehat{\rho}^{-2} t^{-1}),
\end{equation}
which is bounded above by $2 \: \const \: K^\prime c^{-2} \widehat{\rho}^{-2} t^{-1}$
if $t$ is larger than a certain
$w$-dependent constant. Then from length and volume distortion estimates
for the Ricci flow,
we obtain a lower volume bound
$\vol(B_{t^\prime}(x, c^\prime \widehat{\rho} \sqrt{t}))   \: \ge \:
w^{\prime} ( c^\prime \widehat{\rho} \sqrt{t})^3$ on a smaller ball of
controlled radius, for
some $c^\prime = c^\prime(w)$.  Using Theorem \ref{thmI.12.2}, we finally
obtain an upper bound $R \le K^{\prime \prime}(w) (c\widehat{\rho} \sqrt{t})^{-2}$
on  $B_t(x, c\widehat{\rho} \sqrt{t})$ and hence,
by the $\Phi$-almost nonnegative curvature, an upper bound of the form
$|\Rm| \le K(w) \: t^{-1}$ on $B_t(x, c\widehat{\rho} \sqrt{t})$,
provided that $t \ge T$ for an appropriate $T = T(w)$.
Taking $\overline{\rho} = c \widehat{\rho}$, the theorem follows.
\end{proof}

\begin{remark}[estimates for thick--thin geometry]
  \label{rem:kl-thick-thin-decomposition-nonsingular-flows-estimates-thick-thin-geometry}
  \dcref{Remark 56.5}
  \group{Kleiner--Lott Thick--Thin Decomposition}
  \source{kleiner-lott}
  \uses{thm:kl-later-scalar-curvature-bounds-bigger-balls-from-curvature-volume-bounds-positivity-scalar-curvature,thm:kl-thick-thin-decomposition-nonsingular-flows-existence-thick-thin-geometry}

The use of Theorem \ref{thmI.12.2} in the proof of Theorem \ref{thmI.13.1}
also gives an upper bound $|\Rm|(x,t) \le K(A,w) \: t^{-1}$ on
$B_t(x, A \overline{\rho} \sqrt{t})$ if $x \in M_{thick}(w,t)$, for any $A > 0$.
\end{remark}

We now take $w$ sufficiently small. Then for large $t$, $M_{thick}(w,t)$ has a boundary
consisting of tori that are incompressible in $M$ and the interior of
$M_{thick}(w,t)$ admits a complete Riemannian metric with constant
sectional curvature $- \: \frac14$ and finite volume; see
Sections
\ref{hyprig} and \ref{inctori}. In addition,
$M_{thin}(w,t)$ is a graph manifold; see
Section \ref{II.7.4}.
